% !TeX program = xelatex
% 数学 LaTeX 排版 Skill v4.0.0；版式保持 v3 金标准. XeLaTeX 编译至少两次. 不分发字体.
% WZ-LOCKED-CLASS-BEGIN
\begin{filecontents*}[overwrite]{elegantbook-original-adapter.cls}
\NeedsTeXFormat{LaTeX2e}
\ProvidesClass{elegantbook-original-adapter}[2026/09/23 v3.0 fixed mathematical book environments]
\LoadClass[10pt,letterpaper,oneside]{book}
\RequirePackage[UTF8,scheme=plain,fontset=fandol]{ctex}
\RequirePackage{fontspec}
\setmainfont{cmunrm.otf}[BoldFont=cmunbx.otf,ItalicFont=cmunti.otf,BoldItalicFont=cmunbi.otf]
\setCJKfamilyfont{sourceheader}[ItalicFont=FandolKai-Regular.otf]{FandolKai-Regular.otf}
\RequirePackage{amsmath,amssymb,mathrsfs}
\RequirePackage{amsthm,etoolbox}
\RequirePackage{xcolor,geometry,setspace,fancyhdr,enumitem,needspace,xurl}
\definecolor{structurecolor}{HTML}{880E4F}
\definecolor{main}{HTML}{9C27B0}
\geometry{paperwidth=612bp,paperheight=792bp,left=20mm,right=20mm,top=95.04bp,bottom=20mm,headheight=12pt,headsep=25pt,footskip=30pt}
\setstretch{1.6}
\setlength{\parindent}{2em}
\setlength{\parskip}{0pt}
\setlist{nosep,leftmargin=2em}
\AtBeginDocument{\setlength{\abovedisplayskip}{4pt plus 1pt minus 1pt}\setlength{\belowdisplayskip}{4pt plus 1pt minus 1pt}\setlength{\abovedisplayshortskip}{2pt}\setlength{\belowdisplayshortskip}{3pt}}
\fancyhf{}
\fancyhead[L]{{\itshape\CJKfamily{sourceheader}\rightmark}}
\fancyhead[R]{{\itshape\CJKfamily{sourceheader}\leftmark}}
\fancyfoot[C]{\thepage}
\renewcommand{\headrulewidth}{0.4pt}
\renewcommand{\headrule}{{\color{structurecolor}\hrule height\headrulewidth width\headwidth\vskip-\headrulewidth}}
\pagestyle{fancy}
\fancypagestyle{cover}{\fancyhf{}\fancyfoot[C]{\thepage}\renewcommand{\headrulewidth}{0pt}\renewcommand{\headrule}{}}
\RequirePackage{hyperref}
\hypersetup{unicode=true,colorlinks=true,linkcolor=structurecolor,urlcolor=structurecolor,citecolor=structurecolor,linktoc=section,bookmarksnumbered=true,bookmarksopen=true,pdfborder={0 0 0},pdfstartview=FitH}
\setcounter{tocdepth}{1}
\renewcommand*\l@section{\@dottedtocline{1}{1.5em}{2.4em}}
\renewcommand{\thesection}{\thechapter.\arabic{section}}
\newcommand{\originalchapter}[1]{\par\addvspace{14pt}\Needspace{5\baselineskip}\refstepcounter{chapter}\setcounter{section}{0}\addcontentsline{toc}{chapter}{\protect\numberline{\thechapter}#1}\markboth{\thechapter\quad #1}{}\begingroup\centering\color{structurecolor}\bfseries\fontsize{14.4}{20}\selectfont\thechapter\quad #1\par\endgroup\nobreak\vspace{8pt}}
\newcommand{\originalsection}[1]{\par\addvspace{9pt}\Needspace{4\baselineskip}\refstepcounter{section}\addcontentsline{toc}{section}{\protect\hspace*{1.5em}\protect\numberline{\protect\textcolor{structurecolor}{\thesection}}#1}\markright{\thesection\quad #1}\noindent{\fontsize{12}{17}\selectfont\color{structurecolor}\thesection\quad}{\bfseries\fontsize{12}{17}\selectfont #1}\par\nobreak\vspace{4pt}}

% Semantic environments are registered once here, not re-designed in manuscripts.
\newtheoremstyle{wzstatement}{6pt}{5pt}
  {\normalfont\kaishu}{0pt}{\normalfont\bfseries\color{structurecolor}}
  {}{.7em}{\thmname{#1}\thmnumber{ #2}\thmnote{\normalfont\ (#3)}}
\newtheoremstyle{wzdefinition}{6pt}{5pt}
  {\normalfont\songti}{0pt}{\normalfont\bfseries\color{structurecolor}}
  {}{.7em}{\thmname{#1}\thmnumber{ #2}\thmnote{\normalfont\ (#3)}}
\newtheoremstyle{wzproblem}{7pt}{5pt}
  {\normalfont\songti}{2em}{\normalfont\bfseries\color{main}}
  {}{.7em}{\thmname{#1}\thmnumber{ #2}}
\newtheoremstyle{wzremark}{4pt}{4pt}
  {\normalfont\songti}{0pt}{\normalfont\bfseries}
  {}{.7em}{\thmname{#1}}
\theoremstyle{wzstatement}
\newtheorem{theorem}{定理}[chapter]
\newtheorem{lemma}[theorem]{引理}
\newtheorem{proposition}[theorem]{命题}
\newtheorem{corollary}[theorem]{推论}
\theoremstyle{wzdefinition}
\newtheorem{definition}[theorem]{定义}
\theoremstyle{wzproblem}
\newtheorem{example}{例题}[chapter]
\newtheorem{exercise}{习题}[chapter]
\theoremstyle{wzremark}
\newtheorem*{remark}{注}
\renewcommand{\proofname}{证明}
% Retain the native amsthm QED stack; only adjust the fixed typography.
\expandafter\patchcmd\csname\string\proof\endcsname{\normalfont \topsep6\p@\@plus6\p@\relax}
  {\normalfont\songti \topsep5\p@\@plus1\p@\relax}
  {}{\ClassError{elegantbook-original-adapter}{Cannot patch proof spacing}{Check the amsthm version.}}
\expandafter\patchcmd\csname\string\proof\endcsname{\itshape}{\normalfont\bfseries}
  {}{\ClassError{elegantbook-original-adapter}{Cannot patch proof heading}{Check the amsthm version.}}
\expandafter\patchcmd\csname\string\proof\endcsname{\ignorespaces}
  {\setlength{\parindent}{2em}\ignorespaces}
  {}{\ClassError{elegantbook-original-adapter}{Cannot patch proof paragraphs}{Check the amsthm version.}}
\AtBeginEnvironment{proof}{\par\Needspace{3\baselineskip}}
\renewcommand{\qedsymbol}{\ensuremath{\square}}
\newenvironment{solution}{\begin{proof}[解答]}{\end{proof}}
\newcommand{\wzcheckchapter}{\ifnum\value{chapter}<1
  \ClassError{elegantbook-original-adapter}{Numbered mathematics before chapter 1}{Move it after the first originalchapter.}\fi}

\newcommand{\wzregisterguard}[1]{\AtBeginEnvironment{#1}{\wzcheckchapter\par\Needspace{3\baselineskip}}}
\forcsvlist{\wzregisterguard}{theorem,lemma,proposition,corollary,definition,example,exercise}
% Front matter and bibliography are not numbered mathematical chapters.
\newcommand{\originalpreface}{\par\addvspace{10pt}\Needspace{5\baselineskip}
  \noindent{\bfseries\fontsize{12}{17}\selectfont 前言}\par\nobreak\vspace{4pt}}
\newcommand{\originalcontents}{\par\addvspace{14pt}
  \begin{center}{\color{structurecolor}\bfseries\fontsize{14.4}{20}\selectfont 目录}\end{center}
  \vspace{-10pt}\@starttoc{toc}}
\renewcommand{\bibname}{参考文献}
\renewenvironment{thebibliography}[1]{
  \par\addvspace{12pt}\Needspace{3\baselineskip}\phantomsection
  \addcontentsline{toc}{chapter}{\bibname}
  \markboth{\bibname}{}
  \noindent{\color{structurecolor}\bfseries\fontsize{14.4}{20}\selectfont\bibname}\par\nobreak\vspace{6pt}
  \list{\@biblabel{\@arabic\c@enumiv}}
    {\settowidth\labelwidth{\@biblabel{#1}}\leftmargin\labelwidth
     \advance\leftmargin\labelsep\usecounter{enumiv}
     \let\p@enumiv\@empty\renewcommand\theenumiv{\@arabic\c@enumiv}
     \setlength{\itemsep}{3pt}\setlength{\parsep}{0pt}}
  \sloppy\clubpenalty4000\widowpenalty4000\sfcode`\.=1000\relax
}{\def\@noitemerr{\@latex@warning{Empty bibliography}}\endlist}

\numberwithin{equation}{chapter}
\renewcommand{\theequation}{\thechapter.\arabic{equation}}
\allowdisplaybreaks[2]
\emergencystretch=2em
\clubpenalty=6000
\widowpenalty=6000
\raggedbottom
\endinput
\end{filecontents*}
% WZ-LOCKED-CLASS-END
\documentclass{elegantbook-original-adapter}
% 本册有两位章、节号；仅调整目录号栏，不修改锁定类定义。
\makeatletter
\renewcommand*\l@section{\@dottedtocline{1}{1.5em}{3.8em}}
\makeatother

\usepackage{tikz,pgfplots,longtable,booktabs}
\usetikzlibrary{arrows.meta,calc,positioning,patterns}
\pgfplotsset{compat=1.18}
\newcommand{\dd}{\,\mathrm d}
\newcommand{\kb}{k_{\mathrm B}}
\newcommand{\avg}[1]{\langle #1\rangle}
\newcommand{\Tr}{\operatorname{Tr}}
\newcommand{\e}{\mathrm e}
\newcommand{\ii}{\mathrm i}
\newcommand{\vct}[1]{\boldsymbol{#1}}
\newcommand{\ket}[1]{\lvert #1\rangle}
\newcommand{\bra}[1]{\langle #1\rvert}
\newcommand{\source}[1]{\par\noindent\textit{来源与编写：#1}\par}
\newcommand{\psource}[1]{\par\noindent\textit{题源：#1。以下解答为编者独立补写。}\par}
\newcommand{\osource}[1]{\par\noindent\textit{题源：#1。解答依据公开官方解中文重编，补入推导。}\par}
% Keep a chapter heading with its source and opening paragraph.
\pretocmd{\originalchapter}{\Needspace{10\baselineskip}}{}{}
\newcommand{\volpart}[1]{\par\addvspace{16pt}\Needspace{5\baselineskip}\phantomsection\addcontentsline{toc}{part}{#1}\begin{center}{\color{structurecolor}\bfseries\fontsize{16}{22}\selectfont #1}\end{center}}

\hypersetup{pdftitle={凝聚态与量子多体：MIT OCW 中文重编},pdfauthor={中文重编项目}}
\begin{document}
\frontmatter\thispagestyle{cover}
\begin{center}
{\color{structurecolor}\bfseries\fontsize{17}{23}\selectfont 凝聚态与量子多体}\par
\vspace{5pt}{\bfseries 晶体、响应、关联电子与拓扑物态}\par
\vspace{5pt}根据 MIT OpenCourseWare 公开课程中文重编\par
2026年10月12日
\end{center}
\originalpreface
晶体中电子的运动受周期性约束，而电子彼此之间的作用又使单粒子图像出现边界。学习这两个层次，需要把能带、实验响应、低能有效理论和量子多体态接起来。本册读者应已掌握本科量子力学中的算符、微扰与角动量，以及基本热力学、统计分布和傅里叶变换；二次量子化、推迟响应和 Berry 相位在正文重新建立。

第一编以 Patrick A. Lee 的 \href{https://ocw.mit.edu/courses/8-511-theory-of-solids-i-fall-2004/}{8.511，Fall 2004} 为范围，覆盖公开的24讲主题和11套作业。该课原PDF主要是授课摘要，正文的完整计算链由编者补写，不能视为课堂逐字译稿。第二编接续 Lee 的 \href{https://ocw.mit.edu/courses/8-512-theory-of-solids-ii-spring-2009/}{8.512，Spring 2009}，覆盖25讲主题、12套当学期作业及页面所列2004年补充作业；公开讲义1--10还须署名 Mark Spencer Rudner。课程页讲次与个别PDF内部标题存在错位，来源表保留两种定位。第三编依 Xiao-Gang Wen 的 \href{https://ocw.mit.edu/courses/8-513-modern-quantum-many-body-physics-for-condensed-matter-systems-fall-2021/}{8.513，Fall 2021}，四组幻灯片覆盖1--24讲，作业页列11套题。

8.511成绩由作业70\%与一次测验30\%组成，测验原卷未公开；8.512按12次作业评分；8.513要求每周作业与约四页学期论文，未提供论文范本或官方标准解。课堂摘要引用 Marder、Ashcroft--Mermin 等商业书的部分，只保留真实定位，不把书本引用当成已经取得的题面。作业答案若无公开官方解，明确标记为编者独立补写。

原作者为上述教师及获署名的讲义作者，课程平台为 Massachusetts Institute of Technology OpenCourseWare。中文翻译、结构重编、推导补充和可编辑TikZ图均为本项目改编。基于OCW的部分遵守\href{https://creativecommons.org/licenses/by-nc-sa/4.0/}{CC BY-NC-SA 4.0}：署名、非商业、相同方式共享；不表示MIT或原作者认可本项目。8.513幻灯片中明确标注第三方例外的图片不进入本稿，也不随源码包分发。逐件链接、页数、哈希、许可例外及未取得项目见卷末和学科目录中的来源清单。

\clearpage\originalcontents\clearpage\mainmatter
\volpart{第一编\quad 固体理论的单粒子与配对基础（8.511）}
\originalchapter{电子气、晶格与衍射}
\label{ch:crystal}
\source{8.511公开摘要第1--4讲；以下推导为中文重编与补写。}
\originalsection{自由电子气的尺度}
把价电子近似为在体积$V$中运动的自旋$1/2$粒子，离子只提供均匀中和背景。这个模型首先给出可与实验比较的能量尺度。周期边界下$k_i=2\pi n_i/L_i$，每个自旋态占据$k$空间体积$(2\pi)^3/V$。零温时从低能态依次填到费米球边界，故
\begin{equation}
n=\frac NV=2\int_{|\vct k|<k_F}\frac{\dd^3k}{(2\pi)^3}=\frac{k_F^3}{3\pi^2},\qquad E_F=\frac{\hbar^2k_F^2}{2m},\qquad v_F=\frac{\hbar k_F}{m}.
\end{equation}
因$N(E)=V(2mE/\hbar^2)^{3/2}/(3\pi^2)$，两种自旋合计的单位体积态密度为$D(E)=\dd(N/V)/\dd E=3n\sqrt{E/E_F}/(2E_F)$。这里$D$按能量计，不能与每自旋态密度混用。金属的$T_F=E_F/\kb$常远高于室温，因此只有宽度约$\kb T$的费米面薄层能被热激发。

对平滑函数$g$，把费米函数与阶跃函数的差局限于该薄层，令$x=(E-\mu)/\kb T$，奇偶性使一阶修正消失，二阶矩$\int x^2[-f'(x)]\dd x=\pi^2/3$给出 Sommerfeld 展开
\begin{equation}
\int_0^\infty g(E)f(E)\dd E=\int_0^\mu g(E)\dd E+\frac{\pi^2}{6}(\kb T)^2g'(\mu)+O(T^4).
\end{equation}
先把$g=D$代入并固定$N$，得$\mu=E_F[1-\pi^2(\kb T/E_F)^2/12+O(T^4)]$；再把$g=ED$代入，化学势变化与显式温度项合并为$C_V/V=\pi^2D(E_F)\kb^2T/3$。这解释电子比热的线性项，也说明测比热能测费米能附近的态密度。

小磁场使两种自旋能量移为$E\mp\mu_BB$，两支人数差为$D(E_F)\mu_BB$。磁化强度$M=\mu_B(n_\uparrow-n_\downarrow)$，在SI制$\chi_P=\mu_0\mu_B^2D(E_F)$。此处只含自旋顺磁，轨道抗磁另需量子轨道计算。若加上动量弛豫时间$\tau$，稳态$0=-e\vct E-m\vct v/\tau$给$\sigma=ne^2\tau/m$，但$\tau$不是电子气模型内部自动确定的参数。

令每电子球体体积$4\pi(r_sa_B)^3/3=1/n$，其中$a_B=4\pi\epsilon_0\hbar^2/(me^2)$。动能随$r_s^{-2}$，库仑能随$r_s^{-1}$，所以$r_s$越大相互作用越不能忽略。这个估计只判断相对尺度，不能单凭它断言发生哪种有序态。
\originalsection{实空间与倒空间}
\begin{definition}
由三个线性无关向量$\vct a_i$生成的Bravais晶格为$\vct R=\sum_i n_i\vct a_i$，$n_i\in\mathbb Z$。在每个$\vct R$附上相同的原子基底$\vct d_\alpha$，得到一般晶体。基底与原胞是不同对象。
\end{definition}
把实空间基矢作为矩阵$A$的列，定义倒格基矢矩阵$B=2\pi A^{-T}$。这样$\vct a_i\cdot\vct b_j=2\pi\delta_{ij}$，且对任意倒格矢$\vct G=\sum m_i\vct b_i$都有$\e^{\ii\vct G\cdot\vct R}=1$。体积乘积由行列式直接得到$|\det A|\,|\det B|=(2\pi)^3$。第一Brillouin区是倒格点的Wigner--Seitz胞，体积同任意倒格原胞，并非必须画成平行六面体。

入射与散射波的动量转移$\vct Q=\vct k_f-\vct k_i$。弹性散射振幅是
\begin{equation}
\mathcal A(\vct Q)=\sum_{\vct R}\e^{-\ii\vct Q\cdot\vct R}\underbrace{\sum_\alpha f_\alpha(\vct Q)\e^{-\ii\vct Q\cdot\vct d_\alpha}}_{S(\vct Q)}.
\end{equation}
格点和要求$\vct Q=\vct G$，基底结构因子$S$则可使某些允许的倒格点没有衍射峰。不能把消光误判为Bravais晶格变了。有限样品的格点和给出宽约$1/L$的峰。Ewald作图是在$|\vct k_f|=|\vct k_i|$球面上寻找满足动量转移条件的倒格点。

\begin{center}
\begin{tikzpicture}[scale=.85,>=Stealth]
\foreach \x in {-2,...,2}\foreach \y in {-2,...,2}{\fill (\x,\y) circle(1.5pt);}
\draw[structurecolor,thick](-.5,-.5)rectangle(.5,.5);
\draw[->] (0,0)--(1,0)node[below]{$\vct b_1$};\draw[->](0,0)--(0,1)node[left]{$\vct b_2$};
\node at (0,-2.7){方格倒晶格的第一Brillouin区};
\end{tikzpicture}
\end{center}
图中的坐标按倒格常数归一化，方框是离原点最近的区域。真实晶体的长度与形状由$A$决定。
\originalsection{晶体对称与热振动}
保持晶格不变的纯旋转形成点群。若某旋转在二维垂直平面内保持二维晶格，整数基表示的矩阵迹等于$2\cos\theta$，必须是整数且位于$[-2,2]$。因而有限非平凡旋转仅允许$2,3,4,6$重轴，加上恒等的一重轴。三维晶系分类还需比较多条轴的相容性；14种Bravais晶格是对平移晶格的分类，而非所有原子结构的数目。

螺旋与滑移操作是旋转或反射加分数平移，它们可保持含基底晶体，却不必是纯点群操作。研究能带时分数平移会在Bloch态上产生相位，这正是非对称型空间群的简并与纯点群不同的原因。

热振动写成$\vct r_j=\vct R_j+\vct u_j$。若位移服从均值零的高斯分布，则$\avg{\e^{-\ii\vct Q\cdot\vct u}}=\exp[-\avg{(\vct Q\cdot\vct u)^2}/2]$。弹性相干强度因而乘$\e^{-2W}$，$2W=\avg{(\vct Q\cdot\vct u)^2}$。爱因斯坦模型中$\avg{u_x^2}=\hbar\coth(\hbar\omega_E/2\kb T)/(2M\omega_E)$，低温仍有零点振动。随机振动降低Bragg峰权重，同时产生漫散射；它并不把理想周期性的尖峰简单改成同等权重的宽峰。

若位移另有静态调制$\vct u_j=\vct u_0\cos(\vct q\cdot\vct R_j)$，把振幅展开到一阶会出现$\vct G\pm\vct q$卫星峰。$\vct q$与倒格不成有理比时即为非公度结构；这里破坏的是原周期性，不能使用仅含热高斯平均的解释。

\originalchapter{Bloch态与能带计算}
\label{ch:bands}
\source{8.511第5--10讲主题；主推导补写。}
\originalsection{平移对称与Bloch定理}
考虑$H=\vct p^2/(2m)+V(\vct r)$，$V(\vct r+\vct R)=V(\vct r)$。定义$T_{\vct R}\psi(\vct r)=\psi(\vct r+\vct R)$。所有$T_{\vct R}$与$H$对易，且平移彼此对易，可以在每个能量子空间内选共同本征态。因$T_{\vct R}T_{\vct R'}=T_{\vct R+\vct R'}$，本征值是群的字符$\e^{\ii\vct k\cdot\vct R}$。因此
\begin{equation}
\psi_{n\vct k}(\vct r)=\e^{\ii\vct k\cdot\vct r}u_{n\vct k}(\vct r),\qquad u_{n\vct k}(\vct r+\vct R)=u_{n\vct k}(\vct r).
\end{equation}
$\vct k$与$\vct k+\vct G$表示同一个平移字符，故晶体动量只定义到倒格矢。把边长含$N_i$个原胞的周期边界代入，$\vct k\cdot N_i\vct a_i=2\pi n_i$；第一Brillouin区内恰有$N_1N_2N_3=N_c$个$\vct k$。每条不计自旋能带含$N_c$态，每条计双自旋能带含$2N_c$态。完全填满的带可绝缘，部分填充的带通常导电，但多体Mott态不受这个单粒子计数结论保证。

把$V$与$u$作倒格傅里叶展开，Schrödinger方程化为矩阵
\begin{equation}
\sum_{\vct G'}\left[\frac{\hbar^2|\vct k+\vct G|^2}{2m}\delta_{\vct G\vct G'}+V_{\vct G-\vct G'}\right]c_{\vct G'}=E_{n\vct k}c_{\vct G}.
\end{equation}
这既是另一种Bloch证明，也是平面波能带计算的起点。有限截断是近似，应通过增加截断动能检查收敛；不能把截断后矩阵的维数当成真实能带总数。
\originalsection{Bragg简并与能隙}
在$\vct k$与$\vct k-\vct G$自由能近简并处，普通非简并微扰出现小分母，必须先对角化两态子空间。能量为
\begin{equation}
E_\pm=\frac{\epsilon_{\vct k}+\epsilon_{\vct k-\vct G}}2\pm\sqrt{\left(\frac{\epsilon_{\vct k}-\epsilon_{\vct k-\vct G}}2\right)^2+|V_{\vct G}|^2}.
\end{equation}
在Bragg面恰有劈裂$2|V_{\vct G}|$。两个驻波分别把电子密度集中于势能低区与高区，故能隙的物理来源是周期势对两种干涉图样的区别。几个Bragg面相交时必须保留全部简并波，不能逐个独立套两态公式。

若势无磁性且忽略自旋轨道耦合，复共轭把$\vct k$态映到$-\vct k$态，$E_n(\vct k)=E_n(-\vct k)$；若有反演，对应结论也由空间对称得到。含自旋的时间反演满足$\Theta^2=-1$，Kramers配对只保证在时间反演不变动量的简并，不能在任意动量不加条件地称双重简并。
\originalsection{芯态正交与赝势}
真实离子势在核附近很深，价电子为正交于芯态而快速振荡，平面波需要很大截断。令正交芯轨道$\ket{\phi_c}$能量为$E_c$，把光滑试验态投影为$\ket\psi=(1-\sum_c\ket{\phi_c}\bra{\phi_c})\ket\varphi$。将$H\ket\psi=E\ket\psi$代回，得到
\begin{equation}
\left[H+\sum_c(E-E_c)\ket{\phi_c}\bra{\phi_c}\right]\ket\varphi=E\ket\varphi.
\end{equation}
正交化产生排斥的非局域项，部分抵消强吸引芯势。这说明赝势为何可以较弱，也说明它通常依赖能量、角动量且不唯一。可靠赝势应保持所需价电子散射相移或外区对数导数，不是任意削掉深势井。

APW把空间分成原子球与间隙区：球内使用给定能量的径向方程解乘球谐函数，球外使用平面波，并在球面匹配函数。球内径向函数依赖未知能量，使本征方程通常为非线性。线性化方法在参考能量附近同时保留径向函数及其能量导数。OPW降低价电子的振荡负担，APW则把振荡交给适合核附近的基组，两者解决的是同一数值瓶颈的不同方面。
\originalsection{紧束缚与态密度}
若原子轨道空间上局域，先用$\ket{\alpha\vct R}$作基组，构造$\ket{\alpha\vct k}=N_c^{-1/2}\sum_{\vct R}\e^{\ii\vct k\cdot(\vct R+\vct d_\alpha)}\ket{\alpha\vct R}$。非正交基需解$H(\vct k)c=ES(\vct k)c$，不能无声明地令重叠矩阵$S=1$。一维单轨道最近邻且已正交时$E(k)=\epsilon_0-2t\cos ka$，带宽$4|t|$；两个轨道的杂化将交叉变为避免交叉，除非对称性禁止该矩阵元。

一般态密度$D(E)=\int_{\mathrm{BZ}}\dd^dk\,\delta(E-E(\vct k))/(2\pi)^d$。等能面换元给$D(E)=\int_{E(\vct k)=E}\dd S_k/[(2\pi)^d|\nabla_kE|]$，梯度为零的驻点能使态密度奇异。二维鞍点局部$E-E_s=\alpha k_x^2-\beta k_y^2$，令$u=\sqrt\alpha k_x+\sqrt\beta k_y$、$v=\sqrt\alpha k_x-\sqrt\beta k_y$，能量成为$uv$。先用$\delta(E-E_s-uv)$积分$v$，剩余$\int\dd u/|u|$，得到$D(E)\propto\log(\Lambda/|E-E_s|)$。对数来自鞍点附近双曲线的相空间，不是任何极值都具有的形式。

\begin{center}
\begin{tikzpicture}
\begin{axis}[width=10cm,height=4.5cm,xlabel={$ka$},ylabel={$E/t$},xmin=-3.1416,xmax=3.1416,ymin=-2.4,ymax=2.4,xtick={-3.1416,0,3.1416},xticklabels={$-\pi$,0,$\pi$},axis lines=middle]
\addplot[structurecolor,thick,domain=-3.1416:3.1416,samples=150]{-2*cos(deg(x))};
\end{axis}
\end{tikzpicture}
\end{center}
比热、自旋磁化率测低能态密度；光吸收还取决于初末态的联合态密度与跃迁矩阵元；光电子谱测被占据态的谱函数而非单纯能带曲线。周期表中$s,p,d$轨道的不同局域程度决定不同带宽，但完整材料预测还需晶体结构、填充及电子关联，不能用一条自由电子抛物线概括所有元素。

\originalchapter{半导体、半经典运动与量子Hall效应}
\label{ch:transport}
\source{8.511第11--14讲；8.513第1--6讲的先修连接。}
\originalsection{有效质量、掺杂与激子}
在能带极值附近展开$E_n(\vct k_0+\vct q)=E_0+\hbar^2q_i(m^{*-1})_{ij}q_j/2$，故$(m^{*-1})_{ij}=\hbar^{-2}\partial_i\partial_j E$。它是张量，可有不同主轴。价带顶的电子曲率为负；描述缺少的电子时取带正电$+e$的空穴，可获得正的空穴有效质量。

浅施主的长距离吸引势为$-e^2/(4\pi\epsilon_0\epsilon_r r)$。各向同性有效质量近似把氢原子的$m$换成$m^*$、真空介电常数换成$\epsilon_0\epsilon_r$，于是$a^*=\epsilon_r(m/m^*)a_B$、束缚能$E_b=(m^*/m)13.6\,\mathrm{eV}/\epsilon_r^2$。当$a^*$远大于晶格常数时这个包络近似自洽；深杂质的中央胞修正不属于它。电子与空穴形成Wannier激子时把$m^*$换成约化质量$\mu^{-1}=m_e^{*-1}+m_h^{*-1}$，同样要求弱束缚和大半径。间接带隙半导体的光学跃迁还需声子提供晶体动量，这解释电学带隙与光学吸收边之间的区别。
\originalsection{波包方程与轨道量子化}
先忽略Berry曲率，窄波包的中心满足$\dot{\vct r}=\hbar^{-1}\nabla_kE_n$、$\hbar\dot{\vct k}=-e(\vct E+\dot{\vct r}\times\vct B)$。条件是外场在原胞尺度上缓慢变化，且带间跃迁小。均匀电场使$k(t)=k(0)-eEt/\hbar$，一维紧束缚速度随$k$周期变化，所以位移有Bloch振荡，频率$\omega_B=eEa/\hbar$。实际散射使振荡可见性要求$\omega_B\tau\gtrsim1$，强场又可能诱发Zener跨带跃迁。

在$\vct E=0$且$\vct B=B\hat z$时，$k_z$与能量守恒，电子沿$k$空间等能截线运动。对电荷$-e$的电子，$\hbar\dot{\vct k}=-e(\dot{\vct r}\times B\hat z)$；两边左叉乘$\hat z$后消去速度，得$\dd\vct r_\perp=-(\hbar/eB)\hat z\times\dd\vct k_\perp$。这确定了轨道的取向，实空间面积与动量空间面积之比为$(\hbar/eB)^2$。闭合轨道的作用量量子化给Onsager规则
\begin{equation}
A_k(E,k_z)=\frac{2\pi eB}{\hbar}(n+\gamma),\qquad m_c=\frac{\hbar^2}{2\pi}\frac{\partial A_k}{\partial E}.
\end{equation}
无Berry相位的简单抛物线$\gamma=1/2$。$m_c$控制回旋频率$eB/m_c$，不等于任意主轴质量。电子能级经过$E_F$时态密度振荡，$1/B$周期为$\Delta(1/B)=2\pi e/(\hbar A_{\rm ext})$，其中$A_{\rm ext}$是垂直磁场的极值截面积。改变磁场方向便能探测费米面形状。热展宽和有限寿命削弱振荡；磁击穿是强场下相邻轨道之间跨越小能隙的隧穿。
\originalsection{Landau态与边缘态}
二维抛物线电子取$\vct A=-By\hat x$，$p_x$守恒。哈密顿量在$y_0=\hbar k_x/(eB)$附近等价于频率$\omega_c=eB/m$的谐振子，因此$E_n=\hbar\omega_c(n+1/2)$。把导引中心限制在宽$L_y$内，$k_x$间距$2\pi/L_x$给每Landau级每自旋简并$N_\phi=eBL_xL_y/h$。磁长$\ell_B=\sqrt{\hbar/(eB)}$决定波函数宽度，且$N_\phi=A/(2\pi\ell_B^2)$。

若$n$个自旋分辨Landau级填满，体态有能隙。在边界，限制势$V(y)$抬升能级，$E_n(k_x)\simeq\hbar\omega_c(n+1/2)+V(y_0)$，边缘群速度$\hbar^{-1}\dd E/\dd k_x=V'(y_0)/(eB)$。两侧边缘方向相反；一条边同向传播的通道缺少同位置反向态，弱无序不能直接造成后向散射。

把二维电子放在圆柱或环上，沿孔缓慢插入磁通$\Phi_0=h/e$。规范变换把插通量等价成边界相位$2\pi$变化，每条填满带泵送一个电子从一边到另一边。若有$n$条通道，$\Delta Q=ne$；由Faraday电动势积分$\int\mathcal E\dd t=\Phi_0$得$\sigma_{xy}=\Delta Q/\Phi_0=ne^2/h$，符号由电荷与磁场方向约定。此论证要求绝热过程中体能隙或迁移率隙不闭合。实际平台的宽度需要局域态容纳额外粒子，纯净固定密度Landau级计数本身不能给整个平台。
\originalsection{Berry修正的连接}
Bloch本征矢随$\vct k$变化时，定义$\mathcal A_n(\vct k)=\ii\avg{u_{n\vct k}|\nabla_k u_{n\vct k}}$及$\vct\Omega_n=\nabla_k\times\mathcal A_n$。波包作用量含$\hbar\mathcal A_n\cdot\dot{\vct k}$，Euler--Lagrange方程使速度多出$-\dot{\vct k}\times\vct\Omega_n$。该项是异常速度，来源于态空间几何而非额外力。第三编会从绝热演化推导Berry相位与Chern数，再说明没有外磁场的Chern绝缘体如何也有量子化Hall响应。

\originalchapter{交换、密度泛函与Mott绝缘体}
\label{ch:interaction}
\source{8.511第15--18讲。}
\originalsection{二次量子化与Slater行列式}
对于正交单粒子基$\phi_i(\vct r)$，用$c_i^\dagger$创建一个费米子，满足$\{c_i,c_j^\dagger\}=\delta_{ij}$、$\{c_i,c_j\}=0$。占据态$c_{i_1}^\dagger\cdots c_{i_N}^\dagger\ket0$自动实现反对称性。单体与双体算符成为
\begin{equation}
H=\sum_{ij}h_{ij}c_i^\dagger c_j+\frac12\sum_{ijkl}V_{ij;kl}c_i^\dagger c_j^\dagger c_lc_k,
\end{equation}
其中$V_{ij;kl}=\iint\phi_i^*(r)\phi_j^*(r')v(r-r')\phi_k(r)\phi_l(r')\dd r\dd r'$，自旋坐标包含在指标内。$1/2$避免双计粒子对。玻色子改用对易关系，不能把费米子的交换符号直接搬过去。

在一个Slater态中，$\avg{c_i^\dagger c_j^\dagger c_lc_k}=\rho_{ki}\rho_{lj}-\rho_{li}\rho_{kj}$。第二项是交换项，来自反对称性而非额外的经典相互作用。在占据轨道表象，能量为$E=\sum_i h_{ii}+\frac12\sum_{ij}(J_{ij}-K_{ij})$。对轨道变分并以拉格朗日乘子保持正交，得到Hartree--Fock方程
\begin{align}
\left[h_0+\sum_j\int v(r-r')|\phi_j(r')|^2\dd r'\right]\phi_i(r)
-\sum_j\phi_j(r)\int v(r-r')\phi_j^*(r')\phi_i(r')\dd r'=\epsilon_i\phi_i(r).
\end{align}
交换势非局域且只在相同自旋之间出现。乘以$\phi_i^*$积分并求和，$\sum_i\epsilon_i$把粒子对算了两遍；故总能量不是本征值之和。冻结其他轨道时删去轨道$i$使能量差等于$\epsilon_i$，这才是Koopmans关系的适用条件；真实电离后的弛豫与关联会改变它。
\originalsection{均匀电子气的交换能}
均匀正背景抵消Hartree项，平面波仍为HF本征态。暂用库仑形式$v_q=4\pi e_*^2/q^2$，其中$e_*^2=e^2/(4\pi\epsilon_0)$。每自旋交换自能为
\begin{equation}
\Sigma_x(k)=-\int_{p<k_F}\frac{\dd^3p}{(2\pi)^3}\frac{4\pi e_*^2}{|\vct k-\vct p|^2}
=-\frac{e_*^2k_F}{\pi}\left[1+\frac{k_F^2-k^2}{2kk_F}\log\left|\frac{k+k_F}{k-k_F}\right|\right].
\end{equation}
第二式先取极轴沿$\vct k$，角积分给$\log|(k+p)/(k-p)|/(kp)$，再积分$p$。$k=0$极限为$-2e_*^2k_F/\pi$，$k=k_F$为$-e_*^2k_F/\pi$。把自能按已占据态积分并除以2，交换能每电子为$E_x/N=-3e_*^2k_F/(4\pi)$；动能每电子$3E_F/5$。HF在费米面给出对数发散的$\dd\Sigma_x/\dd k$，不是可信的真实费米速度：长程势的屏蔽及相关修正被漏掉了。

交换孔也可从对关联看出。对同自旋粒子，$g_{\uparrow\uparrow}(r)=1-[3j_1(k_Fr)/(k_Fr)]^2$，$g_{\uparrow\uparrow}(0)=0$。反自旋在无关联Slater态中没有交换孔。高$r_s$下电子可因库仑作用形成Wigner晶体，但从均匀HF能量的一个极小值不能证明该相的存在或转变点。
\originalsection{Hohenberg--Kohn与Kohn--Sham}
固定粒子数和电子间作用，非简并基态密度决定外势至多相差常数。证明用反证：若不同外势$v,v'$有相同密度$n$却不同基态，则严格变分原理给$E_v<E_{v'}+\int(v-v')n$及$E_{v'}<E_v+\int(v'-v)n$，相加矛盾。严格不等式依赖不同基态；含简并情形须使用更一般的密度泛函表述，本稿不把此简短反证当作所有数学情形的存在定理。

定义普适泛函$F[n]$，总能量$E_v[n]=F[n]+\int vn$在可表示密度上极小。Kohn--Sham把$F[n]=T_s[n]+E_H[n]+E_{xc}[n]$，$T_s$由具有相同密度的辅助非相互作用轨道计算。变分得到
\begin{equation}
\left[-\frac{\hbar^2\nabla^2}{2m}+v(r)+\int\frac{e_*^2 n(r')}{|r-r'|}\dd r'+v_{xc}(r)\right]\phi_i=\epsilon_i\phi_i,\qquad n=\sum_{i\in\mathrm{occ}}|\phi_i|^2.
\end{equation}
$E_{xc}$包含交换、关联以及真实动能与$T_s$之差。LDA用$E_{xc}=\int n(r)\epsilon_{xc}(n(r))\dd r$，故$v_{xc}=\dd[n\epsilon_{xc}(n)]/\dd n$。总能量可写成
\begin{equation}
E=\sum_i\epsilon_i-E_H+E_{xc}-\int nv_{xc}\dd r.
\end{equation}
两个扣除项去掉本征值中重复算入的平均场。LDA的局域性适于密度变化缓慢区域，不保证正确带隙、强关联或弱范德华结合。Kohn--Sham本征值是辅助系统量，一般不能直接当作所有真实加减电子激发能。
\originalsection{Hubbard模型的强耦合极限}
把局域轨道上的主要作用压为同站排斥，得到$H=-t\sum_{\avg{ij},\sigma}(c_{i\sigma}^\dagger c_{j\sigma}+\mathrm{h.c.})+U\sum_i n_{i\uparrow}n_{i\downarrow}$。半填充且$U\gg t$时，每站一个电子。一次跳跃产生双占据与空站，代价$U$，再跳回才回到低能子空间。二阶简并微扰$H_{\rm eff}=-PH_tQH_tP/U$给
\begin{equation}
H_{\rm eff}=\sum_{\avg{ij}}\frac{4t^2}{U}\left(\vct S_i\cdot\vct S_j-\frac14n_in_j\right).
\end{equation}
两站平行自旋因Pauli原理无法虚跃迁；单重态可降低能量$4t^2/U$，因此交换反铁磁。电荷激发有$U$尺度的隙，低能自旋仍可运动。这说明一个按非相互作用填充本应金属的体系如何成为Mott绝缘体。它不同于晶格周期势在半填充处直接打开的带隙，且强耦合展开要求$|t|/U$小；近转变区域不能只用最低阶Heisenberg模型。

\originalchapter{屏蔽、声子与电子寿命}
\label{ch:phonon}
\source{8.511第18--21讲。}
\originalsection{密度响应与屏蔽}
外部电势扰动使电子密度改变。为避免符号混淆，定义电子能量扰动$\delta v$及静态极化$\delta n=-\Pi\,\delta v$，$\Pi(q,0)>0$。长波极限$\Pi(0,0)=\partial n/\partial\mu=D(E_F)$。总势为$\delta v_{\rm tot}=\delta v_{\rm ext}+v_q\delta n$，自洽解给
\begin{equation}
\delta v_{\rm tot}=\frac{\delta v_{\rm ext}}{\epsilon(q,\omega)},\qquad \epsilon=1+v_q\Pi,\qquad \epsilon(q,0)=1+\frac{\kappa^2}{q^2},\quad\kappa^2=4\pi e_*^2D(E_F).
\end{equation}
点电荷势因此变为Yukawa形式$e_*^2\e^{-\kappa r}/r$。这里Thomas--Fermi近似忽略了$\Pi$的有限$q$结构，不能产生正确的Friedel长距离振荡。

无相互作用费米子的响应来自一个粒子从已占据态跃到未占据态。以$\chi=-\Pi$的另一常用约定写Lindhard函数，
\begin{equation}
\chi_0(q,\omega)=2\int\frac{\dd^3k}{(2\pi)^3}\frac{f(\epsilon_k)-f(\epsilon_{k+q})}{\hbar\omega+\epsilon_k-\epsilon_{k+q}+\ii0^+}.
\end{equation}
分子是可用初末态的占据差，$\ii0^+$保证推迟因果性。三维零温静态积分给
\begin{equation}
\Pi(q,0)=D(E_F)\left[\frac12+\frac{1-x^2}{4x}\log\left|\frac{1+x}{1-x}\right|\right],\qquad x=\frac q{2k_F}.
\end{equation}
$2k_F$处导数非解析，Fourier变换产生远处的$\cos(2k_Fr)/r^3$密度振荡。在高频长波极限，连续性和惯性给$\epsilon\simeq1-\omega_p^2/\omega^2$，$\omega_p^2=ne^2/(m\epsilon_0)$。零点$\epsilon=0$表示自洽电荷集体振荡，区别于单个粒子--空穴激发。
\originalsection{晶格振动的量子化}
在稳定平衡位置附近，势能为$\frac12\sum_{ij,\alpha\beta}u_{i\alpha}K_{i\alpha,j\beta}u_{j\beta}$。质量加权并Fourier变换得到动力学矩阵$D_{\alpha\beta}(q)$，本征值是$\omega_{q\lambda}^2$。正本征值表示谐振稳定；整体平移要求$q=0$存在零频声学模。每个原胞含$r$原子时有$3r$支，通常3支声学与$3r-3$支光学。

正规坐标量子化为
\begin{equation}
\vct u_j=\frac1{\sqrt N}\sum_{q\lambda}\sqrt{\frac{\hbar}{2M\omega_{q\lambda}}}\vct e_{q\lambda}\left(b_{q\lambda}\e^{\ii q\cdot R_j}+b_{q\lambda}^\dagger\e^{-\ii q\cdot R_j}\right),\quad H_{\rm ph}=\sum_{q\lambda}\hbar\omega_{q\lambda}(b^\dagger b+1/2).
\end{equation}
光学模中的不同质量需放回各基底分量。低温只有$\omega\simeq c_\lambda q$声学模重要，态密度$D_\lambda(\omega)/V=\omega^2/(2\pi^2c_\lambda^3)$。能量积分$\int\hbar\omega D_\lambda(\omega)n_B(\omega)\dd\omega$给
\begin{equation}
\frac{C_V}{V}=\frac{2\pi^2\kb^4T^3}{15\hbar^3}\sum_\lambda c_\lambda^{-3}.
\end{equation}
Debye近似用截止频率使总模数等于$3N$；它是低能声学模型，并不描述所有实际高频支。

离子jellium的裸等离子体频率为$\omega_i^2=n_iZ^2e^2/(M\epsilon_0)$。电子静态屏蔽把它改成$\omega^2(q)=\omega_i^2/\epsilon(q,0)\simeq\omega_i^2q^2/\kappa^2$。离子本来有隙的纵电荷振荡因此在长波成为声波，速度$c_L=\omega_i/\kappa$。这是电子可比离子快地调整密度的绝热近似。
\originalsection{电子--声子作用与幂律}
晶格位移改变电子势，保留一阶得到$H_{e\text{-}ph}=\sum_{kq\sigma}g_qc_{k+q,\sigma}^\dagger c_{k\sigma}(b_q+b_{-q}^\dagger)$。声学形变势常有$|g_q|^2\propto q$，因为应变含$qu$，而量子位移$u\propto\omega_q^{-1/2}$。零温费米能上方$E$的电子可发射声子，Golden rule为
\begin{equation}
\frac1\tau=\frac{2\pi}{\hbar}\int\frac{\dd^3q}{(2\pi)^3}|g_q|^2[1-f(\epsilon_{k-q})]\delta(\epsilon_k-\epsilon_{k-q}-\hbar c_sq).
\end{equation}
能量守恒限制$q\lesssim E/(\hbar c_s)$。在$E\ll\hbar\omega_D$且$v_F\gg c_s$时，角积分的$\delta$函数给$1/(\hbar v_Fq)$，结合相空间$q^2\dd q$和$|g_q|^2\propto q$，剩$\int^{q_{\max}}q^2\dd q\propto E^3$。这给单粒子衰减率的$E^3$幂律。输运还须乘$1-\cos\theta_{kk'}\simeq q_\perp^2/(2k_F^2)$，故动量弛豫率多两个$E$次幂。把热激发尺度取$\kb T$得到低温电阻的$T^5$规律，但数值系数、角向限制与材料参数须由具体$g_q$计算。

声子传播子的静态低频贡献为负：对谐振坐标完成平方或积分掉它，得到$V_{\rm eff}(q,\omega)=|g_q|^2\,2\hbar\omega_q/[(\hbar\omega+\ii0^+)^2-(\hbar\omega_q)^2]$。$|\omega|\ll\omega_q$时这是吸引，频率高过声子尺度时符号与强度改变。其迟滞使短时库仑排斥与低频配对吸引的作用窗口不同。对电子实自能展开$\Sigma'(\omega)\simeq-\lambda\hbar\omega$得到质量增强$1+\lambda$；比热相应增强，而自旋响应还含顶角修正，不能仅以质量替换判断所有观测量。

\originalchapter{Cooper不稳定性与超导序}
\label{ch:bcs}
\source{8.511第22--24讲，衔接8.512第9--13讲。}
\originalsection{费米海上方的一对电子}
考虑总动量零、相反自旋的两电子，只允许$0<\xi_k=\epsilon_k-E_F<\hbar\omega_D$。约化吸引势为$-V_0$，配对波函数系数$a_k$满足$(2\xi_k-E)a_k=V_0\sum_{k'}a_{k'}$。消去$a_k$得
\begin{equation}
1=V_0\sum_k\frac1{2\xi_k-E}\simeq N_\uparrow(0)V_0\int_0^{\hbar\omega_D}\frac{\dd\xi}{2\xi+|E|}.
\end{equation}
$N_\uparrow(0)$是每自旋态密度。弱吸引时$|E|\simeq2\hbar\omega_D\exp[-2/(N_\uparrow(0)V_0)]$，任意小吸引都产生束缚。与真空三维两体问题的有限阈值不同，费米海边缘的常态密度导致对数积分发散。两体结果指出正常费米面不稳定，却未给具有有限配对密度的真正基态。
\originalsection{BCS变分与Bogoliubov准粒子}
多体试验态取$\ket\Psi=\prod_k(u_k+v_kc_{k\uparrow}^\dagger c_{-k\downarrow}^\dagger)\ket0$，$|u_k|^2+|v_k|^2=1$。不同$k$的空态与双占据态叠加，相位一致的异常平均$\avg{c_{-k\downarrow}c_{k\uparrow}}=u_k^*v_k$成为序参量。以$K=H-\mu N$变分，令$\Delta=V_0\sum_ku_k^*v_k$，得到$u_kv_k=\Delta/(2E_k)$、$v_k^2=(1-\xi_k/E_k)/2$，其中$E_k=\sqrt{\xi_k^2+|\Delta|^2}$。

同样结果可从Nambu矩阵获得。对$(c_{k\uparrow},c_{-k\downarrow}^\dagger)^T$，单个配对块为
\begin{equation}
\mathcal H_k=\begin{pmatrix}\xi_k&-\Delta\\-\Delta^*&-\xi_k\end{pmatrix},\qquad \mathcal H_k^2=E_k^2I.
\end{equation}
归一化本征矢形成规范变换$\gamma_{k\uparrow}=u_kc_{k\uparrow}-v_kc_{-k\downarrow}^\dagger$，仍满足费米反对易关系。准粒子是电子与空穴的相干混合，创建它至少需能量$|\Delta|$，而打破一对并产生两个准粒子通常需$2|\Delta|$。

有限温度下每个准粒子态占据$f(E_k)$，异常平均乘$1-2f(E_k)=\tanh(E_k/2\kb T)$，故
\begin{equation}
1=N_\uparrow(0)V_0\int_0^{\hbar\omega_D}\frac{\tanh(\sqrt{\xi^2+\Delta^2}/2\kb T)}{\sqrt{\xi^2+\Delta^2}}\dd\xi.
\end{equation}
$T=0$积分为$\operatorname{arsinh}(\hbar\omega_D/\Delta_0)$，弱耦合给$\Delta_0=2\hbar\omega_D\e^{-1/[N_\uparrow(0)V_0]}$。令$\Delta\to0$时积分的有限常数给$\kb T_c=(2\e^\gamma/\pi)\hbar\omega_D\e^{-1/[N_\uparrow(0)V_0]}$，因此$2\Delta_0/(\kb T_c)=2\pi\e^{-\gamma}\simeq3.53$。这是各向同性弱耦合$s$波结论；不同配对对称、强耦合与多带材料可偏离。

正常态密度近似常数时，超导态密度为$D_s(E)=D_n(0)\operatorname{Re}[|E|/\sqrt{E^2-\Delta^2}]$。平方根边缘来自把$\xi$映到$E$的Jacobian。低温热激发按$\e^{-\Delta/\kb T}$减少，节点超导则不能沿用此指数律。
\originalsection{Ginzburg--Landau、London与涡旋}
接近$T_c$，用缓慢变化的复序参量$\Psi$写
\begin{equation}
F=\int\dd^3r\left[\alpha|\Psi|^2+\frac\beta2|\Psi|^4+\frac1{2m_p}|(-\ii\hbar\nabla-q_p\vct A)\Psi|^2+\frac{B^2}{2\mu_0}\right],\quad q_p=-2e.
\end{equation}
稳定性要求$\beta>0$，$\alpha<0$时均匀极小$|\Psi|^2=-\alpha/\beta$。变分对$\Psi^*$给GL方程，对$\vct A$给电流$\vct j=(q_p\hbar/m_p)\operatorname{Im}(\Psi^*\nabla\Psi)-(q_p^2/m_p)|\Psi|^2\vct A$。取$\Psi=|\Psi|\e^{\ii\theta}$并固定幅度，得到$\vct j=(q_p|\Psi|^2/m_p)(\hbar\nabla\theta-q_p\vct A)$。在无涡旋区域取旋度，再用Maxwell方程，$\nabla^2\vct B=\vct B/\lambda^2$，$\lambda^2=m_p/(\mu_0q_p^2|\Psi|^2)$。这是Meissner排磁的空间尺度。

幅度恢复尺度为$\xi=\hbar/\sqrt{2m_p|\alpha|}$。沿远离涡核的闭合回路电流可忽略，单值性要求$\oint\nabla\theta\cdot\dd\vct l=2\pi n$，故通量绝对值$|\Phi|=n h/(2e)$。相位 winding 不能被单值规范变换消掉，涡核内幅度需减小，否则梯度能发散。$\kappa=\lambda/\xi$控制界面能：GL极限$\kappa<1/\sqrt2$为I类，$\kappa>1/\sqrt2$为II类，后者在两个临界场之间形成涡旋态。临界场近似$B_{c2}=\Phi_0/(2\pi\xi^2)$、$B_{c1}\simeq\Phi_0[\log\kappa+O(1)]/(4\pi\lambda^2)$仅适于相应连续理论与大$\kappa$范围。
\originalsection{隧穿与Josephson关系}
两个超导体的配对振幅弱耦合，最低相位依赖自由能为$F_J=-E_J\cos\varphi$，$\varphi$是扣除跨结规范线积分的相位差。相位与转移的对数互为共轭，故$I=(2e/\hbar)\partial_\varphi F_J=I_c\sin\varphi$，$I_c=2eE_J/\hbar$。电压$V$使两侧化学势相差$2eV$，$\dot\varphi=2eV/\hbar$，恒压产生频率$2eV/h$的交流超流。正常准粒子隧穿则为$\int D_1(E)D_2(E+eV)[f(E)-f(E+eV)]\dd E$的谱卷积，与零电压相干对隧穿是不同通道。

\volpart{第二编\quad 响应、无序与关联电子（8.512）}
\originalchapter{因果响应、涨落与散射}
\label{ch:response}
\source{8.512课程页第1--4讲；Mark Spencer Rudner讲义及Lee讲义。}
\originalsection{从微扰演化到Kubo公式}
设未扰动密度矩阵为$\rho_0=\e^{-\beta H_0}/Z$，施加$H'(t)=-f(t)B$。在相互作用绘景，一阶演化算符为
\begin{equation}
U_I(t)\simeq1+\frac\ii\hbar\int_{-\infty}^t f(t')B_I(t')\dd t'.
\end{equation}
将$U_I^\dagger A_IU_I$展开，响应为
\begin{equation}
\delta\avg{A(t)}=\int\dd t'\,\chi_{AB}(t-t')f(t'),\qquad\chi_{AB}(t)=\frac\ii\hbar\theta(t)\avg{[A(t),B(0)]}.
\end{equation}
$\theta(t)$表达原因不能晚于结果。此处采用对$-fB$响应的正静态磁化率约定。OCW若以$+fB$扰动定义响应函数，则其$\chi$为这里的负号，尤其题面中的$S=-2(n_B+1)\chi''$不能与本章符号无说明混用。Fourier约定$\chi(\omega)=\int\e^{\ii\omega t}\chi(t)\dd t$。

插入$H_0\ket n=E_n\ket n$，对$A=B^\dagger$得谱表示
\begin{equation}
\chi_{AA^\dagger}(\omega)=\sum_{nm}\frac{(p_n-p_m)|A_{nm}|^2}{E_m-E_n-\hbar\omega-\ii0^+},\qquad p_n=\frac{\e^{-\beta E_n}}Z.
\end{equation}
分母极点位于复频率下半面，所以$\chi$在上半面解析。对实扰动与Hermitian观测量，$\chi(-\omega)=\chi(\omega)^*$。复平面积分给Kramers--Kronig关系$\chi'(\omega)=\pi^{-1}\mathcal P\int\chi''(\omega')/(\omega'-\omega)\dd\omega'$，以及实虚部互换的负号关系。关系要求响应在无穷远充分衰减；有常数瞬时项时应先扣去该项。
\originalsection{涨落耗散定理}
定义非对称谱$S_{AA^\dagger}(\omega)=\int\e^{\ii\omega t}\avg{A(t)A^\dagger(0)}\dd t$。插入能量态可得$S(\omega)=2\pi\hbar\sum_{nm}p_n|A_{nm}|^2\delta(E_m-E_n-\hbar\omega)$。交换$n,m$后，$p_m=\e^{-\beta\hbar\omega}p_n$，所以$S_{A^\dagger A}(-\omega)=\e^{-\beta\hbar\omega}S_{AA^\dagger}(\omega)$。这正是详细平衡：零温体系能吸收外界能量，却不能从基态自发释放正能量。

利用谱分母的虚部$\pi\delta(E_m-E_n-\hbar\omega)$，得到
\begin{equation}
S(\omega)=\frac{2\hbar}{1-\e^{-\beta\hbar\omega}}\chi''(\omega),\qquad S_{\rm sym}(\omega)=\hbar\coth\frac{\beta\hbar\omega}2\,\chi''(\omega).
\end{equation}
非对称谱对应实验能量吸收与发射，不能无条件换成对称谱。经典低频极限$S_{\rm sym}\simeq2\kb T\chi''/\omega$，热涨落与耗散系数相联。系统必须处于热平衡；在驱动稳态只测到噪声并不能自动套此式。
\originalsection{密度相关与中子散射}
中子通过核的短程散射长度$b$与晶格密度$\rho_Q=\sum_j\e^{-\ii Q\cdot r_j}$耦合。Born近似把散射率写成常数乘$S(Q,\omega)$，其中$Q=k_i-k_f$、$\hbar\omega$是中子失去的能量。结构因子由$\avg{\rho_Q(t)\rho_{-Q}(0)}$给出，既含弹性Bragg峰也含声子等集体激发。

谐振位移满足高斯统计，但不同时间的位移不对易。令$X=-\ii Q u_j(t)$、$Y=\ii Q u_l(0)$，由于$[X,Y]$是数，Baker--Campbell--Hausdorff公式为$\e^X\e^Y=\e^{X+Y}\e^{[X,Y]/2}$。这项量子相位保证声子产生权重$n_q+1$与消灭权重$n_q$不同。展开位移指数到最低非弹性阶得
\begin{align}
S(Q,\omega)\propto\e^{-2W}\bigg[&\sum_G\delta(Q-G)\delta(\omega)\\
+&\sum_{q\lambda}\frac{\hbar|Q\cdot e_{q\lambda}|^2}{2M\omega_{q\lambda}}\big\{(n_q+1)\sum_G\delta(Q-G-q)\delta(\omega-\omega_q)\nonumber\\
&\hspace{35mm}+n_q\sum_G\delta(Q-G+q)\delta(\omega+\omega_q)\big\}\bigg].\nonumber
\end{align}
第一项是弹性峰，后二项是产生与吸收一个声子。选不同方向$Q$可测极化$e_{q\lambda}$，因为强度含投影平方，不能只凭频率判定纵模或横模。
\originalsection{双对易与频率和规则}
对能量本征态$\ket0$与Hermitian算符$A$，插入完备关系得到$\sum_n(E_n-E_0)|A_{n0}|^2=\avg{[A,[H,A]]}/2$。取$A=x$、$H=p^2/(2m)+V(x)$，内对易$[H,x]=-\ii\hbar p/m$，双对易为$\hbar^2/m$，因此$\sum_n(E_n-E_0)|x_{n0}|^2=\hbar^2/(2m)$。这不依赖具体势能。

对密度$\rho_q=\sum_j\e^{-\ii q\cdot r_j}$，位置依赖相互作用与密度对易，动能给$\avg{[\rho_{-q},[H,\rho_q]]}=N\hbar^2q^2/m$。于是密度谱的一阶频率矩由粒子数与裸质量固定，电子相互作用只能重新分配谱权重。连续电子系统的光学和规则$\int_0^\infty\operatorname{Re}\sigma(\omega)\dd\omega=\pi ne^2/(2m)$还包含零频$\delta$项的半轴权重。只截取一条晶格能带时，右侧改由该带动能曲率给出，不能将裸质量规则直接当作有限带模型内部的精确恒等式。

\originalchapter{电导、局域化与变程跳跃}
\label{ch:disorder}
\source{8.512第5--8讲；原PDF中部分讲号与页面讲号不一致，按主题连接。}
\originalsection{电流响应与扩散}
最小耦合$(p+eA)^2/(2m)$使电流含顺磁算符$\vct j_p$及抗磁项$-ne^2\vct A/m$。均匀场下$\vct E(\omega)=\ii\omega\vct A(\omega)$，两项相加才给规范不变电导。仅保留电流--电流相关会在正常金属错误地产生静态排磁响应。Kubo电导可写成$\sigma_{ij}(\omega)=\ii[ne^2\delta_{ij}/m+\Pi^R_{ij}(\omega)]/(\omega+\ii0^+)$，其中$\Pi^R$按电流--电流的相应符号定义。

若粒子长期运动为扩散，密度连续性与Fick定律给$\partial_t\delta n=D\nabla^2\delta n$。对化学势梯度，扩散电流和电势漂移电流在静态平衡应相抵，因而Einstein关系为$\sigma=e^2D\,\partial n/\partial\mu$。对应密度响应有扩散极点$\chi_{nn}(q,\omega)=\chi_{nn}(q,0)Dq^2/(Dq^2-\ii\omega)$。极限顺序重要：先$q\to0$且$\omega\ne0$时总粒子数无法变化，先$\omega\to0$则得到静态压缩率。

弱随机势的一次Born散射给$1/\tau=(2\pi/\hbar)D_\uparrow(E_F)\overline{|V_{kk'}|^2}$，角平均省略在记号中。输运时间乘$1-\cos\theta$，前向散射会明显区分$\tau$与$\tau_t$。在$k_F\ell\gg1$下，$D=v_F^2\tau_t/d$，恢复Drude电导。
\originalsection{边界敏感性与Thouless能量}
在边长$L$的扩散样品中，穿越时间$t_D=L^2/D$，相应能量$E_{\mathrm{Th}}=\hbar D/L^2$。单粒子平均能级间隔$\delta=1/[D(E_F)L^d]$。改变边界相位时，延展态的能量移动可达$E_{\mathrm{Th}}$；局域态的波函数在边界极小，能量移动按$\e^{-L/\xi}$下降。因此$g\sim E_{\mathrm{Th}}/\delta\sim\hbar D(E_F)DL^{d-2}$刻画电导与边界敏感性，比例常数随$e^2/h$的选法而变。

Landauer图像为$G=(e^2/h)\sum_nT_n$，自旋通道另计。理想一维通道也有两端接触电阻；四端测量与两端电导不应混称。无序相干传播需叠加振幅，不能把串联势垒的经典电阻相加当作所有尺度的精确规则。
\originalsection{局域化的尺度流}
定义$\beta(g)=\dd\log g/\dd\log L$。大$g$的经典导体给$\beta\to d-2$；局域态$g\propto\e^{-L/\xi}$给$\beta\simeq\log g$，在小$g$为负。正交对称类且无强自旋轨道与外磁场时，时间反演的一对闭合路径振幅同相，增强回返概率，降低扩散。

量子修正主要来自返回传播子$1/(Dq^2+\tau_\phi^{-1})$，所以$\delta\sigma\propto-\int_{1/L_\phi}^{1/\ell}\dd^dq/q^2$。一维修正随$L_\phi$增长，二维呈负对数，三维红外收敛。二维弱无序可在极大尺度最终局域，一维一般更快；三维可有不稳定临界点分开金属与绝缘体。以上是特定对称类的单参数尺度理论，不能推广为“所有二维体系永远绝缘”：磁场、拓扑项、自旋轨道以及相互作用都可改变流。

\begin{center}
\begin{tikzpicture}[>=Stealth]
\draw[->](0,0)--(8,0)node[right]{$\log g$};\draw[->](0,-2)--(0,1.4)node[above]{$\beta$};
\draw[thick,structurecolor](-.0,-1.8)..controls(2,-.5)and(4,1)..(7.3,1);
\draw[dashed](0,1)--(7.5,1);\node[left]at(0,1){$d-2$};
\fill(3.1,0)circle(2pt);\node[below]at(3.1,0){$g_c$};
\draw[->](1.8,-.8)--(.9,-1.3);\draw[->](4.5,.7)--(5.5,.9);
\end{tikzpicture}
\end{center}
此图仅示意三维正交类的两端极限及一个临界点，不是未经计算的数值$\beta$曲线。
\originalsection{Mott变程跳跃}
在局域绝缘体中，距$r$的两态重叠$\propto\e^{-r/\xi}$，热激活能差$\Delta$的跳跃率$\Gamma\propto\exp[-2r/\xi-\Delta/(\kb T)]$。在能窗$\Delta$与体积$r^d$中找一个态要求$\nu(E_F)\Delta r^d\sim1$。将$\Delta$消去，优化$X(r)=2r/\xi+1/[\kb T\nu r^d]$给$r_{\rm opt}^{d+1}\sim d\xi/[2\kb T\nu]$，因而
\begin{equation}
\rho(T)\sim\rho_0\exp[(T_0/T)^{1/(d+1)}],\qquad \kb T_0\sim\frac1{\nu(E_F)\xi^d}.
\end{equation}
低温最佳跳跃距离比最近邻远，故称变程。该指数要求费米能附近态密度近常数。若库仑相互作用形成软隙，优化能量--距离关系改变，常出现Efros--Shklovskii的$1/2$指数；不能用观察到$1/2$直接反推“一维Mott”。

\originalchapter{超导响应、相干因子与弱连接}
\label{ch:scresponse}
\source{8.512第9--13讲及8.511配对基础。}
\originalsection{横向响应与极限顺序}
正常电子气对静态、空间均匀的纯规范$A$不能产生可观测电流。其顺磁与抗磁响应在$\omega=0,q\to0$抵消；Landau轨道抗磁来自保留空间变化至适当阶的横向响应，不能把常数$A$响应当作磁化。超导体在横向静态极限保留刚度$K_T(0,0)=n_se^2/m$，于是$\vct j_T=-K_T\vct A_T$，恢复London方程。

无散射正常金属在先$q=0$后$\omega\to0$的动态极限也有Drude零频权重。理想电导与Meissner排磁区别在于不同极限的响应，不是只看$\sigma(0)$是否无穷大。规范不变计算还应让序参量相位随纵向扰动调整，否则会违反电荷守恒。
\originalsection{热准粒子与超流密度}
在各向同性洁净$s$波体系，流动使准粒子能量有Doppler位移。准粒子对动量的反向响应给
\begin{equation}
\frac{n_s(T)}n=1+2\int_\Delta^\infty\frac{E\dd E}{\sqrt{E^2-\Delta^2}}\frac{\partial f(E)}{\partial E}.
\end{equation}
零温$\partial f/\partial E$在正能量区为零，$n_s=n$。低温令$E=\Delta+x$，平方根分母$\sqrt{2\Delta x}$，$-f'(E)\simeq\e^{-E/\kb T}/\kb T$，积分给$n-n_s\propto n\sqrt{2\pi\Delta/\kb T}\e^{-\Delta/\kb T}$。穿透深度$\lambda^{-2}(T)=\mu_0n_s(T)e^2/m$。这一低温指数律又依赖无节点能隙。

非磁性弹性杂质会破坏单个$k$态，却不必破坏由时间反演态构成的$s$波配对。若杂质不改变平均态密度且配对近各向同性，能隙方程中频率与能隙的重整化有相同因子，$T_c$在该近似内保持。这是Anderson定理的内容，不是所有无序都无影响。超流刚度与穿透深度仍会变化；脏极限$\ell\ll\xi_0$有$\xi_{\rm dirty}\sim\sqrt{\xi_0\ell}$，$\lambda_{\rm dirty}\sim\lambda_L\sqrt{\xi_0/\ell}$。磁性散射翻转自旋或非$s$波配对中的角向混合通常破坏这项保护。
\originalsection{相干因子的来源}
把电子算符展开为Bogoliubov准粒子，一般双线性观测量包含准粒子散射系数$u_ku_{k'}\mp v_kv_{k'}$和两准粒子创建系数$u_kv_{k'}\pm v_ku_{k'}$。符号由观测量在时间反演下的性质和自旋结构决定，不是任意挑选。

例如密度与自旋观测量的电子、空穴贡献可相消或相加。在费米面附近$u\simeq v$，某些通道被压低，另一些通道增强。因此超导隙上方态密度的平方根峰并不保证所有散射实验都有同样强的峰。核磁弛豫中的相干峰需要特定矩阵元及散射条件；相位改变的配对、寿命展宽与强涨落可削弱它。

做计算时应先写待测算符，再代入$u,v$，再按准粒子初末态取矩阵元。若只写一个未说明适用通道的“相干因子”，难以判断符号是否正确。
\originalsection{隧穿电流与Josephson微观结构}
正常金属--超导体的弱隧穿，在矩阵元近常数时$dI/\dd V\propto D_s(eV)$，低温可测能隙与谱峰。超导--超导体的耗散电流有两边态密度卷积，零温阈值$|eV|\ge\Delta_1+\Delta_2$。配对电流来自隧穿哈密顿量二阶过程：先转移一个电子的中间态有准粒子能量代价，第二个转移恢复配对，两种方向振幅干涉产生$\cos\varphi$能量。

相同$s$波隙、弱隧穿且无磁性时，微观能量积分给Ambegaokar--Baratoff关系
\begin{equation}
I_cR_N=\frac{\pi\Delta(T)}{2e}\tanh\frac{\Delta(T)}{2\kb T}.
\end{equation}
$R_N$是同一结的正常电阻，公式不能用于任意高透明度接触。在超导环的两个结中，相位差相差$2\pi\Phi/\Phi_0$；相同结总临界流$2I_c|\cos(\pi\Phi/\Phi_0)|$，若忽略环自感。这把规范相位与可测磁通直接联系起来。

\originalchapter{磁矩、交换与自旋波}
\label{ch:magnet}
\source{8.512第14--19讲，8.513相互作用玻色子组的磁序主题。}
\originalsection{局域磁矩与交换尺度}
原子部分填满的壳层具有自旋与轨道角动量。Hund规则是在一定相互作用与能级顺序下先使总自旋尽可能大，再判断轨道及自旋轨道耦合；晶场可使轨道角动量淬灭，不能在晶体里无条件套孤立原子的简并。独立局域磁矩的小场配分函数展开给$\chi=C/T$，$C=\mu_0n(g\mu_B)^2S(S+1)/(3\kb)$，与Pauli顺磁近常数的温度行为不同。

相邻局域轨道通过虚跃迁获得超交换。本册两站Hubbard推导给$J=4t^2/U>0$。若包含多个轨道及Hund耦合，虚中间态能量不同，交换可改变符号。金属中的局域矩还可通过传导电子产生RKKY作用$J_{ij}\propto J_K^2\chi_0(R_i-R_j)$；三维长距离随$\cos(2k_FR)/R^3$振荡。该结果来自把电子密度--自旋响应积分到二阶，而非单一近邻超交换的延长。
\originalsection{铁磁自旋波}
铁磁Heisenberg模型$H=-J\sum_{\avg{ij}}\vct S_i\cdot\vct S_j$，$J>0$。基态沿$+z$排列。Holstein--Primakoff精确表示$S^z=S-a^\dagger a$、$S^+=\sqrt{2S-a^\dagger a}\,a$、$S^-=a^\dagger\sqrt{2S-a^\dagger a}$，在低激发密度$\avg{a^\dagger a}\ll2S$时展开根号。保留二次项并Fourier变换得
\begin{equation}
H=E_0+JSz\sum_k(1-\gamma_k)a_k^\dagger a_k,\qquad\gamma_k=\frac1z\sum_\delta\e^{\ii k\cdot\delta}.
\end{equation}
简单立方晶格$\gamma_k\simeq1-a^2k^2/(2d)$，故$\omega_k\propto k^2$。磁化减少$\delta M\propto\int\dd^dk\,n_B(\hbar\omega_k)$，三维给$T^{3/2}$的Bloch律。连续自旋对称下$d\le2$的红外涨落破坏有限温长程序；各向异性或长程作用改变该结论的条件。
\originalsection{反铁磁自旋波与自旋翻转}
双子格反铁磁$H=J\sum_{\avg{ij}}\vct S_i\cdot\vct S_j$，先把B子格自旋坐标旋转，使局域参考态同向，再作HP展开。二次哈密顿量为
\begin{equation}
H_2=JSz\sum_k[a_k^\dagger a_k+b_k^\dagger b_k+\gamma_k(a_kb_{-k}+a_k^\dagger b_{-k}^\dagger)].
\end{equation}
玻色Bogoliubov变换必须满足$u_k^2-v_k^2=1$，于是$\hbar\omega_k=JSz\sqrt{1-\gamma_k^2}$，长波线性。铁磁与反铁磁的Goldstone模色散不同，原因在于低能时间导数结构及未破缺生成元之间的对易，而不是都由“有连续对称”就得相同色散。

加易轴各向异性与沿轴磁场时，Néel态可能转为自旋翻转态：两子格主要反向却横于磁场，并有共同小纵向分量。若定义易轴能密度为$-K\ell_z^2/2$，横向磁化能为$m^2/(2\chi_\perp)-Hm_z$且$\vct m\perp\vct\ell$，极小化$\vct m$给场能$-\chi_\perp H^2(1-\ell_z^2)/2$。比较易轴与横向极小，转变尺度$H_{\rm sf}=\sqrt{K/\chi_\perp}$。系数取决于所定义的$K$，这比不注明哈密顿量就给一个$\sqrt{JK}$更可复核。

中子磁散射测动态自旋结构因子$S^{\alpha\beta}(q,\omega)$。线性自旋波的单磁振子给$\delta(\omega-\omega_q)$峰，多磁振子给连续谱，强度含极化因子与子格干涉。只知道能谱不足以知道某种探针能否看到该模。
\originalsection{巡游磁性与自旋密度波}
在单带Hubbard平均场中，设$n_\uparrow=n/2+m$、$n_\downarrow=n/2-m$。相互作用能$Un_\uparrow n_\downarrow=Un^2/4-Um^2$，自旋分裂带来的动能代价近$ m^2/D_\uparrow(E_F)$。二次系数变负的Stoner条件为$UD_\uparrow(E_F)>1$。采用双自旋总态密度时应写$UD(E_F)/2>1$，避免因约定产生两倍差。

更一般地，密度在波矢$Q$处的自旋不稳定由$1-U\chi_0(Q,0)=0$决定。费米面嵌套若使$\epsilon_{k+Q}\simeq-\epsilon_k$，静态粒子--空穴响应可有大对数，从而优先产生自旋密度波而非均匀铁磁。平均场混合$k\uparrow$与$k+Q\downarrow$，二维矩阵本征能为$\frac12(\epsilon_k+\epsilon_{k+Q})\pm\sqrt{[(\epsilon_k-\epsilon_{k+Q})/2]^2+\Delta_Q^2}$。只有被嵌套部分打开隙时可仍有剩余费米面。平均场忽略强涨落，尤其低维系统不能以Stoner条件当作精确相边界。

\originalchapter{局域矩、Friedel相移与Kondo效应}
\label{ch:kondo}
\source{8.512第20--22讲，含两个文件名不同的第22讲同文版本。}
\originalsection{散射相移与粒子数}
球对称杂质势使每个角动量通道的远区径向波为$\sin(kr-l\pi/2+\delta_l(k))$。在半径$R$盒子中边界量子化$kR-l\pi/2+\delta_l(k)=n\pi$，相移使低于$k_F$的态数变化$\Delta N_l=\delta_l(k_F)/\pi$，包含相应束缚态的计数后成立。每个$l$有$2l+1$角简并及自旋，故$\Delta N=\sum_{l\sigma}(2l+1)\delta_{l\sigma}(E_F)/\pi$。这是Friedel和规则，连接局域屏蔽电荷与散射相移。

例如局域$s$波共振在费米能处$\delta=\pi/2$，每自旋贡献半个电子的相移计数。谱峰的宽度由杂质--传导带杂化决定，不能把半个计数当作杂质可单独处于“半个电子”的本征态。
\originalsection{Anderson模型与局域矩条件}
Anderson杂质模型为
\begin{equation}
H=\sum_{k\sigma}\epsilon_kc_{k\sigma}^\dagger c_{k\sigma}+\epsilon_d\sum_\sigma n_{d\sigma}+Un_{d\uparrow}n_{d\downarrow}+\sum_{k\sigma}(V_kc_{k\sigma}^\dagger d_\sigma+\mathrm{h.c.}).
\end{equation}
无杂化时空、单、双占据能为$0,\epsilon_d,2\epsilon_d+U$。以$E_F=0$计，$\epsilon_d<0<\epsilon_d+U$使单占据最低；若两种电荷代价远大于杂化宽度$\Gamma=\pi D_\uparrow|V|^2$，可形成局域矩区。$U=0$时$G_d^R(\omega)=1/(\hbar\omega-\epsilon_d+\ii\Gamma)$，谱函数Lorentz形，$n_{d\sigma}=1/2-\pi^{-1}\arctan(\epsilon_d/\Gamma)$。

静态Hartree近似用$\epsilon_{d\sigma}=\epsilon_d+Un_{d,-\sigma}$，可预测磁性分裂，但它把量子自旋翻转抑制成静态场，无法得到低温Kondo单重态。判定局域矩形成的中间温区与最终零温屏蔽要分开。
\originalsection{Schrieffer--Wolff与对数流}
在局域矩区，把空占据和双占据作为虚中间态，用二阶有效哈密顿量得到$H_K=J_K\vct S_d\cdot\vct s(0)+W\rho(0)$，
\begin{equation}
J_K=2|V|^2\left(\frac1{|\epsilon_d|}+\frac1{\epsilon_d+U}\right)>0,\quad
W=\frac{|V|^2}{2}\left(\frac1{|\epsilon_d|}-\frac1{\epsilon_d+U}\right),
\end{equation}
这里$\vct s=\frac12c^\dagger\vct\sigma c$，$W$的系数随局域场归一而变，但$J_K$反铁磁性及粒子--空穴对称时$W=0$不变。虚跃迁把两个自旋交换，解释电荷虽冻结而自旋仍与金属耦合。

弱耦合散射的粒子和空穴虚过程给$J_K^2D_\uparrow\log(D/E)$修正。令$l=\log(D_0/D)$，常用约定$j=D_\uparrow J_K$满足$\dd j/\dd l=2j^2+O(j^3)$，解$j(l)=j_0/(1-2j_0l)$。正耦合随降低能量增长，尺度$\kb T_K\sim D_0\exp[-1/(2j_0)]$。当$j\sim1$，该微扰方程已失效；形式发散不是物理量真的无穷大，而是强耦合新描述开始的信号。
\originalsection{单重态与低温共振}
单通道、自旋$1/2$、反铁磁耦合的强耦合极限中，杂质与局域传导自旋形成单重态，能量$-3J_K/4$；三重态为$J_K/4$。低温保留的传导粒子散射有$\pi/2$相移，并有弱的残余相互作用。零温杂质谱在$E_F$附近产生宽约$\kb T_K$的Kondo共振，它不是裸$\epsilon_d$电荷共振的平移。

金属电阻在中间温区有随降温上升的$\log T$修正，低温饱和并有$T^2$修正。磁化率趋有限常数而非继续Curie发散，杂质比热$C_{
m imp}\propto T$。单通道完全屏蔽问题可用局域Fermi液体解释；过屏蔽多通道、自旋大于可用通道能力或伪隙态密度的问题可能不同，本章结论不覆盖这些模型。

\originalchapter{谱函数与Fermi液体}
\label{ch:fermi}
\source{8.512第23--25讲。}
\originalsection{单粒子谱的物理含义}
推迟电子Green函数定义$G^R(k,t)=-(\ii/\hbar)\theta(t)\avg{\{c_k(t),c_k^\dagger(0)\}}$。Fourier变换后$G^R(k,\omega)=[\hbar\omega-\epsilon_k+\mu-\Sigma^R(k,\omega)]^{-1}$，谱函数$A(k,E)=-\pi^{-1}\operatorname{Im}G^R(k,E/\hbar)$满足$\int A(k,E)\dd E=1$。占据数$n_k=\int f(E)A(k,E)\dd E$。光电子发射主要测$fA$并乘矩阵元；隧穿测动量积分，二者并非同一观测量。

若自能在费米面附近可展成$\Sigma^R\simeq\Sigma_0+(1-Z^{-1})E+\delta k\,\partial_k\Sigma-\ii C[E^2+(\pi\kb T)^2]$，则
\begin{equation}
G^R\simeq\frac Z{E-v_F^*\hbar(k-k_F)+\ii\Gamma(E,T)}+G_{
m incoh}.
\end{equation}
$Z=[1-\partial_E\Sigma']^{-1}$是单粒子极点权重，不是总粒子数比例。零温$n_k$在费米面有跳跃$Z$，其余权重在非相干连续谱中。衰减率$\Gamma\propto E^2$使$\Gamma/E\to0$，低能激发才有可分辨的准粒子意义。若自能非解析导致$Z=0$或$\Gamma/E$不消失，不能继续使用此极点展开。
\originalsection{能量泛函与Landau参数}
Fermi液体的状态以准粒子分布变化$\delta n_{k\sigma}$标记，能量展开
\begin{equation}
\delta E=\sum_{k\sigma}\epsilon^*_{k}\delta n_{k\sigma}+\frac1{2V}\sum_{kk'\sigma\sigma'}f_{\sigma\sigma'}(k,k')\delta n_{k\sigma}\delta n_{k'\sigma'}.
\end{equation}
$f$是改变其他准粒子分布时一个准粒子能量的变化，不等同于裸库仑势。各向同性、自旋旋转对称时分成$s,a$通道，再在夹角的Legendre多项式上展开，定义$F_l^{s,a}=D^*(E_F)f_l^{s,a}$，本文$D^*$为双自旋态密度。均匀密度与自旋响应为
\begin{equation}
\frac{\partial n}{\partial\mu}=\frac{D^*(E_F)}{1+F_0^s},\qquad \chi_s=\frac{\mu_0\mu_B^2D^*(E_F)}{1+F_0^a},\qquad\frac{C_V}{V}=\frac{\pi^2}3D^*(E_F)\kb^2T.
\end{equation}
比热只测热态数，压缩率和磁化率还含分布自洽调整，所以质量增强不足以预测它们。

在Galilean不变连续系统，均匀整体boost的电流必须由裸质量决定，准粒子的附带回流给$m^*/m=1+F_1^s/3$。在晶格上Galilean不变性不存在，这条质量关系一般失效。稳定性要求各角动量通道$1+F_l^{s,a}/(2l+1)>0$；某个系数变负意味着相应费米面形变可降低能量，即Pomeranchuk不稳定。
\originalsection{碰撞、零声与Kondo低能态}
线性化无碰撞动力学$\partial_t\delta n+\vct v_F\cdot\nabla_r[\delta n-f'(\epsilon)\delta\epsilon]=0$，相互作用使密度畸变自洽。只保留$F_0^s$的三维纵模令$s=\omega/(qv_F)$，角积分给
\begin{equation}
1+F_0^s\left[1-\frac s2\log\frac{s+1}{s-1}\right]=0.
\end{equation}
排斥$F_0^s>0$有$s>1$解，模在粒子--空穴连续谱外，是零声。碰撞足够频繁$\omega\tau\ll1$时局域平衡建立，得到普通第一声；两个声不是相同极限。粒子--空穴相空间在低温宽度$T$内提供两个独立能量积分，所以准粒子碰撞率$\propto T^2$，有限能量对应$E^2$。

Kondo单重态附近，杂质散射相移可展开$\delta_\sigma(E,\delta n)=\pi/2+\alpha E/(\kb T_K)-\phi\,\delta n_{-\sigma}/[D_\uparrow\kb T_K]+\cdots$。能量项与残余相互作用项共同决定磁化率和比热，局域液体的Wilson比在标准完全屏蔽极限趋2。该结果需要强耦合固定点分析，前章一回路对数流只给$T_K$尺度，不能单独证明精确低温比值。

\volpart{第三编\quad 几何相位与现代量子多体（8.513）}
\originalchapter{绝热几何、Chern数与波包}
\label{ch:berry}
\source{8.513第1--6讲幻灯片，结合8.511能带与Hall效应。}
\originalsection{Berry相位的推导}
缓慢参数$R(t)$使$H(R)\ket{n(R)}=E_n(R)\ket{n(R)}$。若能级孤立且$|\hbar\avg{m|\dot H|n}|\ll|E_m-E_n|^2$，绝热态可写$\ket{\psi(t)}=\e^{-(\ii/\hbar)\int E_n\dd t}\e^{\ii\gamma_n(t)}\ket{n(R(t))}$。代回Schrödinger方程后动力学相位抵消，剩$\dot\gamma_n=\ii\avg{n|\partial_t n}$。因此$\gamma_n=\int\mathcal A\cdot\dd R$，$\mathcal A=\ii\avg{n|\nabla_Rn}$。

本征态相位$\ket n\to\e^{\ii\zeta(R)}\ket n$使$\mathcal A\to\mathcal A-\nabla\zeta$。闭合路的$\e^{\ii\gamma}$不变，开路积分则要同时指定端点相位。Berry联络是选局部相位规范后的几何对象，曲率$\mathcal F_{ij}=\partial_i\mathcal A_j-\partial_j\mathcal A_i$规范不变。用本征方程微分，可把它写成
\begin{equation}
\mathcal F_{ij}=-2\operatorname{Im}\sum_{m\ne n}\frac{\avg{n|\partial_iH|m}\avg{m|\partial_jH|n}}{(E_n-E_m)^2},
\end{equation}
此式适合计算，且说明能隙闭合附近曲率为何可能变大。
\originalsection{两能级系统与拓扑整数}
$H=d_0I+\vct d\cdot\vct\sigma$的本征值$E_\pm=d_0\pm|d|$。用球角表示$\hat d$，上能级可选$\ket+= (\cos(\theta/2),\e^{\ii\phi}\sin(\theta/2))^T$，则$\mathcal A_\phi=-\sin^2(\theta/2)$、$\mathcal F_{\theta\phi}=-\sin\theta/2$；下能级符号相反。闭合路径的相位为带符号立体角的一半。

二维绝缘体占据带的曲率积分定义$C=(2\pi)^{-1}\int_{\mathrm{BZ}}\mathcal F_{k_xk_y}\dd^2k$。为理解整数性，把周期Brillouin环面覆盖为局部规范片；每片上由Stokes定理曲率积分化为边界联络积分，相邻片的差为$-\dd\zeta$。相位过渡函数$\e^{\ii\zeta}$沿闭合边界只能有整数绕数，所以总积分是$2\pi$的整数。$C\ne0$时无法在整个环面选一组既光滑又周期的单带本征矢。

对于两带模型，下带$C=(4\pi)^{-1}\int\hat d\cdot(\partial_{k_x}\hat d\times\partial_{k_y}\hat d)\dd^2k$，在本章联络约定下取这个正号。它是从Brillouin区到球面的带向映射次数。连续改变参数且保持$|d|>0$不能改变整数，因此相变必须闭隙或破坏所要求的对称性。单个连续Dirac锥可贡献半整数的局部曲率积分，完整晶格的其他动量区使总数为整数，不能把未经正则化的一只锥当成独立二维晶格绝缘体。
\originalsection{几何修正的运动方程}
孤立带波包的作用量取
\begin{equation}
L=(\hbar\vct k-e\vct A)\cdot\dot{\vct r}+\hbar\mathcal A_n(\vct k)\cdot\dot{\vct k}-E_n(\vct k)+e\phi.
\end{equation}
对$\vct r,\vct k$变分得到$\hbar\dot{\vct k}=-e(\vct E+\dot{\vct r}\times\vct B)$、$\dot{\vct r}=\hbar^{-1}\nabla_kE_n-\dot{\vct k}\times\vct\Omega_n$。反常速度在电场下给横向电流；填满二维带时普通群速度积分为零，而曲率积分给$\sigma_{xy}=Ce^2/h$，整体符号还随Hall电流与场方向定义。此推导把拓扑整数与可测输运相连。

有磁场时相空间体积也修正为$[1+(e/\hbar)\vct B\cdot\vct\Omega]$倍，因运动坐标不再有简单的正则辛结构。对态密度、磁化与输运做积分时要保留这项；只改速度可能违反连续性。强带间混合、闭隙点及极强外场超出单带波包近似。
\originalsection{拓扑泵与边界}
一维绝缘体若参数随时间周期变化，将$(k,t)$视为二维参数环面，泵送电荷$Q=eC$。体极化在一个周期内改变整数个电荷，却不要求瞬间产生可压缩金属态。Chern绝缘体与平凡绝缘体相邻时，两侧拓扑数差要求净手性边缘通道。可由Dirac质量畴壁直接求出：取$H=k\sigma_z+\ii\sigma_y\partial_y+m(y)\sigma_x$，若$m(y)$由负变正，零横向模满足一支$\psi(y)\propto\exp[-\int^y m(y')\dd y']$，另支指数增长不可归一。唯一可归一支沿边缘有$E=-k$。一支被舍掉的原因是空间归一化，不能将它当成另一个物理反向边缘态。
\originalsection{一般辛结构与相空间守恒}
课程的半经典框架还可脱离正则坐标来理解。设相空间变量为$\xi^I$，一阶作用量$L=a_I(\xi)\dot\xi^I-H(\xi)$，则变分给$b_{IJ}\dot\xi^J=\partial_IH$，$b_{IJ}=\partial_Ia_J-\partial_Ja_I$。$b$反对称且闭合$\partial_Ib_{JK}+\partial_Jb_{KI}+\partial_Kb_{IJ}=0$。若$b$可逆，$b^{IJ}$定义Poisson括号，其Jacobi恒等式来自这条闭合关系；任意指定一个反对称矩阵而不检查闭合性，未必是合法的哈密顿动力学。

相空间体积不是总为$\dd^{2d}\xi$，而是$\dd\Gamma=|\operatorname{Pf}b|\dd^{2d}\xi=\sqrt{\det b}\dd^{2d}\xi$，符号取定向后固定。辛流保持该体积：可由$\mathcal L_Xb=\dd(\iota_Xb)+\iota_X\dd b=\dd(\dd H)=0$推出其$d$次楔积也不变。因此占据函数$g$满足$\partial_tg+\dot\xi^I\partial_Ig=0$，等价的连续性式为$\partial_t(\sqrt{\det b}\,g)+\partial_I(\sqrt{\det b}\,g\dot\xi^I)=0$。Berry曲率导致的态密度修正正是这个测度的具体例子。

与固定温度环境耦合时，若只加扩散流$J^I=-\gamma^{IJ}\partial_Jg$，平衡趋向常数占据，对应无穷温；有限温需漂移配合。经典粒子可取$J^I=-\gamma^{IJ}g\partial_J(\log g+\beta H)$，使$g_0\propto\e^{-\beta H}$无耗散流。费米与玻色占据分别改为$J^I=-\gamma^{IJ}g(1\mp g)\partial_J[\log(g/(1\mp g))+\beta(H-\mu)]$。这里$\gamma$应有非负对称耗散部分，才能使相应自由能下降。这个开放体系模型保持粒子数却不要求体系自身能量守恒；封闭系统的碰撞项另须满足能量与动量守恒。

\originalchapter{相互作用玻色子与相位动力学}
\label{ch:boson}
\source{8.513第7--15讲组：相互作用玻色子、超流与低维有效理论。}
\originalsection{相干态与路径积分}
玻色相干态满足$b\ket z=z\ket z$，完备关系$\int\dd^2z\,\ket z\bra z/\pi=I$（采用归一化态及对应权重）。把$\e^{-\beta H}$分为很多小虚时间片并逐片插入相干态，相邻态重叠产生一阶时间导数。对正规序哈密顿量得到
\begin{equation}
Z=\int\mathcal D(z^*,z)\e^{-S_E/\hbar},\qquad S_E=\int_0^{\beta\hbar}\dd\tau\,[\hbar z^*\partial_\tau z+H(z^*,z)].
\end{equation}
玻色场虚时间周期；费米相干态使用Grassmann变量且反周期，不能用普通复变量的积分规则替代。经典实时间作用量为$L=\ii\hbar z^*\dot z-H$，差一个总导数可把Berry项写为$\hbar n\dot\theta$。本文取$z=\sqrt n\e^{-\ii\theta}$以使该项为正，若改用$\e^{\ii\theta}$，相位定义和共轭动量应同时改号。
\originalsection{Bogoliubov声子}
均匀弱排斥玻色气$H=\int\dd^dx[\hbar^2|\nabla\psi|^2/(2m)+g|\psi|^4/2-\mu|\psi|^2]$。均匀鞍点$n_0=\mu/g$，写$\psi=\sqrt{n_0}+\delta\psi$并保留二次项，$k,-k$块含$(\epsilon_k+gn_0)b_k^\dagger b_k$与$gn_0(b_kb_{-k}+\mathrm{h.c.})/2$。保持对易关系的变换$u_k^2-v_k^2=1$给
\begin{equation}
E_k=\sqrt{\epsilon_k(\epsilon_k+2gn_0)},\quad u_k^2,v_k^2=\frac12\left(\frac{\epsilon_k+gn_0}{E_k}\pm1\right),\quad\epsilon_k=\frac{\hbar^2k^2}{2m}.
\end{equation}
低$k$线性$E_k\simeq\hbar c|k|$，$c=\sqrt{gn_0/m}$；高$k$回到粒子色散。声子是凝聚密度与相位耦合的集体模。$v_k^2$给零温耗尽，计算时若接触势在高维产生紫外问题，应以散射长度重整化$g$，不能无限外推裸接触模型。
\originalsection{密度--相位有效作用量}
写$\psi=\sqrt{n_0+\delta n}\e^{-\ii\theta}$，低梯度保留$L=\int\dd x[\hbar(n_0+\delta n)\dot\theta-g(\delta n)^2/2-\hbar^2n_0(\nabla\theta)^2/(2m)]$。密度涨落无时间导数，是相位的共轭变量。极小化$\delta n=\hbar\dot\theta/g$并代回，得到
\begin{equation}
L_{\rm eff}=\int\dd x\left[\hbar n_0\dot\theta+\frac{\hbar^2}{2g}\dot\theta^2-\frac{\hbar^2n_0}{2m}(\nabla\theta)^2\right].
\end{equation}
第一项在光滑小涨落中是总导数，但对紧致相位的绕数与扭曲边界有作用，不能随意删去。Euler--Lagrange方程给$\ddot\theta-c^2\nabla^2\theta=0$；声速与Bogoliubov结果相同，证明两个描述在低能一致。

一维常写$S_E/\hbar=(K/2\pi)\int\dd x\dd\tau[v^{-1}(\partial_\tau\theta)^2+v(\partial_x\theta)^2]$。高斯反演在二维$(x,v\tau)$给$\avg{[\theta(x)-\theta(0)]^2}=K^{-1}\log(|x|/a)$，于是$\avg{\e^{\ii\theta(x)}\e^{-\ii\theta(0)}}\propto|x|^{-1/(2K)}$。无穷系统的相位方差发散，所以零温一维并非具有常数远距关联的通常凝聚，而是代数序。有限温虚时间有有限周长，长距离关联进一步衰减。
\originalsection{零模、扭曲与LSM约束}
环长$L$的相位写$\theta=2\pi mx/L+\theta_0(t)+\theta_{\ne0}$。取$L_{\rm eff}=\int[A\dot\theta^2/2-B(\partial_x\theta)^2/2+n_0\hbar\dot\theta]$，零模角动量$p_0=AL\dot\theta_0+n_0\hbar L$，因$\theta_0$为$2\pi$周期坐标而$p_0=\hbar N$，$N\in\mathbb Z$。零模能量
\begin{equation}
E_{N,m}=\frac{\hbar^2(N-n_0L)^2}{2AL}+\frac{2\pi^2Bm^2}{L},\qquad P_{N,m}=\frac{2\pi\hbar mN}{L}.
\end{equation}
末式由平移引起相位整体移$2\pi m\Delta x/L$获得。密度Berry项虽不改声子方程，却改变量子数与扭曲后的动量。

格点环上定义$U=\exp[(2\pi\ii/L)\sum_jj n_j]$，整数局域电荷使边界因子为1，平移关系给$TU T^{-1}=U\exp[-2\pi\ii N/L]$（若平移方向反定，符号反转）。故$U\ket0$的动量与基态差$2\pi\nu$模$2\pi$。局域跳跃经$U$只获得$O(1/L)$相位；正负扭曲合并可消去一阶电流项，变分能量增加至多$O(1/L)$。若$\nu=p/q$既约且$q>1$，重复扭曲产生$q$个不同动量的低能态，排除在保持平移与$U(1)$的条件下唯一且有统一能隙的基态。结论可能由简并、破缺或无隙实现，不能把它只说成“必然超流”。该推导要求局域性、适当的局域能量界和守恒电荷，非局域作用不自动满足。

\originalchapter{量子Ising链与Majorana费米子}
\label{ch:ising}
\source{8.513第7--15讲组及作业5--7。}
\originalsection{Jordan--Wigner变换}
横场Ising链$H=-J\sum_j\sigma_j^x\sigma_{j+1}^x-h\sum_j\sigma_j^z$。不同站自旋算符对易，费米算符应反对易，因此不能直接令$c_j=\sigma_j^-$. 加上前方自旋串，定义
\begin{equation}
a_j=\left(\prod_{l<j}\sigma_l^z\right)\sigma_j^x,\qquad b_j=\left(\prod_{l<j}\sigma_l^z\right)\sigma_j^y,\qquad c_j=(a_j+\ii b_j)/2.
\end{equation}
对于$j<k$，$a_j$与$k$串中的$\sigma_j^z$反对易，其他因子对易，所以$\{a_j,a_k\}=0$；同站$a_j^2=b_j^2=1$，完整关系$\{a_j,a_k\}=\{b_j,b_k\}=2\delta_{jk}$、$\{a_j,b_k\}=0$。又有$\sigma_j^z=-\ii a_jb_j$、$\sigma_j^x\sigma_{j+1}^x=-\ii b_ja_{j+1}$，因此开链$H=\ii\sum_j(Jb_ja_{j+1}+ha_jb_j)$。

这里$c_j$为非局域自旋算符，虽然费米哈密顿量局域，但自旋局域观测量可能对应费米字符串。这个区别会直接影响散射谱，不能将单费米激发自动认作一个自旋翻转产生的态。
\originalsection{能谱与量子临界点}
对无穷链或固定边界扇区Fourier变换，再作费米Bogoliubov变换，正能谱为
\begin{equation}
E_k=2\sqrt{J^2+h^2-2Jh\cos k}.
\end{equation}
$|h|<|J|$为自旋有序区，$|h|>|J|$为横场极化区，$|h|=|J|$隙闭合。$h=J>0$时$E_k\simeq2J|k|$，动态指数$z=1$，非临界质量$\Delta=2|J-h|$。连续低能是左右运动的实Majorana场。

在开链且$h=0$时，$b_j$只与$a_{j+1}$配对，$a_1$和$b_L$完全不在哈密顿量中，组合成一个零能费米自由度。$0<|h/J|<1$时边缘模$\gamma_L\propto\sum_j(h/J)^{j-1}a_j$，归一要求$|h/J|<1$。两个端模重叠导致$\Delta E\propto\e^{-L/\xi}$，$\xi^{-1}=|\log|h/J||$。零模不是在任意参数下存在的孤立粒子，体隙和边界是必要条件。
\originalsection{临界零模与四次项的上临界维数}
上述一维链的能谱在临界点线性，但把这一Gaussian结果推广到任意空间维数，需要检查被丢弃的相互作用。8.513第7--15讲组第25--28页用$d$维横场Ising的序参量$q$说明这点。在有限体积$N=L^d$、平均场临界位置$h=2dJ$，均匀模的二次$q^2$项消失，低能相空间拉氏量为
\begin{equation}
L_0=Np\dot q-\frac{Nh}{2}p^2-Ngq^4.
\end{equation}
令$\widetilde p=\sqrt Np,\ \widetilde q=\sqrt Nq$，则$H_0=h\widetilde p^2/2+(g/N)\widetilde q^4$。量子不确定性$\widetilde p\widetilde q\sim1$使基态估算为$h/\widetilde q^2+g\widetilde q^4/N$；两项平衡给$\widetilde q\sim N^{1/6}$及$E_0\sim N^{-1/3}$（省去固定的$h,g$系数）。零模无论多高维都不可简单按自由振子处理。

对非零波矢，Gaussian能量$\omega_k\sim J|k|$，其真空涨落$\widetilde q^2\sim1/|k|$。四次校正$\delta E_k\sim g/(Nk^2)$相对谐振能$E_k\sim J|k|$的比值为$\sim g/(JNk^3)$。最低$k\sim L^{-1}=N^{-1/d}$，所以比值随大小如$N^{3/d-1}$：$d>3$趋零，$d=3$临界，$d<3$增长。这个模式估算与连续场论的幂次计数相符。

取动态指数$z=1$，Euclidean作用写$S=\int\dd\tau\,\dd^dx\,[(\partial_\tau\phi)^2+c^2(\nabla\phi)^2+r\phi^2+u\phi^4]/2$，其中四次项系数的常数$1/2$已吸进$u$。保持梯度项无量纲给$[\phi]=L^{-(d-1)/2}$、$[r]=L^{-2}$、$[u]=L^{d-3}$。粗粒化长度放大$b$时无量纲四次耦合按$u'\simeq b^{3-d}u$变换：$d<3$相关，$d=3$边缘，$d>3$无关。这里的上临界\emph{空间}维数是$d_c=3$，因为量子时间还贡献一维，时空维数$D=d+1=4$；经典静态$\phi^4$模型的上临界空间维数才是4。$d=1$链已知的精确Majorana解不能被高维Gaussian幂次计数替代。

\originalsection{环上的边界与宇称投影}
周期自旋链的最后一条键穿过Jordan--Wigner串，费米边界条件取决于总宇称$P=\prod_j\sigma_j^z=(-1)^{N_f}$。以上约定通常得到$c_{L+1}=-Pc_1$：偶宇称采用反周期，奇宇称采用周期。每个费米边界模型还须投影到匹配宇称，不能把两个完整Fock空间相加，否则会把原$2^L$个自旋态双计。

反周期允许$k=2\pi(n+1/2)/L$，所有$k$可两两配对，Bogoliubov真空偶宇称。周期有$k=0,\pi$自配对点，其占据由$h-J$与$h+J$的符号决定，在跨越临界点时可改变真空宇称。非匹配扇区的最低物理态要额外加入准粒子。有限环的最低能态选择因此不能只对两个动量格点求和后忽略宇称约束。

有序区两种自旋宇称最低态的能量差为指数小而非严格零。能谱$\sqrt{J^2+h^2-2Jh\cos k}$在复$k$平面的最近支点满足$\cosh\kappa=(J^2+h^2)/(2Jh)$，故$\kappa=|\log(J/h)|$。Poisson求和把两种边界格点之差写成高Fourier系数，最近支点给$\e^{-\kappa L}$与可能的代数前因子，这与端模局域长度一致。临界点支点到实轴，指数估计失效，有限尺寸能隙改为$1/L$。
\originalsection{对偶与局域可观测量}
定义对偶键自旋$\mu_{j+1/2}^z=\sigma_j^x\sigma_{j+1}^x$、$\sigma_j^z=\mu_{j-1/2}^x\mu_{j+1/2}^x$，交换$J,h$后哈密顿量形状保持，临界点在自对偶处。环上全局约束与边界仍必须跟踪，对偶不在无条件下把任意边界Hilbert空间一一对应。

自旋序参数$\sigma^x$包含费米字符串；$\sigma^z$是局域费米双线性。因而动态$\sigma^z$相关通常可创建两费米子，谱连续边缘由$E_k+E_{q-k}$决定。若探针改变自旋宇称，需把它作用后的边界扇区一并处理。直接把$E_k$画成中子散射的唯一$\delta$峰会漏掉算符选择规则。

\originalchapter{手性场、中心荷与玻色化}
\label{ch:chiral}
\source{8.513相互作用玻色子组后半、作业6--8。}
\originalsection{低能场与热态计数}
一维线性色散把右行与左行激发分开。单支右行玻色模$k_n=2\pi n/L$，$n>0$，激发能$\sum_n(2\pi\hbar v/L)nN_n$。给定总整数能量$n$的简并数是整数分拆数$p(n)$，生成函数$\prod_{m\ge1}(1-q^m)^{-1}$。高温于有限尺寸能级间隔而低温于微观带宽时，
\begin{equation}
\log Z=-\sum_{m>0}\log(1-\e^{-\beta 2\pi\hbar vm/L})\simeq\frac L{2\pi\beta\hbar v}\int_0^\infty[-\log(1-\e^{-x})]\dd x=\frac{\pi L}{12\beta\hbar v}.
\end{equation}
展开$-\log(1-\e^{-x})=\sum_{r\ge1}\e^{-rx}/r$并逐项积分给$\sum r^{-2}=\pi^2/6$。于是$E=\pi L\kb^2T^2/(12\hbar v)$、$C=\pi L\kb^2T/(6\hbar v)$。一支Majorana费米子只有半个复费米自由度，等价中心荷$1/2$，热能系数为以上一半。

采用$ c_{\rm tot}=c_R+c_L$时，一般同速CFT有$E/L=\pi c_{\rm tot}\kb^2T^2/(12\hbar v)$。非手性教科书若用$c=c_R=c_L$，同一式写$\pi c\kb^2T^2/(6\hbar v)$。先说明中心荷计数，才能避免两倍差。手性净值$c_-=c_R-c_L$控制净能量流与手性异常，而热态总密度取决于$c_{\rm tot}$，两者不同。
\originalsection{态密度与虚时间相关}
由正则热能关系求微正则熵。令$E=aT^2$，$S=\int C\dd T/T=2aT=2\sqrt{aE}$，$a=\pi c_{\rm tot}L\kb^2/(12\hbar v)$。大能量密度但仍在低能场论窗口时，
\begin{equation}
D_{\rm many}(E)\sim\exp\left[\sqrt{\frac{\pi c_{\rm tot}LE}{3\hbar v}}\right],
\end{equation}
指数前的幂因子依具体边界扇区与有限尺寸。分拆数渐近$p(n)\sim\e^{\pi\sqrt{2n/3}}/(4\sqrt3n)$代入$E=2\pi\hbar vn/L$也得到$c_{\rm tot}=1$的结果。

反周期Majorana场写$\lambda(x)=L^{-1/2}\sum_k\lambda_k\e^{\ii kx}$，$k=2\pi(n+1/2)/L$，$\lambda_{-k}=\lambda_k^\dagger$。零温$k>0$为湮灭模，所以$\tau>0$的相关函数是$L^{-1}\sum_{k>0}\e^{\ii kx-vk\tau}$。几何级数给
\begin{equation}
G(x,\tau)=\frac{\ii}{2L\sin[\pi(x+\ii v\tau)/L]}\ \longrightarrow\ \frac{\ii}{2\pi(x+\ii v\tau)}.
\end{equation}
$\tau<0$需按费米时间排序换号，并选择相应解析支。局域反对易关系由完整正负动量和得到$\{\lambda(x),\lambda(y)\}=\delta(x-y)$，不能仅以正动量构造实场。
\originalsection{手性玻色场与顶角算符}
紧致场$\phi\sim\phi+2\pi$的多分量作用量为
\begin{equation}
L=\frac\hbar{4\pi}\int\dd x[K_{IJ}\partial_x\phi_I\partial_t\phi_J-V_{IJ}\partial_x\phi_I\partial_x\phi_J].
\end{equation}
$K$为非退化对称整数矩阵，$V$正定确保能量稳定。由于一阶时间导数，共轭动量受约束，正确的Dirac括号给$[\phi_I(x),\phi_J(y)]=\ii\pi(K^{-1})_{IJ}\operatorname{sgn}(x-y)$，整体方向随手性约定。允许单值算符$V_l=\e^{\ii l_I\phi_I}$，$l\in\mathbb Z^N$。BCH公式给远距交换相位$\exp[-\ii\pi l^TK^{-1}l'\operatorname{sgn}(x-y)]$。

$K=(1)$时$\e^{\ii\phi}$具有费米交换符号，相关按$1/(x+\ii v\tau)$，所以可将其识别为手性费米场。左右两支$K=\operatorname{diag}(1,-1)$及密度--密度相互作用可映射一维相互作用费米子。它们的低能关联通常有相互作用依赖的幂律，单粒子准粒子极点消失；这一维Luttinger液体与三维Landau Fermi液体的用语需区分。
\originalsection{分数化与边缘异常}
顶角$l$的拓扑自旋为$\theta_l=\e^{\ii\pi l^TK^{-1}l}$，两种顶角完整绕行相位$M_{ll'}=\e^{2\pi\ii l^TK^{-1}l'}$。局域物理粒子为$l=Kn$；在玻色体系$K$对角须为偶数，使这些局域粒子自旋为1。一般整数$l$不必局域，可作为二维体任意子的边缘终点。二维Abel拓扑序按$l\sim l+Kn$分类，因此种类数$|\det K|$。

单支手性场无法作为具有相同局域自由度的独立一维晶格低能谱：一个有界周期一维能带的右行与左行交叉整体相抵，不能仅有净手性模式。二维体的边缘允许净手性，因为外场下体响应提供异常流入。异常说明边缘守恒方程需与体一起解释，不意味着整个物理体系不守恒。$c_-\ne0$的手性能量流也阻止在保持相应结构的条件下把边缘全部开隙。

\originalchapter{自由费米拓扑相与质量矩阵}
\label{ch:spt}
\source{8.513第16--19讲幻灯片；分类框架的数学定理按明确范围引用。}
\originalsection{有隙哈密顿量的稳定等价}
非相互作用、保持粒子数的多体哈密顿量$\hat H=\sum_{ab}M_{ab}c_a^\dagger c_b$由Hermitian单体矩阵决定。以费米能为零，有隙指$M$无零本征值。连续变形$M\to Q=\operatorname{sgn}M$保留本征矢与正负能级数，且$Q^2=1$。在零维，负能级数不能在不闭隙下改变；加上完全局域填满或空的平凡轨道后，分类使用稳定差值。

周期体系给每个$k$的占据投影$P(k)=(1-Q(k))/2$。若$P$有$p$维像与$q$维核，其值域为复Grassmann空间$U(p+q)/(U(p)\times U(q))$。高维拓扑分类不是只数某个$k$的负能级，而是研究整个Brillouin区到此空间的连续映射及其对称约束。加入平凡带把小矩阵的偶然障碍与稳定拓扑分开。
\originalsection{对称类与Clifford扩张}
时间反演$T$和粒子--空穴操作$C$为反幺正，分别满足$TH(k)T^{-1}=H(-k)$与$CH(k)C^{-1}=-H(-k)$；其平方可为$\pm1$。两者的乘积可给幺正手征对称$\Gamma H\Gamma^{-1}=-H$。由这些条件产生十种自由费米对称类。BdG中的粒子--空穴关系包含Nambu冗余，不应把它误当独立物理守恒电荷。

连续Dirac模型$H(k)=\sum_{i=1}^d k_i\gamma_i+mM$若要谱满足$E^2=k^2+m^2$，需$\{\gamma_i,\gamma_j\}=2\delta_{ij}$、$\{M,\gamma_i\}=0$、$M^2=1$。将对称操作在实Majorana基中转为与$M$对易或反对易的生成元，寻找允许质量相当于给Clifford代数增加一个生成元。允许质量空间的连通分支决定能否把两种符号质量连接而不闭隙。

Bott周期性给复分类空间两周期、实分类空间八周期。该同伦定理属于所调用的数学输入，本册不声称用有限几个Pauli矩阵例子证明其一般版本。应用时先明确自由、稳定、局域、有隙、保持指定内部对称的假设。常见强拓扑分类如下，晶格额外弱指标及晶体对称指标不含在表内。
\begin{center}\begin{tabular}{llll}\toprule 对称类&对称条件&空间维数&稳定强指标\\\midrule
A&仅粒子数守恒&2&$\mathbb Z$（Chern）\\
AIII&手征，无$T,C$&1&$\mathbb Z$（绕数）\\
D&$C^2=+1$&1&$\mathbb Z_2$\\
D&$C^2=+1$&2&$\mathbb Z$\\
AII&$T^2=-1$&2、3&$\mathbb Z_2$\\
BDI&$T^2=C^2=+1$&1&$\mathbb Z$\\\bottomrule\end{tabular}\end{center}
\originalsection{SSH链与Kitaev链}
SSH链的两个亚晶格跃迁$t_1,t_2$给$H(k)=(t_1+t_2\cos k)\sigma_x+t_2\sin k\sigma_y$。手征算符$\sigma_z$与$H$反对易。复函数$q(k)=t_1+t_2\e^{\ii k}$的闭曲线绕原点整数$\nu=(2\pi\ii)^{-1}\int q^{-1}\partial_kq\dd k$。$|t_2|>|t_1|$绕一圈，反之零；$|t_1|=|t_2|$闭隙。若加$\sigma_z$势差，手征破坏，可在不闭隙下绕开原点，这说明保护对称的作用。

Kitaev链$H=-\sum_j[tc_j^\dagger c_{j+1}+\Delta c_jc_{j+1}+\mathrm{h.c.}]-\mu\sum_j(n_j-1/2)$。取实$t,\Delta$，BdG谱$E_k^2=(-2t\cos k-\mu)^2+4\Delta^2\sin^2k$。在$\Delta\ne0$时闭隙只能于$k=0,\pi$且$\mu=\mp2t$。Majorana质量在这两点符号乘积为$\operatorname{sgn}[(\mu+2t)(\mu-2t)]$，负号区$|\mu|<2|t|$非平凡，开链有两端零模。允许保持$C$但破坏额外实时间反演的扰动时只剩$\mathbb Z_2$保护，不能用实参数特殊模型的整数分类替代一般D类。
\originalsection{体边对应与相互作用边界}
质量畴壁上的零模展示体边对应：允许质量由一相换到另一相，若无法全程保持有隙，界面必须出现低能态或其他相结构。边缘的具体速度、重构和附加平凡通道可变，体指标只固定净异常内容。

自由费米分类不是所有相互作用物态的分类。多体相可有内禀拓扑序，也可能使自由整数指标经相互作用约化；例如特定一维BDI体系在保持相应对称与局域作用下整数可约为八周期。这个例子只说明假设边界，本稿未推导所有相互作用分类。一般SPT是保持对称时不能连续化到平凡态的短程纠缠相，破坏保护对称后可平凡化；内禀拓扑序则不需这些保护对称，并有长程纠缠与任意子。

\originalchapter{拓扑序、弦算符与融合}
\label{ch:topological}
\source{8.513第20--24讲：$\mathbb Z_2$环液体、弦网与拓扑激发；自绘图不复用第三方图片。}
\originalsection{局域不可辨识与环液体}
对称破缺相通常用局域序参量区分，例如磁化方向。拓扑序的不同基态在大系统中对任何有限支撑的局域算符近乎给相同矩阵元，却可由绕系统的非局域环区分。简并依赖空间拓扑，激发的融合与绕行提供进一步数据，不能只以“没有局域序参量”作为完整定义。

方格边上放自旋，toric-code哈密顿量取$H=-J_e\sum_sA_s-J_m\sum_pB_p$，$A_s=\prod_{l\ni s}\sigma_l^x$、$B_p=\prod_{l\in\partial p}\sigma_l^z$。任意星与面共享0或2条边，交换两个Pauli串的符号为$(-1)^2=1$，所以全部对易，且每项平方为1。基态满足$A_s=B_p=+1$。在$\sigma^z$表象，$B_p$约束和$A_s$翻转可形成对闭环图样的相干叠加；具体把环画在直接或对偶格点上取决于表象。

\begin{center}\begin{tikzpicture}[scale=.8]
\draw[step=1,gray] (0,0)grid(4,3);
\foreach \x in {0,...,4}\foreach \y in {0,...,3}{\fill(\x,\y)circle(1.3pt);}
\draw[structurecolor,line width=2pt](1,1)--(2,1)--(2,2)--(1,2)--cycle;
\draw[main,line width=2pt](3,1)--(2,1);\draw[main,line width=2pt](3,1)--(4,1);\draw[main,line width=2pt](3,1)--(3,0);\draw[main,line width=2pt](3,1)--(3,2);
\node[below]at(1.5,-.4){面$B_p$};\node[below]at(3.5,-.4){星$A_s$};
\end{tikzpicture}\end{center}
图只标算符支撑，并非真实材料的显微结构。
\originalsection{简并计数与任意子}
闭合环面上$\prod_sA_s=\prod_pB_p=1$，因为每条边出现两次。边自旋数$N_l$、顶点数$N_s$、面数$N_p$满足Euler关系$N_s-N_l+N_p=0$。独立约束共$N_s+N_p-2$，基态维数$2^{N_l-(N_s+N_p-2)}=4$。球面Euler数为2则维数1；genus为$g$的闭曲面为$4^g$。这是一个实际可算的拓扑简并来源。

沿路径$\gamma$作用$W_e(\gamma)=\prod_{l\in\gamma}\sigma_l^z$，只在两个端点把$A_s$翻为$-1$，产生两个$e$。对偶路径用$\sigma^x$产生两个$m$。移动端点不在路径中部积累额外激发，说明能量只依端点而非弦长。$e$弦与$m$弦交一次时反对易，因此绕行得到$-1$互统计相位。$e$、$m$各自交换为玻色，复合$\epsilon=e\times m$为费米；融合群为$\mathbb Z_2\times\mathbb Z_2$，$e^2=m^2=\epsilon^2=1$。

非收缩Wilson环在基态子空间形成两对反对易逻辑算符，解释4维空间中的两个拓扑量子比特。局域扰动要通过绕系统的高阶过程才连接这些态，有限体系可有$\e^{-L/\xi}$劈裂。有限温下热任意子运动可能产生逻辑错误，因此二维模型的零温拓扑保护不等于任意温度下自纠错存储。
\originalsection{长程纠缠与拓扑熵}
有隙二维拓扑态的区域熵常为$S(A)=\alpha|\partial A|-\gamma+\cdots$，$\alpha$依微观截断，$\gamma=\log\mathcal D$，$\mathcal D=\sqrt{\sum_ad_a^2}$。toric code四类Abel任意子皆$d_a=1$，$\mathcal D=2$，所以$\gamma=\log2$。通过适当三个或更多区域熵的组合消除边界项，可提取这个常数。

对闭环叠加态，跨边界的二值弦端点受总偶宇称约束，允许配置数较独立端点少一因子2，直接解释$-\log2$。一般拓扑熵公式需要有隙、区域尺度大于关联长等条件，临界态和含角点的任意单区拟合不能直接测$\gamma$。有限深度局域幺正不能改变长程拓扑数据；SPT在不要求保持对称时可由此类电路化为平凡态，内禀拓扑序不能。
\originalsection{融合空间与非Abel弦网}
写$a\times b=\sum_cN_{ab}^{c}c$。$N_{ab}^{c}$是两个激发融合为$c$的独立通道数，不是概率。球面上若总拓扑电荷须为真空，$n$个任意子允许态数是乘积融合中真空出现的次数。Abel任意子每次融合唯一，量子维数$d_a=1$；非Abel任意子有多个通道，态数按$d_a^n$增长。

Fibonacci例子只有$1,\tau$，规则$\tau\times\tau=1+\tau$。设$n$个$\tau$融合为$1$与$\tau$的通道数为$D_n,F_n$，则$D_{n+1}=F_n$、$F_{n+1}=D_n+F_n$，$D_1=0,F_1=1$。增长矩阵的最大本征值$d_\tau=(1+\sqrt5)/2$，这是非整数内部态数增长率，而不是单个粒子的普通Hilbert空间维数。四个$\tau$总电荷真空有2个融合态，可编码一量子比特。

融合的不同括号顺序是同一空间的不同基底，由$F$矩阵联系；粒子交换由$R$矩阵表示。四粒子五种括号路线的一致性要求五边形恒等式，融合与编织相容还需六边形恒等式。在Fibonacci规范中可取
\begin{equation}
F_{\tau}^{\tau\tau\tau}=\begin{pmatrix}d^{-1}&d^{-1/2}\\d^{-1/2}&-d^{-1}\end{pmatrix},\qquad d^2=d+1.
\end{equation}
用$d^{-2}+d^{-1}=1$可直接核对$F^2=I$；仅满足融合规则并不足以保证存在一致的量子态模型，还需完整的$F,R$数据。弦网局域重接规则把这些代数相容性转换为晶格波函数的一致性。
\originalsection{范畴语言与有效规范场}
在融合范畴中，对象是激发类型，态空间或局域过程是态射，张量积表示把激发并置，结合子由$F$实现；编织结构在二维描述绕行。这里“高阶”语言组织粒子、弦、界面等不同维的缺陷，不能把它只当新名称的相图。完整高阶范畴分类超出这门公开幻灯片可重建的数学证明范围；本稿用上述可计算融合例子给出课程的物理入口，不宣称证明一般分类。

Abel拓扑序可由$K$矩阵Chern--Simons作用量$S=\hbar(4\pi)^{-1}\int K_{IJ}a_I\wedge\dd a_J$描述。带整数向量$l$的Wilson线积分掉规范场后给$\theta_l=\e^{\ii\pi l^TK^{-1}l}$与互绕相位$\e^{2\pi\ii l^TK^{-1}l'}$，与前章边缘结果一致。$\mathbb Z_2$序用$K=\left(\begin{smallmatrix}0&2\\2&0\end{smallmatrix}\right)$，$l_e=(1,0),l_m=(0,1)$，自旋为1，互绕$-1$，复合自旋$-1$。$|\det K|=4$同环面简并匹配。非AbelFibonacci弦网不能以一个普通Abel整数$K$完整代替，课程所述双Fibonacci对应更一般的非Abel拓扑场论。

\volpart{第四编\quad 公开作业与中文解答}
\originalchapter{8.511作业1--4：晶格与能带}
\psource{Patrick A. Lee，8.511 Fall 2004，公开Problem Sets 1--4；保留原编号。}
\originalsection{作业1（2004年9月21日截止）}
\begin{exercise}[原题1]\label{p:511-1-1}
证明第一Brillouin区体积与实空间Bravais原胞体积的乘积为$(2\pi)^3$。
\end{exercise}
\begin{solution}
令$A=(a_1,a_2,a_3)$，倒格基矩阵$B=2\pi A^{-T}$，则$|\det A||\det B|=(2\pi)^3$。任意倒格原胞与第一Brillouin区各铺砌倒空间，且每个胞恰含一个格点，故体积相同。原胞的任意整数幺模换基不改行列式绝对值，结论与胞的具体形状无关。
\end{solution}
\begin{exercise}[原题2]
证明使Bravais晶格保持不变的绕任意轴旋转只能是一、二、三、四或六重。
\end{exercise}
\begin{solution}
旋转在晶格基中的矩阵为整数矩阵，在三维的本征值是$1,\e^{\ii\theta},\e^{-\ii\theta}$，所以整数迹$1+2\cos\theta$。$2\cos\theta\in\{-2,-1,0,1,2\}$，得到$\theta=\pi,2\pi/3,\pi/2,\pi/3,0$及相应方向的等价旋转。有限阶基本角给所列阶数；五重要求$2\cos(2\pi/5)$非整数，不能保持平移晶格。
\end{solution}
\begin{exercise}[原题3]
CuO$_2$层中Cu形成间距$a$的正方格，O位于每条Cu--Cu最近邻键中点；各层沿$z$简单堆叠，层距$c$，忽略其他原子。(a)给出Bravais晶格、原矢、原胞与基底。(b)原题第二图中O交替向上、向下移位。将该图的条件写清：在Cu坐标$(na,ma)$的右方O，$u_z=\delta(-1)^{n+m}$；在其上方O，$u_z=-\delta(-1)^{n+m}$。给出畸变后原胞、晶格间距和倒晶格；讨论$\delta\to0$时X射线峰如何变化。上下图样翻转不改变平移晶格。
\end{exercise}
\begin{solution}
(a)取$a_1=a\hat x,a_2=a\hat y,a_3=c\hat z$，基底为Cu$(0,0,0)$、O$(a/2,0,0)$、O$(0,a/2,0)$，体积$a^2c$。平面是方格，三维为简单四方Bravais晶格。

(b)平移$(a,0)$改变所有O的移位符号，不能保持畸变；$(a,a)$与$(-a,a)$保持$n+m$偶性。因此新原矢$A_1=a(\hat x+\hat y),A_2=a(-\hat x+\hat y),A_3=c\hat z$，平面原胞边长$\sqrt2a$，体积$2a^2c$，基底含2Cu、4O。倒矢$B_1=(\pi/a)(\hat x+\hat y),B_2=(\pi/a)(-\hat x+\hat y),B_3=(2\pi/c)\hat z$。新增倒格点对应原来禁戒的超结构峰；位移引起振幅一阶$\propto Q_z\delta$，强度$\propto\delta^2$，在畸变消失时逐渐消失。原有主峰位置不需随$\delta$连续漂移，区别在结构因子。
\end{solution}
\originalsection{作业2（2004年9月28日截止）}
原题1开头只有“Marder第38页2.5(b)”指针，该外书题面未取得，不能从后续小问倒造原题。以下保留公开附加(a)、(b)和原题2的完整数学条件。
\begin{exercise}[原题1(a)]
HCP具有平行于包含$a,c$轴的平面的滑移面。说明可把此平面放在哪里，使沿$c/2$平移再反射保持HCP结构。
\end{exercise}
\begin{solution}
用三角原矢$a_1=a(1,0,0),a_2=a(1/2,\sqrt3/2,0),a_3=c\hat z$，两层基底可取$0$和$r_B=(a/2,\sqrt3a/6,c/2)$。滑移面可取竖直平面经过两层点连线的中点且平面法向沿$r_{B\perp}$，即$\vct r_\perp\cdot r_{B\perp}=|r_{B\perp}|^2/2$。其线性反射$R$把三角晶格映回自身：此方向的法向相对$a_1$为$30^\circ$，属于三角格反射对称。仿射反射把原点映到$r_{B\perp}$，再加$c/2$恰为B层；B层映到下一A层。这样明确给出纯$c/2$滑移所需的平面，不能任取竖直晶面。
\end{solution}
\begin{exercise}[原题1(b)(i)--(iii)，公开文件标注Ashcroft--Mermin第109页题5]
NaCl采用边长$a$的FCC晶格，基底$+$离子在原点、$-$离子在$(a/2)\hat x$。倒格矢$K=(4\pi/a)(\nu_1,\nu_2,\nu_3)$，$\nu_i$全为整数或全为半整数，形状因子为$f_\pm$。(i)求结构因子。(ii)闪锌矿的第二基底在$(a/4)(1,1,1)$，求结构因子。(iii)若闭壳层离子的$f_\pm$只依赖$|K|$，如何根据Bragg峰强度区分两种结构？
\end{exercise}
\begin{solution}
(i)$S_K=f_++f_-\e^{-2\pi\ii\nu_1}$，全整数时$f_++f_-$，全半整数时$f_+-f_-$。若$f_+=f_-$，后者消光，是两类原子变得等价导致额外平移对称。

(ii)$S_K=f_++f_-\e^{-\ii\pi(\nu_1+\nu_2+\nu_3)}$。全整数时按和的偶奇为$f_+\pm f_-$；全半整数时相位为$\pm\ii$，得$f_+\pm\ii f_-$。符号依指数约定，强度不变。(iii)在相同$|K|$或已校正原子形状因子的反射族间比较：NaCl的全半整数类按$|f_+-f_-|^2$，闪锌矿对应$|f_+\pm\ii f_-|^2$，实因子时为$f_+^2+f_-^2$；全整数族还按和偶奇分组。须校正几何、多重性、吸收和热振动，不能凭一个峰强度唯一判定。
\end{solution}
\begin{exercise}[原题2(a)--(b)，公开文件标注Ashcroft--Mermin第171页题3]
FCC的$W=(2\pi/a)(1,1/2,0)$处，四个自由波动量为$k,k-(2\pi/a)(1,1,1),k-(2\pi/a)(1,1,-1),k-(2\pi/a)(2,0,0)$。势反演对称，使$U_G$实，令$U_1=U_{111}=U_{11\bar1},U_2=U_{200}$。(a)写四态能量矩阵并求$W$劈裂。(b)在$U=(2\pi/a)(1,1/4,1/4)$用同法求简并能级劈裂。
\end{exercise}
\begin{solution}
(a)四态对角为$\epsilon_i^0$，非对角依次为
$\left(\begin{smallmatrix}0&U_1&U_1&U_2\\U_1&0&U_2&U_1\\U_1&U_2&0&U_1\\U_2&U_1&U_1&0\end{smallmatrix}\right)$。
在$W$对角相同$\epsilon_W$。反对称组$(1,0,0,-1)$与$(0,1,-1,0)$能量$\epsilon_W-U_2$；剩余组$(1,1,1,1)$与$(1,-1,-1,1)$能量$\epsilon_W+U_2\pm2U_1$。

(b)在$U$只有三波简并：第一、第二和第四，矩阵非对角为$(U_1,U_2,U_1)$。第一与第三分量反对称给$\epsilon_U-U_2$；对称组与中间态的$2\times2$块为$\left(\begin{smallmatrix}U_2&\sqrt2U_1\\\sqrt2U_1&0\end{smallmatrix}\right)$，另两根$\epsilon_U+[U_2\pm\sqrt{U_2^2+8U_1^2}]/2$。其余非简并波的一阶作用不包含在该三态公式中。
\end{solution}
\originalsection{作业3（2004年10月5日截止）}
\begin{exercise}[原题1(a)--(e)]
Al为FCC，立方边长$a=4.05\,\text{\AA}$。(a)构造自由电子最低两带，画$\Gamma\to X\to W\to L\to\Gamma$并标简并，$X=(2\pi/a)(1,0,0),L=(2\pi/a)(1/2,1/2,1/2)$。(b)空芯赝势为$r>R_c$时$-Ze_*^2/r$、$r<R_c$时零，$Z=3,R_c=0.6\,\text{\AA}$；$V(q)=-4\pi e_*^2Z\cos(qR_c)/(q^2\Omega)$，$\Omega=a^3/4$，求$X,W,L$劈裂。(c)$W$四态中$S_1$同时交换1与4、2与3，$S_2$同时交换1与2、3与4，证明两者与$H$对易且态可按$\pm1$分类。(d)讨论沿$XW$、$WL$的这些对称与带连接，允许用计算核对。(e)估计费米能并与题面所指商业教材KKR图比较。
\end{exercise}
\begin{solution}
(a)对每个路径点将$\hbar^2|k+G|^2/(2m)$排序，$G=(2\pi/a)(h,k,l)$且$h,k,l$全偶或全奇，取最低两个（遇简并保留重数）。$\Gamma$最低为非简并0；$X$最低2个波简并；$W$最低4个波简并；$L$最低2个波简并。源码中的数值脚本给这一路径的采样图与能量表，$k$点处的简并先排序再合组。

(b)$X$、$L$分别由$G_{200}$与$G_{111}$散射，劈裂$2|V_{200}|$、$2|V_{111}|$；$W$采用上一题四态根，劈裂相对$\epsilon_W$为$-V_{200}$（二重）及$V_{200}\pm2V_{111}$。使用$\hbar^2/(2m)=3.809982\,\mathrm{eV}\,\text{\AA}^{2}$、$e_*^2=14.39965\,\mathrm{eV}\,\text{\AA}$，数值为$V_{111}=0.18768$ eV、$V_{200}=0.97376$ eV，$X$隙1.94752 eV、$L$隙0.37536 eV；$W$偏移$-0.97376$ eV（二重）、0.59840 eV、1.34912 eV。随卷计算证据也保留这些数，核对单位$q^{-2}/\Omega$为长度倒数。

(c)对称矩阵是置换矩阵，$S_i^2=I$且$S_1S_2=S_2S_1$。$W$处相同对角与非对角$V_1,V_2$均不因交换改变，故对易。四个共同本征矢可取$(1,s_2,s_1s_2,s_1)/2$，$s_i=\pm1$，相应偏移$V_2s_1+V_1s_2(1+s_1)$，重现上一题的二重根。

(d)$XW$上动量$(1,y,0)$，自由对角满足$\epsilon_1=\epsilon_4$、$\epsilon_2=\epsilon_3$，所以$S_1$保留，$S_2$除$W$外不保留。$WL$上$k=(1-s/2,1/2,s/2)$，$\epsilon_1=\epsilon_2$、$\epsilon_3=\epsilon_4$，所以$S_2$保留，$S_1$除$W$外不保留。保留不同$S$本征值的支可交叉，同一对称支一般避免交叉。用四态矩阵沿两段对角化即可把端点根按各自本征值连续连接，不能只按排序序号追踪交叉。

(e)每原胞3电子，$n=3/\Omega$给$k_F=(3\pi^2n)^{1/3}$和$E_F=\hbar^2k_F^2/(2m)$，自由模型估计约$11.65\,\mathrm{eV}$。加入弱赝势后需按整个区的占据数积分重新定$\mu$，单靠高对称线不足以精确求填充。所引用Ashcroft--Mermin第205页及Marder第250页KKR图未作为公开OCW文件取得，比较项保留缺口，不伪造书中曲线。
\end{solution}
\begin{center}\begin{tikzpicture}\begin{axis}[width=11cm,height=5.7cm,xlabel={路径},ylabel={能量/eV},title={Al自由电子折叠带},xtick={0,1,2,3,4},xticklabels={$\Gamma$,$X$,$W$,$L$,$\Gamma$},legend pos=north east]
\addplot[thick,color=structurecolor] coordinates {(0,0) (0.016393443,0.00246441) (0.032786885,0.0098576399) (0.049180328,0.02217969) (0.06557377,0.03943056) (0.081967213,0.061610249) (0.098360656,0.088718759) (0.1147541,0.12075609) (0.13114754,0.15772224) (0.14754098,0.19961721) (0.16393443,0.246441) (0.18032787,0.29819361) (0.19672131,0.35487504) (0.21311475,0.41648529) (0.2295082,0.48302435) (0.24590164,0.55449224) (0.26229508,0.63088895) (0.27868852,0.71221448) (0.29508197,0.79846883) (0.31147541,0.889652) (0.32786885,0.98576399) (0.3442623,1.0868048) (0.36065574,1.1927744) (0.37704918,1.3036729) (0.39344262,1.4195001) (0.40983607,1.5402562) (0.42622951,1.6659411) (0.44262295,1.7965549) (0.45901639,1.9320974) (0.47540984,2.0725688) (0.49180328,2.217969) (0.50819672,2.368298) (0.52459016,2.5235558) (0.54098361,2.6837425) (0.55737705,2.8488579) (0.57377049,3.0189022) (0.59016393,3.1938753) (0.60655738,3.3737773) (0.62295082,3.558608) (0.63934426,3.7483676) (0.6557377,3.943056) (0.67213115,4.1426732) (0.68852459,4.3472192) (0.70491803,4.556694) (0.72131148,4.7710977) (0.73770492,4.9904302) (0.75409836,5.2146915) (0.7704918,5.4438816) (0.78688525,5.6780006) (0.80327869,5.9170483) (0.81967213,6.1610249) (0.83606557,6.4099303) (0.85245902,6.6637646) (0.86885246,6.9225276) (0.8852459,7.1862195) (0.90163934,7.4548402) (0.91803279,7.7283897) (0.93442623,8.006868) (0.95081967,8.2902752) (0.96721311,8.5786111) (0.98360656,8.8718759) (1,9.1700695) (1.0163934,9.1706856) (1.0327869,9.1725339) (1.0491803,9.1756144) (1.0655738,9.1799272) (1.0819672,9.1854721) (1.0983607,9.1922492) (1.1147541,9.2002585) (1.1311475,9.2095001) (1.147541,9.2199738) (1.1639344,9.2316798) (1.1803279,9.2446179) (1.1967213,9.2587883) (1.2131148,9.2741908) (1.2295082,9.2908256) (1.2459016,9.3086926) (1.2622951,9.3277918) (1.2786885,9.3481231) (1.295082,9.3696867) (1.3114754,9.3924825) (1.3278689,9.4165105) (1.3442623,9.4417707) (1.3606557,9.4682631) (1.3770492,9.4959877) (1.3934426,9.5249445) (1.4098361,9.5551336) (1.4262295,9.5865548) (1.442623,9.6192082) (1.4590164,9.6530939) (1.4754098,9.6882117) (1.4918033,9.7245618) (1.5081967,9.762144) (1.5245902,9.8009585) (1.5409836,9.8410051) (1.557377,9.882284) (1.5737705,9.9247951) (1.5901639,9.9685383) (1.6065574,10.013514) (1.6229508,10.059722) (1.6393443,10.107161) (1.6557377,10.155834) (1.6721311,10.205738) (1.6885246,10.256874) (1.704918,10.309243) (1.7213115,10.362844) (1.7377049,10.417677) (1.7540984,10.473742) (1.7704918,10.53104) (1.7868852,10.58957) (1.8032787,10.649332) (1.8196721,10.710326) (1.8360656,10.772552) (1.852459,10.836011) (1.8688525,10.900701) (1.8852459,10.966624) (1.9016393,11.03378) (1.9180328,11.102167) (1.9344262,11.171787) (1.9508197,11.242638) (1.9672131,11.314722) (1.9836066,11.388038) (2,11.462587) (2.0163934,11.31349) (2.0327869,11.166858) (2.0491803,11.02269) (2.0655738,10.880986) (2.0819672,10.741747) (2.0983607,10.604972) (2.1147541,10.470662) (2.1311475,10.338816) (2.147541,10.209434) (2.1639344,10.082517) (2.1803279,9.9580646) (2.1967213,9.8360763) (2.2131148,9.7165524) (2.2295082,9.599493) (2.2459016,9.4848979) (2.2622951,9.3727672) (2.2786885,9.263101) (2.295082,9.1558992) (2.3114754,9.0511617) (2.3278689,8.9488887) (2.3442623,8.8490801) (2.3606557,8.7517359) (2.3770492,8.6568561) (2.3934426,8.5644408) (2.4098361,8.4744898) (2.4262295,8.3870032) (2.442623,8.3019811) (2.4590164,8.2194234) (2.4754098,8.13933) (2.4918033,8.0617011) (2.5081967,7.9865366) (2.5245902,7.9138365) (2.5409836,7.8436008) (2.557377,7.7758296) (2.5737705,7.7105227) (2.5901639,7.6476803) (2.6065574,7.5873022) (2.6229508,7.5293886) (2.6393443,7.4739393) (2.6557377,7.4209545) (2.6721311,7.3704341) (2.6885246,7.3223781) (2.704918,7.2767866) (2.7213115,7.2336594) (2.7377049,7.1929966) (2.7540984,7.1547983) (2.7704918,7.1190643) (2.7868852,7.0857948) (2.8032787,7.0549897) (2.8196721,7.0266489) (2.8360656,7.0007726) (2.852459,6.9773607) (2.8688525,6.9564133) (2.8852459,6.9379302) (2.9016393,6.9219115) (2.9180328,6.9083573) (2.9344262,6.8972674) (2.9508197,6.888642) (2.9672131,6.882481) (2.9836066,6.8787843) (3,6.8775521) (3.0166667,6.6502108) (3.0333333,6.4266904) (3.05,6.2069908) (3.0666667,5.9911121) (3.0833333,5.7790542) (3.1,5.5708172) (3.1166667,5.3664011) (3.1333333,5.1658058) (3.15,4.9690314) (3.1666667,4.7760779) (3.1833333,4.5869452) (3.2,4.4016334) (3.2166667,4.2201424) (3.2333333,4.0424723) (3.25,3.8686231) (3.2666667,3.6985947) (3.2833333,3.5323872) (3.3,3.3700005) (3.3166667,3.2114348) (3.3333333,3.0566898) (3.35,2.9057658) (3.3666667,2.7586626) (3.3833333,2.6153802) (3.4,2.4759188) (3.4166667,2.3402782) (3.4333333,2.2084584) (3.45,2.0804595) (3.4666667,1.9562815) (3.4833333,1.8359243) (3.5,1.719388) (3.5166667,1.6066726) (3.5333333,1.497778) (3.55,1.3927043) (3.5666667,1.2914515) (3.5833333,1.1940195) (3.6,1.1004083) (3.6166667,1.0106181) (3.6333333,0.92464868) (3.65,0.84250014) (3.6666667,0.76417246) (3.6833333,0.68966564) (3.7,0.61897969) (3.7166667,0.5521146) (3.7333333,0.48907037) (3.75,0.42984701) (3.7666667,0.37444451) (3.7833333,0.32286286) (3.8,0.27510209) (3.8166667,0.23116217) (3.8333333,0.19104311) (3.85,0.15474492) (3.8666667,0.12226759) (3.8833333,0.093611126) (3.9,0.068775521) (3.9166667,0.047760779) (3.9333333,0.030566898) (3.95,0.01719388) (3.9666667,0.0076417246) (3.9833333,0.0019104311) (4,0)};
\addlegendentry{第一带}
\addplot[thick,color=main] coordinates {(0,27.510209) (0.016393443,27.212015) (0.032786885,26.91875) (0.049180328,26.630414) (0.06557377,26.347007) (0.081967213,26.068529) (0.098360656,25.794979) (0.1147541,25.526359) (0.13114754,25.262667) (0.14754098,25.003904) (0.16393443,24.750069) (0.18032787,24.501164) (0.19672131,24.257187) (0.21311475,24.01814) (0.2295082,23.784021) (0.24590164,23.554831) (0.26229508,23.330569) (0.27868852,23.111237) (0.29508197,22.896833) (0.31147541,22.687358) (0.32786885,22.482812) (0.3442623,22.283195) (0.36065574,22.088507) (0.37704918,21.898747) (0.39344262,21.713916) (0.40983607,21.534014) (0.42622951,21.359041) (0.44262295,21.188997) (0.45901639,21.023881) (0.47540984,20.863695) (0.49180328,20.708437) (0.50819672,20.407779) (0.52459016,19.961721) (0.54098361,19.520591) (0.55737705,19.084391) (0.57377049,18.653119) (0.59016393,18.226776) (0.60655738,17.805362) (0.62295082,17.388877) (0.63934426,16.97732) (0.6557377,16.570693) (0.67213115,16.168994) (0.68852459,15.772224) (0.70491803,15.380383) (0.72131148,14.99347) (0.73770492,14.611487) (0.75409836,14.234432) (0.7704918,13.862306) (0.78688525,13.495109) (0.80327869,13.132841) (0.81967213,12.775501) (0.83606557,12.423091) (0.85245902,12.075609) (0.86885246,11.733056) (0.8852459,11.395432) (0.90163934,11.062736) (0.91803279,10.73497) (0.93442623,10.412132) (0.95081967,10.094223) (0.96721311,9.7812432) (0.98360656,9.4731919) (1,9.1700695) (1.0163934,9.1706856) (1.0327869,9.1725339) (1.0491803,9.1756144) (1.0655738,9.1799272) (1.0819672,9.1854721) (1.0983607,9.1922492) (1.1147541,9.2002585) (1.1311475,9.2095001) (1.147541,9.2199738) (1.1639344,9.2316798) (1.1803279,9.2446179) (1.1967213,9.2587883) (1.2131148,9.2741908) (1.2295082,9.2908256) (1.2459016,9.3086926) (1.2622951,9.3277918) (1.2786885,9.3481231) (1.295082,9.3696867) (1.3114754,9.3924825) (1.3278689,9.4165105) (1.3442623,9.4417707) (1.3606557,9.4682631) (1.3770492,9.4959877) (1.3934426,9.5249445) (1.4098361,9.5551336) (1.4262295,9.5865548) (1.442623,9.6192082) (1.4590164,9.6530939) (1.4754098,9.6882117) (1.4918033,9.7245618) (1.5081967,9.762144) (1.5245902,9.8009585) (1.5409836,9.8410051) (1.557377,9.882284) (1.5737705,9.9247951) (1.5901639,9.9685383) (1.6065574,10.013514) (1.6229508,10.059722) (1.6393443,10.107161) (1.6557377,10.155834) (1.6721311,10.205738) (1.6885246,10.256874) (1.704918,10.309243) (1.7213115,10.362844) (1.7377049,10.417677) (1.7540984,10.473742) (1.7704918,10.53104) (1.7868852,10.58957) (1.8032787,10.649332) (1.8196721,10.710326) (1.8360656,10.772552) (1.852459,10.836011) (1.8688525,10.900701) (1.8852459,10.966624) (1.9016393,11.03378) (1.9180328,11.102167) (1.9344262,11.171787) (1.9508197,11.242638) (1.9672131,11.314722) (1.9836066,11.388038) (2,11.462587) (2.0163934,11.31349) (2.0327869,11.166858) (2.0491803,11.02269) (2.0655738,10.880986) (2.0819672,10.741747) (2.0983607,10.604972) (2.1147541,10.470662) (2.1311475,10.338816) (2.147541,10.209434) (2.1639344,10.082517) (2.1803279,9.9580646) (2.1967213,9.8360763) (2.2131148,9.7165524) (2.2295082,9.599493) (2.2459016,9.4848979) (2.2622951,9.3727672) (2.2786885,9.263101) (2.295082,9.1558992) (2.3114754,9.0511617) (2.3278689,8.9488887) (2.3442623,8.8490801) (2.3606557,8.7517359) (2.3770492,8.6568561) (2.3934426,8.5644408) (2.4098361,8.4744898) (2.4262295,8.3870032) (2.442623,8.3019811) (2.4590164,8.2194234) (2.4754098,8.13933) (2.4918033,8.0617011) (2.5081967,7.9865366) (2.5245902,7.9138365) (2.5409836,7.8436008) (2.557377,7.7758296) (2.5737705,7.7105227) (2.5901639,7.6476803) (2.6065574,7.5873022) (2.6229508,7.5293886) (2.6393443,7.4739393) (2.6557377,7.4209545) (2.6721311,7.3704341) (2.6885246,7.3223781) (2.704918,7.2767866) (2.7213115,7.2336594) (2.7377049,7.1929966) (2.7540984,7.1547983) (2.7704918,7.1190643) (2.7868852,7.0857948) (2.8032787,7.0549897) (2.8196721,7.0266489) (2.8360656,7.0007726) (2.852459,6.9773607) (2.8688525,6.9564133) (2.8852459,6.9379302) (2.9016393,6.9219115) (2.9180328,6.9083573) (2.9344262,6.8972674) (2.9508197,6.888642) (2.9672131,6.882481) (2.9836066,6.8787843) (3,6.8775521) (3.0166667,7.1087143) (3.0333333,7.3436973) (3.05,7.5825012) (3.0666667,7.825126) (3.0833333,8.0715716) (3.1,8.3218381) (3.1166667,8.5759254) (3.1333333,8.8338336) (3.15,9.0955627) (3.1666667,9.3611126) (3.1833333,9.6304834) (3.2,9.9036751) (3.2166667,10.180688) (3.2333333,10.461521) (3.25,10.746175) (3.2666667,11.03465) (3.2833333,11.326946) (3.3,11.623063) (3.3166667,11.923001) (3.3333333,12.226759) (3.35,12.534339) (3.3666667,12.845739) (3.3833333,13.16096) (3.4,13.480002) (3.4166667,13.802865) (3.4333333,14.129549) (3.45,14.460053) (3.4666667,14.794379) (3.4833333,15.132525) (3.5,15.474492) (3.5166667,15.82028) (3.5333333,16.169889) (3.55,16.523319) (3.5666667,16.88057) (3.5833333,17.241641) (3.6,17.606533) (3.6166667,17.975247) (3.6333333,18.347781) (3.65,18.724136) (3.6666667,19.104311) (3.6833333,19.488308) (3.7,19.876126) (3.7166667,20.267764) (3.7333333,20.663223) (3.75,21.062503) (3.7666667,21.465604) (3.7833333,21.872526) (3.8,22.283269) (3.8166667,22.697832) (3.8333333,23.116217) (3.85,23.538422) (3.8666667,23.964448) (3.8833333,24.394295) (3.9,24.827963) (3.9166667,25.265452) (3.9333333,25.706762) (3.95,26.151892) (3.9666667,26.600843) (3.9833333,27.053615) (4,27.510209)};
\addlegendentry{第二带}
\end{axis}\end{tikzpicture}\end{center}

\originalsection{作业4（2004年10月12日截止）}
\begin{exercise}[原题1(a)--(e)]
La$_2$CuO$_4$的平面模型中Cu在方格点、O在键中点，只保留Cu的$d_{x^2-y^2}$轨道（能量$E_d$）与每个O的沿键$p$轨道（能量$E_p$）。忽略轨道重叠及O--O跃迁，只保留Cu--最近O的幅度$t_{pd}$。(a)求三带结构。(b)取La为$3+$、非平面O为$2-$、Cu的$4s$空，证明保留轨道中每胞一空穴并找费米面。(c)找鞍点与态密度对数奇异。(d)若15\%的La位由Sr$^{2+}$替换，画空穴掺杂后的费米面；原题给Yoshida等PRB 63,220501(2001)为实验参考。(e)Mattheiss PRL58,1028(1987)称$E_p=E_d,t_{pd}=-1.85\,\mathrm{eV}$可拟合其APW；检验模型预言，原文三维$X,S$对应二维$(\pi/a,\pi/a)$，$\Gamma,Z$对应$(0,0)$。
\end{exercise}
\begin{solution}
(a)按轨道奇偶性选规范，$d,p_x,p_y$矩阵为
\begin{equation}
H(k)=\begin{pmatrix}E_d&2\ii t\sin(k_xa/2)&-2\ii t\sin(k_ya/2)\\-2\ii t\sin(k_xa/2)&E_p&0\\2\ii t\sin(k_ya/2)&0&E_p\end{pmatrix}.
\end{equation}
一条非键带$E_p$，另外$E_\pm=(E_d+E_p)/2\pm\sqrt{(E_d-E_p)^2/4+4t^2[s_x^2+s_y^2]}$。保留矩阵元相位保证Hermitian，改轨道规范不改能量。

(b)离子平衡给Cu$2+$即$d^9$，相对于满$d^{10}$一个空穴；两平面O的相关$p$轨道近满。因此三轨道6个自旋态有5电子，下两带满，上带半满。对仅最近Cu--O的模型，$s_x^2+s_y^2=1$是半填充等能线，即$\cos k_xa+\cos k_ya=0$，费米面为菱形。真实材料未掺杂时是Mott绝缘体，题目中的费米面是该单粒子模型的预言。

(c)$(\pi/a,0)$与$(0,\pi/a)$处，$s_x^2+s_y^2=1+[(\delta k_y)^2-(\delta k_x)^2]a^2/4+\cdots$，故$E_+-E_F=\alpha[(\delta k_y)^2-(\delta k_x)^2]$，$\alpha>0$。第\ref{ch:bands}章的$uv$积分给对数态密度。

(d)“15\%的La位”按字面每化学式有2个La，新增$0.30$空穴；若实验用La$_{2-x}$Sr$_x$且$x=0.15$，则是$0.15$空穴，两种口径必须分开。电子填充下降，$E_F$向上带较低能量移动，等能条件$s_x^2+s_y^2<1$，电子占据区环绕$\Gamma$收缩，空穴区扩张。精确等能值由面积满足上带电子数$1-p$确定，不能任意挪一条线冒称15\%面积。

(e)$E_p=E_d=E_0$时$E_\pm=E_0\pm2|t|\sqrt{s_x^2+s_y^2}$，$\Gamma$三带同能，$(\pi,0)$分裂$\pm3.70$ eV，$(\pi,\pi)$分裂$\pm5.233$ eV。这是可检验的二维拟合预言；对题面引用论文的完整三维APW曲线未取得可再分发原件，不能据这几个点声称已核验论文拟合。原题第三方晶体结构照片未复用，以上条件已足够重建本模型。
\end{solution}

\begin{center}\begin{tikzpicture}\begin{axis}[width=11cm,height=5.7cm,xlabel={$k_xa$},ylabel={$k_ya$},title={CuO$_2$模型的费米面},axis equal image,xmin=-3.2,xmax=3.2,ymin=-3.2,ymax=3.2,legend pos=outer north east]
\addplot[thick,color=structurecolor] coordinates {(-3.1415927,0.0039220882) (-3.1363567,0.0065420497) (-3.1311207,0.011182367) (-3.1258847,0.016190229) (-3.1206487,0.02130805) (-3.1154127,0.026472131) (-3.1101767,0.031659844) (-3.1049407,0.036861213) (-3.0997048,0.042071173) (-3.0944688,0.047286885) (-3.0892328,0.052506634) (-3.0839968,0.057729325) (-3.0787608,0.062954227) (-3.0735248,0.06818083) (-3.0682888,0.073408772) (-3.0630528,0.078637786) (-3.0578168,0.08386767) (-3.0525809,0.089098273) (-3.0473449,0.094329473) (-3.0421089,0.099561178) (-3.0368729,0.10479331) (-3.0316369,0.11002581) (-3.0264009,0.11525863) (-3.0211649,0.12049172) (-3.0159289,0.12572506) (-3.010693,0.13095861) (-3.005457,0.13619234) (-3.000221,0.14142625) (-2.994985,0.1466603) (-2.989749,0.15189448) (-2.984513,0.15712879) (-2.979277,0.16236321) (-2.974041,0.16759772) (-2.9688051,0.17283233) (-2.9635691,0.17806701) (-2.9583331,0.18330177) (-2.9530971,0.1885366) (-2.9478611,0.19377149) (-2.9426251,0.19900644) (-2.9373891,0.20424145) (-2.9321531,0.2094765) (-2.9269172,0.2147116) (-2.9216812,0.21994674) (-2.9164452,0.22518192) (-2.9112092,0.23041714) (-2.9059732,0.23565239) (-2.9007372,0.24088768) (-2.8955012,0.24612299) (-2.8902652,0.25135834) (-2.8850293,0.25659371) (-2.8797933,0.2618291) (-2.8745573,0.26706452) (-2.8693213,0.27229996) (-2.8640853,0.27753542) (-2.8588493,0.28277091) (-2.8536133,0.28800641) (-2.8483773,0.29324192) (-2.8431414,0.29847746) (-2.8379054,0.30371301) (-2.8326694,0.30894857) (-2.8274334,0.31418415) (-2.8221974,0.31941975) (-2.8169614,0.32465535) (-2.8117254,0.32989097) (-2.8064894,0.3351266) (-2.8012534,0.34036224) (-2.7960175,0.3455979) (-2.7907815,0.35083356) (-2.7855455,0.35606923) (-2.7803095,0.36130491) (-2.7750735,0.3665406) (-2.7698375,0.3717763) (-2.7646015,0.37701201) (-2.7593655,0.38224773) (-2.7541296,0.38748345) (-2.7488936,0.39271918) (-2.7436576,0.39795492) (-2.7384216,0.40319066) (-2.7331856,0.40842641) (-2.7279496,0.41366217) (-2.7227136,0.41889793) (-2.7174776,0.4241337) (-2.7122417,0.42936947) (-2.7070057,0.43460525) (-2.7017697,0.43984104) (-2.6965337,0.44507682) (-2.6912977,0.45031262) (-2.6860617,0.45554842) (-2.6808257,0.46078422) (-2.6755897,0.46602003) (-2.6703538,0.47125584) (-2.6651178,0.47649166) (-2.6598818,0.48172747) (-2.6546458,0.4869633) (-2.6494098,0.49219913) (-2.6441738,0.49743496) (-2.6389378,0.50267079) (-2.6337018,0.50790663) (-2.6284659,0.51314247) (-2.6232299,0.51837831) (-2.6179939,0.52361416) (-2.6127579,0.52885001) (-2.6075219,0.53408586) (-2.6022859,0.53932172) (-2.5970499,0.54455757) (-2.5918139,0.54979343) (-2.586578,0.5550293) (-2.581342,0.56026516) (-2.576106,0.56550103) (-2.57087,0.5707369) (-2.565634,0.57597277) (-2.560398,0.58120865) (-2.555162,0.58644453) (-2.549926,0.59168041) (-2.54469,0.59691629) (-2.5394541,0.60215217) (-2.5342181,0.60738806) (-2.5289821,0.61262394) (-2.5237461,0.61785983) (-2.5185101,0.62309572) (-2.5132741,0.62833162) (-2.5080381,0.63356751) (-2.5028021,0.63880341) (-2.4975662,0.6440393) (-2.4923302,0.6492752) (-2.4870942,0.6545111) (-2.4818582,0.65974701) (-2.4766222,0.66498291) (-2.4713862,0.67021882) (-2.4661502,0.67545472) (-2.4609142,0.68069063) (-2.4556783,0.68592654) (-2.4504423,0.69116245) (-2.4452063,0.69639836) (-2.4399703,0.70163428) (-2.4347343,0.70687019) (-2.4294983,0.71210611) (-2.4242623,0.71734202) (-2.4190263,0.72257794) (-2.4137904,0.72781386) (-2.4085544,0.73304978) (-2.4033184,0.7382857) (-2.3980824,0.74352162) (-2.3928464,0.74875755) (-2.3876104,0.75399347) (-2.3823744,0.7592294) (-2.3771384,0.76446532) (-2.3719025,0.76970125) (-2.3666665,0.77493718) (-2.3614305,0.78017311) (-2.3561945,0.78540904) (-2.3509585,0.79064497) (-2.3457225,0.7958809) (-2.3404865,0.80111684) (-2.3352505,0.80635277) (-2.3300146,0.81158871) (-2.3247786,0.81682464) (-2.3195426,0.82206058) (-2.3143066,0.82729651) (-2.3090706,0.83253245) (-2.3038346,0.83776839) (-2.2985986,0.84300433) (-2.2933626,0.84824027) (-2.2881266,0.85347621) (-2.2828907,0.85871215) (-2.2776547,0.86394809) (-2.2724187,0.86918404) (-2.2671827,0.87441998) (-2.2619467,0.87965593) (-2.2567107,0.88489187) (-2.2514747,0.89012782) (-2.2462387,0.89536376) (-2.2410028,0.90059971) (-2.2357668,0.90583566) (-2.2305308,0.9110716) (-2.2252948,0.91630755) (-2.2200588,0.9215435) (-2.2148228,0.92677945) (-2.2095868,0.9320154) (-2.2043508,0.93725135) (-2.1991149,0.9424873) (-2.1938789,0.94772325) (-2.1886429,0.95295921) (-2.1834069,0.95819516) (-2.1781709,0.96343111) (-2.1729349,0.96866707) (-2.1676989,0.97390302) (-2.1624629,0.97913898) (-2.157227,0.98437493) (-2.151991,0.98961089) (-2.146755,0.99484684) (-2.141519,1.0000828) (-2.136283,1.0053188) (-2.131047,1.0105547) (-2.125811,1.0157907) (-2.120575,1.0210266) (-2.1153391,1.0262626) (-2.1101031,1.0314986) (-2.1048671,1.0367345) (-2.0996311,1.0419705) (-2.0943951,1.0472064) (-2.0891591,1.0524424) (-2.0839231,1.0576784) (-2.0786871,1.0629143) (-2.0734512,1.0681503) (-2.0682152,1.0733862) (-2.0629792,1.0786222) (-2.0577432,1.0838582) (-2.0525072,1.0890941) (-2.0472712,1.0943301) (-2.0420352,1.0995661) (-2.0367992,1.104802) (-2.0315632,1.110038) (-2.0263273,1.115274) (-2.0210913,1.1205099) (-2.0158553,1.1257459) (-2.0106193,1.1309819) (-2.0053833,1.1362178) (-2.0001473,1.1414538) (-1.9949113,1.1466898) (-1.9896753,1.1519257) (-1.9844394,1.1571617) (-1.9792034,1.1623977) (-1.9739674,1.1676336) (-1.9687314,1.1728696) (-1.9634954,1.1781056) (-1.9582594,1.1833415) (-1.9530234,1.1885775) (-1.9477874,1.1938135) (-1.9425515,1.1990495) (-1.9373155,1.2042854) (-1.9320795,1.2095214) (-1.9268435,1.2147574) (-1.9216075,1.2199933) (-1.9163715,1.2252293) (-1.9111355,1.2304653) (-1.9058995,1.2357013) (-1.9006636,1.2409372) (-1.8954276,1.2461732) (-1.8901916,1.2514092) (-1.8849556,1.2566451) (-1.8797196,1.2618811) (-1.8744836,1.2671171) (-1.8692476,1.2723531) (-1.8640116,1.277589) (-1.8587757,1.282825) (-1.8535397,1.288061) (-1.8483037,1.293297) (-1.8430677,1.2985329) (-1.8378317,1.3037689) (-1.8325957,1.3090049) (-1.8273597,1.3142409) (-1.8221237,1.3194769) (-1.8168878,1.3247128) (-1.8116518,1.3299488) (-1.8064158,1.3351848) (-1.8011798,1.3404208) (-1.7959438,1.3456567) (-1.7907078,1.3508927) (-1.7854718,1.3561287) (-1.7802358,1.3613647) (-1.7749998,1.3666007) (-1.7697639,1.3718366) (-1.7645279,1.3770726) (-1.7592919,1.3823086) (-1.7540559,1.3875446) (-1.7488199,1.3927806) (-1.7435839,1.3980165) (-1.7383479,1.4032525) (-1.7331119,1.4084885) (-1.727876,1.4137245) (-1.72264,1.4189605) (-1.717404,1.4241964) (-1.712168,1.4294324) (-1.706932,1.4346684) (-1.701696,1.4399044) (-1.69646,1.4451404) (-1.691224,1.4503764) (-1.6859881,1.4556123) (-1.6807521,1.4608483) (-1.6755161,1.4660843) (-1.6702801,1.4713203) (-1.6650441,1.4765563) (-1.6598081,1.4817923) (-1.6545721,1.4870282) (-1.6493361,1.4922642) (-1.6441002,1.4975002) (-1.6388642,1.5027362) (-1.6336282,1.5079722) (-1.6283922,1.5132082) (-1.6231562,1.5184442) (-1.6179202,1.5236801) (-1.6126842,1.5289161) (-1.6074482,1.5341521) (-1.6022123,1.5393881) (-1.5969763,1.5446241) (-1.5917403,1.5498601) (-1.5865043,1.5550961) (-1.5812683,1.560332) (-1.5760323,1.565568) (-1.5707963,1.570804) (-1.5655603,1.57604) (-1.5603244,1.581276) (-1.5550884,1.586512) (-1.5498524,1.591748) (-1.5446164,1.596984) (-1.5393804,1.6022199) (-1.5341444,1.6074559) (-1.5289084,1.6126919) (-1.5236724,1.6179279) (-1.5184364,1.6231639) (-1.5132005,1.6283999) (-1.5079645,1.6336359) (-1.5027285,1.6388719) (-1.4974925,1.6441079) (-1.4922565,1.6493439) (-1.4870205,1.6545798) (-1.4817845,1.6598158) (-1.4765485,1.6650518) (-1.4713126,1.6702878) (-1.4660766,1.6755238) (-1.4608406,1.6807598) (-1.4556046,1.6859958) (-1.4503686,1.6912318) (-1.4451326,1.6964678) (-1.4398966,1.7017038) (-1.4346606,1.7069398) (-1.4294247,1.7121758) (-1.4241887,1.7174118) (-1.4189527,1.7226478) (-1.4137167,1.7278837) (-1.4084807,1.7331197) (-1.4032447,1.7383557) (-1.3980087,1.7435917) (-1.3927727,1.7488277) (-1.3875368,1.7540637) (-1.3823008,1.7592997) (-1.3770648,1.7645357) (-1.3718288,1.7697717) (-1.3665928,1.7750077) (-1.3613568,1.7802437) (-1.3561208,1.7854797) (-1.3508848,1.7907157) (-1.3456489,1.7959517) (-1.3404129,1.8011877) (-1.3351769,1.8064237) (-1.3299409,1.8116597) (-1.3247049,1.8168957) (-1.3194689,1.8221317) (-1.3142329,1.8273677) (-1.3089969,1.8326037) (-1.303761,1.8378397) (-1.298525,1.8430757) (-1.293289,1.8483117) (-1.288053,1.8535477) (-1.282817,1.8587837) (-1.277581,1.8640197) (-1.272345,1.8692557) (-1.267109,1.8744917) (-1.261873,1.8797277) (-1.2566371,1.8849637) (-1.2514011,1.8901997) (-1.2461651,1.8954357) (-1.2409291,1.9006717) (-1.2356931,1.9059077) (-1.2304571,1.9111437) (-1.2252211,1.9163797) (-1.2199851,1.9216157) (-1.2147492,1.9268517) (-1.2095132,1.9320877) (-1.2042772,1.9373237) (-1.1990412,1.9425597) (-1.1938052,1.9477957) (-1.1885692,1.9530317) (-1.1833332,1.9582677) (-1.1780972,1.9635037) (-1.1728613,1.9687397) (-1.1676253,1.9739757) (-1.1623893,1.9792118) (-1.1571533,1.9844478) (-1.1519173,1.9896838) (-1.1466813,1.9949198) (-1.1414453,2.0001558) (-1.1362093,2.0053918) (-1.1309734,2.0106278) (-1.1257374,2.0158638) (-1.1205014,2.0210998) (-1.1152654,2.0263358) (-1.1100294,2.0315718) (-1.1047934,2.0368078) (-1.0995574,2.0420439) (-1.0943214,2.0472799) (-1.0890855,2.0525159) (-1.0838495,2.0577519) (-1.0786135,2.0629879) (-1.0733775,2.0682239) (-1.0681415,2.0734599) (-1.0629055,2.0786959) (-1.0576695,2.083932) (-1.0524335,2.089168) (-1.0471976,2.094404) (-1.0419616,2.09964) (-1.0367256,2.104876) (-1.0314896,2.110112) (-1.0262536,2.115348) (-1.0210176,2.1205841) (-1.0157816,2.1258201) (-1.0105456,2.1310561) (-1.0053096,2.1362921) (-1.0000737,2.1415281) (-0.99483767,2.1467642) (-0.98960169,2.1520002) (-0.9843657,2.1572362) (-0.97912971,2.1624722) (-0.97389372,2.1677082) (-0.96865773,2.1729443) (-0.96342175,2.1781803) (-0.95818576,2.1834163) (-0.95294977,2.1886523) (-0.94771378,2.1938883) (-0.9424778,2.1991244) (-0.93724181,2.2043604) (-0.93200582,2.2095964) (-0.92676983,2.2148324) (-0.92153385,2.2200685) (-0.91629786,2.2253045) (-0.91106187,2.2305405) (-0.90582588,2.2357765) (-0.90058989,2.2410126) (-0.89535391,2.2462486) (-0.89011792,2.2514846) (-0.88488193,2.2567207) (-0.87964594,2.2619567) (-0.87440996,2.2671927) (-0.86917397,2.2724288) (-0.86393798,2.2776648) (-0.85870199,2.2829008) (-0.853466,2.2881369) (-0.84823002,2.2933729) (-0.84299403,2.2986089) (-0.83775804,2.303845) (-0.83252205,2.309081) (-0.82728607,2.314317) (-0.82205008,2.3195531) (-0.81681409,2.3247891) (-0.8115781,2.3300252) (-0.80634211,2.3352612) (-0.80110613,2.3404972) (-0.79587014,2.3457333) (-0.79063415,2.3509693) (-0.78539816,2.3562054) (-0.78016218,2.3614414) (-0.77492619,2.3666775) (-0.7696902,2.3719135) (-0.76445421,2.3771496) (-0.75921822,2.3823856) (-0.75398224,2.3876217) (-0.74874625,2.3928577) (-0.74351026,2.3980938) (-0.73827427,2.4033298) (-0.73303829,2.4085659) (-0.7278023,2.4138019) (-0.72256631,2.419038) (-0.71733032,2.424274) (-0.71209433,2.4295101) (-0.70685835,2.4347461) (-0.70162236,2.4399822) (-0.69638637,2.4452183) (-0.69115038,2.4504543) (-0.6859144,2.4556904) (-0.68067841,2.4609265) (-0.67544242,2.4661625) (-0.67020643,2.4713986) (-0.66497045,2.4766347) (-0.65973446,2.4818707) (-0.65449847,2.4871068) (-0.64926248,2.4923429) (-0.64402649,2.497579) (-0.63879051,2.502815) (-0.63355452,2.5080511) (-0.62831853,2.5132872) (-0.62308254,2.5185233) (-0.61784656,2.5237594) (-0.61261057,2.5289955) (-0.60737458,2.5342316) (-0.60213859,2.5394676) (-0.5969026,2.5447037) (-0.59166662,2.5499398) (-0.58643063,2.5551759) (-0.58119464,2.560412) (-0.57595865,2.5656481) (-0.57072267,2.5708842) (-0.56548668,2.5761203) (-0.56025069,2.5813564) (-0.5550147,2.5865925) (-0.54977871,2.5918287) (-0.54454273,2.5970648) (-0.53930674,2.6023009) (-0.53407075,2.607537) (-0.52883476,2.6127731) (-0.52359878,2.6180093) (-0.51836279,2.6232454) (-0.5131268,2.6284815) (-0.50789081,2.6337177) (-0.50265482,2.6389538) (-0.49741884,2.6441899) (-0.49218285,2.6494261) (-0.48694686,2.6546622) (-0.48171087,2.6598984) (-0.47647489,2.6651345) (-0.4712389,2.6703707) (-0.46600291,2.6756069) (-0.46076692,2.680843) (-0.45553093,2.6860792) (-0.45029495,2.6913154) (-0.44505896,2.6965516) (-0.43982297,2.7017877) (-0.43458698,2.7070239) (-0.429351,2.7122601) (-0.42411501,2.7174963) (-0.41887902,2.7227325) (-0.41364303,2.7279688) (-0.40840704,2.733205) (-0.40317106,2.7384412) (-0.39793507,2.7436774) (-0.39269908,2.7489137) (-0.38746309,2.7541499) (-0.38222711,2.7593862) (-0.37699112,2.7646224) (-0.37175513,2.7698587) (-0.36651914,2.775095) (-0.36128316,2.7803313) (-0.35604717,2.7855676) (-0.35081118,2.7908039) (-0.34557519,2.7960402) (-0.3403392,2.8012765) (-0.33510322,2.8065128) (-0.32986723,2.8117492) (-0.32463124,2.8169855) (-0.31939525,2.8222219) (-0.31415927,2.8274583) (-0.30892328,2.8326947) (-0.30368729,2.8379311) (-0.2984513,2.8431675) (-0.29321531,2.848404) (-0.28797933,2.8536404) (-0.28274334,2.8588769) (-0.27750735,2.8641134) (-0.27227136,2.8693499) (-0.26703538,2.8745864) (-0.26179939,2.879823) (-0.2565634,2.8850596) (-0.25132741,2.8902962) (-0.24609142,2.8955328) (-0.24085544,2.9007695) (-0.23561945,2.9060062) (-0.23038346,2.9112429) (-0.22514747,2.9164796) (-0.21991149,2.9217164) (-0.2146755,2.9269533) (-0.20943951,2.9321901) (-0.20420352,2.9374271) (-0.19896753,2.942664) (-0.19373155,2.9479011) (-0.18849556,2.9531381) (-0.18325957,2.9583753) (-0.17802358,2.9636125) (-0.1727876,2.9688498) (-0.16755161,2.9740872) (-0.16231562,2.9793246) (-0.15707963,2.9845622) (-0.15184364,2.9897999) (-0.14660766,2.9950377) (-0.14137167,3.0002756) (-0.13613568,3.0055137) (-0.13089969,3.0107519) (-0.12566371,3.0159903) (-0.12042772,3.021229) (-0.11519173,3.0264679) (-0.10995574,3.031707) (-0.10471976,3.0369465) (-0.099483767,3.0421864) (-0.09424778,3.0474266) (-0.089011792,3.0526674) (-0.083775804,3.0579088) (-0.078539816,3.0631509) (-0.073303829,3.0683939) (-0.068067841,3.073638) (-0.062831853,3.0788834) (-0.057595865,3.0841306) (-0.052359878,3.0893799) (-0.04712389,3.0946323) (-0.041887902,3.0998888) (-0.036651914,3.1051512) (-0.031415927,3.1104226) (-0.026179939,3.1157082) (-0.020943951,3.1210192) (-0.015707963,3.1263822) (-0.010471976,3.1318829) (-0.0052359878,3.1381238) (0.0052359878,3.1381238) (0.010471976,3.1318829) (0.015707963,3.1263822) (0.020943951,3.1210192) (0.026179939,3.1157082) (0.031415927,3.1104226) (0.036651914,3.1051512) (0.041887902,3.0998888) (0.04712389,3.0946323) (0.052359878,3.0893799) (0.057595865,3.0841306) (0.062831853,3.0788834) (0.068067841,3.073638) (0.073303829,3.0683939) (0.078539816,3.0631509) (0.083775804,3.0579088) (0.089011792,3.0526674) (0.09424778,3.0474266) (0.099483767,3.0421864) (0.10471976,3.0369465) (0.10995574,3.031707) (0.11519173,3.0264679) (0.12042772,3.021229) (0.12566371,3.0159903) (0.13089969,3.0107519) (0.13613568,3.0055137) (0.14137167,3.0002756) (0.14660766,2.9950377) (0.15184364,2.9897999) (0.15707963,2.9845622) (0.16231562,2.9793246) (0.16755161,2.9740872) (0.1727876,2.9688498) (0.17802358,2.9636125) (0.18325957,2.9583753) (0.18849556,2.9531381) (0.19373155,2.9479011) (0.19896753,2.942664) (0.20420352,2.9374271) (0.20943951,2.9321901) (0.2146755,2.9269533) (0.21991149,2.9217164) (0.22514747,2.9164796) (0.23038346,2.9112429) (0.23561945,2.9060062) (0.24085544,2.9007695) (0.24609142,2.8955328) (0.25132741,2.8902962) (0.2565634,2.8850596) (0.26179939,2.879823) (0.26703538,2.8745864) (0.27227136,2.8693499) (0.27750735,2.8641134) (0.28274334,2.8588769) (0.28797933,2.8536404) (0.29321531,2.848404) (0.2984513,2.8431675) (0.30368729,2.8379311) (0.30892328,2.8326947) (0.31415927,2.8274583) (0.31939525,2.8222219) (0.32463124,2.8169855) (0.32986723,2.8117492) (0.33510322,2.8065128) (0.3403392,2.8012765) (0.34557519,2.7960402) (0.35081118,2.7908039) (0.35604717,2.7855676) (0.36128316,2.7803313) (0.36651914,2.775095) (0.37175513,2.7698587) (0.37699112,2.7646224) (0.38222711,2.7593862) (0.38746309,2.7541499) (0.39269908,2.7489137) (0.39793507,2.7436774) (0.40317106,2.7384412) (0.40840704,2.733205) (0.41364303,2.7279688) (0.41887902,2.7227325) (0.42411501,2.7174963) (0.429351,2.7122601) (0.43458698,2.7070239) (0.43982297,2.7017877) (0.44505896,2.6965516) (0.45029495,2.6913154) (0.45553093,2.6860792) (0.46076692,2.680843) (0.46600291,2.6756069) (0.4712389,2.6703707) (0.47647489,2.6651345) (0.48171087,2.6598984) (0.48694686,2.6546622) (0.49218285,2.6494261) (0.49741884,2.6441899) (0.50265482,2.6389538) (0.50789081,2.6337177) (0.5131268,2.6284815) (0.51836279,2.6232454) (0.52359878,2.6180093) (0.52883476,2.6127731) (0.53407075,2.607537) (0.53930674,2.6023009) (0.54454273,2.5970648) (0.54977871,2.5918287) (0.5550147,2.5865925) (0.56025069,2.5813564) (0.56548668,2.5761203) (0.57072267,2.5708842) (0.57595865,2.5656481) (0.58119464,2.560412) (0.58643063,2.5551759) (0.59166662,2.5499398) (0.5969026,2.5447037) (0.60213859,2.5394676) (0.60737458,2.5342316) (0.61261057,2.5289955) (0.61784656,2.5237594) (0.62308254,2.5185233) (0.62831853,2.5132872) (0.63355452,2.5080511) (0.63879051,2.502815) (0.64402649,2.497579) (0.64926248,2.4923429) (0.65449847,2.4871068) (0.65973446,2.4818707) (0.66497045,2.4766347) (0.67020643,2.4713986) (0.67544242,2.4661625) (0.68067841,2.4609265) (0.6859144,2.4556904) (0.69115038,2.4504543) (0.69638637,2.4452183) (0.70162236,2.4399822) (0.70685835,2.4347461) (0.71209433,2.4295101) (0.71733032,2.424274) (0.72256631,2.419038) (0.7278023,2.4138019) (0.73303829,2.4085659) (0.73827427,2.4033298) (0.74351026,2.3980938) (0.74874625,2.3928577) (0.75398224,2.3876217) (0.75921822,2.3823856) (0.76445421,2.3771496) (0.7696902,2.3719135) (0.77492619,2.3666775) (0.78016218,2.3614414) (0.78539816,2.3562054) (0.79063415,2.3509693) (0.79587014,2.3457333) (0.80110613,2.3404972) (0.80634211,2.3352612) (0.8115781,2.3300252) (0.81681409,2.3247891) (0.82205008,2.3195531) (0.82728607,2.314317) (0.83252205,2.309081) (0.83775804,2.303845) (0.84299403,2.2986089) (0.84823002,2.2933729) (0.853466,2.2881369) (0.85870199,2.2829008) (0.86393798,2.2776648) (0.86917397,2.2724288) (0.87440996,2.2671927) (0.87964594,2.2619567) (0.88488193,2.2567207) (0.89011792,2.2514846) (0.89535391,2.2462486) (0.90058989,2.2410126) (0.90582588,2.2357765) (0.91106187,2.2305405) (0.91629786,2.2253045) (0.92153385,2.2200685) (0.92676983,2.2148324) (0.93200582,2.2095964) (0.93724181,2.2043604) (0.9424778,2.1991244) (0.94771378,2.1938883) (0.95294977,2.1886523) (0.95818576,2.1834163) (0.96342175,2.1781803) (0.96865773,2.1729443) (0.97389372,2.1677082) (0.97912971,2.1624722) (0.9843657,2.1572362) (0.98960169,2.1520002) (0.99483767,2.1467642) (1.0000737,2.1415281) (1.0053096,2.1362921) (1.0105456,2.1310561) (1.0157816,2.1258201) (1.0210176,2.1205841) (1.0262536,2.115348) (1.0314896,2.110112) (1.0367256,2.104876) (1.0419616,2.09964) (1.0471976,2.094404) (1.0524335,2.089168) (1.0576695,2.083932) (1.0629055,2.0786959) (1.0681415,2.0734599) (1.0733775,2.0682239) (1.0786135,2.0629879) (1.0838495,2.0577519) (1.0890855,2.0525159) (1.0943214,2.0472799) (1.0995574,2.0420439) (1.1047934,2.0368078) (1.1100294,2.0315718) (1.1152654,2.0263358) (1.1205014,2.0210998) (1.1257374,2.0158638) (1.1309734,2.0106278) (1.1362093,2.0053918) (1.1414453,2.0001558) (1.1466813,1.9949198) (1.1519173,1.9896838) (1.1571533,1.9844478) (1.1623893,1.9792118) (1.1676253,1.9739757) (1.1728613,1.9687397) (1.1780972,1.9635037) (1.1833332,1.9582677) (1.1885692,1.9530317) (1.1938052,1.9477957) (1.1990412,1.9425597) (1.2042772,1.9373237) (1.2095132,1.9320877) (1.2147492,1.9268517) (1.2199851,1.9216157) (1.2252211,1.9163797) (1.2304571,1.9111437) (1.2356931,1.9059077) (1.2409291,1.9006717) (1.2461651,1.8954357) (1.2514011,1.8901997) (1.2566371,1.8849637) (1.261873,1.8797277) (1.267109,1.8744917) (1.272345,1.8692557) (1.277581,1.8640197) (1.282817,1.8587837) (1.288053,1.8535477) (1.293289,1.8483117) (1.298525,1.8430757) (1.303761,1.8378397) (1.3089969,1.8326037) (1.3142329,1.8273677) (1.3194689,1.8221317) (1.3247049,1.8168957) (1.3299409,1.8116597) (1.3351769,1.8064237) (1.3404129,1.8011877) (1.3456489,1.7959517) (1.3508848,1.7907157) (1.3561208,1.7854797) (1.3613568,1.7802437) (1.3665928,1.7750077) (1.3718288,1.7697717) (1.3770648,1.7645357) (1.3823008,1.7592997) (1.3875368,1.7540637) (1.3927727,1.7488277) (1.3980087,1.7435917) (1.4032447,1.7383557) (1.4084807,1.7331197) (1.4137167,1.7278837) (1.4189527,1.7226478) (1.4241887,1.7174118) (1.4294247,1.7121758) (1.4346606,1.7069398) (1.4398966,1.7017038) (1.4451326,1.6964678) (1.4503686,1.6912318) (1.4556046,1.6859958) (1.4608406,1.6807598) (1.4660766,1.6755238) (1.4713126,1.6702878) (1.4765485,1.6650518) (1.4817845,1.6598158) (1.4870205,1.6545798) (1.4922565,1.6493439) (1.4974925,1.6441079) (1.5027285,1.6388719) (1.5079645,1.6336359) (1.5132005,1.6283999) (1.5184364,1.6231639) (1.5236724,1.6179279) (1.5289084,1.6126919) (1.5341444,1.6074559) (1.5393804,1.6022199) (1.5446164,1.596984) (1.5498524,1.591748) (1.5550884,1.586512) (1.5603244,1.581276) (1.5655603,1.57604) (1.5707963,1.570804) (1.5760323,1.565568) (1.5812683,1.560332) (1.5865043,1.5550961) (1.5917403,1.5498601) (1.5969763,1.5446241) (1.6022123,1.5393881) (1.6074482,1.5341521) (1.6126842,1.5289161) (1.6179202,1.5236801) (1.6231562,1.5184442) (1.6283922,1.5132082) (1.6336282,1.5079722) (1.6388642,1.5027362) (1.6441002,1.4975002) (1.6493361,1.4922642) (1.6545721,1.4870282) (1.6598081,1.4817923) (1.6650441,1.4765563) (1.6702801,1.4713203) (1.6755161,1.4660843) (1.6807521,1.4608483) (1.6859881,1.4556123) (1.691224,1.4503764) (1.69646,1.4451404) (1.701696,1.4399044) (1.706932,1.4346684) (1.712168,1.4294324) (1.717404,1.4241964) (1.72264,1.4189605) (1.727876,1.4137245) (1.7331119,1.4084885) (1.7383479,1.4032525) (1.7435839,1.3980165) (1.7488199,1.3927806) (1.7540559,1.3875446) (1.7592919,1.3823086) (1.7645279,1.3770726) (1.7697639,1.3718366) (1.7749998,1.3666007) (1.7802358,1.3613647) (1.7854718,1.3561287) (1.7907078,1.3508927) (1.7959438,1.3456567) (1.8011798,1.3404208) (1.8064158,1.3351848) (1.8116518,1.3299488) (1.8168878,1.3247128) (1.8221237,1.3194769) (1.8273597,1.3142409) (1.8325957,1.3090049) (1.8378317,1.3037689) (1.8430677,1.2985329) (1.8483037,1.293297) (1.8535397,1.288061) (1.8587757,1.282825) (1.8640116,1.277589) (1.8692476,1.2723531) (1.8744836,1.2671171) (1.8797196,1.2618811) (1.8849556,1.2566451) (1.8901916,1.2514092) (1.8954276,1.2461732) (1.9006636,1.2409372) (1.9058995,1.2357013) (1.9111355,1.2304653) (1.9163715,1.2252293) (1.9216075,1.2199933) (1.9268435,1.2147574) (1.9320795,1.2095214) (1.9373155,1.2042854) (1.9425515,1.1990495) (1.9477874,1.1938135) (1.9530234,1.1885775) (1.9582594,1.1833415) (1.9634954,1.1781056) (1.9687314,1.1728696) (1.9739674,1.1676336) (1.9792034,1.1623977) (1.9844394,1.1571617) (1.9896753,1.1519257) (1.9949113,1.1466898) (2.0001473,1.1414538) (2.0053833,1.1362178) (2.0106193,1.1309819) (2.0158553,1.1257459) (2.0210913,1.1205099) (2.0263273,1.115274) (2.0315632,1.110038) (2.0367992,1.104802) (2.0420352,1.0995661) (2.0472712,1.0943301) (2.0525072,1.0890941) (2.0577432,1.0838582) (2.0629792,1.0786222) (2.0682152,1.0733862) (2.0734512,1.0681503) (2.0786871,1.0629143) (2.0839231,1.0576784) (2.0891591,1.0524424) (2.0943951,1.0472064) (2.0996311,1.0419705) (2.1048671,1.0367345) (2.1101031,1.0314986) (2.1153391,1.0262626) (2.120575,1.0210266) (2.125811,1.0157907) (2.131047,1.0105547) (2.136283,1.0053188) (2.141519,1.0000828) (2.146755,0.99484684) (2.151991,0.98961089) (2.157227,0.98437493) (2.1624629,0.97913898) (2.1676989,0.97390302) (2.1729349,0.96866707) (2.1781709,0.96343111) (2.1834069,0.95819516) (2.1886429,0.95295921) (2.1938789,0.94772325) (2.1991149,0.9424873) (2.2043508,0.93725135) (2.2095868,0.9320154) (2.2148228,0.92677945) (2.2200588,0.9215435) (2.2252948,0.91630755) (2.2305308,0.9110716) (2.2357668,0.90583566) (2.2410028,0.90059971) (2.2462387,0.89536376) (2.2514747,0.89012782) (2.2567107,0.88489187) (2.2619467,0.87965593) (2.2671827,0.87441998) (2.2724187,0.86918404) (2.2776547,0.86394809) (2.2828907,0.85871215) (2.2881266,0.85347621) (2.2933626,0.84824027) (2.2985986,0.84300433) (2.3038346,0.83776839) (2.3090706,0.83253245) (2.3143066,0.82729651) (2.3195426,0.82206058) (2.3247786,0.81682464) (2.3300146,0.81158871) (2.3352505,0.80635277) (2.3404865,0.80111684) (2.3457225,0.7958809) (2.3509585,0.79064497) (2.3561945,0.78540904) (2.3614305,0.78017311) (2.3666665,0.77493718) (2.3719025,0.76970125) (2.3771384,0.76446532) (2.3823744,0.7592294) (2.3876104,0.75399347) (2.3928464,0.74875755) (2.3980824,0.74352162) (2.4033184,0.7382857) (2.4085544,0.73304978) (2.4137904,0.72781386) (2.4190263,0.72257794) (2.4242623,0.71734202) (2.4294983,0.71210611) (2.4347343,0.70687019) (2.4399703,0.70163428) (2.4452063,0.69639836) (2.4504423,0.69116245) (2.4556783,0.68592654) (2.4609142,0.68069063) (2.4661502,0.67545472) (2.4713862,0.67021882) (2.4766222,0.66498291) (2.4818582,0.65974701) (2.4870942,0.6545111) (2.4923302,0.6492752) (2.4975662,0.6440393) (2.5028021,0.63880341) (2.5080381,0.63356751) (2.5132741,0.62833162) (2.5185101,0.62309572) (2.5237461,0.61785983) (2.5289821,0.61262394) (2.5342181,0.60738806) (2.5394541,0.60215217) (2.54469,0.59691629) (2.549926,0.59168041) (2.555162,0.58644453) (2.560398,0.58120865) (2.565634,0.57597277) (2.57087,0.5707369) (2.576106,0.56550103) (2.581342,0.56026516) (2.586578,0.5550293) (2.5918139,0.54979343) (2.5970499,0.54455757) (2.6022859,0.53932172) (2.6075219,0.53408586) (2.6127579,0.52885001) (2.6179939,0.52361416) (2.6232299,0.51837831) (2.6284659,0.51314247) (2.6337018,0.50790663) (2.6389378,0.50267079) (2.6441738,0.49743496) (2.6494098,0.49219913) (2.6546458,0.4869633) (2.6598818,0.48172747) (2.6651178,0.47649166) (2.6703538,0.47125584) (2.6755897,0.46602003) (2.6808257,0.46078422) (2.6860617,0.45554842) (2.6912977,0.45031262) (2.6965337,0.44507682) (2.7017697,0.43984104) (2.7070057,0.43460525) (2.7122417,0.42936947) (2.7174776,0.4241337) (2.7227136,0.41889793) (2.7279496,0.41366217) (2.7331856,0.40842641) (2.7384216,0.40319066) (2.7436576,0.39795492) (2.7488936,0.39271918) (2.7541296,0.38748345) (2.7593655,0.38224773) (2.7646015,0.37701201) (2.7698375,0.3717763) (2.7750735,0.3665406) (2.7803095,0.36130491) (2.7855455,0.35606923) (2.7907815,0.35083356) (2.7960175,0.3455979) (2.8012534,0.34036224) (2.8064894,0.3351266) (2.8117254,0.32989097) (2.8169614,0.32465535) (2.8221974,0.31941975) (2.8274334,0.31418415) (2.8326694,0.30894857) (2.8379054,0.30371301) (2.8431414,0.29847746) (2.8483773,0.29324192) (2.8536133,0.28800641) (2.8588493,0.28277091) (2.8640853,0.27753542) (2.8693213,0.27229996) (2.8745573,0.26706452) (2.8797933,0.2618291) (2.8850293,0.25659371) (2.8902652,0.25135834) (2.8955012,0.24612299) (2.9007372,0.24088768) (2.9059732,0.23565239) (2.9112092,0.23041714) (2.9164452,0.22518192) (2.9216812,0.21994674) (2.9269172,0.2147116) (2.9321531,0.2094765) (2.9373891,0.20424145) (2.9426251,0.19900644) (2.9478611,0.19377149) (2.9530971,0.1885366) (2.9583331,0.18330177) (2.9635691,0.17806701) (2.9688051,0.17283233) (2.974041,0.16759772) (2.979277,0.16236321) (2.984513,0.15712879) (2.989749,0.15189448) (2.994985,0.1466603) (3.000221,0.14142625) (3.005457,0.13619234) (3.010693,0.13095861) (3.0159289,0.12572506) (3.0211649,0.12049172) (3.0264009,0.11525863) (3.0316369,0.11002581) (3.0368729,0.10479331) (3.0421089,0.099561178) (3.0473449,0.094329473) (3.0525809,0.089098273) (3.0578168,0.08386767) (3.0630528,0.078637786) (3.0682888,0.073408772) (3.0735248,0.06818083) (3.0787608,0.062954227) (3.0839968,0.057729325) (3.0892328,0.052506634) (3.0944688,0.047286885) (3.0997048,0.042071173) (3.1049407,0.036861213) (3.1101767,0.031659844) (3.1154127,0.026472131) (3.1206487,0.02130805) (3.1258847,0.016190229) (3.1311207,0.011182367) (3.1363567,0.0065420497) (3.1415927,0.0039220882) (3.1415927,-0.0039220882) (3.1363567,-0.0065420497) (3.1311207,-0.011182367) (3.1258847,-0.016190229) (3.1206487,-0.02130805) (3.1154127,-0.026472131) (3.1101767,-0.031659844) (3.1049407,-0.036861213) (3.0997048,-0.042071173) (3.0944688,-0.047286885) (3.0892328,-0.052506634) (3.0839968,-0.057729325) (3.0787608,-0.062954227) (3.0735248,-0.06818083) (3.0682888,-0.073408772) (3.0630528,-0.078637786) (3.0578168,-0.08386767) (3.0525809,-0.089098273) (3.0473449,-0.094329473) (3.0421089,-0.099561178) (3.0368729,-0.10479331) (3.0316369,-0.11002581) (3.0264009,-0.11525863) (3.0211649,-0.12049172) (3.0159289,-0.12572506) (3.010693,-0.13095861) (3.005457,-0.13619234) (3.000221,-0.14142625) (2.994985,-0.1466603) (2.989749,-0.15189448) (2.984513,-0.15712879) (2.979277,-0.16236321) (2.974041,-0.16759772) (2.9688051,-0.17283233) (2.9635691,-0.17806701) (2.9583331,-0.18330177) (2.9530971,-0.1885366) (2.9478611,-0.19377149) (2.9426251,-0.19900644) (2.9373891,-0.20424145) (2.9321531,-0.2094765) (2.9269172,-0.2147116) (2.9216812,-0.21994674) (2.9164452,-0.22518192) (2.9112092,-0.23041714) (2.9059732,-0.23565239) (2.9007372,-0.24088768) (2.8955012,-0.24612299) (2.8902652,-0.25135834) (2.8850293,-0.25659371) (2.8797933,-0.2618291) (2.8745573,-0.26706452) (2.8693213,-0.27229996) (2.8640853,-0.27753542) (2.8588493,-0.28277091) (2.8536133,-0.28800641) (2.8483773,-0.29324192) (2.8431414,-0.29847746) (2.8379054,-0.30371301) (2.8326694,-0.30894857) (2.8274334,-0.31418415) (2.8221974,-0.31941975) (2.8169614,-0.32465535) (2.8117254,-0.32989097) (2.8064894,-0.3351266) (2.8012534,-0.34036224) (2.7960175,-0.3455979) (2.7907815,-0.35083356) (2.7855455,-0.35606923) (2.7803095,-0.36130491) (2.7750735,-0.3665406) (2.7698375,-0.3717763) (2.7646015,-0.37701201) (2.7593655,-0.38224773) (2.7541296,-0.38748345) (2.7488936,-0.39271918) (2.7436576,-0.39795492) (2.7384216,-0.40319066) (2.7331856,-0.40842641) (2.7279496,-0.41366217) (2.7227136,-0.41889793) (2.7174776,-0.4241337) (2.7122417,-0.42936947) (2.7070057,-0.43460525) (2.7017697,-0.43984104) (2.6965337,-0.44507682) (2.6912977,-0.45031262) (2.6860617,-0.45554842) (2.6808257,-0.46078422) (2.6755897,-0.46602003) (2.6703538,-0.47125584) (2.6651178,-0.47649166) (2.6598818,-0.48172747) (2.6546458,-0.4869633) (2.6494098,-0.49219913) (2.6441738,-0.49743496) (2.6389378,-0.50267079) (2.6337018,-0.50790663) (2.6284659,-0.51314247) (2.6232299,-0.51837831) (2.6179939,-0.52361416) (2.6127579,-0.52885001) (2.6075219,-0.53408586) (2.6022859,-0.53932172) (2.5970499,-0.54455757) (2.5918139,-0.54979343) (2.586578,-0.5550293) (2.581342,-0.56026516) (2.576106,-0.56550103) (2.57087,-0.5707369) (2.565634,-0.57597277) (2.560398,-0.58120865) (2.555162,-0.58644453) (2.549926,-0.59168041) (2.54469,-0.59691629) (2.5394541,-0.60215217) (2.5342181,-0.60738806) (2.5289821,-0.61262394) (2.5237461,-0.61785983) (2.5185101,-0.62309572) (2.5132741,-0.62833162) (2.5080381,-0.63356751) (2.5028021,-0.63880341) (2.4975662,-0.6440393) (2.4923302,-0.6492752) (2.4870942,-0.6545111) (2.4818582,-0.65974701) (2.4766222,-0.66498291) (2.4713862,-0.67021882) (2.4661502,-0.67545472) (2.4609142,-0.68069063) (2.4556783,-0.68592654) (2.4504423,-0.69116245) (2.4452063,-0.69639836) (2.4399703,-0.70163428) (2.4347343,-0.70687019) (2.4294983,-0.71210611) (2.4242623,-0.71734202) (2.4190263,-0.72257794) (2.4137904,-0.72781386) (2.4085544,-0.73304978) (2.4033184,-0.7382857) (2.3980824,-0.74352162) (2.3928464,-0.74875755) (2.3876104,-0.75399347) (2.3823744,-0.7592294) (2.3771384,-0.76446532) (2.3719025,-0.76970125) (2.3666665,-0.77493718) (2.3614305,-0.78017311) (2.3561945,-0.78540904) (2.3509585,-0.79064497) (2.3457225,-0.7958809) (2.3404865,-0.80111684) (2.3352505,-0.80635277) (2.3300146,-0.81158871) (2.3247786,-0.81682464) (2.3195426,-0.82206058) (2.3143066,-0.82729651) (2.3090706,-0.83253245) (2.3038346,-0.83776839) (2.2985986,-0.84300433) (2.2933626,-0.84824027) (2.2881266,-0.85347621) (2.2828907,-0.85871215) (2.2776547,-0.86394809) (2.2724187,-0.86918404) (2.2671827,-0.87441998) (2.2619467,-0.87965593) (2.2567107,-0.88489187) (2.2514747,-0.89012782) (2.2462387,-0.89536376) (2.2410028,-0.90059971) (2.2357668,-0.90583566) (2.2305308,-0.9110716) (2.2252948,-0.91630755) (2.2200588,-0.9215435) (2.2148228,-0.92677945) (2.2095868,-0.9320154) (2.2043508,-0.93725135) (2.1991149,-0.9424873) (2.1938789,-0.94772325) (2.1886429,-0.95295921) (2.1834069,-0.95819516) (2.1781709,-0.96343111) (2.1729349,-0.96866707) (2.1676989,-0.97390302) (2.1624629,-0.97913898) (2.157227,-0.98437493) (2.151991,-0.98961089) (2.146755,-0.99484684) (2.141519,-1.0000828) (2.136283,-1.0053188) (2.131047,-1.0105547) (2.125811,-1.0157907) (2.120575,-1.0210266) (2.1153391,-1.0262626) (2.1101031,-1.0314986) (2.1048671,-1.0367345) (2.0996311,-1.0419705) (2.0943951,-1.0472064) (2.0891591,-1.0524424) (2.0839231,-1.0576784) (2.0786871,-1.0629143) (2.0734512,-1.0681503) (2.0682152,-1.0733862) (2.0629792,-1.0786222) (2.0577432,-1.0838582) (2.0525072,-1.0890941) (2.0472712,-1.0943301) (2.0420352,-1.0995661) (2.0367992,-1.104802) (2.0315632,-1.110038) (2.0263273,-1.115274) (2.0210913,-1.1205099) (2.0158553,-1.1257459) (2.0106193,-1.1309819) (2.0053833,-1.1362178) (2.0001473,-1.1414538) (1.9949113,-1.1466898) (1.9896753,-1.1519257) (1.9844394,-1.1571617) (1.9792034,-1.1623977) (1.9739674,-1.1676336) (1.9687314,-1.1728696) (1.9634954,-1.1781056) (1.9582594,-1.1833415) (1.9530234,-1.1885775) (1.9477874,-1.1938135) (1.9425515,-1.1990495) (1.9373155,-1.2042854) (1.9320795,-1.2095214) (1.9268435,-1.2147574) (1.9216075,-1.2199933) (1.9163715,-1.2252293) (1.9111355,-1.2304653) (1.9058995,-1.2357013) (1.9006636,-1.2409372) (1.8954276,-1.2461732) (1.8901916,-1.2514092) (1.8849556,-1.2566451) (1.8797196,-1.2618811) (1.8744836,-1.2671171) (1.8692476,-1.2723531) (1.8640116,-1.277589) (1.8587757,-1.282825) (1.8535397,-1.288061) (1.8483037,-1.293297) (1.8430677,-1.2985329) (1.8378317,-1.3037689) (1.8325957,-1.3090049) (1.8273597,-1.3142409) (1.8221237,-1.3194769) (1.8168878,-1.3247128) (1.8116518,-1.3299488) (1.8064158,-1.3351848) (1.8011798,-1.3404208) (1.7959438,-1.3456567) (1.7907078,-1.3508927) (1.7854718,-1.3561287) (1.7802358,-1.3613647) (1.7749998,-1.3666007) (1.7697639,-1.3718366) (1.7645279,-1.3770726) (1.7592919,-1.3823086) (1.7540559,-1.3875446) (1.7488199,-1.3927806) (1.7435839,-1.3980165) (1.7383479,-1.4032525) (1.7331119,-1.4084885) (1.727876,-1.4137245) (1.72264,-1.4189605) (1.717404,-1.4241964) (1.712168,-1.4294324) (1.706932,-1.4346684) (1.701696,-1.4399044) (1.69646,-1.4451404) (1.691224,-1.4503764) (1.6859881,-1.4556123) (1.6807521,-1.4608483) (1.6755161,-1.4660843) (1.6702801,-1.4713203) (1.6650441,-1.4765563) (1.6598081,-1.4817923) (1.6545721,-1.4870282) (1.6493361,-1.4922642) (1.6441002,-1.4975002) (1.6388642,-1.5027362) (1.6336282,-1.5079722) (1.6283922,-1.5132082) (1.6231562,-1.5184442) (1.6179202,-1.5236801) (1.6126842,-1.5289161) (1.6074482,-1.5341521) (1.6022123,-1.5393881) (1.5969763,-1.5446241) (1.5917403,-1.5498601) (1.5865043,-1.5550961) (1.5812683,-1.560332) (1.5760323,-1.565568) (1.5707963,-1.570804) (1.5655603,-1.57604) (1.5603244,-1.581276) (1.5550884,-1.586512) (1.5498524,-1.591748) (1.5446164,-1.596984) (1.5393804,-1.6022199) (1.5341444,-1.6074559) (1.5289084,-1.6126919) (1.5236724,-1.6179279) (1.5184364,-1.6231639) (1.5132005,-1.6283999) (1.5079645,-1.6336359) (1.5027285,-1.6388719) (1.4974925,-1.6441079) (1.4922565,-1.6493439) (1.4870205,-1.6545798) (1.4817845,-1.6598158) (1.4765485,-1.6650518) (1.4713126,-1.6702878) (1.4660766,-1.6755238) (1.4608406,-1.6807598) (1.4556046,-1.6859958) (1.4503686,-1.6912318) (1.4451326,-1.6964678) (1.4398966,-1.7017038) (1.4346606,-1.7069398) (1.4294247,-1.7121758) (1.4241887,-1.7174118) (1.4189527,-1.7226478) (1.4137167,-1.7278837) (1.4084807,-1.7331197) (1.4032447,-1.7383557) (1.3980087,-1.7435917) (1.3927727,-1.7488277) (1.3875368,-1.7540637) (1.3823008,-1.7592997) (1.3770648,-1.7645357) (1.3718288,-1.7697717) (1.3665928,-1.7750077) (1.3613568,-1.7802437) (1.3561208,-1.7854797) (1.3508848,-1.7907157) (1.3456489,-1.7959517) (1.3404129,-1.8011877) (1.3351769,-1.8064237) (1.3299409,-1.8116597) (1.3247049,-1.8168957) (1.3194689,-1.8221317) (1.3142329,-1.8273677) (1.3089969,-1.8326037) (1.303761,-1.8378397) (1.298525,-1.8430757) (1.293289,-1.8483117) (1.288053,-1.8535477) (1.282817,-1.8587837) (1.277581,-1.8640197) (1.272345,-1.8692557) (1.267109,-1.8744917) (1.261873,-1.8797277) (1.2566371,-1.8849637) (1.2514011,-1.8901997) (1.2461651,-1.8954357) (1.2409291,-1.9006717) (1.2356931,-1.9059077) (1.2304571,-1.9111437) (1.2252211,-1.9163797) (1.2199851,-1.9216157) (1.2147492,-1.9268517) (1.2095132,-1.9320877) (1.2042772,-1.9373237) (1.1990412,-1.9425597) (1.1938052,-1.9477957) (1.1885692,-1.9530317) (1.1833332,-1.9582677) (1.1780972,-1.9635037) (1.1728613,-1.9687397) (1.1676253,-1.9739757) (1.1623893,-1.9792118) (1.1571533,-1.9844478) (1.1519173,-1.9896838) (1.1466813,-1.9949198) (1.1414453,-2.0001558) (1.1362093,-2.0053918) (1.1309734,-2.0106278) (1.1257374,-2.0158638) (1.1205014,-2.0210998) (1.1152654,-2.0263358) (1.1100294,-2.0315718) (1.1047934,-2.0368078) (1.0995574,-2.0420439) (1.0943214,-2.0472799) (1.0890855,-2.0525159) (1.0838495,-2.0577519) (1.0786135,-2.0629879) (1.0733775,-2.0682239) (1.0681415,-2.0734599) (1.0629055,-2.0786959) (1.0576695,-2.083932) (1.0524335,-2.089168) (1.0471976,-2.094404) (1.0419616,-2.09964) (1.0367256,-2.104876) (1.0314896,-2.110112) (1.0262536,-2.115348) (1.0210176,-2.1205841) (1.0157816,-2.1258201) (1.0105456,-2.1310561) (1.0053096,-2.1362921) (1.0000737,-2.1415281) (0.99483767,-2.1467642) (0.98960169,-2.1520002) (0.9843657,-2.1572362) (0.97912971,-2.1624722) (0.97389372,-2.1677082) (0.96865773,-2.1729443) (0.96342175,-2.1781803) (0.95818576,-2.1834163) (0.95294977,-2.1886523) (0.94771378,-2.1938883) (0.9424778,-2.1991244) (0.93724181,-2.2043604) (0.93200582,-2.2095964) (0.92676983,-2.2148324) (0.92153385,-2.2200685) (0.91629786,-2.2253045) (0.91106187,-2.2305405) (0.90582588,-2.2357765) (0.90058989,-2.2410126) (0.89535391,-2.2462486) (0.89011792,-2.2514846) (0.88488193,-2.2567207) (0.87964594,-2.2619567) (0.87440996,-2.2671927) (0.86917397,-2.2724288) (0.86393798,-2.2776648) (0.85870199,-2.2829008) (0.853466,-2.2881369) (0.84823002,-2.2933729) (0.84299403,-2.2986089) (0.83775804,-2.303845) (0.83252205,-2.309081) (0.82728607,-2.314317) (0.82205008,-2.3195531) (0.81681409,-2.3247891) (0.8115781,-2.3300252) (0.80634211,-2.3352612) (0.80110613,-2.3404972) (0.79587014,-2.3457333) (0.79063415,-2.3509693) (0.78539816,-2.3562054) (0.78016218,-2.3614414) (0.77492619,-2.3666775) (0.7696902,-2.3719135) (0.76445421,-2.3771496) (0.75921822,-2.3823856) (0.75398224,-2.3876217) (0.74874625,-2.3928577) (0.74351026,-2.3980938) (0.73827427,-2.4033298) (0.73303829,-2.4085659) (0.7278023,-2.4138019) (0.72256631,-2.419038) (0.71733032,-2.424274) (0.71209433,-2.4295101) (0.70685835,-2.4347461) (0.70162236,-2.4399822) (0.69638637,-2.4452183) (0.69115038,-2.4504543) (0.6859144,-2.4556904) (0.68067841,-2.4609265) (0.67544242,-2.4661625) (0.67020643,-2.4713986) (0.66497045,-2.4766347) (0.65973446,-2.4818707) (0.65449847,-2.4871068) (0.64926248,-2.4923429) (0.64402649,-2.497579) (0.63879051,-2.502815) (0.63355452,-2.5080511) (0.62831853,-2.5132872) (0.62308254,-2.5185233) (0.61784656,-2.5237594) (0.61261057,-2.5289955) (0.60737458,-2.5342316) (0.60213859,-2.5394676) (0.5969026,-2.5447037) (0.59166662,-2.5499398) (0.58643063,-2.5551759) (0.58119464,-2.560412) (0.57595865,-2.5656481) (0.57072267,-2.5708842) (0.56548668,-2.5761203) (0.56025069,-2.5813564) (0.5550147,-2.5865925) (0.54977871,-2.5918287) (0.54454273,-2.5970648) (0.53930674,-2.6023009) (0.53407075,-2.607537) (0.52883476,-2.6127731) (0.52359878,-2.6180093) (0.51836279,-2.6232454) (0.5131268,-2.6284815) (0.50789081,-2.6337177) (0.50265482,-2.6389538) (0.49741884,-2.6441899) (0.49218285,-2.6494261) (0.48694686,-2.6546622) (0.48171087,-2.6598984) (0.47647489,-2.6651345) (0.4712389,-2.6703707) (0.46600291,-2.6756069) (0.46076692,-2.680843) (0.45553093,-2.6860792) (0.45029495,-2.6913154) (0.44505896,-2.6965516) (0.43982297,-2.7017877) (0.43458698,-2.7070239) (0.429351,-2.7122601) (0.42411501,-2.7174963) (0.41887902,-2.7227325) (0.41364303,-2.7279688) (0.40840704,-2.733205) (0.40317106,-2.7384412) (0.39793507,-2.7436774) (0.39269908,-2.7489137) (0.38746309,-2.7541499) (0.38222711,-2.7593862) (0.37699112,-2.7646224) (0.37175513,-2.7698587) (0.36651914,-2.775095) (0.36128316,-2.7803313) (0.35604717,-2.7855676) (0.35081118,-2.7908039) (0.34557519,-2.7960402) (0.3403392,-2.8012765) (0.33510322,-2.8065128) (0.32986723,-2.8117492) (0.32463124,-2.8169855) (0.31939525,-2.8222219) (0.31415927,-2.8274583) (0.30892328,-2.8326947) (0.30368729,-2.8379311) (0.2984513,-2.8431675) (0.29321531,-2.848404) (0.28797933,-2.8536404) (0.28274334,-2.8588769) (0.27750735,-2.8641134) (0.27227136,-2.8693499) (0.26703538,-2.8745864) (0.26179939,-2.879823) (0.2565634,-2.8850596) (0.25132741,-2.8902962) (0.24609142,-2.8955328) (0.24085544,-2.9007695) (0.23561945,-2.9060062) (0.23038346,-2.9112429) (0.22514747,-2.9164796) (0.21991149,-2.9217164) (0.2146755,-2.9269533) (0.20943951,-2.9321901) (0.20420352,-2.9374271) (0.19896753,-2.942664) (0.19373155,-2.9479011) (0.18849556,-2.9531381) (0.18325957,-2.9583753) (0.17802358,-2.9636125) (0.1727876,-2.9688498) (0.16755161,-2.9740872) (0.16231562,-2.9793246) (0.15707963,-2.9845622) (0.15184364,-2.9897999) (0.14660766,-2.9950377) (0.14137167,-3.0002756) (0.13613568,-3.0055137) (0.13089969,-3.0107519) (0.12566371,-3.0159903) (0.12042772,-3.021229) (0.11519173,-3.0264679) (0.10995574,-3.031707) (0.10471976,-3.0369465) (0.099483767,-3.0421864) (0.09424778,-3.0474266) (0.089011792,-3.0526674) (0.083775804,-3.0579088) (0.078539816,-3.0631509) (0.073303829,-3.0683939) (0.068067841,-3.073638) (0.062831853,-3.0788834) (0.057595865,-3.0841306) (0.052359878,-3.0893799) (0.04712389,-3.0946323) (0.041887902,-3.0998888) (0.036651914,-3.1051512) (0.031415927,-3.1104226) (0.026179939,-3.1157082) (0.020943951,-3.1210192) (0.015707963,-3.1263822) (0.010471976,-3.1318829) (0.0052359878,-3.1381238) (-0.0052359878,-3.1381238) (-0.010471976,-3.1318829) (-0.015707963,-3.1263822) (-0.020943951,-3.1210192) (-0.026179939,-3.1157082) (-0.031415927,-3.1104226) (-0.036651914,-3.1051512) (-0.041887902,-3.0998888) (-0.04712389,-3.0946323) (-0.052359878,-3.0893799) (-0.057595865,-3.0841306) (-0.062831853,-3.0788834) (-0.068067841,-3.073638) (-0.073303829,-3.0683939) (-0.078539816,-3.0631509) (-0.083775804,-3.0579088) (-0.089011792,-3.0526674) (-0.09424778,-3.0474266) (-0.099483767,-3.0421864) (-0.10471976,-3.0369465) (-0.10995574,-3.031707) (-0.11519173,-3.0264679) (-0.12042772,-3.021229) (-0.12566371,-3.0159903) (-0.13089969,-3.0107519) (-0.13613568,-3.0055137) (-0.14137167,-3.0002756) (-0.14660766,-2.9950377) (-0.15184364,-2.9897999) (-0.15707963,-2.9845622) (-0.16231562,-2.9793246) (-0.16755161,-2.9740872) (-0.1727876,-2.9688498) (-0.17802358,-2.9636125) (-0.18325957,-2.9583753) (-0.18849556,-2.9531381) (-0.19373155,-2.9479011) (-0.19896753,-2.942664) (-0.20420352,-2.9374271) (-0.20943951,-2.9321901) (-0.2146755,-2.9269533) (-0.21991149,-2.9217164) (-0.22514747,-2.9164796) (-0.23038346,-2.9112429) (-0.23561945,-2.9060062) (-0.24085544,-2.9007695) (-0.24609142,-2.8955328) (-0.25132741,-2.8902962) (-0.2565634,-2.8850596) (-0.26179939,-2.879823) (-0.26703538,-2.8745864) (-0.27227136,-2.8693499) (-0.27750735,-2.8641134) (-0.28274334,-2.8588769) (-0.28797933,-2.8536404) (-0.29321531,-2.848404) (-0.2984513,-2.8431675) (-0.30368729,-2.8379311) (-0.30892328,-2.8326947) (-0.31415927,-2.8274583) (-0.31939525,-2.8222219) (-0.32463124,-2.8169855) (-0.32986723,-2.8117492) (-0.33510322,-2.8065128) (-0.3403392,-2.8012765) (-0.34557519,-2.7960402) (-0.35081118,-2.7908039) (-0.35604717,-2.7855676) (-0.36128316,-2.7803313) (-0.36651914,-2.775095) (-0.37175513,-2.7698587) (-0.37699112,-2.7646224) (-0.38222711,-2.7593862) (-0.38746309,-2.7541499) (-0.39269908,-2.7489137) (-0.39793507,-2.7436774) (-0.40317106,-2.7384412) (-0.40840704,-2.733205) (-0.41364303,-2.7279688) (-0.41887902,-2.7227325) (-0.42411501,-2.7174963) (-0.429351,-2.7122601) (-0.43458698,-2.7070239) (-0.43982297,-2.7017877) (-0.44505896,-2.6965516) (-0.45029495,-2.6913154) (-0.45553093,-2.6860792) (-0.46076692,-2.680843) (-0.46600291,-2.6756069) (-0.4712389,-2.6703707) (-0.47647489,-2.6651345) (-0.48171087,-2.6598984) (-0.48694686,-2.6546622) (-0.49218285,-2.6494261) (-0.49741884,-2.6441899) (-0.50265482,-2.6389538) (-0.50789081,-2.6337177) (-0.5131268,-2.6284815) (-0.51836279,-2.6232454) (-0.52359878,-2.6180093) (-0.52883476,-2.6127731) (-0.53407075,-2.607537) (-0.53930674,-2.6023009) (-0.54454273,-2.5970648) (-0.54977871,-2.5918287) (-0.5550147,-2.5865925) (-0.56025069,-2.5813564) (-0.56548668,-2.5761203) (-0.57072267,-2.5708842) (-0.57595865,-2.5656481) (-0.58119464,-2.560412) (-0.58643063,-2.5551759) (-0.59166662,-2.5499398) (-0.5969026,-2.5447037) (-0.60213859,-2.5394676) (-0.60737458,-2.5342316) (-0.61261057,-2.5289955) (-0.61784656,-2.5237594) (-0.62308254,-2.5185233) (-0.62831853,-2.5132872) (-0.63355452,-2.5080511) (-0.63879051,-2.502815) (-0.64402649,-2.497579) (-0.64926248,-2.4923429) (-0.65449847,-2.4871068) (-0.65973446,-2.4818707) (-0.66497045,-2.4766347) (-0.67020643,-2.4713986) (-0.67544242,-2.4661625) (-0.68067841,-2.4609265) (-0.6859144,-2.4556904) (-0.69115038,-2.4504543) (-0.69638637,-2.4452183) (-0.70162236,-2.4399822) (-0.70685835,-2.4347461) (-0.71209433,-2.4295101) (-0.71733032,-2.424274) (-0.72256631,-2.419038) (-0.7278023,-2.4138019) (-0.73303829,-2.4085659) (-0.73827427,-2.4033298) (-0.74351026,-2.3980938) (-0.74874625,-2.3928577) (-0.75398224,-2.3876217) (-0.75921822,-2.3823856) (-0.76445421,-2.3771496) (-0.7696902,-2.3719135) (-0.77492619,-2.3666775) (-0.78016218,-2.3614414) (-0.78539816,-2.3562054) (-0.79063415,-2.3509693) (-0.79587014,-2.3457333) (-0.80110613,-2.3404972) (-0.80634211,-2.3352612) (-0.8115781,-2.3300252) (-0.81681409,-2.3247891) (-0.82205008,-2.3195531) (-0.82728607,-2.314317) (-0.83252205,-2.309081) (-0.83775804,-2.303845) (-0.84299403,-2.2986089) (-0.84823002,-2.2933729) (-0.853466,-2.2881369) (-0.85870199,-2.2829008) (-0.86393798,-2.2776648) (-0.86917397,-2.2724288) (-0.87440996,-2.2671927) (-0.87964594,-2.2619567) (-0.88488193,-2.2567207) (-0.89011792,-2.2514846) (-0.89535391,-2.2462486) (-0.90058989,-2.2410126) (-0.90582588,-2.2357765) (-0.91106187,-2.2305405) (-0.91629786,-2.2253045) (-0.92153385,-2.2200685) (-0.92676983,-2.2148324) (-0.93200582,-2.2095964) (-0.93724181,-2.2043604) (-0.9424778,-2.1991244) (-0.94771378,-2.1938883) (-0.95294977,-2.1886523) (-0.95818576,-2.1834163) (-0.96342175,-2.1781803) (-0.96865773,-2.1729443) (-0.97389372,-2.1677082) (-0.97912971,-2.1624722) (-0.9843657,-2.1572362) (-0.98960169,-2.1520002) (-0.99483767,-2.1467642) (-1.0000737,-2.1415281) (-1.0053096,-2.1362921) (-1.0105456,-2.1310561) (-1.0157816,-2.1258201) (-1.0210176,-2.1205841) (-1.0262536,-2.115348) (-1.0314896,-2.110112) (-1.0367256,-2.104876) (-1.0419616,-2.09964) (-1.0471976,-2.094404) (-1.0524335,-2.089168) (-1.0576695,-2.083932) (-1.0629055,-2.0786959) (-1.0681415,-2.0734599) (-1.0733775,-2.0682239) (-1.0786135,-2.0629879) (-1.0838495,-2.0577519) (-1.0890855,-2.0525159) (-1.0943214,-2.0472799) (-1.0995574,-2.0420439) (-1.1047934,-2.0368078) (-1.1100294,-2.0315718) (-1.1152654,-2.0263358) (-1.1205014,-2.0210998) (-1.1257374,-2.0158638) (-1.1309734,-2.0106278) (-1.1362093,-2.0053918) (-1.1414453,-2.0001558) (-1.1466813,-1.9949198) (-1.1519173,-1.9896838) (-1.1571533,-1.9844478) (-1.1623893,-1.9792118) (-1.1676253,-1.9739757) (-1.1728613,-1.9687397) (-1.1780972,-1.9635037) (-1.1833332,-1.9582677) (-1.1885692,-1.9530317) (-1.1938052,-1.9477957) (-1.1990412,-1.9425597) (-1.2042772,-1.9373237) (-1.2095132,-1.9320877) (-1.2147492,-1.9268517) (-1.2199851,-1.9216157) (-1.2252211,-1.9163797) (-1.2304571,-1.9111437) (-1.2356931,-1.9059077) (-1.2409291,-1.9006717) (-1.2461651,-1.8954357) (-1.2514011,-1.8901997) (-1.2566371,-1.8849637) (-1.261873,-1.8797277) (-1.267109,-1.8744917) (-1.272345,-1.8692557) (-1.277581,-1.8640197) (-1.282817,-1.8587837) (-1.288053,-1.8535477) (-1.293289,-1.8483117) (-1.298525,-1.8430757) (-1.303761,-1.8378397) (-1.3089969,-1.8326037) (-1.3142329,-1.8273677) (-1.3194689,-1.8221317) (-1.3247049,-1.8168957) (-1.3299409,-1.8116597) (-1.3351769,-1.8064237) (-1.3404129,-1.8011877) (-1.3456489,-1.7959517) (-1.3508848,-1.7907157) (-1.3561208,-1.7854797) (-1.3613568,-1.7802437) (-1.3665928,-1.7750077) (-1.3718288,-1.7697717) (-1.3770648,-1.7645357) (-1.3823008,-1.7592997) (-1.3875368,-1.7540637) (-1.3927727,-1.7488277) (-1.3980087,-1.7435917) (-1.4032447,-1.7383557) (-1.4084807,-1.7331197) (-1.4137167,-1.7278837) (-1.4189527,-1.7226478) (-1.4241887,-1.7174118) (-1.4294247,-1.7121758) (-1.4346606,-1.7069398) (-1.4398966,-1.7017038) (-1.4451326,-1.6964678) (-1.4503686,-1.6912318) (-1.4556046,-1.6859958) (-1.4608406,-1.6807598) (-1.4660766,-1.6755238) (-1.4713126,-1.6702878) (-1.4765485,-1.6650518) (-1.4817845,-1.6598158) (-1.4870205,-1.6545798) (-1.4922565,-1.6493439) (-1.4974925,-1.6441079) (-1.5027285,-1.6388719) (-1.5079645,-1.6336359) (-1.5132005,-1.6283999) (-1.5184364,-1.6231639) (-1.5236724,-1.6179279) (-1.5289084,-1.6126919) (-1.5341444,-1.6074559) (-1.5393804,-1.6022199) (-1.5446164,-1.596984) (-1.5498524,-1.591748) (-1.5550884,-1.586512) (-1.5603244,-1.581276) (-1.5655603,-1.57604) (-1.5707963,-1.570804) (-1.5760323,-1.565568) (-1.5812683,-1.560332) (-1.5865043,-1.5550961) (-1.5917403,-1.5498601) (-1.5969763,-1.5446241) (-1.6022123,-1.5393881) (-1.6074482,-1.5341521) (-1.6126842,-1.5289161) (-1.6179202,-1.5236801) (-1.6231562,-1.5184442) (-1.6283922,-1.5132082) (-1.6336282,-1.5079722) (-1.6388642,-1.5027362) (-1.6441002,-1.4975002) (-1.6493361,-1.4922642) (-1.6545721,-1.4870282) (-1.6598081,-1.4817923) (-1.6650441,-1.4765563) (-1.6702801,-1.4713203) (-1.6755161,-1.4660843) (-1.6807521,-1.4608483) (-1.6859881,-1.4556123) (-1.691224,-1.4503764) (-1.69646,-1.4451404) (-1.701696,-1.4399044) (-1.706932,-1.4346684) (-1.712168,-1.4294324) (-1.717404,-1.4241964) (-1.72264,-1.4189605) (-1.727876,-1.4137245) (-1.7331119,-1.4084885) (-1.7383479,-1.4032525) (-1.7435839,-1.3980165) (-1.7488199,-1.3927806) (-1.7540559,-1.3875446) (-1.7592919,-1.3823086) (-1.7645279,-1.3770726) (-1.7697639,-1.3718366) (-1.7749998,-1.3666007) (-1.7802358,-1.3613647) (-1.7854718,-1.3561287) (-1.7907078,-1.3508927) (-1.7959438,-1.3456567) (-1.8011798,-1.3404208) (-1.8064158,-1.3351848) (-1.8116518,-1.3299488) (-1.8168878,-1.3247128) (-1.8221237,-1.3194769) (-1.8273597,-1.3142409) (-1.8325957,-1.3090049) (-1.8378317,-1.3037689) (-1.8430677,-1.2985329) (-1.8483037,-1.293297) (-1.8535397,-1.288061) (-1.8587757,-1.282825) (-1.8640116,-1.277589) (-1.8692476,-1.2723531) (-1.8744836,-1.2671171) (-1.8797196,-1.2618811) (-1.8849556,-1.2566451) (-1.8901916,-1.2514092) (-1.8954276,-1.2461732) (-1.9006636,-1.2409372) (-1.9058995,-1.2357013) (-1.9111355,-1.2304653) (-1.9163715,-1.2252293) (-1.9216075,-1.2199933) (-1.9268435,-1.2147574) (-1.9320795,-1.2095214) (-1.9373155,-1.2042854) (-1.9425515,-1.1990495) (-1.9477874,-1.1938135) (-1.9530234,-1.1885775) (-1.9582594,-1.1833415) (-1.9634954,-1.1781056) (-1.9687314,-1.1728696) (-1.9739674,-1.1676336) (-1.9792034,-1.1623977) (-1.9844394,-1.1571617) (-1.9896753,-1.1519257) (-1.9949113,-1.1466898) (-2.0001473,-1.1414538) (-2.0053833,-1.1362178) (-2.0106193,-1.1309819) (-2.0158553,-1.1257459) (-2.0210913,-1.1205099) (-2.0263273,-1.115274) (-2.0315632,-1.110038) (-2.0367992,-1.104802) (-2.0420352,-1.0995661) (-2.0472712,-1.0943301) (-2.0525072,-1.0890941) (-2.0577432,-1.0838582) (-2.0629792,-1.0786222) (-2.0682152,-1.0733862) (-2.0734512,-1.0681503) (-2.0786871,-1.0629143) (-2.0839231,-1.0576784) (-2.0891591,-1.0524424) (-2.0943951,-1.0472064) (-2.0996311,-1.0419705) (-2.1048671,-1.0367345) (-2.1101031,-1.0314986) (-2.1153391,-1.0262626) (-2.120575,-1.0210266) (-2.125811,-1.0157907) (-2.131047,-1.0105547) (-2.136283,-1.0053188) (-2.141519,-1.0000828) (-2.146755,-0.99484684) (-2.151991,-0.98961089) (-2.157227,-0.98437493) (-2.1624629,-0.97913898) (-2.1676989,-0.97390302) (-2.1729349,-0.96866707) (-2.1781709,-0.96343111) (-2.1834069,-0.95819516) (-2.1886429,-0.95295921) (-2.1938789,-0.94772325) (-2.1991149,-0.9424873) (-2.2043508,-0.93725135) (-2.2095868,-0.9320154) (-2.2148228,-0.92677945) (-2.2200588,-0.9215435) (-2.2252948,-0.91630755) (-2.2305308,-0.9110716) (-2.2357668,-0.90583566) (-2.2410028,-0.90059971) (-2.2462387,-0.89536376) (-2.2514747,-0.89012782) (-2.2567107,-0.88489187) (-2.2619467,-0.87965593) (-2.2671827,-0.87441998) (-2.2724187,-0.86918404) (-2.2776547,-0.86394809) (-2.2828907,-0.85871215) (-2.2881266,-0.85347621) (-2.2933626,-0.84824027) (-2.2985986,-0.84300433) (-2.3038346,-0.83776839) (-2.3090706,-0.83253245) (-2.3143066,-0.82729651) (-2.3195426,-0.82206058) (-2.3247786,-0.81682464) (-2.3300146,-0.81158871) (-2.3352505,-0.80635277) (-2.3404865,-0.80111684) (-2.3457225,-0.7958809) (-2.3509585,-0.79064497) (-2.3561945,-0.78540904) (-2.3614305,-0.78017311) (-2.3666665,-0.77493718) (-2.3719025,-0.76970125) (-2.3771384,-0.76446532) (-2.3823744,-0.7592294) (-2.3876104,-0.75399347) (-2.3928464,-0.74875755) (-2.3980824,-0.74352162) (-2.4033184,-0.7382857) (-2.4085544,-0.73304978) (-2.4137904,-0.72781386) (-2.4190263,-0.72257794) (-2.4242623,-0.71734202) (-2.4294983,-0.71210611) (-2.4347343,-0.70687019) (-2.4399703,-0.70163428) (-2.4452063,-0.69639836) (-2.4504423,-0.69116245) (-2.4556783,-0.68592654) (-2.4609142,-0.68069063) (-2.4661502,-0.67545472) (-2.4713862,-0.67021882) (-2.4766222,-0.66498291) (-2.4818582,-0.65974701) (-2.4870942,-0.6545111) (-2.4923302,-0.6492752) (-2.4975662,-0.6440393) (-2.5028021,-0.63880341) (-2.5080381,-0.63356751) (-2.5132741,-0.62833162) (-2.5185101,-0.62309572) (-2.5237461,-0.61785983) (-2.5289821,-0.61262394) (-2.5342181,-0.60738806) (-2.5394541,-0.60215217) (-2.54469,-0.59691629) (-2.549926,-0.59168041) (-2.555162,-0.58644453) (-2.560398,-0.58120865) (-2.565634,-0.57597277) (-2.57087,-0.5707369) (-2.576106,-0.56550103) (-2.581342,-0.56026516) (-2.586578,-0.5550293) (-2.5918139,-0.54979343) (-2.5970499,-0.54455757) (-2.6022859,-0.53932172) (-2.6075219,-0.53408586) (-2.6127579,-0.52885001) (-2.6179939,-0.52361416) (-2.6232299,-0.51837831) (-2.6284659,-0.51314247) (-2.6337018,-0.50790663) (-2.6389378,-0.50267079) (-2.6441738,-0.49743496) (-2.6494098,-0.49219913) (-2.6546458,-0.4869633) (-2.6598818,-0.48172747) (-2.6651178,-0.47649166) (-2.6703538,-0.47125584) (-2.6755897,-0.46602003) (-2.6808257,-0.46078422) (-2.6860617,-0.45554842) (-2.6912977,-0.45031262) (-2.6965337,-0.44507682) (-2.7017697,-0.43984104) (-2.7070057,-0.43460525) (-2.7122417,-0.42936947) (-2.7174776,-0.4241337) (-2.7227136,-0.41889793) (-2.7279496,-0.41366217) (-2.7331856,-0.40842641) (-2.7384216,-0.40319066) (-2.7436576,-0.39795492) (-2.7488936,-0.39271918) (-2.7541296,-0.38748345) (-2.7593655,-0.38224773) (-2.7646015,-0.37701201) (-2.7698375,-0.3717763) (-2.7750735,-0.3665406) (-2.7803095,-0.36130491) (-2.7855455,-0.35606923) (-2.7907815,-0.35083356) (-2.7960175,-0.3455979) (-2.8012534,-0.34036224) (-2.8064894,-0.3351266) (-2.8117254,-0.32989097) (-2.8169614,-0.32465535) (-2.8221974,-0.31941975) (-2.8274334,-0.31418415) (-2.8326694,-0.30894857) (-2.8379054,-0.30371301) (-2.8431414,-0.29847746) (-2.8483773,-0.29324192) (-2.8536133,-0.28800641) (-2.8588493,-0.28277091) (-2.8640853,-0.27753542) (-2.8693213,-0.27229996) (-2.8745573,-0.26706452) (-2.8797933,-0.2618291) (-2.8850293,-0.25659371) (-2.8902652,-0.25135834) (-2.8955012,-0.24612299) (-2.9007372,-0.24088768) (-2.9059732,-0.23565239) (-2.9112092,-0.23041714) (-2.9164452,-0.22518192) (-2.9216812,-0.21994674) (-2.9269172,-0.2147116) (-2.9321531,-0.2094765) (-2.9373891,-0.20424145) (-2.9426251,-0.19900644) (-2.9478611,-0.19377149) (-2.9530971,-0.1885366) (-2.9583331,-0.18330177) (-2.9635691,-0.17806701) (-2.9688051,-0.17283233) (-2.974041,-0.16759772) (-2.979277,-0.16236321) (-2.984513,-0.15712879) (-2.989749,-0.15189448) (-2.994985,-0.1466603) (-3.000221,-0.14142625) (-3.005457,-0.13619234) (-3.010693,-0.13095861) (-3.0159289,-0.12572506) (-3.0211649,-0.12049172) (-3.0264009,-0.11525863) (-3.0316369,-0.11002581) (-3.0368729,-0.10479331) (-3.0421089,-0.099561178) (-3.0473449,-0.094329473) (-3.0525809,-0.089098273) (-3.0578168,-0.08386767) (-3.0630528,-0.078637786) (-3.0682888,-0.073408772) (-3.0735248,-0.06818083) (-3.0787608,-0.062954227) (-3.0839968,-0.057729325) (-3.0892328,-0.052506634) (-3.0944688,-0.047286885) (-3.0997048,-0.042071173) (-3.1049407,-0.036861213) (-3.1101767,-0.031659844) (-3.1154127,-0.026472131) (-3.1206487,-0.02130805) (-3.1258847,-0.016190229) (-3.1311207,-0.011182367) (-3.1363567,-0.0065420497) (-3.1415927,-0.0039220882)};
\addlegendentry{$p=0$}
\addplot[thick,color=main] coordinates {(-2.586578,0.057843896) (-2.581342,0.094303316) (-2.576106,0.12038277) (-2.57087,0.1419246) (-2.565634,0.16076333) (-2.560398,0.17776042) (-2.555162,0.193402) (-2.549926,0.20799408) (-2.54469,0.22174392) (-2.5394541,0.23479951) (-2.5342181,0.24727081) (-2.5289821,0.25924214) (-2.5237461,0.27077978) (-2.5185101,0.28193697) (-2.5132741,0.29275719) (-2.5080381,0.30327648) (-2.5028021,0.31352511) (-2.4975662,0.32352878) (-2.4923302,0.33330953) (-2.4870942,0.34288642) (-2.4818582,0.35227605) (-2.4766222,0.36149299) (-2.4713862,0.37055012) (-2.4661502,0.37945885) (-2.4609142,0.38822939) (-2.4556783,0.39687087) (-2.4504423,0.40539153) (-2.4452063,0.41379881) (-2.4399703,0.42209948) (-2.4347343,0.43029968) (-2.4294983,0.43840503) (-2.4242623,0.44642069) (-2.4190263,0.45435138) (-2.4137904,0.46220145) (-2.4085544,0.46997493) (-2.4033184,0.47767555) (-2.3980824,0.48530674) (-2.3928464,0.49287173) (-2.3876104,0.50037348) (-2.3823744,0.5078148) (-2.3771384,0.51519828) (-2.3719025,0.52252636) (-2.3666665,0.52980132) (-2.3614305,0.53702529) (-2.3561945,0.5442003) (-2.3509585,0.55132824) (-2.3457225,0.55841089) (-2.3404865,0.56544994) (-2.3352505,0.57244698) (-2.3300146,0.57940351) (-2.3247786,0.58632097) (-2.3195426,0.59320069) (-2.3143066,0.60004397) (-2.3090706,0.60685202) (-2.3038346,0.61362599) (-2.2985986,0.62036699) (-2.2933626,0.62707605) (-2.2881266,0.63375418) (-2.2828907,0.64040234) (-2.2776547,0.64702142) (-2.2724187,0.65361229) (-2.2671827,0.66017578) (-2.2619467,0.66671268) (-2.2567107,0.67322376) (-2.2514747,0.67970973) (-2.2462387,0.68617129) (-2.2410028,0.69260911) (-2.2357668,0.69902383) (-2.2305308,0.70541605) (-2.2252948,0.71178638) (-2.2200588,0.71813537) (-2.2148228,0.72446357) (-2.2095868,0.7307715) (-2.2043508,0.73705967) (-2.1991149,0.74332856) (-2.1938789,0.74957863) (-2.1886429,0.75581034) (-2.1834069,0.76202412) (-2.1781709,0.76822039) (-2.1729349,0.77439955) (-2.1676989,0.78056199) (-2.1624629,0.78670808) (-2.157227,0.79283819) (-2.151991,0.79895267) (-2.146755,0.80505186) (-2.141519,0.81113609) (-2.136283,0.81720566) (-2.131047,0.82326089) (-2.125811,0.82930208) (-2.120575,0.83532951) (-2.1153391,0.84134346) (-2.1101031,0.8473442) (-2.1048671,0.85333199) (-2.0996311,0.85930708) (-2.0943951,0.86526971) (-2.0891591,0.87122013) (-2.0839231,0.87715857) (-2.0786871,0.88308525) (-2.0734512,0.88900038) (-2.0682152,0.89490418) (-2.0629792,0.90079685) (-2.0577432,0.90667858) (-2.0525072,0.91254958) (-2.0472712,0.91841003) (-2.0420352,0.92426012) (-2.0367992,0.9301) (-2.0315632,0.93592987) (-2.0263273,0.94174989) (-2.0210913,0.94756022) (-2.0158553,0.95336101) (-2.0106193,0.95915242) (-2.0053833,0.96493461) (-2.0001473,0.97070771) (-1.9949113,0.97647187) (-1.9896753,0.98222723) (-1.9844394,0.98797392) (-1.9792034,0.99371206) (-1.9739674,0.9994418) (-1.9687314,1.0051632) (-1.9634954,1.0108765) (-1.9582594,1.0165818) (-1.9530234,1.022279) (-1.9477874,1.0279685) (-1.9425515,1.0336503) (-1.9373155,1.0393244) (-1.9320795,1.0449911) (-1.9268435,1.0506503) (-1.9216075,1.0563022) (-1.9163715,1.0619469) (-1.9111355,1.0675845) (-1.9058995,1.0732151) (-1.9006636,1.0788387) (-1.8954276,1.0844555) (-1.8901916,1.0900655) (-1.8849556,1.0956688) (-1.8797196,1.1012656) (-1.8744836,1.1068558) (-1.8692476,1.1124395) (-1.8640116,1.1180169) (-1.8587757,1.1235879) (-1.8535397,1.1291527) (-1.8483037,1.1347114) (-1.8430677,1.140264) (-1.8378317,1.1458105) (-1.8325957,1.1513511) (-1.8273597,1.1568857) (-1.8221237,1.1624145) (-1.8168878,1.1679376) (-1.8116518,1.1734549) (-1.8064158,1.1789665) (-1.8011798,1.1844725) (-1.7959438,1.189973) (-1.7907078,1.195468) (-1.7854718,1.2009575) (-1.7802358,1.2064415) (-1.7749998,1.2119203) (-1.7697639,1.2173937) (-1.7645279,1.2228618) (-1.7592919,1.2283248) (-1.7540559,1.2337825) (-1.7488199,1.2392352) (-1.7435839,1.2446827) (-1.7383479,1.2501252) (-1.7331119,1.2555627) (-1.727876,1.2609952) (-1.72264,1.2664227) (-1.717404,1.2718454) (-1.712168,1.2772632) (-1.706932,1.2826762) (-1.701696,1.2880843) (-1.69646,1.2934877) (-1.691224,1.2988864) (-1.6859881,1.3042804) (-1.6807521,1.3096697) (-1.6755161,1.3150544) (-1.6702801,1.3204345) (-1.6650441,1.3258099) (-1.6598081,1.3311808) (-1.6545721,1.3365472) (-1.6493361,1.3419091) (-1.6441002,1.3472665) (-1.6388642,1.3526194) (-1.6336282,1.3579679) (-1.6283922,1.363312) (-1.6231562,1.3686517) (-1.6179202,1.373987) (-1.6126842,1.3793179) (-1.6074482,1.3846446) (-1.6022123,1.3899669) (-1.5969763,1.3952849) (-1.5917403,1.4005986) (-1.5865043,1.405908) (-1.5812683,1.4112132) (-1.5760323,1.4165142) (-1.5707963,1.421811) (-1.5655603,1.4271035) (-1.5603244,1.4323918) (-1.5550884,1.437676) (-1.5498524,1.442956) (-1.5446164,1.4482318) (-1.5393804,1.4535035) (-1.5341444,1.4587711) (-1.5289084,1.4640345) (-1.5236724,1.4692938) (-1.5184364,1.4745489) (-1.5132005,1.4798) (-1.5079645,1.485047) (-1.5027285,1.4902899) (-1.4974925,1.4955287) (-1.4922565,1.5007634) (-1.4870205,1.505994) (-1.4817845,1.5112206) (-1.4765485,1.5164431) (-1.4713126,1.5216615) (-1.4660766,1.5268759) (-1.4608406,1.5320862) (-1.4556046,1.5372924) (-1.4503686,1.5424946) (-1.4451326,1.5476927) (-1.4398966,1.5528867) (-1.4346606,1.5580767) (-1.4294247,1.5632627) (-1.4241887,1.5684445) (-1.4189527,1.5736223) (-1.4137167,1.5787961) (-1.4084807,1.5839657) (-1.4032447,1.5891313) (-1.3980087,1.5942928) (-1.3927727,1.5994502) (-1.3875368,1.6046035) (-1.3823008,1.6097527) (-1.3770648,1.6148978) (-1.3718288,1.6200388) (-1.3665928,1.6251756) (-1.3613568,1.6303083) (-1.3561208,1.6354369) (-1.3508848,1.6405613) (-1.3456489,1.6456816) (-1.3404129,1.6507977) (-1.3351769,1.6559096) (-1.3299409,1.6610173) (-1.3247049,1.6661208) (-1.3194689,1.6712201) (-1.3142329,1.6763152) (-1.3089969,1.681406) (-1.303761,1.6864925) (-1.298525,1.6915748) (-1.293289,1.6966527) (-1.288053,1.7017264) (-1.282817,1.7067957) (-1.277581,1.7118607) (-1.272345,1.7169213) (-1.267109,1.7219775) (-1.261873,1.7270294) (-1.2566371,1.7320768) (-1.2514011,1.7371198) (-1.2461651,1.7421583) (-1.2409291,1.7471923) (-1.2356931,1.7522218) (-1.2304571,1.7572468) (-1.2252211,1.7622672) (-1.2199851,1.7672831) (-1.2147492,1.7722943) (-1.2095132,1.7773009) (-1.2042772,1.7823029) (-1.1990412,1.7873002) (-1.1938052,1.7922928) (-1.1885692,1.7972806) (-1.1833332,1.8022637) (-1.1780972,1.8072419) (-1.1728613,1.8122154) (-1.1676253,1.817184) (-1.1623893,1.8221477) (-1.1571533,1.8271065) (-1.1519173,1.8320603) (-1.1466813,1.8370091) (-1.1414453,1.8419529) (-1.1362093,1.8468917) (-1.1309734,1.8518254) (-1.1257374,1.8567539) (-1.1205014,1.8616773) (-1.1152654,1.8665954) (-1.1100294,1.8715084) (-1.1047934,1.876416) (-1.0995574,1.8813183) (-1.0943214,1.8862152) (-1.0890855,1.8911068) (-1.0838495,1.8959928) (-1.0786135,1.9008734) (-1.0733775,1.9057484) (-1.0681415,1.9106179) (-1.0629055,1.9154817) (-1.0576695,1.9203398) (-1.0524335,1.9251922) (-1.0471976,1.9300388) (-1.0419616,1.9348796) (-1.0367256,1.9397145) (-1.0314896,1.9445435) (-1.0262536,1.9493664) (-1.0210176,1.9541834) (-1.0157816,1.9589942) (-1.0105456,1.9637989) (-1.0053096,1.9685973) (-1.0000737,1.9733895) (-0.99483767,1.9781754) (-0.98960169,1.9829548) (-0.9843657,1.9877278) (-0.97912971,1.9924943) (-0.97389372,1.9972542) (-0.96865773,2.0020075) (-0.96342175,2.0067541) (-0.95818576,2.0114938) (-0.95294977,2.0162268) (-0.94771378,2.0209528) (-0.9424778,2.0256718) (-0.93724181,2.0303837) (-0.93200582,2.0350885) (-0.92676983,2.0397861) (-0.92153385,2.0444765) (-0.91629786,2.0491594) (-0.91106187,2.0538349) (-0.90582588,2.0585029) (-0.90058989,2.0631632) (-0.89535391,2.0678159) (-0.89011792,2.0724608) (-0.88488193,2.0770978) (-0.87964594,2.0817269) (-0.87440996,2.0863479) (-0.86917397,2.0909608) (-0.86393798,2.0955655) (-0.85870199,2.1001618) (-0.853466,2.1047498) (-0.84823002,2.1093292) (-0.84299403,2.1139) (-0.83775804,2.1184621) (-0.83252205,2.1230153) (-0.82728607,2.1275597) (-0.82205008,2.132095) (-0.81681409,2.1366212) (-0.8115781,2.1411381) (-0.80634211,2.1456457) (-0.80110613,2.1501438) (-0.79587014,2.1546323) (-0.79063415,2.1591111) (-0.78539816,2.16358) (-0.78016218,2.1680391) (-0.77492619,2.172488) (-0.7696902,2.1769267) (-0.76445421,2.1813552) (-0.75921822,2.1857731) (-0.75398224,2.1901805) (-0.74874625,2.1945771) (-0.74351026,2.1989629) (-0.73827427,2.2033376) (-0.73303829,2.2077012) (-0.7278023,2.2120535) (-0.72256631,2.2163944) (-0.71733032,2.2207237) (-0.71209433,2.2250412) (-0.70685835,2.2293469) (-0.70162236,2.2336405) (-0.69638637,2.2379218) (-0.69115038,2.2421908) (-0.6859144,2.2464473) (-0.68067841,2.250691) (-0.67544242,2.2549219) (-0.67020643,2.2591397) (-0.66497045,2.2633443) (-0.65973446,2.2675354) (-0.65449847,2.271713) (-0.64926248,2.2758768) (-0.64402649,2.2800266) (-0.63879051,2.2841622) (-0.63355452,2.2882835) (-0.62831853,2.2923903) (-0.62308254,2.2964823) (-0.61784656,2.3005593) (-0.61261057,2.3046211) (-0.60737458,2.3086676) (-0.60213859,2.3126985) (-0.5969026,2.3167135) (-0.59166662,2.3207126) (-0.58643063,2.3246953) (-0.58119464,2.3286616) (-0.57595865,2.3326111) (-0.57072267,2.3365437) (-0.56548668,2.3404591) (-0.56025069,2.3443571) (-0.5550147,2.3482373) (-0.54977871,2.3520996) (-0.54454273,2.3559437) (-0.53930674,2.3597694) (-0.53407075,2.3635763) (-0.52883476,2.3673643) (-0.52359878,2.371133) (-0.51836279,2.3748821) (-0.5131268,2.3786115) (-0.50789081,2.3823207) (-0.50265482,2.3860096) (-0.49741884,2.3896778) (-0.49218285,2.393325) (-0.48694686,2.396951) (-0.48171087,2.4005554) (-0.47647489,2.4041379) (-0.4712389,2.4076983) (-0.46600291,2.4112361) (-0.46076692,2.4147511) (-0.45553093,2.4182429) (-0.45029495,2.4217113) (-0.44505896,2.4251559) (-0.43982297,2.4285763) (-0.43458698,2.4319722) (-0.429351,2.4353432) (-0.42411501,2.4386891) (-0.41887902,2.4420095) (-0.41364303,2.4453039) (-0.40840704,2.4485721) (-0.40317106,2.4518137) (-0.39793507,2.4550282) (-0.39269908,2.4582154) (-0.38746309,2.4613749) (-0.38222711,2.4645062) (-0.37699112,2.467609) (-0.37175513,2.4706829) (-0.36651914,2.4737275) (-0.36128316,2.4767425) (-0.35604717,2.4797273) (-0.35081118,2.4826816) (-0.34557519,2.4856051) (-0.3403392,2.4884972) (-0.33510322,2.4913576) (-0.32986723,2.4941859) (-0.32463124,2.4969817) (-0.31939525,2.4997445) (-0.31415927,2.5024739) (-0.30892328,2.5051694) (-0.30368729,2.5078308) (-0.2984513,2.5104575) (-0.29321531,2.5130491) (-0.28797933,2.5156051) (-0.28274334,2.5181252) (-0.27750735,2.5206089) (-0.27227136,2.5230558) (-0.26703538,2.5254655) (-0.26179939,2.5278375) (-0.2565634,2.5301713) (-0.25132741,2.5324666) (-0.24609142,2.5347229) (-0.24085544,2.5369398) (-0.23561945,2.5391168) (-0.23038346,2.5412535) (-0.22514747,2.5433496) (-0.21991149,2.5454045) (-0.2146755,2.5474179) (-0.20943951,2.5493893) (-0.20420352,2.5513183) (-0.19896753,2.5532045) (-0.19373155,2.5550475) (-0.18849556,2.5568469) (-0.18325957,2.5586022) (-0.17802358,2.5603132) (-0.1727876,2.5619794) (-0.16755161,2.5636003) (-0.16231562,2.5651758) (-0.15707963,2.5667052) (-0.15184364,2.5681884) (-0.14660766,2.569625) (-0.14137167,2.5710145) (-0.13613568,2.5723567) (-0.13089969,2.5736512) (-0.12566371,2.5748977) (-0.12042772,2.5760959) (-0.11519173,2.5772455) (-0.10995574,2.5783461) (-0.10471976,2.5793976) (-0.099483767,2.5803996) (-0.09424778,2.5813518) (-0.089011792,2.582254) (-0.083775804,2.5831061) (-0.078539816,2.5839076) (-0.073303829,2.5846585) (-0.068067841,2.5853585) (-0.062831853,2.5860074) (-0.057595865,2.5866051) (-0.052359878,2.5871514) (-0.04712389,2.5876461) (-0.041887902,2.5880891) (-0.036651914,2.5884802) (-0.031415927,2.5888195) (-0.026179939,2.5891067) (-0.020943951,2.5893418) (-0.015707963,2.5895247) (-0.010471976,2.5896554) (-0.0052359878,2.5897339) (-4.4408921e-16,2.58976) (0.0052359878,2.5897339) (0.010471976,2.5896554) (0.015707963,2.5895247) (0.020943951,2.5893418) (0.026179939,2.5891067) (0.031415927,2.5888195) (0.036651914,2.5884802) (0.041887902,2.5880891) (0.04712389,2.5876461) (0.052359878,2.5871514) (0.057595865,2.5866051) (0.062831853,2.5860074) (0.068067841,2.5853585) (0.073303829,2.5846585) (0.078539816,2.5839076) (0.083775804,2.5831061) (0.089011792,2.582254) (0.09424778,2.5813518) (0.099483767,2.5803996) (0.10471976,2.5793976) (0.10995574,2.5783461) (0.11519173,2.5772455) (0.12042772,2.5760959) (0.12566371,2.5748977) (0.13089969,2.5736512) (0.13613568,2.5723567) (0.14137167,2.5710145) (0.14660766,2.569625) (0.15184364,2.5681884) (0.15707963,2.5667052) (0.16231562,2.5651758) (0.16755161,2.5636003) (0.1727876,2.5619794) (0.17802358,2.5603132) (0.18325957,2.5586022) (0.18849556,2.5568469) (0.19373155,2.5550475) (0.19896753,2.5532045) (0.20420352,2.5513183) (0.20943951,2.5493893) (0.2146755,2.5474179) (0.21991149,2.5454045) (0.22514747,2.5433496) (0.23038346,2.5412535) (0.23561945,2.5391168) (0.24085544,2.5369398) (0.24609142,2.5347229) (0.25132741,2.5324666) (0.2565634,2.5301713) (0.26179939,2.5278375) (0.26703538,2.5254655) (0.27227136,2.5230558) (0.27750735,2.5206089) (0.28274334,2.5181252) (0.28797933,2.5156051) (0.29321531,2.5130491) (0.2984513,2.5104575) (0.30368729,2.5078308) (0.30892328,2.5051694) (0.31415927,2.5024739) (0.31939525,2.4997445) (0.32463124,2.4969817) (0.32986723,2.4941859) (0.33510322,2.4913576) (0.3403392,2.4884972) (0.34557519,2.4856051) (0.35081118,2.4826816) (0.35604717,2.4797273) (0.36128316,2.4767425) (0.36651914,2.4737275) (0.37175513,2.4706829) (0.37699112,2.467609) (0.38222711,2.4645062) (0.38746309,2.4613749) (0.39269908,2.4582154) (0.39793507,2.4550282) (0.40317106,2.4518137) (0.40840704,2.4485721) (0.41364303,2.4453039) (0.41887902,2.4420095) (0.42411501,2.4386891) (0.429351,2.4353432) (0.43458698,2.4319722) (0.43982297,2.4285763) (0.44505896,2.4251559) (0.45029495,2.4217113) (0.45553093,2.4182429) (0.46076692,2.4147511) (0.46600291,2.4112361) (0.4712389,2.4076983) (0.47647489,2.4041379) (0.48171087,2.4005554) (0.48694686,2.396951) (0.49218285,2.393325) (0.49741884,2.3896778) (0.50265482,2.3860096) (0.50789081,2.3823207) (0.5131268,2.3786115) (0.51836279,2.3748821) (0.52359878,2.371133) (0.52883476,2.3673643) (0.53407075,2.3635763) (0.53930674,2.3597694) (0.54454273,2.3559437) (0.54977871,2.3520996) (0.5550147,2.3482373) (0.56025069,2.3443571) (0.56548668,2.3404591) (0.57072267,2.3365437) (0.57595865,2.3326111) (0.58119464,2.3286616) (0.58643063,2.3246953) (0.59166662,2.3207126) (0.5969026,2.3167135) (0.60213859,2.3126985) (0.60737458,2.3086676) (0.61261057,2.3046211) (0.61784656,2.3005593) (0.62308254,2.2964823) (0.62831853,2.2923903) (0.63355452,2.2882835) (0.63879051,2.2841622) (0.64402649,2.2800266) (0.64926248,2.2758768) (0.65449847,2.271713) (0.65973446,2.2675354) (0.66497045,2.2633443) (0.67020643,2.2591397) (0.67544242,2.2549219) (0.68067841,2.250691) (0.6859144,2.2464473) (0.69115038,2.2421908) (0.69638637,2.2379218) (0.70162236,2.2336405) (0.70685835,2.2293469) (0.71209433,2.2250412) (0.71733032,2.2207237) (0.72256631,2.2163944) (0.7278023,2.2120535) (0.73303829,2.2077012) (0.73827427,2.2033376) (0.74351026,2.1989629) (0.74874625,2.1945771) (0.75398224,2.1901805) (0.75921822,2.1857731) (0.76445421,2.1813552) (0.7696902,2.1769267) (0.77492619,2.172488) (0.78016218,2.1680391) (0.78539816,2.16358) (0.79063415,2.1591111) (0.79587014,2.1546323) (0.80110613,2.1501438) (0.80634211,2.1456457) (0.8115781,2.1411381) (0.81681409,2.1366212) (0.82205008,2.132095) (0.82728607,2.1275597) (0.83252205,2.1230153) (0.83775804,2.1184621) (0.84299403,2.1139) (0.84823002,2.1093292) (0.853466,2.1047498) (0.85870199,2.1001618) (0.86393798,2.0955655) (0.86917397,2.0909608) (0.87440996,2.0863479) (0.87964594,2.0817269) (0.88488193,2.0770978) (0.89011792,2.0724608) (0.89535391,2.0678159) (0.90058989,2.0631632) (0.90582588,2.0585029) (0.91106187,2.0538349) (0.91629786,2.0491594) (0.92153385,2.0444765) (0.92676983,2.0397861) (0.93200582,2.0350885) (0.93724181,2.0303837) (0.9424778,2.0256718) (0.94771378,2.0209528) (0.95294977,2.0162268) (0.95818576,2.0114938) (0.96342175,2.0067541) (0.96865773,2.0020075) (0.97389372,1.9972542) (0.97912971,1.9924943) (0.9843657,1.9877278) (0.98960169,1.9829548) (0.99483767,1.9781754) (1.0000737,1.9733895) (1.0053096,1.9685973) (1.0105456,1.9637989) (1.0157816,1.9589942) (1.0210176,1.9541834) (1.0262536,1.9493664) (1.0314896,1.9445435) (1.0367256,1.9397145) (1.0419616,1.9348796) (1.0471976,1.9300388) (1.0524335,1.9251922) (1.0576695,1.9203398) (1.0629055,1.9154817) (1.0681415,1.9106179) (1.0733775,1.9057484) (1.0786135,1.9008734) (1.0838495,1.8959928) (1.0890855,1.8911068) (1.0943214,1.8862152) (1.0995574,1.8813183) (1.1047934,1.876416) (1.1100294,1.8715084) (1.1152654,1.8665954) (1.1205014,1.8616773) (1.1257374,1.8567539) (1.1309734,1.8518254) (1.1362093,1.8468917) (1.1414453,1.8419529) (1.1466813,1.8370091) (1.1519173,1.8320603) (1.1571533,1.8271065) (1.1623893,1.8221477) (1.1676253,1.817184) (1.1728613,1.8122154) (1.1780972,1.8072419) (1.1833332,1.8022637) (1.1885692,1.7972806) (1.1938052,1.7922928) (1.1990412,1.7873002) (1.2042772,1.7823029) (1.2095132,1.7773009) (1.2147492,1.7722943) (1.2199851,1.7672831) (1.2252211,1.7622672) (1.2304571,1.7572468) (1.2356931,1.7522218) (1.2409291,1.7471923) (1.2461651,1.7421583) (1.2514011,1.7371198) (1.2566371,1.7320768) (1.261873,1.7270294) (1.267109,1.7219775) (1.272345,1.7169213) (1.277581,1.7118607) (1.282817,1.7067957) (1.288053,1.7017264) (1.293289,1.6966527) (1.298525,1.6915748) (1.303761,1.6864925) (1.3089969,1.681406) (1.3142329,1.6763152) (1.3194689,1.6712201) (1.3247049,1.6661208) (1.3299409,1.6610173) (1.3351769,1.6559096) (1.3404129,1.6507977) (1.3456489,1.6456816) (1.3508848,1.6405613) (1.3561208,1.6354369) (1.3613568,1.6303083) (1.3665928,1.6251756) (1.3718288,1.6200388) (1.3770648,1.6148978) (1.3823008,1.6097527) (1.3875368,1.6046035) (1.3927727,1.5994502) (1.3980087,1.5942928) (1.4032447,1.5891313) (1.4084807,1.5839657) (1.4137167,1.5787961) (1.4189527,1.5736223) (1.4241887,1.5684445) (1.4294247,1.5632627) (1.4346606,1.5580767) (1.4398966,1.5528867) (1.4451326,1.5476927) (1.4503686,1.5424946) (1.4556046,1.5372924) (1.4608406,1.5320862) (1.4660766,1.5268759) (1.4713126,1.5216615) (1.4765485,1.5164431) (1.4817845,1.5112206) (1.4870205,1.505994) (1.4922565,1.5007634) (1.4974925,1.4955287) (1.5027285,1.4902899) (1.5079645,1.485047) (1.5132005,1.4798) (1.5184364,1.4745489) (1.5236724,1.4692938) (1.5289084,1.4640345) (1.5341444,1.4587711) (1.5393804,1.4535035) (1.5446164,1.4482318) (1.5498524,1.442956) (1.5550884,1.437676) (1.5603244,1.4323918) (1.5655603,1.4271035) (1.5707963,1.421811) (1.5760323,1.4165142) (1.5812683,1.4112132) (1.5865043,1.405908) (1.5917403,1.4005986) (1.5969763,1.3952849) (1.6022123,1.3899669) (1.6074482,1.3846446) (1.6126842,1.3793179) (1.6179202,1.373987) (1.6231562,1.3686517) (1.6283922,1.363312) (1.6336282,1.3579679) (1.6388642,1.3526194) (1.6441002,1.3472665) (1.6493361,1.3419091) (1.6545721,1.3365472) (1.6598081,1.3311808) (1.6650441,1.3258099) (1.6702801,1.3204345) (1.6755161,1.3150544) (1.6807521,1.3096697) (1.6859881,1.3042804) (1.691224,1.2988864) (1.69646,1.2934877) (1.701696,1.2880843) (1.706932,1.2826762) (1.712168,1.2772632) (1.717404,1.2718454) (1.72264,1.2664227) (1.727876,1.2609952) (1.7331119,1.2555627) (1.7383479,1.2501252) (1.7435839,1.2446827) (1.7488199,1.2392352) (1.7540559,1.2337825) (1.7592919,1.2283248) (1.7645279,1.2228618) (1.7697639,1.2173937) (1.7749998,1.2119203) (1.7802358,1.2064415) (1.7854718,1.2009575) (1.7907078,1.195468) (1.7959438,1.189973) (1.8011798,1.1844725) (1.8064158,1.1789665) (1.8116518,1.1734549) (1.8168878,1.1679376) (1.8221237,1.1624145) (1.8273597,1.1568857) (1.8325957,1.1513511) (1.8378317,1.1458105) (1.8430677,1.140264) (1.8483037,1.1347114) (1.8535397,1.1291527) (1.8587757,1.1235879) (1.8640116,1.1180169) (1.8692476,1.1124395) (1.8744836,1.1068558) (1.8797196,1.1012656) (1.8849556,1.0956688) (1.8901916,1.0900655) (1.8954276,1.0844555) (1.9006636,1.0788387) (1.9058995,1.0732151) (1.9111355,1.0675845) (1.9163715,1.0619469) (1.9216075,1.0563022) (1.9268435,1.0506503) (1.9320795,1.0449911) (1.9373155,1.0393244) (1.9425515,1.0336503) (1.9477874,1.0279685) (1.9530234,1.022279) (1.9582594,1.0165818) (1.9634954,1.0108765) (1.9687314,1.0051632) (1.9739674,0.9994418) (1.9792034,0.99371206) (1.9844394,0.98797392) (1.9896753,0.98222723) (1.9949113,0.97647187) (2.0001473,0.97070771) (2.0053833,0.96493461) (2.0106193,0.95915242) (2.0158553,0.95336101) (2.0210913,0.94756022) (2.0263273,0.94174989) (2.0315632,0.93592987) (2.0367992,0.9301) (2.0420352,0.92426012) (2.0472712,0.91841003) (2.0525072,0.91254958) (2.0577432,0.90667858) (2.0629792,0.90079685) (2.0682152,0.89490418) (2.0734512,0.88900038) (2.0786871,0.88308525) (2.0839231,0.87715857) (2.0891591,0.87122013) (2.0943951,0.86526971) (2.0996311,0.85930708) (2.1048671,0.85333199) (2.1101031,0.8473442) (2.1153391,0.84134346) (2.120575,0.83532951) (2.125811,0.82930208) (2.131047,0.82326089) (2.136283,0.81720566) (2.141519,0.81113609) (2.146755,0.80505186) (2.151991,0.79895267) (2.157227,0.79283819) (2.1624629,0.78670808) (2.1676989,0.78056199) (2.1729349,0.77439955) (2.1781709,0.76822039) (2.1834069,0.76202412) (2.1886429,0.75581034) (2.1938789,0.74957863) (2.1991149,0.74332856) (2.2043508,0.73705967) (2.2095868,0.7307715) (2.2148228,0.72446357) (2.2200588,0.71813537) (2.2252948,0.71178638) (2.2305308,0.70541605) (2.2357668,0.69902383) (2.2410028,0.69260911) (2.2462387,0.68617129) (2.2514747,0.67970973) (2.2567107,0.67322376) (2.2619467,0.66671268) (2.2671827,0.66017578) (2.2724187,0.65361229) (2.2776547,0.64702142) (2.2828907,0.64040234) (2.2881266,0.63375418) (2.2933626,0.62707605) (2.2985986,0.62036699) (2.3038346,0.61362599) (2.3090706,0.60685202) (2.3143066,0.60004397) (2.3195426,0.59320069) (2.3247786,0.58632097) (2.3300146,0.57940351) (2.3352505,0.57244698) (2.3404865,0.56544994) (2.3457225,0.55841089) (2.3509585,0.55132824) (2.3561945,0.5442003) (2.3614305,0.53702529) (2.3666665,0.52980132) (2.3719025,0.52252636) (2.3771384,0.51519828) (2.3823744,0.5078148) (2.3876104,0.50037348) (2.3928464,0.49287173) (2.3980824,0.48530674) (2.4033184,0.47767555) (2.4085544,0.46997493) (2.4137904,0.46220145) (2.4190263,0.45435138) (2.4242623,0.44642069) (2.4294983,0.43840503) (2.4347343,0.43029968) (2.4399703,0.42209948) (2.4452063,0.41379881) (2.4504423,0.40539153) (2.4556783,0.39687087) (2.4609142,0.38822939) (2.4661502,0.37945885) (2.4713862,0.37055012) (2.4766222,0.36149299) (2.4818582,0.35227605) (2.4870942,0.34288642) (2.4923302,0.33330953) (2.4975662,0.32352878) (2.5028021,0.31352511) (2.5080381,0.30327648) (2.5132741,0.29275719) (2.5185101,0.28193697) (2.5237461,0.27077978) (2.5289821,0.25924214) (2.5342181,0.24727081) (2.5394541,0.23479951) (2.54469,0.22174392) (2.549926,0.20799408) (2.555162,0.193402) (2.560398,0.17776042) (2.565634,0.16076333) (2.57087,0.1419246) (2.576106,0.12038277) (2.581342,0.094303316) (2.586578,0.057843896) (2.586578,-0.057843896) (2.581342,-0.094303316) (2.576106,-0.12038277) (2.57087,-0.1419246) (2.565634,-0.16076333) (2.560398,-0.17776042) (2.555162,-0.193402) (2.549926,-0.20799408) (2.54469,-0.22174392) (2.5394541,-0.23479951) (2.5342181,-0.24727081) (2.5289821,-0.25924214) (2.5237461,-0.27077978) (2.5185101,-0.28193697) (2.5132741,-0.29275719) (2.5080381,-0.30327648) (2.5028021,-0.31352511) (2.4975662,-0.32352878) (2.4923302,-0.33330953) (2.4870942,-0.34288642) (2.4818582,-0.35227605) (2.4766222,-0.36149299) (2.4713862,-0.37055012) (2.4661502,-0.37945885) (2.4609142,-0.38822939) (2.4556783,-0.39687087) (2.4504423,-0.40539153) (2.4452063,-0.41379881) (2.4399703,-0.42209948) (2.4347343,-0.43029968) (2.4294983,-0.43840503) (2.4242623,-0.44642069) (2.4190263,-0.45435138) (2.4137904,-0.46220145) (2.4085544,-0.46997493) (2.4033184,-0.47767555) (2.3980824,-0.48530674) (2.3928464,-0.49287173) (2.3876104,-0.50037348) (2.3823744,-0.5078148) (2.3771384,-0.51519828) (2.3719025,-0.52252636) (2.3666665,-0.52980132) (2.3614305,-0.53702529) (2.3561945,-0.5442003) (2.3509585,-0.55132824) (2.3457225,-0.55841089) (2.3404865,-0.56544994) (2.3352505,-0.57244698) (2.3300146,-0.57940351) (2.3247786,-0.58632097) (2.3195426,-0.59320069) (2.3143066,-0.60004397) (2.3090706,-0.60685202) (2.3038346,-0.61362599) (2.2985986,-0.62036699) (2.2933626,-0.62707605) (2.2881266,-0.63375418) (2.2828907,-0.64040234) (2.2776547,-0.64702142) (2.2724187,-0.65361229) (2.2671827,-0.66017578) (2.2619467,-0.66671268) (2.2567107,-0.67322376) (2.2514747,-0.67970973) (2.2462387,-0.68617129) (2.2410028,-0.69260911) (2.2357668,-0.69902383) (2.2305308,-0.70541605) (2.2252948,-0.71178638) (2.2200588,-0.71813537) (2.2148228,-0.72446357) (2.2095868,-0.7307715) (2.2043508,-0.73705967) (2.1991149,-0.74332856) (2.1938789,-0.74957863) (2.1886429,-0.75581034) (2.1834069,-0.76202412) (2.1781709,-0.76822039) (2.1729349,-0.77439955) (2.1676989,-0.78056199) (2.1624629,-0.78670808) (2.157227,-0.79283819) (2.151991,-0.79895267) (2.146755,-0.80505186) (2.141519,-0.81113609) (2.136283,-0.81720566) (2.131047,-0.82326089) (2.125811,-0.82930208) (2.120575,-0.83532951) (2.1153391,-0.84134346) (2.1101031,-0.8473442) (2.1048671,-0.85333199) (2.0996311,-0.85930708) (2.0943951,-0.86526971) (2.0891591,-0.87122013) (2.0839231,-0.87715857) (2.0786871,-0.88308525) (2.0734512,-0.88900038) (2.0682152,-0.89490418) (2.0629792,-0.90079685) (2.0577432,-0.90667858) (2.0525072,-0.91254958) (2.0472712,-0.91841003) (2.0420352,-0.92426012) (2.0367992,-0.9301) (2.0315632,-0.93592987) (2.0263273,-0.94174989) (2.0210913,-0.94756022) (2.0158553,-0.95336101) (2.0106193,-0.95915242) (2.0053833,-0.96493461) (2.0001473,-0.97070771) (1.9949113,-0.97647187) (1.9896753,-0.98222723) (1.9844394,-0.98797392) (1.9792034,-0.99371206) (1.9739674,-0.9994418) (1.9687314,-1.0051632) (1.9634954,-1.0108765) (1.9582594,-1.0165818) (1.9530234,-1.022279) (1.9477874,-1.0279685) (1.9425515,-1.0336503) (1.9373155,-1.0393244) (1.9320795,-1.0449911) (1.9268435,-1.0506503) (1.9216075,-1.0563022) (1.9163715,-1.0619469) (1.9111355,-1.0675845) (1.9058995,-1.0732151) (1.9006636,-1.0788387) (1.8954276,-1.0844555) (1.8901916,-1.0900655) (1.8849556,-1.0956688) (1.8797196,-1.1012656) (1.8744836,-1.1068558) (1.8692476,-1.1124395) (1.8640116,-1.1180169) (1.8587757,-1.1235879) (1.8535397,-1.1291527) (1.8483037,-1.1347114) (1.8430677,-1.140264) (1.8378317,-1.1458105) (1.8325957,-1.1513511) (1.8273597,-1.1568857) (1.8221237,-1.1624145) (1.8168878,-1.1679376) (1.8116518,-1.1734549) (1.8064158,-1.1789665) (1.8011798,-1.1844725) (1.7959438,-1.189973) (1.7907078,-1.195468) (1.7854718,-1.2009575) (1.7802358,-1.2064415) (1.7749998,-1.2119203) (1.7697639,-1.2173937) (1.7645279,-1.2228618) (1.7592919,-1.2283248) (1.7540559,-1.2337825) (1.7488199,-1.2392352) (1.7435839,-1.2446827) (1.7383479,-1.2501252) (1.7331119,-1.2555627) (1.727876,-1.2609952) (1.72264,-1.2664227) (1.717404,-1.2718454) (1.712168,-1.2772632) (1.706932,-1.2826762) (1.701696,-1.2880843) (1.69646,-1.2934877) (1.691224,-1.2988864) (1.6859881,-1.3042804) (1.6807521,-1.3096697) (1.6755161,-1.3150544) (1.6702801,-1.3204345) (1.6650441,-1.3258099) (1.6598081,-1.3311808) (1.6545721,-1.3365472) (1.6493361,-1.3419091) (1.6441002,-1.3472665) (1.6388642,-1.3526194) (1.6336282,-1.3579679) (1.6283922,-1.363312) (1.6231562,-1.3686517) (1.6179202,-1.373987) (1.6126842,-1.3793179) (1.6074482,-1.3846446) (1.6022123,-1.3899669) (1.5969763,-1.3952849) (1.5917403,-1.4005986) (1.5865043,-1.405908) (1.5812683,-1.4112132) (1.5760323,-1.4165142) (1.5707963,-1.421811) (1.5655603,-1.4271035) (1.5603244,-1.4323918) (1.5550884,-1.437676) (1.5498524,-1.442956) (1.5446164,-1.4482318) (1.5393804,-1.4535035) (1.5341444,-1.4587711) (1.5289084,-1.4640345) (1.5236724,-1.4692938) (1.5184364,-1.4745489) (1.5132005,-1.4798) (1.5079645,-1.485047) (1.5027285,-1.4902899) (1.4974925,-1.4955287) (1.4922565,-1.5007634) (1.4870205,-1.505994) (1.4817845,-1.5112206) (1.4765485,-1.5164431) (1.4713126,-1.5216615) (1.4660766,-1.5268759) (1.4608406,-1.5320862) (1.4556046,-1.5372924) (1.4503686,-1.5424946) (1.4451326,-1.5476927) (1.4398966,-1.5528867) (1.4346606,-1.5580767) (1.4294247,-1.5632627) (1.4241887,-1.5684445) (1.4189527,-1.5736223) (1.4137167,-1.5787961) (1.4084807,-1.5839657) (1.4032447,-1.5891313) (1.3980087,-1.5942928) (1.3927727,-1.5994502) (1.3875368,-1.6046035) (1.3823008,-1.6097527) (1.3770648,-1.6148978) (1.3718288,-1.6200388) (1.3665928,-1.6251756) (1.3613568,-1.6303083) (1.3561208,-1.6354369) (1.3508848,-1.6405613) (1.3456489,-1.6456816) (1.3404129,-1.6507977) (1.3351769,-1.6559096) (1.3299409,-1.6610173) (1.3247049,-1.6661208) (1.3194689,-1.6712201) (1.3142329,-1.6763152) (1.3089969,-1.681406) (1.303761,-1.6864925) (1.298525,-1.6915748) (1.293289,-1.6966527) (1.288053,-1.7017264) (1.282817,-1.7067957) (1.277581,-1.7118607) (1.272345,-1.7169213) (1.267109,-1.7219775) (1.261873,-1.7270294) (1.2566371,-1.7320768) (1.2514011,-1.7371198) (1.2461651,-1.7421583) (1.2409291,-1.7471923) (1.2356931,-1.7522218) (1.2304571,-1.7572468) (1.2252211,-1.7622672) (1.2199851,-1.7672831) (1.2147492,-1.7722943) (1.2095132,-1.7773009) (1.2042772,-1.7823029) (1.1990412,-1.7873002) (1.1938052,-1.7922928) (1.1885692,-1.7972806) (1.1833332,-1.8022637) (1.1780972,-1.8072419) (1.1728613,-1.8122154) (1.1676253,-1.817184) (1.1623893,-1.8221477) (1.1571533,-1.8271065) (1.1519173,-1.8320603) (1.1466813,-1.8370091) (1.1414453,-1.8419529) (1.1362093,-1.8468917) (1.1309734,-1.8518254) (1.1257374,-1.8567539) (1.1205014,-1.8616773) (1.1152654,-1.8665954) (1.1100294,-1.8715084) (1.1047934,-1.876416) (1.0995574,-1.8813183) (1.0943214,-1.8862152) (1.0890855,-1.8911068) (1.0838495,-1.8959928) (1.0786135,-1.9008734) (1.0733775,-1.9057484) (1.0681415,-1.9106179) (1.0629055,-1.9154817) (1.0576695,-1.9203398) (1.0524335,-1.9251922) (1.0471976,-1.9300388) (1.0419616,-1.9348796) (1.0367256,-1.9397145) (1.0314896,-1.9445435) (1.0262536,-1.9493664) (1.0210176,-1.9541834) (1.0157816,-1.9589942) (1.0105456,-1.9637989) (1.0053096,-1.9685973) (1.0000737,-1.9733895) (0.99483767,-1.9781754) (0.98960169,-1.9829548) (0.9843657,-1.9877278) (0.97912971,-1.9924943) (0.97389372,-1.9972542) (0.96865773,-2.0020075) (0.96342175,-2.0067541) (0.95818576,-2.0114938) (0.95294977,-2.0162268) (0.94771378,-2.0209528) (0.9424778,-2.0256718) (0.93724181,-2.0303837) (0.93200582,-2.0350885) (0.92676983,-2.0397861) (0.92153385,-2.0444765) (0.91629786,-2.0491594) (0.91106187,-2.0538349) (0.90582588,-2.0585029) (0.90058989,-2.0631632) (0.89535391,-2.0678159) (0.89011792,-2.0724608) (0.88488193,-2.0770978) (0.87964594,-2.0817269) (0.87440996,-2.0863479) (0.86917397,-2.0909608) (0.86393798,-2.0955655) (0.85870199,-2.1001618) (0.853466,-2.1047498) (0.84823002,-2.1093292) (0.84299403,-2.1139) (0.83775804,-2.1184621) (0.83252205,-2.1230153) (0.82728607,-2.1275597) (0.82205008,-2.132095) (0.81681409,-2.1366212) (0.8115781,-2.1411381) (0.80634211,-2.1456457) (0.80110613,-2.1501438) (0.79587014,-2.1546323) (0.79063415,-2.1591111) (0.78539816,-2.16358) (0.78016218,-2.1680391) (0.77492619,-2.172488) (0.7696902,-2.1769267) (0.76445421,-2.1813552) (0.75921822,-2.1857731) (0.75398224,-2.1901805) (0.74874625,-2.1945771) (0.74351026,-2.1989629) (0.73827427,-2.2033376) (0.73303829,-2.2077012) (0.7278023,-2.2120535) (0.72256631,-2.2163944) (0.71733032,-2.2207237) (0.71209433,-2.2250412) (0.70685835,-2.2293469) (0.70162236,-2.2336405) (0.69638637,-2.2379218) (0.69115038,-2.2421908) (0.6859144,-2.2464473) (0.68067841,-2.250691) (0.67544242,-2.2549219) (0.67020643,-2.2591397) (0.66497045,-2.2633443) (0.65973446,-2.2675354) (0.65449847,-2.271713) (0.64926248,-2.2758768) (0.64402649,-2.2800266) (0.63879051,-2.2841622) (0.63355452,-2.2882835) (0.62831853,-2.2923903) (0.62308254,-2.2964823) (0.61784656,-2.3005593) (0.61261057,-2.3046211) (0.60737458,-2.3086676) (0.60213859,-2.3126985) (0.5969026,-2.3167135) (0.59166662,-2.3207126) (0.58643063,-2.3246953) (0.58119464,-2.3286616) (0.57595865,-2.3326111) (0.57072267,-2.3365437) (0.56548668,-2.3404591) (0.56025069,-2.3443571) (0.5550147,-2.3482373) (0.54977871,-2.3520996) (0.54454273,-2.3559437) (0.53930674,-2.3597694) (0.53407075,-2.3635763) (0.52883476,-2.3673643) (0.52359878,-2.371133) (0.51836279,-2.3748821) (0.5131268,-2.3786115) (0.50789081,-2.3823207) (0.50265482,-2.3860096) (0.49741884,-2.3896778) (0.49218285,-2.393325) (0.48694686,-2.396951) (0.48171087,-2.4005554) (0.47647489,-2.4041379) (0.4712389,-2.4076983) (0.46600291,-2.4112361) (0.46076692,-2.4147511) (0.45553093,-2.4182429) (0.45029495,-2.4217113) (0.44505896,-2.4251559) (0.43982297,-2.4285763) (0.43458698,-2.4319722) (0.429351,-2.4353432) (0.42411501,-2.4386891) (0.41887902,-2.4420095) (0.41364303,-2.4453039) (0.40840704,-2.4485721) (0.40317106,-2.4518137) (0.39793507,-2.4550282) (0.39269908,-2.4582154) (0.38746309,-2.4613749) (0.38222711,-2.4645062) (0.37699112,-2.467609) (0.37175513,-2.4706829) (0.36651914,-2.4737275) (0.36128316,-2.4767425) (0.35604717,-2.4797273) (0.35081118,-2.4826816) (0.34557519,-2.4856051) (0.3403392,-2.4884972) (0.33510322,-2.4913576) (0.32986723,-2.4941859) (0.32463124,-2.4969817) (0.31939525,-2.4997445) (0.31415927,-2.5024739) (0.30892328,-2.5051694) (0.30368729,-2.5078308) (0.2984513,-2.5104575) (0.29321531,-2.5130491) (0.28797933,-2.5156051) (0.28274334,-2.5181252) (0.27750735,-2.5206089) (0.27227136,-2.5230558) (0.26703538,-2.5254655) (0.26179939,-2.5278375) (0.2565634,-2.5301713) (0.25132741,-2.5324666) (0.24609142,-2.5347229) (0.24085544,-2.5369398) (0.23561945,-2.5391168) (0.23038346,-2.5412535) (0.22514747,-2.5433496) (0.21991149,-2.5454045) (0.2146755,-2.5474179) (0.20943951,-2.5493893) (0.20420352,-2.5513183) (0.19896753,-2.5532045) (0.19373155,-2.5550475) (0.18849556,-2.5568469) (0.18325957,-2.5586022) (0.17802358,-2.5603132) (0.1727876,-2.5619794) (0.16755161,-2.5636003) (0.16231562,-2.5651758) (0.15707963,-2.5667052) (0.15184364,-2.5681884) (0.14660766,-2.569625) (0.14137167,-2.5710145) (0.13613568,-2.5723567) (0.13089969,-2.5736512) (0.12566371,-2.5748977) (0.12042772,-2.5760959) (0.11519173,-2.5772455) (0.10995574,-2.5783461) (0.10471976,-2.5793976) (0.099483767,-2.5803996) (0.09424778,-2.5813518) (0.089011792,-2.582254) (0.083775804,-2.5831061) (0.078539816,-2.5839076) (0.073303829,-2.5846585) (0.068067841,-2.5853585) (0.062831853,-2.5860074) (0.057595865,-2.5866051) (0.052359878,-2.5871514) (0.04712389,-2.5876461) (0.041887902,-2.5880891) (0.036651914,-2.5884802) (0.031415927,-2.5888195) (0.026179939,-2.5891067) (0.020943951,-2.5893418) (0.015707963,-2.5895247) (0.010471976,-2.5896554) (0.0052359878,-2.5897339) (-4.4408921e-16,-2.58976) (-0.0052359878,-2.5897339) (-0.010471976,-2.5896554) (-0.015707963,-2.5895247) (-0.020943951,-2.5893418) (-0.026179939,-2.5891067) (-0.031415927,-2.5888195) (-0.036651914,-2.5884802) (-0.041887902,-2.5880891) (-0.04712389,-2.5876461) (-0.052359878,-2.5871514) (-0.057595865,-2.5866051) (-0.062831853,-2.5860074) (-0.068067841,-2.5853585) (-0.073303829,-2.5846585) (-0.078539816,-2.5839076) (-0.083775804,-2.5831061) (-0.089011792,-2.582254) (-0.09424778,-2.5813518) (-0.099483767,-2.5803996) (-0.10471976,-2.5793976) (-0.10995574,-2.5783461) (-0.11519173,-2.5772455) (-0.12042772,-2.5760959) (-0.12566371,-2.5748977) (-0.13089969,-2.5736512) (-0.13613568,-2.5723567) (-0.14137167,-2.5710145) (-0.14660766,-2.569625) (-0.15184364,-2.5681884) (-0.15707963,-2.5667052) (-0.16231562,-2.5651758) (-0.16755161,-2.5636003) (-0.1727876,-2.5619794) (-0.17802358,-2.5603132) (-0.18325957,-2.5586022) (-0.18849556,-2.5568469) (-0.19373155,-2.5550475) (-0.19896753,-2.5532045) (-0.20420352,-2.5513183) (-0.20943951,-2.5493893) (-0.2146755,-2.5474179) (-0.21991149,-2.5454045) (-0.22514747,-2.5433496) (-0.23038346,-2.5412535) (-0.23561945,-2.5391168) (-0.24085544,-2.5369398) (-0.24609142,-2.5347229) (-0.25132741,-2.5324666) (-0.2565634,-2.5301713) (-0.26179939,-2.5278375) (-0.26703538,-2.5254655) (-0.27227136,-2.5230558) (-0.27750735,-2.5206089) (-0.28274334,-2.5181252) (-0.28797933,-2.5156051) (-0.29321531,-2.5130491) (-0.2984513,-2.5104575) (-0.30368729,-2.5078308) (-0.30892328,-2.5051694) (-0.31415927,-2.5024739) (-0.31939525,-2.4997445) (-0.32463124,-2.4969817) (-0.32986723,-2.4941859) (-0.33510322,-2.4913576) (-0.3403392,-2.4884972) (-0.34557519,-2.4856051) (-0.35081118,-2.4826816) (-0.35604717,-2.4797273) (-0.36128316,-2.4767425) (-0.36651914,-2.4737275) (-0.37175513,-2.4706829) (-0.37699112,-2.467609) (-0.38222711,-2.4645062) (-0.38746309,-2.4613749) (-0.39269908,-2.4582154) (-0.39793507,-2.4550282) (-0.40317106,-2.4518137) (-0.40840704,-2.4485721) (-0.41364303,-2.4453039) (-0.41887902,-2.4420095) (-0.42411501,-2.4386891) (-0.429351,-2.4353432) (-0.43458698,-2.4319722) (-0.43982297,-2.4285763) (-0.44505896,-2.4251559) (-0.45029495,-2.4217113) (-0.45553093,-2.4182429) (-0.46076692,-2.4147511) (-0.46600291,-2.4112361) (-0.4712389,-2.4076983) (-0.47647489,-2.4041379) (-0.48171087,-2.4005554) (-0.48694686,-2.396951) (-0.49218285,-2.393325) (-0.49741884,-2.3896778) (-0.50265482,-2.3860096) (-0.50789081,-2.3823207) (-0.5131268,-2.3786115) (-0.51836279,-2.3748821) (-0.52359878,-2.371133) (-0.52883476,-2.3673643) (-0.53407075,-2.3635763) (-0.53930674,-2.3597694) (-0.54454273,-2.3559437) (-0.54977871,-2.3520996) (-0.5550147,-2.3482373) (-0.56025069,-2.3443571) (-0.56548668,-2.3404591) (-0.57072267,-2.3365437) (-0.57595865,-2.3326111) (-0.58119464,-2.3286616) (-0.58643063,-2.3246953) (-0.59166662,-2.3207126) (-0.5969026,-2.3167135) (-0.60213859,-2.3126985) (-0.60737458,-2.3086676) (-0.61261057,-2.3046211) (-0.61784656,-2.3005593) (-0.62308254,-2.2964823) (-0.62831853,-2.2923903) (-0.63355452,-2.2882835) (-0.63879051,-2.2841622) (-0.64402649,-2.2800266) (-0.64926248,-2.2758768) (-0.65449847,-2.271713) (-0.65973446,-2.2675354) (-0.66497045,-2.2633443) (-0.67020643,-2.2591397) (-0.67544242,-2.2549219) (-0.68067841,-2.250691) (-0.6859144,-2.2464473) (-0.69115038,-2.2421908) (-0.69638637,-2.2379218) (-0.70162236,-2.2336405) (-0.70685835,-2.2293469) (-0.71209433,-2.2250412) (-0.71733032,-2.2207237) (-0.72256631,-2.2163944) (-0.7278023,-2.2120535) (-0.73303829,-2.2077012) (-0.73827427,-2.2033376) (-0.74351026,-2.1989629) (-0.74874625,-2.1945771) (-0.75398224,-2.1901805) (-0.75921822,-2.1857731) (-0.76445421,-2.1813552) (-0.7696902,-2.1769267) (-0.77492619,-2.172488) (-0.78016218,-2.1680391) (-0.78539816,-2.16358) (-0.79063415,-2.1591111) (-0.79587014,-2.1546323) (-0.80110613,-2.1501438) (-0.80634211,-2.1456457) (-0.8115781,-2.1411381) (-0.81681409,-2.1366212) (-0.82205008,-2.132095) (-0.82728607,-2.1275597) (-0.83252205,-2.1230153) (-0.83775804,-2.1184621) (-0.84299403,-2.1139) (-0.84823002,-2.1093292) (-0.853466,-2.1047498) (-0.85870199,-2.1001618) (-0.86393798,-2.0955655) (-0.86917397,-2.0909608) (-0.87440996,-2.0863479) (-0.87964594,-2.0817269) (-0.88488193,-2.0770978) (-0.89011792,-2.0724608) (-0.89535391,-2.0678159) (-0.90058989,-2.0631632) (-0.90582588,-2.0585029) (-0.91106187,-2.0538349) (-0.91629786,-2.0491594) (-0.92153385,-2.0444765) (-0.92676983,-2.0397861) (-0.93200582,-2.0350885) (-0.93724181,-2.0303837) (-0.9424778,-2.0256718) (-0.94771378,-2.0209528) (-0.95294977,-2.0162268) (-0.95818576,-2.0114938) (-0.96342175,-2.0067541) (-0.96865773,-2.0020075) (-0.97389372,-1.9972542) (-0.97912971,-1.9924943) (-0.9843657,-1.9877278) (-0.98960169,-1.9829548) (-0.99483767,-1.9781754) (-1.0000737,-1.9733895) (-1.0053096,-1.9685973) (-1.0105456,-1.9637989) (-1.0157816,-1.9589942) (-1.0210176,-1.9541834) (-1.0262536,-1.9493664) (-1.0314896,-1.9445435) (-1.0367256,-1.9397145) (-1.0419616,-1.9348796) (-1.0471976,-1.9300388) (-1.0524335,-1.9251922) (-1.0576695,-1.9203398) (-1.0629055,-1.9154817) (-1.0681415,-1.9106179) (-1.0733775,-1.9057484) (-1.0786135,-1.9008734) (-1.0838495,-1.8959928) (-1.0890855,-1.8911068) (-1.0943214,-1.8862152) (-1.0995574,-1.8813183) (-1.1047934,-1.876416) (-1.1100294,-1.8715084) (-1.1152654,-1.8665954) (-1.1205014,-1.8616773) (-1.1257374,-1.8567539) (-1.1309734,-1.8518254) (-1.1362093,-1.8468917) (-1.1414453,-1.8419529) (-1.1466813,-1.8370091) (-1.1519173,-1.8320603) (-1.1571533,-1.8271065) (-1.1623893,-1.8221477) (-1.1676253,-1.817184) (-1.1728613,-1.8122154) (-1.1780972,-1.8072419) (-1.1833332,-1.8022637) (-1.1885692,-1.7972806) (-1.1938052,-1.7922928) (-1.1990412,-1.7873002) (-1.2042772,-1.7823029) (-1.2095132,-1.7773009) (-1.2147492,-1.7722943) (-1.2199851,-1.7672831) (-1.2252211,-1.7622672) (-1.2304571,-1.7572468) (-1.2356931,-1.7522218) (-1.2409291,-1.7471923) (-1.2461651,-1.7421583) (-1.2514011,-1.7371198) (-1.2566371,-1.7320768) (-1.261873,-1.7270294) (-1.267109,-1.7219775) (-1.272345,-1.7169213) (-1.277581,-1.7118607) (-1.282817,-1.7067957) (-1.288053,-1.7017264) (-1.293289,-1.6966527) (-1.298525,-1.6915748) (-1.303761,-1.6864925) (-1.3089969,-1.681406) (-1.3142329,-1.6763152) (-1.3194689,-1.6712201) (-1.3247049,-1.6661208) (-1.3299409,-1.6610173) (-1.3351769,-1.6559096) (-1.3404129,-1.6507977) (-1.3456489,-1.6456816) (-1.3508848,-1.6405613) (-1.3561208,-1.6354369) (-1.3613568,-1.6303083) (-1.3665928,-1.6251756) (-1.3718288,-1.6200388) (-1.3770648,-1.6148978) (-1.3823008,-1.6097527) (-1.3875368,-1.6046035) (-1.3927727,-1.5994502) (-1.3980087,-1.5942928) (-1.4032447,-1.5891313) (-1.4084807,-1.5839657) (-1.4137167,-1.5787961) (-1.4189527,-1.5736223) (-1.4241887,-1.5684445) (-1.4294247,-1.5632627) (-1.4346606,-1.5580767) (-1.4398966,-1.5528867) (-1.4451326,-1.5476927) (-1.4503686,-1.5424946) (-1.4556046,-1.5372924) (-1.4608406,-1.5320862) (-1.4660766,-1.5268759) (-1.4713126,-1.5216615) (-1.4765485,-1.5164431) (-1.4817845,-1.5112206) (-1.4870205,-1.505994) (-1.4922565,-1.5007634) (-1.4974925,-1.4955287) (-1.5027285,-1.4902899) (-1.5079645,-1.485047) (-1.5132005,-1.4798) (-1.5184364,-1.4745489) (-1.5236724,-1.4692938) (-1.5289084,-1.4640345) (-1.5341444,-1.4587711) (-1.5393804,-1.4535035) (-1.5446164,-1.4482318) (-1.5498524,-1.442956) (-1.5550884,-1.437676) (-1.5603244,-1.4323918) (-1.5655603,-1.4271035) (-1.5707963,-1.421811) (-1.5760323,-1.4165142) (-1.5812683,-1.4112132) (-1.5865043,-1.405908) (-1.5917403,-1.4005986) (-1.5969763,-1.3952849) (-1.6022123,-1.3899669) (-1.6074482,-1.3846446) (-1.6126842,-1.3793179) (-1.6179202,-1.373987) (-1.6231562,-1.3686517) (-1.6283922,-1.363312) (-1.6336282,-1.3579679) (-1.6388642,-1.3526194) (-1.6441002,-1.3472665) (-1.6493361,-1.3419091) (-1.6545721,-1.3365472) (-1.6598081,-1.3311808) (-1.6650441,-1.3258099) (-1.6702801,-1.3204345) (-1.6755161,-1.3150544) (-1.6807521,-1.3096697) (-1.6859881,-1.3042804) (-1.691224,-1.2988864) (-1.69646,-1.2934877) (-1.701696,-1.2880843) (-1.706932,-1.2826762) (-1.712168,-1.2772632) (-1.717404,-1.2718454) (-1.72264,-1.2664227) (-1.727876,-1.2609952) (-1.7331119,-1.2555627) (-1.7383479,-1.2501252) (-1.7435839,-1.2446827) (-1.7488199,-1.2392352) (-1.7540559,-1.2337825) (-1.7592919,-1.2283248) (-1.7645279,-1.2228618) (-1.7697639,-1.2173937) (-1.7749998,-1.2119203) (-1.7802358,-1.2064415) (-1.7854718,-1.2009575) (-1.7907078,-1.195468) (-1.7959438,-1.189973) (-1.8011798,-1.1844725) (-1.8064158,-1.1789665) (-1.8116518,-1.1734549) (-1.8168878,-1.1679376) (-1.8221237,-1.1624145) (-1.8273597,-1.1568857) (-1.8325957,-1.1513511) (-1.8378317,-1.1458105) (-1.8430677,-1.140264) (-1.8483037,-1.1347114) (-1.8535397,-1.1291527) (-1.8587757,-1.1235879) (-1.8640116,-1.1180169) (-1.8692476,-1.1124395) (-1.8744836,-1.1068558) (-1.8797196,-1.1012656) (-1.8849556,-1.0956688) (-1.8901916,-1.0900655) (-1.8954276,-1.0844555) (-1.9006636,-1.0788387) (-1.9058995,-1.0732151) (-1.9111355,-1.0675845) (-1.9163715,-1.0619469) (-1.9216075,-1.0563022) (-1.9268435,-1.0506503) (-1.9320795,-1.0449911) (-1.9373155,-1.0393244) (-1.9425515,-1.0336503) (-1.9477874,-1.0279685) (-1.9530234,-1.022279) (-1.9582594,-1.0165818) (-1.9634954,-1.0108765) (-1.9687314,-1.0051632) (-1.9739674,-0.9994418) (-1.9792034,-0.99371206) (-1.9844394,-0.98797392) (-1.9896753,-0.98222723) (-1.9949113,-0.97647187) (-2.0001473,-0.97070771) (-2.0053833,-0.96493461) (-2.0106193,-0.95915242) (-2.0158553,-0.95336101) (-2.0210913,-0.94756022) (-2.0263273,-0.94174989) (-2.0315632,-0.93592987) (-2.0367992,-0.9301) (-2.0420352,-0.92426012) (-2.0472712,-0.91841003) (-2.0525072,-0.91254958) (-2.0577432,-0.90667858) (-2.0629792,-0.90079685) (-2.0682152,-0.89490418) (-2.0734512,-0.88900038) (-2.0786871,-0.88308525) (-2.0839231,-0.87715857) (-2.0891591,-0.87122013) (-2.0943951,-0.86526971) (-2.0996311,-0.85930708) (-2.1048671,-0.85333199) (-2.1101031,-0.8473442) (-2.1153391,-0.84134346) (-2.120575,-0.83532951) (-2.125811,-0.82930208) (-2.131047,-0.82326089) (-2.136283,-0.81720566) (-2.141519,-0.81113609) (-2.146755,-0.80505186) (-2.151991,-0.79895267) (-2.157227,-0.79283819) (-2.1624629,-0.78670808) (-2.1676989,-0.78056199) (-2.1729349,-0.77439955) (-2.1781709,-0.76822039) (-2.1834069,-0.76202412) (-2.1886429,-0.75581034) (-2.1938789,-0.74957863) (-2.1991149,-0.74332856) (-2.2043508,-0.73705967) (-2.2095868,-0.7307715) (-2.2148228,-0.72446357) (-2.2200588,-0.71813537) (-2.2252948,-0.71178638) (-2.2305308,-0.70541605) (-2.2357668,-0.69902383) (-2.2410028,-0.69260911) (-2.2462387,-0.68617129) (-2.2514747,-0.67970973) (-2.2567107,-0.67322376) (-2.2619467,-0.66671268) (-2.2671827,-0.66017578) (-2.2724187,-0.65361229) (-2.2776547,-0.64702142) (-2.2828907,-0.64040234) (-2.2881266,-0.63375418) (-2.2933626,-0.62707605) (-2.2985986,-0.62036699) (-2.3038346,-0.61362599) (-2.3090706,-0.60685202) (-2.3143066,-0.60004397) (-2.3195426,-0.59320069) (-2.3247786,-0.58632097) (-2.3300146,-0.57940351) (-2.3352505,-0.57244698) (-2.3404865,-0.56544994) (-2.3457225,-0.55841089) (-2.3509585,-0.55132824) (-2.3561945,-0.5442003) (-2.3614305,-0.53702529) (-2.3666665,-0.52980132) (-2.3719025,-0.52252636) (-2.3771384,-0.51519828) (-2.3823744,-0.5078148) (-2.3876104,-0.50037348) (-2.3928464,-0.49287173) (-2.3980824,-0.48530674) (-2.4033184,-0.47767555) (-2.4085544,-0.46997493) (-2.4137904,-0.46220145) (-2.4190263,-0.45435138) (-2.4242623,-0.44642069) (-2.4294983,-0.43840503) (-2.4347343,-0.43029968) (-2.4399703,-0.42209948) (-2.4452063,-0.41379881) (-2.4504423,-0.40539153) (-2.4556783,-0.39687087) (-2.4609142,-0.38822939) (-2.4661502,-0.37945885) (-2.4713862,-0.37055012) (-2.4766222,-0.36149299) (-2.4818582,-0.35227605) (-2.4870942,-0.34288642) (-2.4923302,-0.33330953) (-2.4975662,-0.32352878) (-2.5028021,-0.31352511) (-2.5080381,-0.30327648) (-2.5132741,-0.29275719) (-2.5185101,-0.28193697) (-2.5237461,-0.27077978) (-2.5289821,-0.25924214) (-2.5342181,-0.24727081) (-2.5394541,-0.23479951) (-2.54469,-0.22174392) (-2.549926,-0.20799408) (-2.555162,-0.193402) (-2.560398,-0.17776042) (-2.565634,-0.16076333) (-2.57087,-0.1419246) (-2.576106,-0.12038277) (-2.581342,-0.094303316) (-2.586578,-0.057843896)};
\addlegendentry{$p=0.15$}
\addplot[thick,color=blue] coordinates {(-2.2619467,0.060750487) (-2.2567107,0.10856931) (-2.2514747,0.14117281) (-2.2462387,0.16768416) (-2.2410028,0.19066173) (-2.2357668,0.21126161) (-2.2305308,0.23012311) (-2.2252948,0.24764367) (-2.2200588,0.26409022) (-2.2148228,0.27965222) (-2.2095868,0.29446986) (-2.2043508,0.30865028) (-2.1991149,0.32227754) (-2.1938789,0.33541898) (-2.1886429,0.34812956) (-2.1834069,0.36045478) (-2.1781709,0.37243284) (-2.1729349,0.38409615) (-2.1676989,0.39547249) (-2.1624629,0.40658588) (-2.157227,0.41745726) (-2.151991,0.42810501) (-2.146755,0.43854534) (-2.141519,0.44879268) (-2.136283,0.4588599) (-2.131047,0.46875853) (-2.125811,0.47849898) (-2.120575,0.48809067) (-2.1153391,0.49754213) (-2.1101031,0.50686115) (-2.1048671,0.51605486) (-2.0996311,0.52512977) (-2.0943951,0.53409189) (-2.0891591,0.54294674) (-2.0839231,0.55169943) (-2.0786871,0.5603547) (-2.0734512,0.56891693) (-2.0682152,0.57739021) (-2.0629792,0.58577835) (-2.0577432,0.59408489) (-2.0525072,0.60231316) (-2.0472712,0.61046626) (-2.0420352,0.61854713) (-2.0367992,0.62655849) (-2.0315632,0.63450293) (-2.0263273,0.64238288) (-2.0210913,0.65020063) (-2.0158553,0.65795833) (-2.0106193,0.66565805) (-2.0053833,0.67330171) (-2.0001473,0.68089115) (-1.9949113,0.68842811) (-1.9896753,0.69591425) (-1.9844394,0.70335112) (-1.9792034,0.71074024) (-1.9739674,0.71808301) (-1.9687314,0.72538079) (-1.9634954,0.73263488) (-1.9582594,0.7398465) (-1.9530234,0.74701684) (-1.9477874,0.75414701) (-1.9425515,0.76123811) (-1.9373155,0.76829114) (-1.9320795,0.77530711) (-1.9268435,0.78228695) (-1.9216075,0.78923157) (-1.9163715,0.79614185) (-1.9111355,0.8030186) (-1.9058995,0.80986264) (-1.9006636,0.81667473) (-1.8954276,0.82345562) (-1.8901916,0.83020601) (-1.8849556,0.83692659) (-1.8797196,0.84361802) (-1.8744836,0.85028093) (-1.8692476,0.85691594) (-1.8640116,0.86352363) (-1.8587757,0.87010458) (-1.8535397,0.87665932) (-1.8483037,0.8831884) (-1.8430677,0.88969232) (-1.8378317,0.89617157) (-1.8325957,0.90262664) (-1.8273597,0.90905797) (-1.8221237,0.91546603) (-1.8168878,0.92185123) (-1.8116518,0.928214) (-1.8064158,0.93455473) (-1.8011798,0.94087382) (-1.7959438,0.94717165) (-1.7907078,0.95344858) (-1.7854718,0.95970496) (-1.7802358,0.96594114) (-1.7749998,0.97215746) (-1.7697639,0.97835422) (-1.7645279,0.98453175) (-1.7592919,0.99069035) (-1.7540559,0.99683031) (-1.7488199,1.0029519) (-1.7435839,1.0090554) (-1.7383479,1.0151412) (-1.7331119,1.0212093) (-1.727876,1.0272602) (-1.72264,1.033294) (-1.717404,1.039311) (-1.712168,1.0453115) (-1.706932,1.0512956) (-1.701696,1.0572635) (-1.69646,1.0632155) (-1.691224,1.0691518) (-1.6859881,1.0750726) (-1.6807521,1.080978) (-1.6755161,1.0868683) (-1.6702801,1.0927437) (-1.6650441,1.0986043) (-1.6598081,1.1044502) (-1.6545721,1.1102818) (-1.6493361,1.1160991) (-1.6441002,1.1219022) (-1.6388642,1.1276914) (-1.6336282,1.1334669) (-1.6283922,1.1392286) (-1.6231562,1.1449768) (-1.6179202,1.1507117) (-1.6126842,1.1564333) (-1.6074482,1.1621418) (-1.6022123,1.1678374) (-1.5969763,1.17352) (-1.5917403,1.1791899) (-1.5865043,1.1848472) (-1.5812683,1.190492) (-1.5760323,1.1961244) (-1.5707963,1.2017445) (-1.5655603,1.2073524) (-1.5603244,1.2129482) (-1.5550884,1.218532) (-1.5498524,1.2241039) (-1.5446164,1.2296641) (-1.5393804,1.2352125) (-1.5341444,1.2407493) (-1.5289084,1.2462745) (-1.5236724,1.2517883) (-1.5184364,1.2572907) (-1.5132005,1.2627818) (-1.5079645,1.2682616) (-1.5027285,1.2737303) (-1.4974925,1.279188) (-1.4922565,1.2846346) (-1.4870205,1.2900702) (-1.4817845,1.295495) (-1.4765485,1.3009089) (-1.4713126,1.306312) (-1.4660766,1.3117045) (-1.4608406,1.3170863) (-1.4556046,1.3224575) (-1.4503686,1.3278181) (-1.4451326,1.3331682) (-1.4398966,1.3385079) (-1.4346606,1.3438372) (-1.4294247,1.3491561) (-1.4241887,1.3544647) (-1.4189527,1.359763) (-1.4137167,1.365051) (-1.4084807,1.3703289) (-1.4032447,1.3755966) (-1.3980087,1.3808541) (-1.3927727,1.3861015) (-1.3875368,1.3913389) (-1.3823008,1.3965662) (-1.3770648,1.4017835) (-1.3718288,1.4069908) (-1.3665928,1.4121882) (-1.3613568,1.4173756) (-1.3561208,1.4225531) (-1.3508848,1.4277207) (-1.3456489,1.4328784) (-1.3404129,1.4380262) (-1.3351769,1.4431642) (-1.3299409,1.4482923) (-1.3247049,1.4534107) (-1.3194689,1.4585192) (-1.3142329,1.4636179) (-1.3089969,1.4687069) (-1.303761,1.473786) (-1.298525,1.4788554) (-1.293289,1.483915) (-1.288053,1.4889649) (-1.282817,1.494005) (-1.277581,1.4990353) (-1.272345,1.5040558) (-1.267109,1.5090667) (-1.261873,1.5140677) (-1.2566371,1.519059) (-1.2514011,1.5240405) (-1.2461651,1.5290122) (-1.2409291,1.5339742) (-1.2356931,1.5389263) (-1.2304571,1.5438687) (-1.2252211,1.5488012) (-1.2199851,1.5537239) (-1.2147492,1.5586368) (-1.2095132,1.5635398) (-1.2042772,1.568433) (-1.1990412,1.5733163) (-1.1938052,1.5781897) (-1.1885692,1.5830531) (-1.1833332,1.5879067) (-1.1780972,1.5927503) (-1.1728613,1.5975838) (-1.1676253,1.6024074) (-1.1623893,1.607221) (-1.1571533,1.6120245) (-1.1519173,1.6168179) (-1.1466813,1.6216013) (-1.1414453,1.6263745) (-1.1362093,1.6311375) (-1.1309734,1.6358903) (-1.1257374,1.6406329) (-1.1205014,1.6453652) (-1.1152654,1.6500873) (-1.1100294,1.654799) (-1.1047934,1.6595003) (-1.0995574,1.6641913) (-1.0943214,1.6688718) (-1.0890855,1.6735418) (-1.0838495,1.6782013) (-1.0786135,1.6828502) (-1.0733775,1.6874885) (-1.0681415,1.6921161) (-1.0629055,1.6967331) (-1.0576695,1.7013393) (-1.0524335,1.7059347) (-1.0471976,1.7105192) (-1.0419616,1.7150929) (-1.0367256,1.7196556) (-1.0314896,1.7242074) (-1.0262536,1.728748) (-1.0210176,1.7332776) (-1.0157816,1.737796) (-1.0105456,1.7423032) (-1.0053096,1.7467991) (-1.0000737,1.7512837) (-0.99483767,1.7557569) (-0.98960169,1.7602186) (-0.9843657,1.7646688) (-0.97912971,1.7691075) (-0.97389372,1.7735344) (-0.96865773,1.7779497) (-0.96342175,1.7823532) (-0.95818576,1.7867448) (-0.95294977,1.7911245) (-0.94771378,1.7954923) (-0.9424778,1.7998479) (-0.93724181,1.8041914) (-0.93200582,1.8085228) (-0.92676983,1.8128418) (-0.92153385,1.8171485) (-0.91629786,1.8214427) (-0.91106187,1.8257244) (-0.90582588,1.8299935) (-0.90058989,1.8342499) (-0.89535391,1.8384936) (-0.89011792,1.8427244) (-0.88488193,1.8469422) (-0.87964594,1.8511471) (-0.87440996,1.8553388) (-0.86917397,1.8595174) (-0.86393798,1.8636826) (-0.85870199,1.8678345) (-0.853466,1.8719729) (-0.84823002,1.8760977) (-0.84299403,1.8802088) (-0.83775804,1.8843062) (-0.83252205,1.8883898) (-0.82728607,1.8924593) (-0.82205008,1.8965148) (-0.81681409,1.9005562) (-0.8115781,1.9045833) (-0.80634211,1.908596) (-0.80110613,1.9125942) (-0.79587014,1.9165779) (-0.79063415,1.9205468) (-0.78539816,1.924501) (-0.78016218,1.9284402) (-0.77492619,1.9323645) (-0.7696902,1.9362735) (-0.76445421,1.9401674) (-0.75921822,1.9440458) (-0.75398224,1.9479088) (-0.74874625,1.9517561) (-0.74351026,1.9555878) (-0.73827427,1.9594035) (-0.73303829,1.9632033) (-0.7278023,1.966987) (-0.72256631,1.9707545) (-0.71733032,1.9745056) (-0.71209433,1.9782402) (-0.70685835,1.9819583) (-0.70162236,1.9856596) (-0.69638637,1.989344) (-0.69115038,1.9930114) (-0.6859144,1.9966616) (-0.68067841,2.0002946) (-0.67544242,2.0039102) (-0.67020643,2.0075082) (-0.66497045,2.0110885) (-0.65973446,2.0146509) (-0.65449847,2.0181954) (-0.64926248,2.0217218) (-0.64402649,2.0252298) (-0.63879051,2.0287195) (-0.63355452,2.0321906) (-0.62831853,2.035643) (-0.62308254,2.0390765) (-0.61784656,2.042491) (-0.61261057,2.0458863) (-0.60737458,2.0492623) (-0.60213859,2.0526188) (-0.5969026,2.0559557) (-0.59166662,2.0592727) (-0.58643063,2.0625698) (-0.58119464,2.0658468) (-0.57595865,2.0691036) (-0.57072267,2.0723398) (-0.56548668,2.0755555) (-0.56025069,2.0787504) (-0.5550147,2.0819244) (-0.54977871,2.0850772) (-0.54454273,2.0882088) (-0.53930674,2.091319) (-0.53407075,2.0944075) (-0.52883476,2.0974743) (-0.52359878,2.1005191) (-0.51836279,2.1035418) (-0.5131268,2.1065422) (-0.50789081,2.1095201) (-0.50265482,2.1124754) (-0.49741884,2.1154078) (-0.49218285,2.1183173) (-0.48694686,2.1212036) (-0.48171087,2.1240665) (-0.47647489,2.1269059) (-0.4712389,2.1297215) (-0.46600291,2.1325133) (-0.46076692,2.135281) (-0.45553093,2.1380244) (-0.45029495,2.1407434) (-0.44505896,2.1434378) (-0.43982297,2.1461074) (-0.43458698,2.1487519) (-0.429351,2.1513714) (-0.42411501,2.1539654) (-0.41887902,2.1565339) (-0.41364303,2.1590767) (-0.40840704,2.1615936) (-0.40317106,2.1640844) (-0.39793507,2.166549) (-0.39269908,2.168987) (-0.38746309,2.1713984) (-0.38222711,2.173783) (-0.37699112,2.1761406) (-0.37175513,2.178471) (-0.36651914,2.180774) (-0.36128316,2.1830494) (-0.35604717,2.1852971) (-0.35081118,2.1875169) (-0.34557519,2.1897086) (-0.3403392,2.1918719) (-0.33510322,2.1940068) (-0.32986723,2.196113) (-0.32463124,2.1981905) (-0.31939525,2.2002389) (-0.31415927,2.2022581) (-0.30892328,2.2042479) (-0.30368729,2.2062082) (-0.2984513,2.2081388) (-0.29321531,2.2100395) (-0.28797933,2.2119102) (-0.28274334,2.2137506) (-0.27750735,2.2155606) (-0.27227136,2.21734) (-0.26703538,2.2190888) (-0.26179939,2.2208066) (-0.2565634,2.2224933) (-0.25132741,2.2241489) (-0.24609142,2.225773) (-0.24085544,2.2273656) (-0.23561945,2.2289265) (-0.23038346,2.2304556) (-0.22514747,2.2319526) (-0.21991149,2.2334176) (-0.2146755,2.2348502) (-0.20943951,2.2362503) (-0.20420352,2.2376179) (-0.19896753,2.2389528) (-0.19373155,2.2402548) (-0.18849556,2.2415237) (-0.18325957,2.2427596) (-0.17802358,2.2439622) (-0.1727876,2.2451314) (-0.16755161,2.2462671) (-0.16231562,2.2473691) (-0.15707963,2.2484374) (-0.15184364,2.2494719) (-0.14660766,2.2504723) (-0.14137167,2.2514387) (-0.13613568,2.2523709) (-0.13089969,2.2532688) (-0.12566371,2.2541322) (-0.12042772,2.2549612) (-0.11519173,2.2557556) (-0.10995574,2.2565154) (-0.10471976,2.2572403) (-0.099483767,2.2579305) (-0.09424778,2.2585857) (-0.089011792,2.2592059) (-0.083775804,2.2597911) (-0.078539816,2.2603412) (-0.073303829,2.260856) (-0.068067841,2.2613356) (-0.062831853,2.2617799) (-0.057595865,2.2621889) (-0.052359878,2.2625624) (-0.04712389,2.2629005) (-0.041887902,2.2632031) (-0.036651914,2.2634702) (-0.031415927,2.2637018) (-0.026179939,2.2638977) (-0.020943951,2.2640581) (-0.015707963,2.2641828) (-0.010471976,2.264272) (-0.0052359878,2.2643254) (-4.4408921e-16,2.2643433) (0.0052359878,2.2643254) (0.010471976,2.264272) (0.015707963,2.2641828) (0.020943951,2.2640581) (0.026179939,2.2638977) (0.031415927,2.2637018) (0.036651914,2.2634702) (0.041887902,2.2632031) (0.04712389,2.2629005) (0.052359878,2.2625624) (0.057595865,2.2621889) (0.062831853,2.2617799) (0.068067841,2.2613356) (0.073303829,2.260856) (0.078539816,2.2603412) (0.083775804,2.2597911) (0.089011792,2.2592059) (0.09424778,2.2585857) (0.099483767,2.2579305) (0.10471976,2.2572403) (0.10995574,2.2565154) (0.11519173,2.2557556) (0.12042772,2.2549612) (0.12566371,2.2541322) (0.13089969,2.2532688) (0.13613568,2.2523709) (0.14137167,2.2514387) (0.14660766,2.2504723) (0.15184364,2.2494719) (0.15707963,2.2484374) (0.16231562,2.2473691) (0.16755161,2.2462671) (0.1727876,2.2451314) (0.17802358,2.2439622) (0.18325957,2.2427596) (0.18849556,2.2415237) (0.19373155,2.2402548) (0.19896753,2.2389528) (0.20420352,2.2376179) (0.20943951,2.2362503) (0.2146755,2.2348502) (0.21991149,2.2334176) (0.22514747,2.2319526) (0.23038346,2.2304556) (0.23561945,2.2289265) (0.24085544,2.2273656) (0.24609142,2.225773) (0.25132741,2.2241489) (0.2565634,2.2224933) (0.26179939,2.2208066) (0.26703538,2.2190888) (0.27227136,2.21734) (0.27750735,2.2155606) (0.28274334,2.2137506) (0.28797933,2.2119102) (0.29321531,2.2100395) (0.2984513,2.2081388) (0.30368729,2.2062082) (0.30892328,2.2042479) (0.31415927,2.2022581) (0.31939525,2.2002389) (0.32463124,2.1981905) (0.32986723,2.196113) (0.33510322,2.1940068) (0.3403392,2.1918719) (0.34557519,2.1897086) (0.35081118,2.1875169) (0.35604717,2.1852971) (0.36128316,2.1830494) (0.36651914,2.180774) (0.37175513,2.178471) (0.37699112,2.1761406) (0.38222711,2.173783) (0.38746309,2.1713984) (0.39269908,2.168987) (0.39793507,2.166549) (0.40317106,2.1640844) (0.40840704,2.1615936) (0.41364303,2.1590767) (0.41887902,2.1565339) (0.42411501,2.1539654) (0.429351,2.1513714) (0.43458698,2.1487519) (0.43982297,2.1461074) (0.44505896,2.1434378) (0.45029495,2.1407434) (0.45553093,2.1380244) (0.46076692,2.135281) (0.46600291,2.1325133) (0.4712389,2.1297215) (0.47647489,2.1269059) (0.48171087,2.1240665) (0.48694686,2.1212036) (0.49218285,2.1183173) (0.49741884,2.1154078) (0.50265482,2.1124754) (0.50789081,2.1095201) (0.5131268,2.1065422) (0.51836279,2.1035418) (0.52359878,2.1005191) (0.52883476,2.0974743) (0.53407075,2.0944075) (0.53930674,2.091319) (0.54454273,2.0882088) (0.54977871,2.0850772) (0.5550147,2.0819244) (0.56025069,2.0787504) (0.56548668,2.0755555) (0.57072267,2.0723398) (0.57595865,2.0691036) (0.58119464,2.0658468) (0.58643063,2.0625698) (0.59166662,2.0592727) (0.5969026,2.0559557) (0.60213859,2.0526188) (0.60737458,2.0492623) (0.61261057,2.0458863) (0.61784656,2.042491) (0.62308254,2.0390765) (0.62831853,2.035643) (0.63355452,2.0321906) (0.63879051,2.0287195) (0.64402649,2.0252298) (0.64926248,2.0217218) (0.65449847,2.0181954) (0.65973446,2.0146509) (0.66497045,2.0110885) (0.67020643,2.0075082) (0.67544242,2.0039102) (0.68067841,2.0002946) (0.6859144,1.9966616) (0.69115038,1.9930114) (0.69638637,1.989344) (0.70162236,1.9856596) (0.70685835,1.9819583) (0.71209433,1.9782402) (0.71733032,1.9745056) (0.72256631,1.9707545) (0.7278023,1.966987) (0.73303829,1.9632033) (0.73827427,1.9594035) (0.74351026,1.9555878) (0.74874625,1.9517561) (0.75398224,1.9479088) (0.75921822,1.9440458) (0.76445421,1.9401674) (0.7696902,1.9362735) (0.77492619,1.9323645) (0.78016218,1.9284402) (0.78539816,1.924501) (0.79063415,1.9205468) (0.79587014,1.9165779) (0.80110613,1.9125942) (0.80634211,1.908596) (0.8115781,1.9045833) (0.81681409,1.9005562) (0.82205008,1.8965148) (0.82728607,1.8924593) (0.83252205,1.8883898) (0.83775804,1.8843062) (0.84299403,1.8802088) (0.84823002,1.8760977) (0.853466,1.8719729) (0.85870199,1.8678345) (0.86393798,1.8636826) (0.86917397,1.8595174) (0.87440996,1.8553388) (0.87964594,1.8511471) (0.88488193,1.8469422) (0.89011792,1.8427244) (0.89535391,1.8384936) (0.90058989,1.8342499) (0.90582588,1.8299935) (0.91106187,1.8257244) (0.91629786,1.8214427) (0.92153385,1.8171485) (0.92676983,1.8128418) (0.93200582,1.8085228) (0.93724181,1.8041914) (0.9424778,1.7998479) (0.94771378,1.7954923) (0.95294977,1.7911245) (0.95818576,1.7867448) (0.96342175,1.7823532) (0.96865773,1.7779497) (0.97389372,1.7735344) (0.97912971,1.7691075) (0.9843657,1.7646688) (0.98960169,1.7602186) (0.99483767,1.7557569) (1.0000737,1.7512837) (1.0053096,1.7467991) (1.0105456,1.7423032) (1.0157816,1.737796) (1.0210176,1.7332776) (1.0262536,1.728748) (1.0314896,1.7242074) (1.0367256,1.7196556) (1.0419616,1.7150929) (1.0471976,1.7105192) (1.0524335,1.7059347) (1.0576695,1.7013393) (1.0629055,1.6967331) (1.0681415,1.6921161) (1.0733775,1.6874885) (1.0786135,1.6828502) (1.0838495,1.6782013) (1.0890855,1.6735418) (1.0943214,1.6688718) (1.0995574,1.6641913) (1.1047934,1.6595003) (1.1100294,1.654799) (1.1152654,1.6500873) (1.1205014,1.6453652) (1.1257374,1.6406329) (1.1309734,1.6358903) (1.1362093,1.6311375) (1.1414453,1.6263745) (1.1466813,1.6216013) (1.1519173,1.6168179) (1.1571533,1.6120245) (1.1623893,1.607221) (1.1676253,1.6024074) (1.1728613,1.5975838) (1.1780972,1.5927503) (1.1833332,1.5879067) (1.1885692,1.5830531) (1.1938052,1.5781897) (1.1990412,1.5733163) (1.2042772,1.568433) (1.2095132,1.5635398) (1.2147492,1.5586368) (1.2199851,1.5537239) (1.2252211,1.5488012) (1.2304571,1.5438687) (1.2356931,1.5389263) (1.2409291,1.5339742) (1.2461651,1.5290122) (1.2514011,1.5240405) (1.2566371,1.519059) (1.261873,1.5140677) (1.267109,1.5090667) (1.272345,1.5040558) (1.277581,1.4990353) (1.282817,1.494005) (1.288053,1.4889649) (1.293289,1.483915) (1.298525,1.4788554) (1.303761,1.473786) (1.3089969,1.4687069) (1.3142329,1.4636179) (1.3194689,1.4585192) (1.3247049,1.4534107) (1.3299409,1.4482923) (1.3351769,1.4431642) (1.3404129,1.4380262) (1.3456489,1.4328784) (1.3508848,1.4277207) (1.3561208,1.4225531) (1.3613568,1.4173756) (1.3665928,1.4121882) (1.3718288,1.4069908) (1.3770648,1.4017835) (1.3823008,1.3965662) (1.3875368,1.3913389) (1.3927727,1.3861015) (1.3980087,1.3808541) (1.4032447,1.3755966) (1.4084807,1.3703289) (1.4137167,1.365051) (1.4189527,1.359763) (1.4241887,1.3544647) (1.4294247,1.3491561) (1.4346606,1.3438372) (1.4398966,1.3385079) (1.4451326,1.3331682) (1.4503686,1.3278181) (1.4556046,1.3224575) (1.4608406,1.3170863) (1.4660766,1.3117045) (1.4713126,1.306312) (1.4765485,1.3009089) (1.4817845,1.295495) (1.4870205,1.2900702) (1.4922565,1.2846346) (1.4974925,1.279188) (1.5027285,1.2737303) (1.5079645,1.2682616) (1.5132005,1.2627818) (1.5184364,1.2572907) (1.5236724,1.2517883) (1.5289084,1.2462745) (1.5341444,1.2407493) (1.5393804,1.2352125) (1.5446164,1.2296641) (1.5498524,1.2241039) (1.5550884,1.218532) (1.5603244,1.2129482) (1.5655603,1.2073524) (1.5707963,1.2017445) (1.5760323,1.1961244) (1.5812683,1.190492) (1.5865043,1.1848472) (1.5917403,1.1791899) (1.5969763,1.17352) (1.6022123,1.1678374) (1.6074482,1.1621418) (1.6126842,1.1564333) (1.6179202,1.1507117) (1.6231562,1.1449768) (1.6283922,1.1392286) (1.6336282,1.1334669) (1.6388642,1.1276914) (1.6441002,1.1219022) (1.6493361,1.1160991) (1.6545721,1.1102818) (1.6598081,1.1044502) (1.6650441,1.0986043) (1.6702801,1.0927437) (1.6755161,1.0868683) (1.6807521,1.080978) (1.6859881,1.0750726) (1.691224,1.0691518) (1.69646,1.0632155) (1.701696,1.0572635) (1.706932,1.0512956) (1.712168,1.0453115) (1.717404,1.039311) (1.72264,1.033294) (1.727876,1.0272602) (1.7331119,1.0212093) (1.7383479,1.0151412) (1.7435839,1.0090554) (1.7488199,1.0029519) (1.7540559,0.99683031) (1.7592919,0.99069035) (1.7645279,0.98453175) (1.7697639,0.97835422) (1.7749998,0.97215746) (1.7802358,0.96594114) (1.7854718,0.95970496) (1.7907078,0.95344858) (1.7959438,0.94717165) (1.8011798,0.94087382) (1.8064158,0.93455473) (1.8116518,0.928214) (1.8168878,0.92185123) (1.8221237,0.91546603) (1.8273597,0.90905797) (1.8325957,0.90262664) (1.8378317,0.89617157) (1.8430677,0.88969232) (1.8483037,0.8831884) (1.8535397,0.87665932) (1.8587757,0.87010458) (1.8640116,0.86352363) (1.8692476,0.85691594) (1.8744836,0.85028093) (1.8797196,0.84361802) (1.8849556,0.83692659) (1.8901916,0.83020601) (1.8954276,0.82345562) (1.9006636,0.81667473) (1.9058995,0.80986264) (1.9111355,0.8030186) (1.9163715,0.79614185) (1.9216075,0.78923157) (1.9268435,0.78228695) (1.9320795,0.77530711) (1.9373155,0.76829114) (1.9425515,0.76123811) (1.9477874,0.75414701) (1.9530234,0.74701684) (1.9582594,0.7398465) (1.9634954,0.73263488) (1.9687314,0.72538079) (1.9739674,0.71808301) (1.9792034,0.71074024) (1.9844394,0.70335112) (1.9896753,0.69591425) (1.9949113,0.68842811) (2.0001473,0.68089115) (2.0053833,0.67330171) (2.0106193,0.66565805) (2.0158553,0.65795833) (2.0210913,0.65020063) (2.0263273,0.64238288) (2.0315632,0.63450293) (2.0367992,0.62655849) (2.0420352,0.61854713) (2.0472712,0.61046626) (2.0525072,0.60231316) (2.0577432,0.59408489) (2.0629792,0.58577835) (2.0682152,0.57739021) (2.0734512,0.56891693) (2.0786871,0.5603547) (2.0839231,0.55169943) (2.0891591,0.54294674) (2.0943951,0.53409189) (2.0996311,0.52512977) (2.1048671,0.51605486) (2.1101031,0.50686115) (2.1153391,0.49754213) (2.120575,0.48809067) (2.125811,0.47849898) (2.131047,0.46875853) (2.136283,0.4588599) (2.141519,0.44879268) (2.146755,0.43854534) (2.151991,0.42810501) (2.157227,0.41745726) (2.1624629,0.40658588) (2.1676989,0.39547249) (2.1729349,0.38409615) (2.1781709,0.37243284) (2.1834069,0.36045478) (2.1886429,0.34812956) (2.1938789,0.33541898) (2.1991149,0.32227754) (2.2043508,0.30865028) (2.2095868,0.29446986) (2.2148228,0.27965222) (2.2200588,0.26409022) (2.2252948,0.24764367) (2.2305308,0.23012311) (2.2357668,0.21126161) (2.2410028,0.19066173) (2.2462387,0.16768416) (2.2514747,0.14117281) (2.2567107,0.10856931) (2.2619467,0.060750487) (2.2619467,-0.060750487) (2.2567107,-0.10856931) (2.2514747,-0.14117281) (2.2462387,-0.16768416) (2.2410028,-0.19066173) (2.2357668,-0.21126161) (2.2305308,-0.23012311) (2.2252948,-0.24764367) (2.2200588,-0.26409022) (2.2148228,-0.27965222) (2.2095868,-0.29446986) (2.2043508,-0.30865028) (2.1991149,-0.32227754) (2.1938789,-0.33541898) (2.1886429,-0.34812956) (2.1834069,-0.36045478) (2.1781709,-0.37243284) (2.1729349,-0.38409615) (2.1676989,-0.39547249) (2.1624629,-0.40658588) (2.157227,-0.41745726) (2.151991,-0.42810501) (2.146755,-0.43854534) (2.141519,-0.44879268) (2.136283,-0.4588599) (2.131047,-0.46875853) (2.125811,-0.47849898) (2.120575,-0.48809067) (2.1153391,-0.49754213) (2.1101031,-0.50686115) (2.1048671,-0.51605486) (2.0996311,-0.52512977) (2.0943951,-0.53409189) (2.0891591,-0.54294674) (2.0839231,-0.55169943) (2.0786871,-0.5603547) (2.0734512,-0.56891693) (2.0682152,-0.57739021) (2.0629792,-0.58577835) (2.0577432,-0.59408489) (2.0525072,-0.60231316) (2.0472712,-0.61046626) (2.0420352,-0.61854713) (2.0367992,-0.62655849) (2.0315632,-0.63450293) (2.0263273,-0.64238288) (2.0210913,-0.65020063) (2.0158553,-0.65795833) (2.0106193,-0.66565805) (2.0053833,-0.67330171) (2.0001473,-0.68089115) (1.9949113,-0.68842811) (1.9896753,-0.69591425) (1.9844394,-0.70335112) (1.9792034,-0.71074024) (1.9739674,-0.71808301) (1.9687314,-0.72538079) (1.9634954,-0.73263488) (1.9582594,-0.7398465) (1.9530234,-0.74701684) (1.9477874,-0.75414701) (1.9425515,-0.76123811) (1.9373155,-0.76829114) (1.9320795,-0.77530711) (1.9268435,-0.78228695) (1.9216075,-0.78923157) (1.9163715,-0.79614185) (1.9111355,-0.8030186) (1.9058995,-0.80986264) (1.9006636,-0.81667473) (1.8954276,-0.82345562) (1.8901916,-0.83020601) (1.8849556,-0.83692659) (1.8797196,-0.84361802) (1.8744836,-0.85028093) (1.8692476,-0.85691594) (1.8640116,-0.86352363) (1.8587757,-0.87010458) (1.8535397,-0.87665932) (1.8483037,-0.8831884) (1.8430677,-0.88969232) (1.8378317,-0.89617157) (1.8325957,-0.90262664) (1.8273597,-0.90905797) (1.8221237,-0.91546603) (1.8168878,-0.92185123) (1.8116518,-0.928214) (1.8064158,-0.93455473) (1.8011798,-0.94087382) (1.7959438,-0.94717165) (1.7907078,-0.95344858) (1.7854718,-0.95970496) (1.7802358,-0.96594114) (1.7749998,-0.97215746) (1.7697639,-0.97835422) (1.7645279,-0.98453175) (1.7592919,-0.99069035) (1.7540559,-0.99683031) (1.7488199,-1.0029519) (1.7435839,-1.0090554) (1.7383479,-1.0151412) (1.7331119,-1.0212093) (1.727876,-1.0272602) (1.72264,-1.033294) (1.717404,-1.039311) (1.712168,-1.0453115) (1.706932,-1.0512956) (1.701696,-1.0572635) (1.69646,-1.0632155) (1.691224,-1.0691518) (1.6859881,-1.0750726) (1.6807521,-1.080978) (1.6755161,-1.0868683) (1.6702801,-1.0927437) (1.6650441,-1.0986043) (1.6598081,-1.1044502) (1.6545721,-1.1102818) (1.6493361,-1.1160991) (1.6441002,-1.1219022) (1.6388642,-1.1276914) (1.6336282,-1.1334669) (1.6283922,-1.1392286) (1.6231562,-1.1449768) (1.6179202,-1.1507117) (1.6126842,-1.1564333) (1.6074482,-1.1621418) (1.6022123,-1.1678374) (1.5969763,-1.17352) (1.5917403,-1.1791899) (1.5865043,-1.1848472) (1.5812683,-1.190492) (1.5760323,-1.1961244) (1.5707963,-1.2017445) (1.5655603,-1.2073524) (1.5603244,-1.2129482) (1.5550884,-1.218532) (1.5498524,-1.2241039) (1.5446164,-1.2296641) (1.5393804,-1.2352125) (1.5341444,-1.2407493) (1.5289084,-1.2462745) (1.5236724,-1.2517883) (1.5184364,-1.2572907) (1.5132005,-1.2627818) (1.5079645,-1.2682616) (1.5027285,-1.2737303) (1.4974925,-1.279188) (1.4922565,-1.2846346) (1.4870205,-1.2900702) (1.4817845,-1.295495) (1.4765485,-1.3009089) (1.4713126,-1.306312) (1.4660766,-1.3117045) (1.4608406,-1.3170863) (1.4556046,-1.3224575) (1.4503686,-1.3278181) (1.4451326,-1.3331682) (1.4398966,-1.3385079) (1.4346606,-1.3438372) (1.4294247,-1.3491561) (1.4241887,-1.3544647) (1.4189527,-1.359763) (1.4137167,-1.365051) (1.4084807,-1.3703289) (1.4032447,-1.3755966) (1.3980087,-1.3808541) (1.3927727,-1.3861015) (1.3875368,-1.3913389) (1.3823008,-1.3965662) (1.3770648,-1.4017835) (1.3718288,-1.4069908) (1.3665928,-1.4121882) (1.3613568,-1.4173756) (1.3561208,-1.4225531) (1.3508848,-1.4277207) (1.3456489,-1.4328784) (1.3404129,-1.4380262) (1.3351769,-1.4431642) (1.3299409,-1.4482923) (1.3247049,-1.4534107) (1.3194689,-1.4585192) (1.3142329,-1.4636179) (1.3089969,-1.4687069) (1.303761,-1.473786) (1.298525,-1.4788554) (1.293289,-1.483915) (1.288053,-1.4889649) (1.282817,-1.494005) (1.277581,-1.4990353) (1.272345,-1.5040558) (1.267109,-1.5090667) (1.261873,-1.5140677) (1.2566371,-1.519059) (1.2514011,-1.5240405) (1.2461651,-1.5290122) (1.2409291,-1.5339742) (1.2356931,-1.5389263) (1.2304571,-1.5438687) (1.2252211,-1.5488012) (1.2199851,-1.5537239) (1.2147492,-1.5586368) (1.2095132,-1.5635398) (1.2042772,-1.568433) (1.1990412,-1.5733163) (1.1938052,-1.5781897) (1.1885692,-1.5830531) (1.1833332,-1.5879067) (1.1780972,-1.5927503) (1.1728613,-1.5975838) (1.1676253,-1.6024074) (1.1623893,-1.607221) (1.1571533,-1.6120245) (1.1519173,-1.6168179) (1.1466813,-1.6216013) (1.1414453,-1.6263745) (1.1362093,-1.6311375) (1.1309734,-1.6358903) (1.1257374,-1.6406329) (1.1205014,-1.6453652) (1.1152654,-1.6500873) (1.1100294,-1.654799) (1.1047934,-1.6595003) (1.0995574,-1.6641913) (1.0943214,-1.6688718) (1.0890855,-1.6735418) (1.0838495,-1.6782013) (1.0786135,-1.6828502) (1.0733775,-1.6874885) (1.0681415,-1.6921161) (1.0629055,-1.6967331) (1.0576695,-1.7013393) (1.0524335,-1.7059347) (1.0471976,-1.7105192) (1.0419616,-1.7150929) (1.0367256,-1.7196556) (1.0314896,-1.7242074) (1.0262536,-1.728748) (1.0210176,-1.7332776) (1.0157816,-1.737796) (1.0105456,-1.7423032) (1.0053096,-1.7467991) (1.0000737,-1.7512837) (0.99483767,-1.7557569) (0.98960169,-1.7602186) (0.9843657,-1.7646688) (0.97912971,-1.7691075) (0.97389372,-1.7735344) (0.96865773,-1.7779497) (0.96342175,-1.7823532) (0.95818576,-1.7867448) (0.95294977,-1.7911245) (0.94771378,-1.7954923) (0.9424778,-1.7998479) (0.93724181,-1.8041914) (0.93200582,-1.8085228) (0.92676983,-1.8128418) (0.92153385,-1.8171485) (0.91629786,-1.8214427) (0.91106187,-1.8257244) (0.90582588,-1.8299935) (0.90058989,-1.8342499) (0.89535391,-1.8384936) (0.89011792,-1.8427244) (0.88488193,-1.8469422) (0.87964594,-1.8511471) (0.87440996,-1.8553388) (0.86917397,-1.8595174) (0.86393798,-1.8636826) (0.85870199,-1.8678345) (0.853466,-1.8719729) (0.84823002,-1.8760977) (0.84299403,-1.8802088) (0.83775804,-1.8843062) (0.83252205,-1.8883898) (0.82728607,-1.8924593) (0.82205008,-1.8965148) (0.81681409,-1.9005562) (0.8115781,-1.9045833) (0.80634211,-1.908596) (0.80110613,-1.9125942) (0.79587014,-1.9165779) (0.79063415,-1.9205468) (0.78539816,-1.924501) (0.78016218,-1.9284402) (0.77492619,-1.9323645) (0.7696902,-1.9362735) (0.76445421,-1.9401674) (0.75921822,-1.9440458) (0.75398224,-1.9479088) (0.74874625,-1.9517561) (0.74351026,-1.9555878) (0.73827427,-1.9594035) (0.73303829,-1.9632033) (0.7278023,-1.966987) (0.72256631,-1.9707545) (0.71733032,-1.9745056) (0.71209433,-1.9782402) (0.70685835,-1.9819583) (0.70162236,-1.9856596) (0.69638637,-1.989344) (0.69115038,-1.9930114) (0.6859144,-1.9966616) (0.68067841,-2.0002946) (0.67544242,-2.0039102) (0.67020643,-2.0075082) (0.66497045,-2.0110885) (0.65973446,-2.0146509) (0.65449847,-2.0181954) (0.64926248,-2.0217218) (0.64402649,-2.0252298) (0.63879051,-2.0287195) (0.63355452,-2.0321906) (0.62831853,-2.035643) (0.62308254,-2.0390765) (0.61784656,-2.042491) (0.61261057,-2.0458863) (0.60737458,-2.0492623) (0.60213859,-2.0526188) (0.5969026,-2.0559557) (0.59166662,-2.0592727) (0.58643063,-2.0625698) (0.58119464,-2.0658468) (0.57595865,-2.0691036) (0.57072267,-2.0723398) (0.56548668,-2.0755555) (0.56025069,-2.0787504) (0.5550147,-2.0819244) (0.54977871,-2.0850772) (0.54454273,-2.0882088) (0.53930674,-2.091319) (0.53407075,-2.0944075) (0.52883476,-2.0974743) (0.52359878,-2.1005191) (0.51836279,-2.1035418) (0.5131268,-2.1065422) (0.50789081,-2.1095201) (0.50265482,-2.1124754) (0.49741884,-2.1154078) (0.49218285,-2.1183173) (0.48694686,-2.1212036) (0.48171087,-2.1240665) (0.47647489,-2.1269059) (0.4712389,-2.1297215) (0.46600291,-2.1325133) (0.46076692,-2.135281) (0.45553093,-2.1380244) (0.45029495,-2.1407434) (0.44505896,-2.1434378) (0.43982297,-2.1461074) (0.43458698,-2.1487519) (0.429351,-2.1513714) (0.42411501,-2.1539654) (0.41887902,-2.1565339) (0.41364303,-2.1590767) (0.40840704,-2.1615936) (0.40317106,-2.1640844) (0.39793507,-2.166549) (0.39269908,-2.168987) (0.38746309,-2.1713984) (0.38222711,-2.173783) (0.37699112,-2.1761406) (0.37175513,-2.178471) (0.36651914,-2.180774) (0.36128316,-2.1830494) (0.35604717,-2.1852971) (0.35081118,-2.1875169) (0.34557519,-2.1897086) (0.3403392,-2.1918719) (0.33510322,-2.1940068) (0.32986723,-2.196113) (0.32463124,-2.1981905) (0.31939525,-2.2002389) (0.31415927,-2.2022581) (0.30892328,-2.2042479) (0.30368729,-2.2062082) (0.2984513,-2.2081388) (0.29321531,-2.2100395) (0.28797933,-2.2119102) (0.28274334,-2.2137506) (0.27750735,-2.2155606) (0.27227136,-2.21734) (0.26703538,-2.2190888) (0.26179939,-2.2208066) (0.2565634,-2.2224933) (0.25132741,-2.2241489) (0.24609142,-2.225773) (0.24085544,-2.2273656) (0.23561945,-2.2289265) (0.23038346,-2.2304556) (0.22514747,-2.2319526) (0.21991149,-2.2334176) (0.2146755,-2.2348502) (0.20943951,-2.2362503) (0.20420352,-2.2376179) (0.19896753,-2.2389528) (0.19373155,-2.2402548) (0.18849556,-2.2415237) (0.18325957,-2.2427596) (0.17802358,-2.2439622) (0.1727876,-2.2451314) (0.16755161,-2.2462671) (0.16231562,-2.2473691) (0.15707963,-2.2484374) (0.15184364,-2.2494719) (0.14660766,-2.2504723) (0.14137167,-2.2514387) (0.13613568,-2.2523709) (0.13089969,-2.2532688) (0.12566371,-2.2541322) (0.12042772,-2.2549612) (0.11519173,-2.2557556) (0.10995574,-2.2565154) (0.10471976,-2.2572403) (0.099483767,-2.2579305) (0.09424778,-2.2585857) (0.089011792,-2.2592059) (0.083775804,-2.2597911) (0.078539816,-2.2603412) (0.073303829,-2.260856) (0.068067841,-2.2613356) (0.062831853,-2.2617799) (0.057595865,-2.2621889) (0.052359878,-2.2625624) (0.04712389,-2.2629005) (0.041887902,-2.2632031) (0.036651914,-2.2634702) (0.031415927,-2.2637018) (0.026179939,-2.2638977) (0.020943951,-2.2640581) (0.015707963,-2.2641828) (0.010471976,-2.264272) (0.0052359878,-2.2643254) (-4.4408921e-16,-2.2643433) (-0.0052359878,-2.2643254) (-0.010471976,-2.264272) (-0.015707963,-2.2641828) (-0.020943951,-2.2640581) (-0.026179939,-2.2638977) (-0.031415927,-2.2637018) (-0.036651914,-2.2634702) (-0.041887902,-2.2632031) (-0.04712389,-2.2629005) (-0.052359878,-2.2625624) (-0.057595865,-2.2621889) (-0.062831853,-2.2617799) (-0.068067841,-2.2613356) (-0.073303829,-2.260856) (-0.078539816,-2.2603412) (-0.083775804,-2.2597911) (-0.089011792,-2.2592059) (-0.09424778,-2.2585857) (-0.099483767,-2.2579305) (-0.10471976,-2.2572403) (-0.10995574,-2.2565154) (-0.11519173,-2.2557556) (-0.12042772,-2.2549612) (-0.12566371,-2.2541322) (-0.13089969,-2.2532688) (-0.13613568,-2.2523709) (-0.14137167,-2.2514387) (-0.14660766,-2.2504723) (-0.15184364,-2.2494719) (-0.15707963,-2.2484374) (-0.16231562,-2.2473691) (-0.16755161,-2.2462671) (-0.1727876,-2.2451314) (-0.17802358,-2.2439622) (-0.18325957,-2.2427596) (-0.18849556,-2.2415237) (-0.19373155,-2.2402548) (-0.19896753,-2.2389528) (-0.20420352,-2.2376179) (-0.20943951,-2.2362503) (-0.2146755,-2.2348502) (-0.21991149,-2.2334176) (-0.22514747,-2.2319526) (-0.23038346,-2.2304556) (-0.23561945,-2.2289265) (-0.24085544,-2.2273656) (-0.24609142,-2.225773) (-0.25132741,-2.2241489) (-0.2565634,-2.2224933) (-0.26179939,-2.2208066) (-0.26703538,-2.2190888) (-0.27227136,-2.21734) (-0.27750735,-2.2155606) (-0.28274334,-2.2137506) (-0.28797933,-2.2119102) (-0.29321531,-2.2100395) (-0.2984513,-2.2081388) (-0.30368729,-2.2062082) (-0.30892328,-2.2042479) (-0.31415927,-2.2022581) (-0.31939525,-2.2002389) (-0.32463124,-2.1981905) (-0.32986723,-2.196113) (-0.33510322,-2.1940068) (-0.3403392,-2.1918719) (-0.34557519,-2.1897086) (-0.35081118,-2.1875169) (-0.35604717,-2.1852971) (-0.36128316,-2.1830494) (-0.36651914,-2.180774) (-0.37175513,-2.178471) (-0.37699112,-2.1761406) (-0.38222711,-2.173783) (-0.38746309,-2.1713984) (-0.39269908,-2.168987) (-0.39793507,-2.166549) (-0.40317106,-2.1640844) (-0.40840704,-2.1615936) (-0.41364303,-2.1590767) (-0.41887902,-2.1565339) (-0.42411501,-2.1539654) (-0.429351,-2.1513714) (-0.43458698,-2.1487519) (-0.43982297,-2.1461074) (-0.44505896,-2.1434378) (-0.45029495,-2.1407434) (-0.45553093,-2.1380244) (-0.46076692,-2.135281) (-0.46600291,-2.1325133) (-0.4712389,-2.1297215) (-0.47647489,-2.1269059) (-0.48171087,-2.1240665) (-0.48694686,-2.1212036) (-0.49218285,-2.1183173) (-0.49741884,-2.1154078) (-0.50265482,-2.1124754) (-0.50789081,-2.1095201) (-0.5131268,-2.1065422) (-0.51836279,-2.1035418) (-0.52359878,-2.1005191) (-0.52883476,-2.0974743) (-0.53407075,-2.0944075) (-0.53930674,-2.091319) (-0.54454273,-2.0882088) (-0.54977871,-2.0850772) (-0.5550147,-2.0819244) (-0.56025069,-2.0787504) (-0.56548668,-2.0755555) (-0.57072267,-2.0723398) (-0.57595865,-2.0691036) (-0.58119464,-2.0658468) (-0.58643063,-2.0625698) (-0.59166662,-2.0592727) (-0.5969026,-2.0559557) (-0.60213859,-2.0526188) (-0.60737458,-2.0492623) (-0.61261057,-2.0458863) (-0.61784656,-2.042491) (-0.62308254,-2.0390765) (-0.62831853,-2.035643) (-0.63355452,-2.0321906) (-0.63879051,-2.0287195) (-0.64402649,-2.0252298) (-0.64926248,-2.0217218) (-0.65449847,-2.0181954) (-0.65973446,-2.0146509) (-0.66497045,-2.0110885) (-0.67020643,-2.0075082) (-0.67544242,-2.0039102) (-0.68067841,-2.0002946) (-0.6859144,-1.9966616) (-0.69115038,-1.9930114) (-0.69638637,-1.989344) (-0.70162236,-1.9856596) (-0.70685835,-1.9819583) (-0.71209433,-1.9782402) (-0.71733032,-1.9745056) (-0.72256631,-1.9707545) (-0.7278023,-1.966987) (-0.73303829,-1.9632033) (-0.73827427,-1.9594035) (-0.74351026,-1.9555878) (-0.74874625,-1.9517561) (-0.75398224,-1.9479088) (-0.75921822,-1.9440458) (-0.76445421,-1.9401674) (-0.7696902,-1.9362735) (-0.77492619,-1.9323645) (-0.78016218,-1.9284402) (-0.78539816,-1.924501) (-0.79063415,-1.9205468) (-0.79587014,-1.9165779) (-0.80110613,-1.9125942) (-0.80634211,-1.908596) (-0.8115781,-1.9045833) (-0.81681409,-1.9005562) (-0.82205008,-1.8965148) (-0.82728607,-1.8924593) (-0.83252205,-1.8883898) (-0.83775804,-1.8843062) (-0.84299403,-1.8802088) (-0.84823002,-1.8760977) (-0.853466,-1.8719729) (-0.85870199,-1.8678345) (-0.86393798,-1.8636826) (-0.86917397,-1.8595174) (-0.87440996,-1.8553388) (-0.87964594,-1.8511471) (-0.88488193,-1.8469422) (-0.89011792,-1.8427244) (-0.89535391,-1.8384936) (-0.90058989,-1.8342499) (-0.90582588,-1.8299935) (-0.91106187,-1.8257244) (-0.91629786,-1.8214427) (-0.92153385,-1.8171485) (-0.92676983,-1.8128418) (-0.93200582,-1.8085228) (-0.93724181,-1.8041914) (-0.9424778,-1.7998479) (-0.94771378,-1.7954923) (-0.95294977,-1.7911245) (-0.95818576,-1.7867448) (-0.96342175,-1.7823532) (-0.96865773,-1.7779497) (-0.97389372,-1.7735344) (-0.97912971,-1.7691075) (-0.9843657,-1.7646688) (-0.98960169,-1.7602186) (-0.99483767,-1.7557569) (-1.0000737,-1.7512837) (-1.0053096,-1.7467991) (-1.0105456,-1.7423032) (-1.0157816,-1.737796) (-1.0210176,-1.7332776) (-1.0262536,-1.728748) (-1.0314896,-1.7242074) (-1.0367256,-1.7196556) (-1.0419616,-1.7150929) (-1.0471976,-1.7105192) (-1.0524335,-1.7059347) (-1.0576695,-1.7013393) (-1.0629055,-1.6967331) (-1.0681415,-1.6921161) (-1.0733775,-1.6874885) (-1.0786135,-1.6828502) (-1.0838495,-1.6782013) (-1.0890855,-1.6735418) (-1.0943214,-1.6688718) (-1.0995574,-1.6641913) (-1.1047934,-1.6595003) (-1.1100294,-1.654799) (-1.1152654,-1.6500873) (-1.1205014,-1.6453652) (-1.1257374,-1.6406329) (-1.1309734,-1.6358903) (-1.1362093,-1.6311375) (-1.1414453,-1.6263745) (-1.1466813,-1.6216013) (-1.1519173,-1.6168179) (-1.1571533,-1.6120245) (-1.1623893,-1.607221) (-1.1676253,-1.6024074) (-1.1728613,-1.5975838) (-1.1780972,-1.5927503) (-1.1833332,-1.5879067) (-1.1885692,-1.5830531) (-1.1938052,-1.5781897) (-1.1990412,-1.5733163) (-1.2042772,-1.568433) (-1.2095132,-1.5635398) (-1.2147492,-1.5586368) (-1.2199851,-1.5537239) (-1.2252211,-1.5488012) (-1.2304571,-1.5438687) (-1.2356931,-1.5389263) (-1.2409291,-1.5339742) (-1.2461651,-1.5290122) (-1.2514011,-1.5240405) (-1.2566371,-1.519059) (-1.261873,-1.5140677) (-1.267109,-1.5090667) (-1.272345,-1.5040558) (-1.277581,-1.4990353) (-1.282817,-1.494005) (-1.288053,-1.4889649) (-1.293289,-1.483915) (-1.298525,-1.4788554) (-1.303761,-1.473786) (-1.3089969,-1.4687069) (-1.3142329,-1.4636179) (-1.3194689,-1.4585192) (-1.3247049,-1.4534107) (-1.3299409,-1.4482923) (-1.3351769,-1.4431642) (-1.3404129,-1.4380262) (-1.3456489,-1.4328784) (-1.3508848,-1.4277207) (-1.3561208,-1.4225531) (-1.3613568,-1.4173756) (-1.3665928,-1.4121882) (-1.3718288,-1.4069908) (-1.3770648,-1.4017835) (-1.3823008,-1.3965662) (-1.3875368,-1.3913389) (-1.3927727,-1.3861015) (-1.3980087,-1.3808541) (-1.4032447,-1.3755966) (-1.4084807,-1.3703289) (-1.4137167,-1.365051) (-1.4189527,-1.359763) (-1.4241887,-1.3544647) (-1.4294247,-1.3491561) (-1.4346606,-1.3438372) (-1.4398966,-1.3385079) (-1.4451326,-1.3331682) (-1.4503686,-1.3278181) (-1.4556046,-1.3224575) (-1.4608406,-1.3170863) (-1.4660766,-1.3117045) (-1.4713126,-1.306312) (-1.4765485,-1.3009089) (-1.4817845,-1.295495) (-1.4870205,-1.2900702) (-1.4922565,-1.2846346) (-1.4974925,-1.279188) (-1.5027285,-1.2737303) (-1.5079645,-1.2682616) (-1.5132005,-1.2627818) (-1.5184364,-1.2572907) (-1.5236724,-1.2517883) (-1.5289084,-1.2462745) (-1.5341444,-1.2407493) (-1.5393804,-1.2352125) (-1.5446164,-1.2296641) (-1.5498524,-1.2241039) (-1.5550884,-1.218532) (-1.5603244,-1.2129482) (-1.5655603,-1.2073524) (-1.5707963,-1.2017445) (-1.5760323,-1.1961244) (-1.5812683,-1.190492) (-1.5865043,-1.1848472) (-1.5917403,-1.1791899) (-1.5969763,-1.17352) (-1.6022123,-1.1678374) (-1.6074482,-1.1621418) (-1.6126842,-1.1564333) (-1.6179202,-1.1507117) (-1.6231562,-1.1449768) (-1.6283922,-1.1392286) (-1.6336282,-1.1334669) (-1.6388642,-1.1276914) (-1.6441002,-1.1219022) (-1.6493361,-1.1160991) (-1.6545721,-1.1102818) (-1.6598081,-1.1044502) (-1.6650441,-1.0986043) (-1.6702801,-1.0927437) (-1.6755161,-1.0868683) (-1.6807521,-1.080978) (-1.6859881,-1.0750726) (-1.691224,-1.0691518) (-1.69646,-1.0632155) (-1.701696,-1.0572635) (-1.706932,-1.0512956) (-1.712168,-1.0453115) (-1.717404,-1.039311) (-1.72264,-1.033294) (-1.727876,-1.0272602) (-1.7331119,-1.0212093) (-1.7383479,-1.0151412) (-1.7435839,-1.0090554) (-1.7488199,-1.0029519) (-1.7540559,-0.99683031) (-1.7592919,-0.99069035) (-1.7645279,-0.98453175) (-1.7697639,-0.97835422) (-1.7749998,-0.97215746) (-1.7802358,-0.96594114) (-1.7854718,-0.95970496) (-1.7907078,-0.95344858) (-1.7959438,-0.94717165) (-1.8011798,-0.94087382) (-1.8064158,-0.93455473) (-1.8116518,-0.928214) (-1.8168878,-0.92185123) (-1.8221237,-0.91546603) (-1.8273597,-0.90905797) (-1.8325957,-0.90262664) (-1.8378317,-0.89617157) (-1.8430677,-0.88969232) (-1.8483037,-0.8831884) (-1.8535397,-0.87665932) (-1.8587757,-0.87010458) (-1.8640116,-0.86352363) (-1.8692476,-0.85691594) (-1.8744836,-0.85028093) (-1.8797196,-0.84361802) (-1.8849556,-0.83692659) (-1.8901916,-0.83020601) (-1.8954276,-0.82345562) (-1.9006636,-0.81667473) (-1.9058995,-0.80986264) (-1.9111355,-0.8030186) (-1.9163715,-0.79614185) (-1.9216075,-0.78923157) (-1.9268435,-0.78228695) (-1.9320795,-0.77530711) (-1.9373155,-0.76829114) (-1.9425515,-0.76123811) (-1.9477874,-0.75414701) (-1.9530234,-0.74701684) (-1.9582594,-0.7398465) (-1.9634954,-0.73263488) (-1.9687314,-0.72538079) (-1.9739674,-0.71808301) (-1.9792034,-0.71074024) (-1.9844394,-0.70335112) (-1.9896753,-0.69591425) (-1.9949113,-0.68842811) (-2.0001473,-0.68089115) (-2.0053833,-0.67330171) (-2.0106193,-0.66565805) (-2.0158553,-0.65795833) (-2.0210913,-0.65020063) (-2.0263273,-0.64238288) (-2.0315632,-0.63450293) (-2.0367992,-0.62655849) (-2.0420352,-0.61854713) (-2.0472712,-0.61046626) (-2.0525072,-0.60231316) (-2.0577432,-0.59408489) (-2.0629792,-0.58577835) (-2.0682152,-0.57739021) (-2.0734512,-0.56891693) (-2.0786871,-0.5603547) (-2.0839231,-0.55169943) (-2.0891591,-0.54294674) (-2.0943951,-0.53409189) (-2.0996311,-0.52512977) (-2.1048671,-0.51605486) (-2.1101031,-0.50686115) (-2.1153391,-0.49754213) (-2.120575,-0.48809067) (-2.125811,-0.47849898) (-2.131047,-0.46875853) (-2.136283,-0.4588599) (-2.141519,-0.44879268) (-2.146755,-0.43854534) (-2.151991,-0.42810501) (-2.157227,-0.41745726) (-2.1624629,-0.40658588) (-2.1676989,-0.39547249) (-2.1729349,-0.38409615) (-2.1781709,-0.37243284) (-2.1834069,-0.36045478) (-2.1886429,-0.34812956) (-2.1938789,-0.33541898) (-2.1991149,-0.32227754) (-2.2043508,-0.30865028) (-2.2095868,-0.29446986) (-2.2148228,-0.27965222) (-2.2200588,-0.26409022) (-2.2252948,-0.24764367) (-2.2305308,-0.23012311) (-2.2357668,-0.21126161) (-2.2410028,-0.19066173) (-2.2462387,-0.16768416) (-2.2514747,-0.14117281) (-2.2567107,-0.10856931) (-2.2619467,-0.060750487)};
\addlegendentry{$p=0.3$}
\end{axis}\end{tikzpicture}\end{center}


\originalchapter{8.511作业5--11：运动、关联与配对}
\psource{Patrick A. Lee，8.511 Fall 2004，公开Problem Sets5--11。}
\originalsection{作业5（2004年10月19日截止）}
\begin{exercise}[原题1(a)--(c)]
一维周期势的带隙为$\Delta$，在均匀电场$\mathcal E$中价带与导带倾斜，等能的两点A、B可发生Zener隧穿。(a)证明在禁带中仍可作$\Psi_k=\sum_Ga_{k-G}\e^{\ii(k-G)x}$展开而$k$可复。令$G_0=2\pi/a,k=G_0/2+\kappa,E=E_0+\varepsilon,E_0=\hbar^2G_0^2/(8m)$，在两波近似下求$\kappa^2$。(b)若$\varepsilon(x)=e\mathcal E x$，求A、B距离。(c)按WKB求$P=\exp[-2\int_B^A|\kappa(x)|\dd x]$，给出半经典单带图像失效条件。
\end{exercise}
\begin{solution}
(a)Bloch平移关系可解析延拓到复$k$，禁带中复波矢描述衰减而不代表周期盒的可归一本征态。两波矩阵在区界附近对角$E_0\pm\hbar^2G_0\kappa/(2m)$，略去共同$O(\kappa^2)$项，得到$\varepsilon^2=|V_{G_0}|^2+\hbar^2(2E_0/m)\kappa^2$，即$\kappa^2=2m(\varepsilon^2-|V_{G_0}|^2)/(4\hbar^2E_0)$。$|\varepsilon|<|V|$时$\kappa$虚，$\Delta=2|V|$。

(b)$x_\pm=\pm|V|/(e\mathcal E)$，距离$\Delta/(e\mathcal E)$。(c)积分是椭圆面积：$\int_{-|V|}^{|V|}\sqrt{|V|^2-\varepsilon^2}\dd\varepsilon=\pi|V|^2/2$。因此
\begin{equation}
P=\exp\left[-\frac{\pi\sqrt{m/(2E_0)}|V|^2}{\hbar e\mathcal E}\right]=\exp\left[-\frac{\pi^2\Delta^2}{8E_0e\mathcal E a}\right].
\end{equation}
指数不再远大于1时带间隧穿不可忽略。WKB本身要求慢变场与局部带描述，不能把该近似指数当作任意强场的精确概率。
\end{solution}
\begin{exercise}[原题2(a)--(b)]
电子位于边长$L$的立方盒中，磁场$B\hat z$，采用原题Gaussian制$H=(\hbar\nabla/\ii+eA/c)^2/(2m)$与$A=-By\hat x$。(a)求能量。(b)求每级简并，并与Onsager规则比较。
\end{exercise}
\begin{solution}
平面波$\e^{\ii(k_xx+k_zz)}$使$y$方向化为中心$y_0=\hbar ck_x/(eB)$的谐振子，$E_{n,k_z}=\hbar\omega_c(n+1/2)+\hbar^2k_z^2/(2m)$，$\omega_c=eB/(mc)$。$k_x$间距$2\pi/L$，中心落在$0<y_0<L$的数量$eBL^2/(hc)$，不含自旋。固定$k_z$时$A_k=\pi k_\perp^2$，Onsager$A_k=(2\pi eB/\hbar c)(n+1/2)$给完全相同能量和一单位磁通对应一态的计数。
\end{solution}
\originalsection{作业6（2004年10月26日截止，修正版）}
\begin{exercise}[原题1(a)--(b)]
将K近似为自由电子金属。(a)求Shubnikov--de Haas振荡的$1/B$周期。(b)求$B=1\,\mathrm T$时实空间极值轨道面积。原题未给晶格常数或密度，给出依$n$的答案并明确数值假设。
\end{exercise}
\begin{solution}
$A_k=\pi k_F^2$，SI制$\Delta(1/B)=2e/(\hbar k_F^2)$，$A_r=(\hbar/eB)^2A_k$。若取常见室温近似密度$0.862\,\mathrm{g/cm^3}$、摩尔质量$39.0983\,\mathrm{g/mol}$且每原子一价电子，$n=1.328\times10^{28}\,\mathrm{m^{-3}}$，$k_F\simeq7.325\times10^9\,\mathrm{m^{-1}}$，周期约$5.67\times10^{-5}\,\mathrm{T^{-1}}$，$A_r\simeq7.30\times10^{-11}\,\mathrm{m^2}$。这些数字是所列模型输入的计算，不是源题给定数据；实材料参数随温度而变。
\end{solution}
\begin{exercise}[原题2(a)--(d)]
能带$E=\hbar^2[k_x^2/m_x+k_y^2/m_y+k_z^2/m_z]/2$。在半经典方程左侧加$k/\tau$，令$E,B$恒定、$B=B\hat z$、$j=-env$。(a)求$3\times3$DC电导矩阵，以$\omega_{cx}=eB/(m_xc),\omega_{cy}=eB/(m_yc)$表示。(b)求电阻矩阵。(c)讨论低高场Hall电导与Hall电阻。(d)画对角电导和电阻随场变化。
\end{exercise}
\begin{solution}
令$D=1+\omega_{cx}\omega_{cy}\tau^2$。稳态方程为$v_x+\tau\omega_{cx}v_y=-e\tau E_x/m_x$、$v_y-\tau\omega_{cy}v_x=-e\tau E_y/m_y$。解得
\begin{equation}
\sigma=ne^2\tau\begin{pmatrix}(m_xD)^{-1}&-eB\tau/(cm_xm_yD)&0\\eB\tau/(cm_xm_yD)&(m_yD)^{-1}&0\\0&0&m_z^{-1}\end{pmatrix},\quad
\rho=\begin{pmatrix}m_x/(ne^2\tau)&B/(nec)&0\\-B/(nec)&m_y/(ne^2\tau)&0\\0&0&m_z/(ne^2\tau)\end{pmatrix}.
\end{equation}
低场$\sigma_{xy}\propto-B$，高场$\sigma_{xy}\simeq-nec/B$；$\rho_{xy}=B/(nec)$始终线性。$\sigma_{xx},\sigma_{yy}$按$1/(1+\mathrm{const}\,B^2)$下降，$\sigma_{zz}$常数；所有对角$\rho$均常数，该单椭球、常$\tau$模型没有磁电阻，虽有磁电导。SI制把每个$B/c$换成$B$。正文与数值图采用SI，原题表达保留Gaussian换算说明。
\end{solution}
\originalsection{作业7（2004年11月2日截止）}
原题1仅指“Marder第17章题6”，另给$ j_Q=\Pi j-\kappa\partial_xT$、$\Pi_p>0,\Pi_n<0$。完整装置、边界与要求不在公开题面中，故保留引用缺口；不能仅凭热流式猜题。
\begin{exercise}[原题2(a)--(b)]
(a)写N个相互作用电子HF基态能$E_N$，说明为何不等于一电子能级之和。(b)移去第$i$轨道电子并冻结其余轨道，证明$E_N-E_{N-1}(i)=\epsilon_i$，即原题定义的有符号电离能Koopmans关系。
\end{exercise}
\begin{solution}
$E_N=\sum_ih_{ii}+\frac12\sum_{ij}(J_{ij}-K_{ij})$，而$\epsilon_i=h_{ii}+\sum_j(J_{ij}-K_{ij})$。双和每对两次，故$E_N=\sum_i\epsilon_i-\frac12\sum_{ij}(J_{ij}-K_{ij})$。删除$i$后差为$h_{ii}+\sum_j(J_{ij}-K_{ij})=\epsilon_i$，自作用$J_{ii}-K_{ii}=0$。实际移走电子所需的正功常定义为$E_{N-1}-E_N=-\epsilon_i$，原题“ionization energy”的符号应按其明确定义理解。冻结轨道是定理必要的近似。
\end{solution}
\originalsection{作业8（2004年11月16日截止）}
\begin{exercise}[原题1(a)--(c)]
(a)均匀正背景中的电子气，求HF本征值与基态能。(b)LDA中$E_{xc}=\int n\epsilon_{xc}(n)\dd r$，$\mu_{xc}=\dd(n\epsilon_{xc})/\dd n$，证明外势$v(r)$中的总能为$\sum_i\epsilon_i-E_H+\int n(\epsilon_{xc}-\mu_{xc})\dd r$。(c)把$\epsilon_{xc}$只取(a)的交换能，求均匀体系KS能级，并比较HF的费米速度、带宽和基态能；背景自能也须计入。
\end{exercise}
\begin{solution}
(a)第\ref{ch:interaction}章给$\epsilon_k=\hbar^2k^2/(2m)+\Sigma_x(k)$，$E/N=3E_F/5-3e_*^2k_F/(4\pi)$。正--正、正--负和电子Hartree长程项相互抵消，若漏背景自能则会剩无限体积发散。

(b)将KS方程按占据轨道求和得$\sum_i\epsilon_i=T_s+\int vn+2E_H+\int n\mu_{xc}$；与$E=T_s+\int vn+E_H+\int n\epsilon_{xc}$比较即得公式。

(c)$\epsilon_x=-3e_*^2k_F/(4\pi)$且$k_F\propto n^{1/3}$，所以$\mu_x=4\epsilon_x/3=-e_*^2k_F/\pi$。均匀KS能$\epsilon_k^{KS}=\hbar^2k^2/(2m)-e_*^2k_F/\pi$，交换是常数，速度为裸$\hbar k_F/m$，带宽$E_F$。HF自能在0与$k_F$间上升$e_*^2k_F/\pi$，故带宽$E_F+e_*^2k_F/\pi$，费米速度有对数发散。KS按(b)计算总能仍与HF交换总能相同：辅助能级差不等于总能泛函差。
\end{solution}
原题2要求阅读Marder附录C学习二次量子化，没有独立计算题。本册第\ref{ch:interaction}章提供所需先修，阅读指针不计为已取得的书本题。
\originalsection{作业9（2004年11月23日截止）}
\begin{exercise}[原题1--2]
两站Hubbard模型$H=t\sum_\sigma(c_{1\sigma}^\dagger c_{2\sigma}+\mathrm{h.c.})+U\sum_{i=1,2}n_{i\uparrow}n_{i\downarrow}$。(1)用创建算符写1、2、3、4电子Hilbert基及矩阵，给1、3、4电子本征能与态。(2)精确求2电子本征能与简并，并在$U\gg|t|$与自旋$1/2$的$J\vct S_1\cdot\vct S_2$比较，求$J$。
\end{exercise}
\begin{solution}
固定单轨道顺序$(1\uparrow,1\downarrow,2\uparrow,2\downarrow)$，各粒子数基为从四个$c_\alpha^\dagger$中选$N$个按顺序乘真空，维数依次$4,6,4,1$。矩阵元由$tc_\alpha^\dagger c_\beta$按反对易换序得，源码另以此规则实际构矩阵复核。

$N=1$每自旋矩阵$\left(\begin{smallmatrix}0&t\\t&0\end{smallmatrix}\right)$，态$(c_{1\sigma}^\dagger\pm c_{2\sigma}^\dagger)\ket0/\sqrt2$，能量$\pm t$。$N=3$按一个空穴基等价于对角$U$的两个$2\times2$块，能量$U\pm t$各二重，本征矢是两种相同剩余自旋配置的双占据位置叠加，相位由固定基顺序决定；严格向量与完整矩阵随计算证据给出。$N=4$唯一全占据态，跃迁因Pauli原理阻塞，能量$2U$。

$N=2$三重态$T_+=c_{1\uparrow}^\dagger c_{2\uparrow}^\dagger\ket0,T_-=c_{1\downarrow}^\dagger c_{2\downarrow}^\dagger\ket0,T_0=(c_{1\uparrow}^\dagger c_{2\downarrow}^\dagger+c_{1\downarrow}^\dagger c_{2\uparrow}^\dagger)\ket0/\sqrt2$皆能量0。单重态$S=(c_{1\uparrow}^\dagger c_{2\downarrow}^\dagger-c_{1\downarrow}^\dagger c_{2\uparrow}^\dagger)\ket0/\sqrt2$与$D_+=(c_{1\uparrow}^\dagger c_{1\downarrow}^\dagger+c_{2\uparrow}^\dagger c_{2\downarrow}^\dagger)\ket0/\sqrt2$构成矩阵$\left(\begin{smallmatrix}0&2t\\2t&U\end{smallmatrix}\right)$，本征根$E_\pm=(U\pm\sqrt{U^2+16t^2})/2$，态归一化为$S+(E_\pm/2t)D_+$。$D_-=(D_1-D_2)/\sqrt2$能量$U$。$t=0$用连续极限或直接原基处理。

强耦合$E_-=-4t^2/U+O(t^4/U^3)$，三重态仍0，匹配$J(\vct S_1\cdot\vct S_2-1/4)$，故$J=4t^2/U$。比较普通Heisenberg时需加常数$-J/4$，否则本征能绝对值不会相同。
\end{solution}
\originalsection{作业10（2004年11月30日截止）}
\begin{exercise}[原题1(a)--(b)]
零温费米能上方$E$的电子：(a)证明低能声子散射单粒子宽度量级$\hbar/\tau\sim E^3/(\hbar\omega_D)^2$；(b)输运时间在跃迁概率中乘$1-\cos\theta$，证明$\hbar/\tau_t\sim E^5/(\hbar\omega_D)^4$。只求幂律及自然尺度，材料系数可省略。
\end{exercise}
\begin{solution}
声学声子$\omega_q=c_sq$且$|g_q|^2\propto q$。三维相空间$q^2\dd q$，能量守恒角积分给$1/(v_Fq)$，零温Pauli因子限制$\hbar c_sq<E$，所以率$\propto\int_0^{E/\hbar c_s}q^2\dd q\propto E^3$。低温大角输运由小声子动量造成的电子转角$\theta\sim q/k_F$，故多$\int q^2$权重，率$\propto E^5$。以Debye能归一后量纲给题示尺度，但系数仍含耦合及$k_F/q_D$比；替换$E\to\kb T$只给有限温幂律，不给精确系数。该推导要求三维、声学形变耦合与$E\ll\hbar\omega_D$。
\end{solution}
\originalsection{作业11（2004年12月9日截止）}
\begin{exercise}[原题1(a)--(c)]
约化BCS哈密顿量$H=\sum_{k\sigma}\epsilon_kc_{k\sigma}^\dagger c_{k\sigma}-V_0\sum_{kk'}c_{k\uparrow}^\dagger c_{-k\downarrow}^\dagger c_{-k'\downarrow}c_{k'\uparrow}$与BCS态。(a)求$W=\avg{H-\mu N}$。(b)在$u_k^2+|v_k|^2=1$下变分，令$u=\cos\theta,v=\sin\theta$，求$\tan2\theta=\Delta/\xi$与零温隙方程。(c)从有限温方程推低温$\Delta(T)=\Delta_0[1-a(T)\e^{-\Delta_0/\kb T}]$，证明$a(T)\propto\sqrt T$。
\end{exercise}
\begin{solution}
(a)$\avg{n_{k\sigma}}=|v_k|^2$，配对平均$u_k^*v_k$，得$W=2\sum_k\xi_k|v_k|^2-V_0|\sum_ku_k^*v_k|^2$（约化模型的热力学极限，重复$k$自项的$O(1/V)$约定不改隙方程）。(b)统一相位令$\Delta$实，$\partial_{\theta_k}W=2\xi_k\sin2\theta_k-2\Delta\cos2\theta_k=0$，所以$\sin2\theta=\Delta/\sqrt{\xi^2+\Delta^2}$，$\Delta=V_0\sum\Delta/(2E_k)$。近费米面每自旋态密度常数给$1=N_\uparrow(0)V_0\int_0^{\hbar\omega_D}\dd\xi/\sqrt{\xi^2+\Delta^2}$。

(c)从有限温积分扣除零温积分，弱耦合下$\log(\Delta_0/\Delta)=2\int_0^\infty f(\sqrt{\xi^2+\Delta^2})\dd\xi/\sqrt{\xi^2+\Delta^2}$。低温以$E\simeq\Delta_0+\xi^2/(2\Delta_0)$估计高斯积分，右侧$\simeq\sqrt{2\pi\kb T/\Delta_0}\e^{-\Delta_0/\kb T}$，故
$\Delta_0-\Delta\simeq\sqrt{2\pi\Delta_0\kb T}\e^{-\Delta_0/\kb T}$。这里$a(T)=\sqrt{2\pi\kb T/\Delta_0}$，不是与温度无关常数；原题也说明其含$\sqrt T$。
\end{solution}

\originalchapter{8.512作业1--3：涨落、介电与Peierls不稳定性}
\psource{Patrick A. Lee，8.512 Spring 2009，公开作业1--3；题页保留2008底稿和2009改期痕迹，不把印刷日期混作独立课次。}
\originalsection{作业1（公开修订题页，2009年2月）}
\begin{exercise}[原题1(a)--(b)]\label{p:512-1-1}
原题取$\hbar=1$，$S(q,\omega)=\int\dd x\dd t\,\e^{-\ii qx+\ii\omega t}\avg{\rho(x,t)\rho(0,0)}$，$n_B=(\e^{\beta\omega}-1)^{-1}$，响应采用外势$+f\rho$约定。(a)证明$\chi''=(\e^{-\beta\omega}-1)S/2$与$S=-2(n_B+1)\chi''$。(b)证明$\chi''(q,\omega)=-\chi''(-q,-\omega)$、$S(-q,-\omega)=\e^{-\beta\omega}S(q,\omega)$，解释详细平衡。
\end{exercise}
\begin{solution}
插入能量态，$S=2\pi\sum_{nm}p_n|\rho_{nm}(q)|^2\delta(\omega-E_m+E_n)$。交换$n,m$使$q,\omega$均反转，权重比$p_m/p_n=\e^{-\beta\omega}$，得到详细平衡。该响应约定的谱虚部为$-\pi\sum(p_n-p_m)|\rho_{nm}|^2\delta(\omega-E_m+E_n)$，故$\chi''=-(1-\e^{-\beta\omega})S/2$。利用$n_B+1=(1-\e^{-\beta\omega})^{-1}$得第二式。两次反转令虚部反号。探针失能$\omega>0$的概率相对于得能大$\e^{\beta\omega}$倍，是样品吸收与释放能量的热平衡比例。
\end{solution}
\begin{exercise}[原题2(a)--(e)]\label{p:512-1-2}
晶体核位置$r_j=R_j+u_j$，$\hbar=1$，$u_j=\sum_q(2NM\omega_q)^{-1/2}(a_q\e^{\ii(qR_j-\omega_qt)}+\mathrm{h.c.})$；先假设三种极化简并，选择沿$Q$的一种并将$Q\cdot u_j$写为$Qu_j$。定义$F(Q,t)=\sum_{jl}\e^{-\ii Q(R_j-R_l)}F_{jl}(t)$，$F_{jl}=\avg{\e^{-\ii Qu_j(t)}\e^{\ii Qu_l(0)}}$，$S=\int\e^{\ii\omega t}F\dd t$。(a)证明$F_{jl}=\avg{\e^{-\ii Q[u_j(t)-u_l(0)]}}\e^{[Qu_j(t),Qu_l(0)]/2}$，可调用高斯平均公式。(b)令$\theta_{jl}=q(R_j-R_l)-\omega_qt$，证明
\begin{equation}
F_{jl}=\e^{-2W}\exp\left\{\frac{Q^2}{2NM}\sum_q\frac{(2n_q+1)\cos\theta_{jl}+\ii\sin\theta_{jl}}{\omega_q}\right\},\quad2W=\frac{Q^2}{2NM}\sum_q\frac{2n_q+1}{\omega_q}.
\end{equation}
(c)展开指数至一个声子阶，写出弹性及$Q=G\pm q,\omega=\pm\omega_q$项，$V^*$为倒格原胞体积。(d)解释各项。(e)怎样测声子极化向量？
\end{exercise}
\begin{solution}
(a)谐振位移的不同时间对易子为数，BCH恰截断，因$(-\ii)(\ii)=1$得到所给正号。高斯平均为$\exp[-Q^2\avg{(u_j-u_l)^2}/2]$。(b)位移自相关的两个本征权重为$(n_q+1)\e^{\ii\theta}+n_q\e^{-\ii\theta}=(2n_q+1)\cos\theta+\ii\sin\theta$；自方差项组成$2W$。

(c)把该恒等式代入指数的一阶，格点和给倒格$\delta$，时间积分给能量$\delta$。若用通常$\int\e^{\ii\omega t}\dd t=2\pi\delta(\omega)$，完整表达为
\begin{align}
S=2\pi NV^*\e^{-2W}\bigg[&\sum_G\delta(Q-G)\delta(\omega)\\
+&\frac{Q^2}{2NM}\sum_q\frac1{\omega_q}\Big\{(n_q+1)\sum_G\delta(Q-G-q)\delta(\omega-\omega_q)\nonumber\\
&\hspace{31mm}+n_q\sum_G\delta(Q-G+q)\delta(\omega+\omega_q)\Big\}\bigg].\nonumber
\end{align}
原题把Fourier归一的$2\pi$吸收到其$\delta$记号中；须同时保持$S$定义与实验散射率的约定。(d)Bragg弹性峰、创造一个声子、吸收一个声子依次对应三项；零温吸收项为零。(e)恢复非简并极化后强度为$|Q\cdot\lambda_{q\alpha}|^2$，改变$Q$方向并在相同$q,\omega_{q\alpha}$测强度可重建投影，纵横模需几何与多方向数据区分。
\end{solution}
\originalsection{作业2（题页修订为2009年2月23日）}
\begin{exercise}[原题1(a)--(c)]\label{p:512-2-1}
高速电子在金属内因发射等离子体激元而失能。(a)小$q$时$\operatorname{Im}[-1/\epsilon(q,\omega)]\simeq A(q)\delta(\omega-\omega_p)$，由和规则求$A$。(b)写出发射一个激元进入固角的概率表达式。(c)取电子能100 keV等典型值估计平均自由程，说明选取参数。
\end{exercise}
\begin{solution}
(a)逆介电和规则$\int_0^\infty\omega\operatorname{Im}[-\epsilon^{-1}]\dd\omega=\pi\omega_p^2/2$，所以$A=\pi\omega_p/2$。同样可由$\epsilon=1-\omega_p^2/(\omega+\ii0)^2$直接取极点检验。

(b)选择平面波归一体积$V$与能量变量$E=\hbar\omega$，外电子的屏蔽库仑势谱给微分率
\begin{equation}
\dd\Gamma=\frac2\hbar\frac{\dd^3q}{(2\pi)^3}\,v_q\operatorname{Im}[-\epsilon^{-1}(q,\omega_q)]\,,\qquad \hbar\omega_q=\epsilon_k-\epsilon_{k-q},\quad v_q=\frac{4\pi e_*^2}{q^2}.
\end{equation}
代入$A\delta(\omega_q-\omega_p)$再由$q=k-k'$换到$k'^2\dd k'\dd\Omega$，并对$k'$积分，就得定角发射率；可写为$\dd\Gamma/\dd\Omega=(2/\hbar)(2\pi)^{-3}v_{|k-k'|}A\,k'^2/|\partial_{k'}\omega_q|$，$k'$由能量守恒确定。该表达指定了归一与Jacobian，避免把速度与波矢单位混在一起。

(c)在$v$远大于电子费米速度时忽略反冲，$\omega_q\simeq\vct v\cdot\vct q$，角积分给$1/(vq)$，于是$\Gamma\simeq(e_*^2\omega_p/\hbar v)\log(q_{\max}/q_{\min})$，$q_{\min}=\omega_p/v$，$\lambda=v/\Gamma$。取$\hbar\omega_p=15$ eV，100 keV电子的相对论速度$v=c\sqrt{1-(1+100/511)^{-2}}\simeq0.548c$；取激元进入粒子--空穴连续谱的截断$q_{\max}\sim\omega_p/v_F$、$v_F=1.5\times10^6$ m/s，得到对数约4.7、平均自由程约$1.15\times10^{-7}$ m量级。若采用其他截断，数字可变化；大$q$时单极点近似失效，这不是材料专属精确寿命。
\end{solution}
\begin{exercise}[原题2(a)--(c)]\label{p:512-2-2}
(a)证明周期固体的RPA介电函数为$1+v_q\sum_{k\alpha\beta}|\avg{k+q,\beta|\e^{\ii q\cdot r}|k,\alpha}|^2(f_{k\alpha}-f_{k+q,\beta})/(E_{k+q,\beta}-E_{k\alpha}-\hbar\omega-\ii0)$，各自旋与体积计数要一致。(b)静态长波半导体以平均带隙$\Delta$替代跃迁能；证明$\sum_b(E_b-E_a)|\avg{b|\e^{\ii q r}|a}|^2=\hbar^2q^2/(2m)$，说明所需平衡态条件。(c)再把该和中的跃迁能取$\Delta$，求$\epsilon(0)=1+(\hbar\omega_p/\Delta)^2$，估计Si、Ge并讨论与实验比较的限制。
\end{exercise}
\begin{solution}
(a)先在Bloch能量基对外势作一阶响应，矩阵元加占据差给粒子--空穴泡图，再自洽加Hartree势形成$\epsilon=1+v_q\Pi$，每个$k$应带$1/V$，自旋若不在$\alpha$中再乘2。

(b)$O=\e^{\ii q r}$满足$O^\dagger HO=H+\hbar\vct q\cdot\vct p/m+\hbar^2q^2/(2m)$。在无流的时间反演对称态或成对平均中$\avg{p}=0$，完备关系给所需式。对一个具有非零流的单Bloch态，还应保留$\hbar q\cdot\avg p/m$；原题提时间反演是为使该项取消，并不保证任意单个$k$态本身无流。

(c)绝缘体的带间贡献包含正反跃迁，对占据态求和后$\Pi\simeq n\hbar^2q^2/(m\Delta^2)$，故$\epsilon-1=\hbar^2\omega_p^2/\Delta^2$。若用价电子等离子体能约16 eV与有效平均跃迁尺度Si约4.9 eV、Ge约4.1 eV，则预测约11.7、16.2。这里选择平均光学跃迁尺度是模型输入，不能把间接最低带隙Si约1.1 eV、Ge约0.7 eV代入并宣称成功：那会给大得多的介电常数。与实验做精确比较需核对温度、频率及实测光谱；本题仅给模型量级及失效原因。
\end{solution}
\originalsection{作业3（题页修订为2009年3月2日）}
\begin{exercise}[原题1(a)--(e)]\label{p:512-3-1}
一维H链间距$a$，每原子一个$1s$轨道、一电子，只保留最近邻$V(a)<0$并忽略重叠。(a)求带与费米能。(b)原子位移$(-1)^n\delta$，令$\Delta=2\delta\dd V/\dd a$，证明跃迁交替$V\pm\Delta$并求新带、金属或绝缘性质。(c)求小$|\Delta/V|$时电子能改变的$\Delta^2\log|\Delta/V|$项及系数。(d)晶格代价每站$b\delta^2$，证明均匀链不稳定。(e)求一维自由费米气静态极化，证明$2k_F$对数奇异。
\end{exercise}
\begin{solution}
(a)$E(k)=\epsilon_{1s}+2V\cos ka$，双自旋半填充，$k_F=\pi/(2a)$、$E_F=\epsilon_{1s}$。(b)相邻距为$a\pm2\delta$，Taylor展开给交替跳跃。两站胞$q\in[-\pi/(2a),\pi/(2a)]$，$E_\pm=\epsilon_{1s}\pm2\sqrt{V^2\cos^2qa+\Delta^2\sin^2qa}$。下带填满，边界隙$4|\Delta|$，变成带绝缘体。

(c)按每原子并含自旋，电子能$E_{\mathrm{el}}/N=-(4/\pi)\int_0^{\pi/2}\sqrt{V^2\cos^2x+\Delta^2\sin^2x}\dd x$，忽略共同$\epsilon_{1s}$。边界$x=\pi/2-y$处差$\sqrt{V^2y^2+\Delta^2}-|V|y$，积分的对数部分$[\Delta^2/(2|V|)]\log(|V|/|\Delta|)$，所以$\delta E_{\mathrm{el}}/N=[2\Delta^2/(\pi|V|)]\log|\Delta/V|+O(\Delta^2)$。(d)$\Delta\propto\delta$，负$\delta^2\log(1/|\delta|)$压过任何固定正$b\delta^2$，在该零温无限链模型中任意小畸变可降能。有限温、有限长或电子关联改变红外截止。

(e)取双自旋态与$\Pi=-\chi$，一维积分直接得
\begin{equation}
\Pi(q,0)=\frac{2m}{\pi\hbar^2|q|}\log\left|\frac{|q|+2k_F}{|q|-2k_F}\right|.
\end{equation}
$q\to0$极限$2m/(\pi\hbar^2k_F)$，等于双自旋态密度；$q=2k_F$有对数，来源于两个费米点被同一波矢连接。
\end{solution}
\begin{exercise}[原题2]\label{p:512-3-2}
电子仅在$x,y$平面运动而电磁场在三维，求长波二维等离子体色散并证明$\omega_p\propto\sqrt{|q|}$。
\end{exercise}
\begin{solution}
平面库仑Fourier为$v_q=e^2/(2\epsilon_0\epsilon_r|q|)$；惯性密度响应高频为$\chi_0\simeq n_{2D}q^2/(m\omega^2)$，自洽条件$1-v_q\chi_0=0$给$\omega_p^2=n_{2D}e^2|q|/(2m\epsilon_0\epsilon_r)$。若错误把场也限成二维使势Fourier不同，就会改变色散。该式忽略迟滞、有限厚度及强Landau阻尼。
\end{solution}

\originalchapter{8.512作业4--8：电导、无序与超导}
\psource{Patrick A. Lee，8.512 Spring 2009，按网页作业编号；PDF04内部底稿写作业5、PDF05内部底稿写作业4，修订标记恢复网页编号。}
\originalsection{作业4（修订为2009年3月9日）}
\begin{exercise}[原题1(a)--(c)]\label{p:512-4-1}
电子气受$E(r,t)=E_0\e^{\ii(qr-\omega t)}$，弛豫时间Boltzmann方程$\partial_tf+v\cdot\nabla_rf-eE\cdot\nabla_pf=-(f-f_0)/\tau$。(a)取$f=f_0+\Phi(k)\e^{\ii(qr-\omega t)}$，一阶求沿$\hat e\parallel E_0$的电导。(b)由光子吸收与发射的能量耗散导出Kubo谱表达，并用Kramers--Kronig恢复复电导。(c)假设电子态是平面波，$|q|\ll k_F$且展宽$\eta=\hbar/\tau$，说明谱表达如何化为(a)。
\end{exercise}
\begin{solution}
(a)一阶式$[\tau^{-1}-\ii\omega+\ii q\cdot v]\Phi=e(E_0\cdot v)f_0'$，再用$j=-2e\int v\Phi\dd^3k/(2\pi)^3$，得
\begin{equation}
\sigma(q,\omega)=2e^2\int\frac{\dd^3k}{(2\pi)^3}\frac{\tau(\hat e\cdot v_k)^2[-f_0'(\epsilon_k)]}{1-\ii\tau(\omega-q\cdot v_k)}.
\end{equation}
(b)以实场$E_0\cos\omega t$，$A$幅度$E_0/\omega$，跃迁率的正频项由Golden rule给出。吸收减发射使权重$f_a-f_b$，平均耗散$\sigma'E_0^2/2$给
\begin{equation}
\sigma'(q,\omega)=\frac{2\pi e^2}{m^2V\omega}\sum_{ab}(f_a-f_b)|\avg{b|\e^{\ii qr}\hat e\cdot p|a}|^2\delta(E_b-E_a-\hbar\omega).
\end{equation}
这里2为自旋，$\pi\delta$采用通常归一。将$\pi\delta(x)$视为$\operatorname{Im}[-(x+\ii0)^{-1}]$，恢复满足因果性且含合适零频项的复响应；顺、抗磁两部分的结合见第\ref{ch:disorder}章。若用原题将$\pi$吸入谱$\delta$的记号，两式必须同时换约定。

(c)矩阵元要求$k_b=k_a+q$，小$q$下$\Delta E=\hbar q\cdot v$、$f_k-f_{k+q}=-\hbar q\cdot v f_0'$。将谱极点展宽、考虑$\omega=0$所需的静态抵消项后，得到$\tau/[1-\ii\tau(\omega-q\cdot v)]$。严格任意纵向$q$的守恒碰撞核还需局域化学势修正；原题的单$\tau$模型在此按给定近似使用。
\end{solution}
\begin{exercise}[原题2(a)--(b)]\label{p:512-4-2}
各向同性低温金属中取$E\perp q$。(a)从上式证明$\sigma_\perp=\sigma_0(1-\ii\omega\tau)^{-1}(3/4)\int_{-1}^1(1-x^2)/(1+sx)\dd x$，$s=\ii qv_F\tau/(1-\ii\omega\tau)$，$\sigma_0=ne^2\tau/m$。(b)求小$|s|$和大$|s|$的极限，并用粒子运动解释低频大$q\ell$时$1/(q\ell)$压低因子。
\end{exercise}
\begin{solution}
(a)令$q\parallel z,E\parallel x$，$v_x^2=v_F^2(1-x^2)\cos^2\phi$，对方位角积分$\pi$，径向积分由$-f'=\delta(E-E_F)$固定费米面。用$n/m=D(E_F)v_F^2/3$，得到所给归一。

(b)小$s$展开$1/(1+sx)$，奇次积分为0，$\int(1-x^2)\dd x=4/3$、$\int(x^2-x^4)\dd x=4/15$，故$\sigma_\perp=\sigma_0[1+s^2/5+O(s^4)]/(1-\ii\omega\tau)$。低频$s=\ii Q,Q=q\ell\gg1$时，实积分$\int(1-x^2)/(1+Q^2x^2)\dd x\simeq\pi/Q$，所以$\sigma_\perp\simeq3\pi\sigma_0/(4q\ell)$。只有$|v_z|/v_F\lesssim1/(q\ell)$的电子在一段平均自由程内不跨越许多反向电场周期，其速度主要平行电场；这一薄角带占总角空间$O(1/q\ell)$，给相同压低尺度。任意复杂$s$的大值展开应保留解析支，不能直接将这个实低频极限推广到所有频率。
\end{solution}
\originalsection{作业5（网页编号5，底稿题页写2008年3月10日）}
\begin{exercise}[原题1(a)--(f)，(f)选做]\label{p:512-5-1}
一维紧束缚跃迁$t$，随机站势$V_i\sim U[-W,W]$。原题把站$N$与站1认作同一站，所以独立站数$M=N-1$。(a)写周期$M\times M$矩阵。(b)改为$\psi_N=\e^{\ii\phi}\psi_1$。(c)证明与每键跃迁$t\e^{\ii\phi/M}$的周期模型等价，求穿环磁通。(d)取$W/t=2,N=20$，数值画零能附近10级随$\phi$变化，增大$N$比较。(e)求Thouless量$g=E_T/\delta$，$E_T=\partial_\phi^2E_\alpha|_0$，按附近10级与多随机样本平均，检查$N$依赖。(f)选做二维方格，比较$W/t=4,9$。
\end{exercise}
\begin{solution}
(a)$H_{ii}=V_i$、$H_{i,i+1}=H_{i+1,i}=t$，角元$H_{M,1}=H_{1,M}=t$。(b)取$H_{M,1}=t\e^{\ii\phi},H_{1,M}=t\e^{-\ii\phi}$，总绕环相位$\phi$。(c)局域变换$\psi_j=\e^{\ii j\phi/M}\widetilde\psi_j$把相位均匀分配到每键，净相位不变，$\Phi/\Phi_0=\phi/(2\pi)$，此处电子通量量子$\Phi_0=hc/e$。

(d)--(e)随卷脚本固定种子851205，直接对Hermitian矩阵本征值求解，取$N=20,40,80$等尺寸及多个独立样本。按能量选零附近10级，曲率用$[E(\delta\phi)-2E(0)+E(-\delta\phi)]/\delta\phi^2$，另用更小步长检查稳定性。能级曲率有正负；若直接有符号平均可因相消趋零，通常的正电导尺度须平均$|E''|$或RMS。证据同时保留有符号值与$|E''|$，$\delta$由所选能级跨度除9估计，$g=\avg{|E''|}/\avg\delta$。图和数据为该有限样本实验，局域化趋势由$g$随$L$衰减判断，单一$N=20$不能精确给无限系统指数。

(f)二维采用$L_x\times L_y$环面，一方向穿相位，另一方向周期；每站4邻居且每键Hermitian配对。源码计算多个小尺寸在$W/t=4,9$的曲率分布与$g$，强无序边界敏感性更弱。有限尺寸对照支持方法复现，不冒称数值证明二维无穷尺寸局域化。
\end{solution}
\begin{center}\begin{tikzpicture}\begin{axis}[width=11cm,height=5.7cm,xlabel={$\phi$},ylabel={$E/t$},title={一维无序环：$N=20$，10个能级},]
\addplot[thick,color=structurecolor] coordinates {(-3.1415927,-1.2419138) (-3.0630528,-1.2419166) (-2.984513,-1.2419247) (-2.9059732,-1.2419381) (-2.8274334,-1.2419568) (-2.7488936,-1.2419807) (-2.6703538,-1.2420095) (-2.5918139,-1.2420432) (-2.5132741,-1.2420815) (-2.4347343,-1.2421241) (-2.3561945,-1.2421708) (-2.2776547,-1.2422213) (-2.1991149,-1.2422753) (-2.120575,-1.2423324) (-2.0420352,-1.2423923) (-1.9634954,-1.2424547) (-1.8849556,-1.242519) (-1.8064158,-1.242585) (-1.727876,-1.2426522) (-1.6493361,-1.2427202) (-1.5707963,-1.2427885) (-1.4922565,-1.2428569) (-1.4137167,-1.2429247) (-1.3351769,-1.2429917) (-1.2566371,-1.2430574) (-1.1780972,-1.2431214) (-1.0995574,-1.2431833) (-1.0210176,-1.2432427) (-0.9424778,-1.2432993) (-0.86393798,-1.2433528) (-0.78539816,-1.2434027) (-0.70685835,-1.2434489) (-0.62831853,-1.243491) (-0.54977871,-1.2435287) (-0.4712389,-1.2435619) (-0.39269908,-1.2435903) (-0.31415927,-1.2436138) (-0.23561945,-1.2436323) (-0.15707963,-1.2436455) (-0.078539816,-1.2436535) (0,-1.2436561) (0.078539816,-1.2436535) (0.15707963,-1.2436455) (0.23561945,-1.2436323) (0.31415927,-1.2436138) (0.39269908,-1.2435903) (0.4712389,-1.2435619) (0.54977871,-1.2435287) (0.62831853,-1.243491) (0.70685835,-1.2434489) (0.78539816,-1.2434027) (0.86393798,-1.2433528) (0.9424778,-1.2432993) (1.0210176,-1.2432427) (1.0995574,-1.2431833) (1.1780972,-1.2431214) (1.2566371,-1.2430574) (1.3351769,-1.2429917) (1.4137167,-1.2429247) (1.4922565,-1.2428569) (1.5707963,-1.2427885) (1.6493361,-1.2427202) (1.727876,-1.2426522) (1.8064158,-1.242585) (1.8849556,-1.242519) (1.9634954,-1.2424547) (2.0420352,-1.2423923) (2.120575,-1.2423324) (2.1991149,-1.2422753) (2.2776547,-1.2422213) (2.3561945,-1.2421708) (2.4347343,-1.2421241) (2.5132741,-1.2420815) (2.5918139,-1.2420432) (2.6703538,-1.2420095) (2.7488936,-1.2419807) (2.8274334,-1.2419568) (2.9059732,-1.2419381) (2.984513,-1.2419247) (3.0630528,-1.2419166) (3.1415927,-1.2419138)};
\addplot[thick,color=main] coordinates {(-3.1415927,-0.90317756) (-3.0630528,-0.90316488) (-2.984513,-0.90312691) (-2.9059732,-0.90306385) (-2.8274334,-0.90297604) (-2.7488936,-0.90286394) (-2.6703538,-0.90272814) (-2.5918139,-0.90256939) (-2.5132741,-0.90238852) (-2.4347343,-0.90218652) (-2.3561945,-0.9019645) (-2.2776547,-0.90172365) (-2.1991149,-0.90146532) (-2.120575,-0.90119095) (-2.0420352,-0.90090206) (-1.9634954,-0.90060031) (-1.8849556,-0.90028742) (-1.8064158,-0.8999652) (-1.727876,-0.89963555) (-1.6493361,-0.8993004) (-1.5707963,-0.89896177) (-1.4922565,-0.89862172) (-1.4137167,-0.89828233) (-1.3351769,-0.89794572) (-1.2566371,-0.89761401) (-1.1780972,-0.89728931) (-1.0995574,-0.89697373) (-1.0210176,-0.89666933) (-0.9424778,-0.89637813) (-0.86393798,-0.89610208) (-0.78539816,-0.89584306) (-0.70685835,-0.89560285) (-0.62831853,-0.89538309) (-0.54977871,-0.89518535) (-0.4712389,-0.89501101) (-0.39269908,-0.8948613) (-0.31415927,-0.89473732) (-0.23561945,-0.89463994) (-0.15707963,-0.89456987) (-0.078539816,-0.89452762) (0,-0.89451351) (0.078539816,-0.89452762) (0.15707963,-0.89456987) (0.23561945,-0.89463994) (0.31415927,-0.89473732) (0.39269908,-0.8948613) (0.4712389,-0.89501101) (0.54977871,-0.89518535) (0.62831853,-0.89538309) (0.70685835,-0.89560285) (0.78539816,-0.89584306) (0.86393798,-0.89610208) (0.9424778,-0.89637813) (1.0210176,-0.89666933) (1.0995574,-0.89697373) (1.1780972,-0.89728931) (1.2566371,-0.89761401) (1.3351769,-0.89794572) (1.4137167,-0.89828233) (1.4922565,-0.89862172) (1.5707963,-0.89896177) (1.6493361,-0.8993004) (1.727876,-0.89963555) (1.8064158,-0.8999652) (1.8849556,-0.90028742) (1.9634954,-0.90060031) (2.0420352,-0.90090206) (2.120575,-0.90119095) (2.1991149,-0.90146532) (2.2776547,-0.90172365) (2.3561945,-0.9019645) (2.4347343,-0.90218652) (2.5132741,-0.90238852) (2.5918139,-0.90256939) (2.6703538,-0.90272814) (2.7488936,-0.90286394) (2.8274334,-0.90297604) (2.9059732,-0.90306385) (2.984513,-0.90312691) (3.0630528,-0.90316488) (3.1415927,-0.90317756)};
\addplot[thick,color=blue] coordinates {(-3.1415927,-0.7379716) (-3.0630528,-0.73798646) (-2.984513,-0.73803096) (-2.9059732,-0.73810485) (-2.8274334,-0.73820773) (-2.7488936,-0.73833903) (-2.6703538,-0.73849804) (-2.5918139,-0.73868387) (-2.5132741,-0.7388955) (-2.4347343,-0.73913175) (-2.3561945,-0.7393913) (-2.2776547,-0.7396727) (-2.1991149,-0.73997436) (-2.120575,-0.74029456) (-2.0420352,-0.74063147) (-1.9634954,-0.74098315) (-1.8849556,-0.74134756) (-1.8064158,-0.74172255) (-1.727876,-0.74210593) (-1.6493361,-0.7424954) (-1.5707963,-0.7428886) (-1.4922565,-0.74328316) (-1.4137167,-0.74367665) (-1.3351769,-0.74406662) (-1.2566371,-0.74445063) (-1.1780972,-0.74482624) (-1.0995574,-0.74519104) (-1.0210176,-0.74554268) (-0.9424778,-0.74587884) (-0.86393798,-0.7461973) (-0.78539816,-0.74649595) (-0.70685835,-0.74677277) (-0.62831853,-0.74702586) (-0.54977871,-0.74725351) (-0.4712389,-0.74745414) (-0.39269908,-0.74762634) (-0.31415927,-0.74776893) (-0.23561945,-0.74788088) (-0.15707963,-0.74796143) (-0.078539816,-0.74800998) (0,-0.7480262) (0.078539816,-0.74800998) (0.15707963,-0.74796143) (0.23561945,-0.74788088) (0.31415927,-0.74776893) (0.39269908,-0.74762634) (0.4712389,-0.74745414) (0.54977871,-0.74725351) (0.62831853,-0.74702586) (0.70685835,-0.74677277) (0.78539816,-0.74649595) (0.86393798,-0.7461973) (0.9424778,-0.74587884) (1.0210176,-0.74554268) (1.0995574,-0.74519104) (1.1780972,-0.74482624) (1.2566371,-0.74445063) (1.3351769,-0.74406662) (1.4137167,-0.74367665) (1.4922565,-0.74328316) (1.5707963,-0.7428886) (1.6493361,-0.7424954) (1.727876,-0.74210593) (1.8064158,-0.74172255) (1.8849556,-0.74134756) (1.9634954,-0.74098315) (2.0420352,-0.74063147) (2.120575,-0.74029456) (2.1991149,-0.73997436) (2.2776547,-0.7396727) (2.3561945,-0.7393913) (2.4347343,-0.73913175) (2.5132741,-0.7388955) (2.5918139,-0.73868387) (2.6703538,-0.73849804) (2.7488936,-0.73833903) (2.8274334,-0.73820773) (2.9059732,-0.73810485) (2.984513,-0.73803096) (3.0630528,-0.73798646) (3.1415927,-0.7379716)};
\addplot[thick,color=gray] coordinates {(-3.1415927,-0.16642575) (-3.0630528,-0.16640191) (-2.984513,-0.16633051) (-2.9059732,-0.1662119) (-2.8274334,-0.16604666) (-2.7488936,-0.16583563) (-2.6703538,-0.16557984) (-2.5918139,-0.16528058) (-2.5132741,-0.16493936) (-2.4347343,-0.16455791) (-2.3561945,-0.16413819) (-2.2776547,-0.16368236) (-2.1991149,-0.16319281) (-2.120575,-0.16267213) (-2.0420352,-0.16212311) (-1.9634954,-0.16154872) (-1.8849556,-0.16095213) (-1.8064158,-0.16033668) (-1.727876,-0.15970586) (-1.6493361,-0.15906332) (-1.5707963,-0.15841284) (-1.4922565,-0.1577583) (-1.4137167,-0.15710371) (-1.3351769,-0.15645313) (-1.2566371,-0.15581067) (-1.1780972,-0.15518049) (-1.0995574,-0.15456674) (-1.0210176,-0.15397353) (-0.9424778,-0.15340491) (-0.86393798,-0.15286484) (-0.78539816,-0.15235716) (-0.70685835,-0.1518855) (-0.62831853,-0.15145333) (-0.54977871,-0.15106386) (-0.4712389,-0.15072001) (-0.39269908,-0.15042442) (-0.31415927,-0.15017935) (-0.23561945,-0.14998672) (-0.15707963,-0.14984802) (-0.078539816,-0.14976437) (0,-0.14973641) (0.078539816,-0.14976437) (0.15707963,-0.14984802) (0.23561945,-0.14998672) (0.31415927,-0.15017935) (0.39269908,-0.15042442) (0.4712389,-0.15072001) (0.54977871,-0.15106386) (0.62831853,-0.15145333) (0.70685835,-0.1518855) (0.78539816,-0.15235716) (0.86393798,-0.15286484) (0.9424778,-0.15340491) (1.0210176,-0.15397353) (1.0995574,-0.15456674) (1.1780972,-0.15518049) (1.2566371,-0.15581067) (1.3351769,-0.15645313) (1.4137167,-0.15710371) (1.4922565,-0.1577583) (1.5707963,-0.15841284) (1.6493361,-0.15906332) (1.727876,-0.15970586) (1.8064158,-0.16033668) (1.8849556,-0.16095213) (1.9634954,-0.16154872) (2.0420352,-0.16212311) (2.120575,-0.16267213) (2.1991149,-0.16319281) (2.2776547,-0.16368236) (2.3561945,-0.16413819) (2.4347343,-0.16455791) (2.5132741,-0.16493936) (2.5918139,-0.16528058) (2.6703538,-0.16557984) (2.7488936,-0.16583563) (2.8274334,-0.16604666) (2.9059732,-0.1662119) (2.984513,-0.16633051) (3.0630528,-0.16640191) (3.1415927,-0.16642575)};
\addplot[thick,color=orange] coordinates {(-3.1415927,0.1822732) (-3.0630528,0.18148306) (-2.984513,0.1795113) (-2.9059732,0.17694124) (-2.8274334,0.17409312) (-2.7488936,0.17111178) (-2.6703538,0.16806639) (-2.5918139,0.16499386) (-2.5132741,0.16191624) (-2.4347343,0.15884816) (-2.3561945,0.15580027) (-2.2776547,0.15278097) (-2.1991149,0.14979732) (-2.120575,0.14685554) (-2.0420352,0.14396134) (-1.9634954,0.1411201) (-1.8849556,0.13833698) (-1.8064158,0.135617) (-1.727876,0.1329651) (-1.6493361,0.13038616) (-1.5707963,0.12788501) (-1.4922565,0.12546647) (-1.4137167,0.1231353) (-1.3351769,0.12089625) (-1.2566371,0.11875404) (-1.1780972,0.11671329) (-1.0995574,0.11477859) (-1.0210176,0.11295442) (-0.9424778,0.11124513) (-0.86393798,0.10965494) (-0.78539816,0.10818789) (-0.70685835,0.10684782) (-0.62831853,0.10563834) (-0.54977871,0.10456276) (-0.4712389,0.10362413) (-0.39269908,0.10282513) (-0.31415927,0.10216812) (-0.23561945,0.10165504) (-0.15707963,0.10128743) (-0.078539816,0.10106642) (0,0.10099267) (0.078539816,0.10106642) (0.15707963,0.10128743) (0.23561945,0.10165504) (0.31415927,0.10216812) (0.39269908,0.10282513) (0.4712389,0.10362413) (0.54977871,0.10456276) (0.62831853,0.10563834) (0.70685835,0.10684782) (0.78539816,0.10818789) (0.86393798,0.10965494) (0.9424778,0.11124513) (1.0210176,0.11295442) (1.0995574,0.11477859) (1.1780972,0.11671329) (1.2566371,0.11875404) (1.3351769,0.12089625) (1.4137167,0.1231353) (1.4922565,0.12546647) (1.5707963,0.12788501) (1.6493361,0.13038616) (1.727876,0.1329651) (1.8064158,0.135617) (1.8849556,0.13833698) (1.9634954,0.1411201) (2.0420352,0.14396134) (2.120575,0.14685554) (2.1991149,0.14979732) (2.2776547,0.15278097) (2.3561945,0.15580027) (2.4347343,0.15884816) (2.5132741,0.16191624) (2.5918139,0.16499386) (2.6703538,0.16806639) (2.7488936,0.17111178) (2.8274334,0.17409312) (2.9059732,0.17694124) (2.984513,0.1795113) (3.0630528,0.18148306) (3.1415927,0.1822732)};
\addplot[thick,color=teal] coordinates {(-3.1415927,0.19531149) (-3.0630528,0.19610381) (-2.984513,0.19808209) (-2.9059732,0.20066302) (-2.8274334,0.20352633) (-2.7488936,0.20652716) (-2.6703538,0.20959631) (-2.5918139,0.21269684) (-2.5132741,0.21580665) (-2.4347343,0.21891104) (-2.3561945,0.22199928) (-2.2776547,0.2250629) (-2.1991149,0.22809473) (-2.120575,0.23108841) (-2.0420352,0.23403809) (-1.9634954,0.23693823) (-1.8849556,0.23978345) (-1.8064158,0.24256851) (-1.727876,0.24528821) (-1.6493361,0.24793737) (-1.5707963,0.25051083) (-1.4922565,0.25300341) (-1.4137167,0.25540992) (-1.3351769,0.25772517) (-1.2566371,0.25994396) (-1.1780972,0.26206112) (-1.0995574,0.26407151) (-1.0210176,0.26597007) (-0.9424778,0.26775181) (-0.86393798,0.26941187) (-0.78539816,0.27094557) (-0.70685835,0.27234841) (-0.62831853,0.27361615) (-0.54977871,0.27474484) (-0.4712389,0.27573086) (-0.39269908,0.27657096) (-0.31415927,0.27726231) (-0.23561945,0.27780256) (-0.15707963,0.27818982) (-0.078539816,0.27842272) (0,0.27850045) (0.078539816,0.27842272) (0.15707963,0.27818982) (0.23561945,0.27780256) (0.31415927,0.27726231) (0.39269908,0.27657096) (0.4712389,0.27573086) (0.54977871,0.27474484) (0.62831853,0.27361615) (0.70685835,0.27234841) (0.78539816,0.27094557) (0.86393798,0.26941187) (0.9424778,0.26775181) (1.0210176,0.26597007) (1.0995574,0.26407151) (1.1780972,0.26206112) (1.2566371,0.25994396) (1.3351769,0.25772517) (1.4137167,0.25540992) (1.4922565,0.25300341) (1.5707963,0.25051083) (1.6493361,0.24793737) (1.727876,0.24528821) (1.8064158,0.24256851) (1.8849556,0.23978345) (1.9634954,0.23693823) (2.0420352,0.23403809) (2.120575,0.23108841) (2.1991149,0.22809473) (2.2776547,0.2250629) (2.3561945,0.22199928) (2.4347343,0.21891104) (2.5132741,0.21580665) (2.5918139,0.21269684) (2.6703538,0.20959631) (2.7488936,0.20652716) (2.8274334,0.20352633) (2.9059732,0.20066302) (2.984513,0.19808209) (3.0630528,0.19610381) (3.1415927,0.19531149)};
\addplot[thick,color=black] coordinates {(-3.1415927,0.51642451) (-3.0630528,0.51639826) (-2.984513,0.51631965) (-2.9059732,0.51618903) (-2.8274334,0.51600701) (-2.7488936,0.51577446) (-2.6703538,0.51549247) (-2.5918139,0.51516238) (-2.5132741,0.51478576) (-2.4347343,0.51436444) (-2.3561945,0.51390047) (-2.2776547,0.51339614) (-2.1991149,0.51285397) (-2.120575,0.51227671) (-2.0420352,0.51166732) (-1.9634954,0.51102899) (-1.8849556,0.51036513) (-1.8064158,0.50967932) (-1.727876,0.50897537) (-1.6493361,0.50825725) (-1.5707963,0.5075291) (-1.4922565,0.50679522) (-1.4137167,0.50606004) (-1.3351769,0.50532812) (-1.2566371,0.5046041) (-1.1780972,0.50389266) (-1.0995574,0.50319856) (-1.0210176,0.50252653) (-0.9424778,0.50188126) (-0.86393798,0.50126738) (-0.78539816,0.50068937) (-0.70685835,0.50015157) (-0.62831853,0.49965809) (-0.54977871,0.49921277) (-0.4712389,0.49881915) (-0.39269908,0.49848041) (-0.31415927,0.49819931) (-0.23561945,0.4979782) (-0.15707963,0.49781891) (-0.078539816,0.4977228) (0,0.49769067) (0.078539816,0.4977228) (0.15707963,0.49781891) (0.23561945,0.4979782) (0.31415927,0.49819931) (0.39269908,0.49848041) (0.4712389,0.49881915) (0.54977871,0.49921277) (0.62831853,0.49965809) (0.70685835,0.50015157) (0.78539816,0.50068937) (0.86393798,0.50126738) (0.9424778,0.50188126) (1.0210176,0.50252653) (1.0995574,0.50319856) (1.1780972,0.50389266) (1.2566371,0.5046041) (1.3351769,0.50532812) (1.4137167,0.50606004) (1.4922565,0.50679522) (1.5707963,0.5075291) (1.6493361,0.50825725) (1.727876,0.50897537) (1.8064158,0.50967932) (1.8849556,0.51036513) (1.9634954,0.51102899) (2.0420352,0.51166732) (2.120575,0.51227671) (2.1991149,0.51285397) (2.2776547,0.51339614) (2.3561945,0.51390047) (2.4347343,0.51436444) (2.5132741,0.51478576) (2.5918139,0.51516238) (2.6703538,0.51549247) (2.7488936,0.51577446) (2.8274334,0.51600701) (2.9059732,0.51618903) (2.984513,0.51631965) (3.0630528,0.51639826) (3.1415927,0.51642451)};
\addplot[thick,color=structurecolor] coordinates {(-3.1415927,0.92147543) (-3.0630528,0.9214817) (-2.984513,0.92150047) (-2.9059732,0.92153163) (-2.8274334,0.92157496) (-2.7488936,0.9216302) (-2.6703538,0.921697) (-2.5918139,0.92177493) (-2.5132741,0.92186348) (-2.4347343,0.92196212) (-2.3561945,0.92207019) (-2.2776547,0.92218704) (-2.1991149,0.9223119) (-2.120575,0.92244401) (-2.0420352,0.92258253) (-1.9634954,0.92272659) (-1.8849556,0.92287529) (-1.8064158,0.9230277) (-1.727876,0.92318288) (-1.6493361,0.92333986) (-1.5707963,0.92349767) (-1.4922565,0.92365534) (-1.4137167,0.92381189) (-1.3351769,0.92396638) (-1.2566371,0.92411784) (-1.1780972,0.92426536) (-1.0995574,0.92440804) (-1.0210176,0.92454501) (-0.9424778,0.92467543) (-0.86393798,0.92479854) (-0.78539816,0.92491357) (-0.70685835,0.92501983) (-0.62831853,0.9251167) (-0.54977871,0.92520358) (-0.4712389,0.92527996) (-0.39269908,0.92534537) (-0.31415927,0.92539943) (-0.23561945,0.92544182) (-0.15707963,0.92547228) (-0.078539816,0.92549063) (0,0.92549675) (0.078539816,0.92549063) (0.15707963,0.92547228) (0.23561945,0.92544182) (0.31415927,0.92539943) (0.39269908,0.92534537) (0.4712389,0.92527996) (0.54977871,0.92520358) (0.62831853,0.9251167) (0.70685835,0.92501983) (0.78539816,0.92491357) (0.86393798,0.92479854) (0.9424778,0.92467543) (1.0210176,0.92454501) (1.0995574,0.92440804) (1.1780972,0.92426536) (1.2566371,0.92411784) (1.3351769,0.92396638) (1.4137167,0.92381189) (1.4922565,0.92365534) (1.5707963,0.92349767) (1.6493361,0.92333986) (1.727876,0.92318288) (1.8064158,0.9230277) (1.8849556,0.92287529) (1.9634954,0.92272659) (2.0420352,0.92258253) (2.120575,0.92244401) (2.1991149,0.9223119) (2.2776547,0.92218704) (2.3561945,0.92207019) (2.4347343,0.92196212) (2.5132741,0.92186348) (2.5918139,0.92177493) (2.6703538,0.921697) (2.7488936,0.9216302) (2.8274334,0.92157496) (2.9059732,0.92153163) (2.984513,0.92150047) (3.0630528,0.9214817) (3.1415927,0.92147543)};
\addplot[thick,color=main] coordinates {(-3.1415927,1.3685779) (-3.0630528,1.3685752) (-2.984513,1.368567) (-2.9059732,1.3685535) (-2.8274334,1.3685347) (-2.7488936,1.3685107) (-2.6703538,1.3684817) (-2.5918139,1.3684478) (-2.5132741,1.3684093) (-2.4347343,1.3683664) (-2.3561945,1.3683193) (-2.2776547,1.3682684) (-2.1991149,1.3682139) (-2.120575,1.3681563) (-2.0420352,1.3680958) (-1.9634954,1.3680328) (-1.8849556,1.3679678) (-1.8064158,1.3679011) (-1.727876,1.367833) (-1.6493361,1.3677642) (-1.5707963,1.3676949) (-1.4922565,1.3676256) (-1.4137167,1.3675567) (-1.3351769,1.3674887) (-1.2566371,1.367422) (-1.1780972,1.3673569) (-1.0995574,1.3672939) (-1.0210176,1.3672334) (-0.9424778,1.3671757) (-0.86393798,1.3671213) (-0.78539816,1.3670703) (-0.70685835,1.3670233) (-0.62831853,1.3669803) (-0.54977871,1.3669418) (-0.4712389,1.3669079) (-0.39269908,1.3668789) (-0.31415927,1.3668548) (-0.23561945,1.366836) (-0.15707963,1.3668225) (-0.078539816,1.3668143) (0,1.3668116) (0.078539816,1.3668143) (0.15707963,1.3668225) (0.23561945,1.366836) (0.31415927,1.3668548) (0.39269908,1.3668789) (0.4712389,1.3669079) (0.54977871,1.3669418) (0.62831853,1.3669803) (0.70685835,1.3670233) (0.78539816,1.3670703) (0.86393798,1.3671213) (0.9424778,1.3671757) (1.0210176,1.3672334) (1.0995574,1.3672939) (1.1780972,1.3673569) (1.2566371,1.367422) (1.3351769,1.3674887) (1.4137167,1.3675567) (1.4922565,1.3676256) (1.5707963,1.3676949) (1.6493361,1.3677642) (1.727876,1.367833) (1.8064158,1.3679011) (1.8849556,1.3679678) (1.9634954,1.3680328) (2.0420352,1.3680958) (2.120575,1.3681563) (2.1991149,1.3682139) (2.2776547,1.3682684) (2.3561945,1.3683193) (2.4347343,1.3683664) (2.5132741,1.3684093) (2.5918139,1.3684478) (2.6703538,1.3684817) (2.7488936,1.3685107) (2.8274334,1.3685347) (2.9059732,1.3685535) (2.984513,1.368567) (3.0630528,1.3685752) (3.1415927,1.3685779)};
\addplot[thick,color=blue] coordinates {(-3.1415927,1.5897831) (-3.0630528,1.5897844) (-2.984513,1.5897883) (-2.9059732,1.5897948) (-2.8274334,1.5898039) (-2.7488936,1.5898155) (-2.6703538,1.5898294) (-2.5918139,1.5898458) (-2.5132741,1.5898643) (-2.4347343,1.589885) (-2.3561945,1.5899076) (-2.2776547,1.5899321) (-2.1991149,1.5899583) (-2.120575,1.589986) (-2.0420352,1.590015) (-1.9634954,1.5900453) (-1.8849556,1.5900765) (-1.8064158,1.5901085) (-1.727876,1.5901411) (-1.6493361,1.5901742) (-1.5707963,1.5902074) (-1.4922565,1.5902406) (-1.4137167,1.5902735) (-1.3351769,1.5903061) (-1.2566371,1.590338) (-1.1780972,1.5903691) (-1.0995574,1.5903992) (-1.0210176,1.5904281) (-0.9424778,1.5904557) (-0.86393798,1.5904817) (-0.78539816,1.590506) (-0.70685835,1.5905284) (-0.62831853,1.5905489) (-0.54977871,1.5905673) (-0.4712389,1.5905835) (-0.39269908,1.5905973) (-0.31415927,1.5906087) (-0.23561945,1.5906177) (-0.15707963,1.5906242) (-0.078539816,1.590628) (0,1.5906293) (0.078539816,1.590628) (0.15707963,1.5906242) (0.23561945,1.5906177) (0.31415927,1.5906087) (0.39269908,1.5905973) (0.4712389,1.5905835) (0.54977871,1.5905673) (0.62831853,1.5905489) (0.70685835,1.5905284) (0.78539816,1.590506) (0.86393798,1.5904817) (0.9424778,1.5904557) (1.0210176,1.5904281) (1.0995574,1.5903992) (1.1780972,1.5903691) (1.2566371,1.590338) (1.3351769,1.5903061) (1.4137167,1.5902735) (1.4922565,1.5902406) (1.5707963,1.5902074) (1.6493361,1.5901742) (1.727876,1.5901411) (1.8064158,1.5901085) (1.8849556,1.5900765) (1.9634954,1.5900453) (2.0420352,1.590015) (2.120575,1.589986) (2.1991149,1.5899583) (2.2776547,1.5899321) (2.3561945,1.5899076) (2.4347343,1.589885) (2.5132741,1.5898643) (2.5918139,1.5898458) (2.6703538,1.5898294) (2.7488936,1.5898155) (2.8274334,1.5898039) (2.9059732,1.5897948) (2.984513,1.5897883) (3.0630528,1.5897844) (3.1415927,1.5897831)};
\end{axis}\end{tikzpicture}\end{center}

\begin{center}\begin{tikzpicture}\begin{axis}[width=11cm,height=5.7cm,xlabel={$N$},ylabel={$g$},title={有限样本Thouless电导尺度},ymode=log,mark=*]
\addplot[thick,color=structurecolor] coordinates {(20,0.033114471) (40,0.0060952993) (80,7.5557181e-05) (140,2.5499047e-07)};
\addlegendentry{$W/t=2$}
\end{axis}\end{tikzpicture}\end{center}

\begin{center}\begin{tabular}{rrrr}\toprule $N$&独立站数&样本数&$g$\\\midrule
20&19&40&$0.03311$\\
40&39&40&$0.006095$\\
80&79&40&$7.56\times10^{-5}$\\\bottomrule\end{tabular}\end{center}
大尺寸的极小曲率逐渐接近浮点分辨率；图中更大尺寸只作带步长比较的数值边界记录，不用于拟合精确局域长度。
\begin{center}\begin{tikzpicture}\begin{axis}[width=11cm,height=5.7cm,xlabel={$L$},ylabel={$g$},title={二维有限尺寸对照},ymode=log,mark=*]
\addplot[thick,color=structurecolor] coordinates {(4,0.25939206) (6,0.13353791) (8,0.099962056)};
\addlegendentry{$W/t=4$}
\addplot[thick,color=main] coordinates {(4,0.023333335) (6,0.012444557) (8,0.0062368641)};
\addlegendentry{$W/t=9$}
\end{axis}\end{tikzpicture}\end{center}

\originalsection{作业6（底稿截止2008年3月31日）}
\begin{exercise}[原题1(a)--(b)]\label{p:512-6-1}
定义$J_\mu=-K_{\mu\nu}A_\nu$，将$K$分为横向与纵向投影。(a)由线性响应求自由费米气的静态$K_\perp(q,0)$，可取$q\parallel z$算$K_{xx}$。(b)证明双自旋Landau抗磁Gaussian式$\chi_D=-e^2k_F/(12\pi^2mc^2)$，是Pauli值$-1/3$。给原题$\hbar=1$之外的单位恢复。
\end{exercise}
\begin{solution}
(a)顺磁泡图与抗磁项相加，SI制可写
\begin{equation}
K_{xx}(q,0)=ne^2/m+(2e^2\hbar^2/m^2)\int\dd^3k\,k_x^2[f(\epsilon_{k+q/2})-f(\epsilon_{k-q/2})]/[(2\pi)^3(\epsilon_{k+q/2}-\epsilon_{k-q/2})]
\end{equation}
$q=0$项以积分分部$\int k_x\partial_{k_x}f=-\int f$恰抵消。为求$q^2$，令$b=\hbar^2q^2/(8m)$、$a=\hbar^2k_zq/(2m)$，比值为$f'(\epsilon+b)+a^2f'''(\epsilon+b)/6+\cdots$。把$b f''+a^2f'''/6$在球面作角积分，$\avg{k_x^2}=k^2/3$、$\avg{k_x^2k_z^2}=k^4/15$，径向的$\delta$导数积分给
\begin{equation}
K_\perp(q,0)=\frac{e^2k_F}{12\pi^2m}q^2+O(q^4)\qquad(\text{SI电流--}A\text{核}).
\end{equation}
(b)静态磁化电流$j=\nabla\times M$，线性$M=\chi_DB/\mu_0$，而$B=\nabla\times A$，横向Fourier给$j=(\chi_D/\mu_0)q^2A$，故$\chi_D=-\mu_0e^2k_F/(12\pi^2m)$。Pauli值$\chi_P=\mu_0\mu_B^2D(E_F)=\mu_0e^2k_F/(4\pi^2m)$，比值$-1/3$。Gaussian把$\mu_0$的磁化关系换成相应$c^{-2}$得到原式；裸质量与非相互作用假设不可省。
\end{solution}
\originalsection{作业7（2009年4月6日截止）}
\begin{exercise}[原题1(a)--(d)]\label{p:512-7-1}
无序超导体用正常态精确轨道$\phi_\alpha$表示$\int j_x^P=e\sum_{\alpha\beta\sigma}v_{\alpha\beta}c_{\beta\sigma}^\dagger c_{\alpha\sigma}$，$v_{\alpha\beta}=\int\phi_\beta^*p_x\phi_\alpha/m$，$\xi_\alpha=\epsilon_\alpha-\mu$，$E_\alpha=\sqrt{\xi_\alpha^2+\Delta^2}$，$u_\alpha^2=(1+\xi_\alpha/E_\alpha)/2$，$v_\alpha^2=(1-\xi_\alpha/E_\alpha)/2$，$\hbar=1$。(a)证明零温光吸收只创建两个准粒子，权重$(u_\alpha v_\beta-v_\alpha u_\beta)^2|v_{\alpha\beta}|^2\delta(\omega-E_\alpha-E_\beta)$，并以$f(\xi,\xi')=\Omega^{-1}\sum|v_{\alpha\beta}|^2\delta(\xi-\xi_\alpha)\delta(\xi'-\xi_\beta)$写积分。(b)证明相干因子$[1-\xi\xi'/(EE')-\Delta^2/(EE')]/2$。(c)说明近费米面脏极限$f$可近常数，$\Delta\to0$恢复正常电导与自旋因子。(d)化简至Mattis--Bardeen积分、数值画$\sigma_1/\sigma_N$并讨论；要求$\Delta\tau\ll1$。
\end{exercise}
\begin{solution}
(a)Bogoliubov展开电流，作用于零准粒子真空仅留下双创建项；速度矩阵在时间反演下变号，电子与空穴振幅相减，得$\sigma_1=(\pi e^2/\omega)\iint(uv'-vu')^2f(\xi,\xi')\delta(\omega-E-E')\dd\xi\dd\xi'$。自旋及双计数依原归一计入$f$和两创建通道。

(b)展开平方，使用$u^2,v^2$与$uv=\Delta/(2E)$，直接得所列结果。(c)正常态轨道的速度相关能量宽度约$1/\tau$，若$|\xi-\xi'|\ll1/\tau$，脏金属近似可将$f$固定在费米面值。$\Delta=0$时相干因子只保留$\xi,\xi'$异号的初末态，积分长度为$\omega$，恰恢复有限$\sigma_N$。四象限对称求和后$\xi\xi'$项抵消，剩题示$1-\Delta^2/(EE')$。

(d)改变量$\dd\xi=E\dd E/\sqrt{E^2-\Delta^2}$，$\delta$积分给
\begin{equation}
\frac{\sigma_1(\omega)}{\sigma_N}=\frac{\theta(\omega-2\Delta)}{\omega}\int_\Delta^{\omega-\Delta}\frac{E(\omega-E)-\Delta^2}{\sqrt{E^2-\Delta^2}\sqrt{(\omega-E)^2-\Delta^2}}\dd E.
\end{equation}
取$E=\Delta+(\omega-2\Delta)\sin^2\theta$可消端点平方根奇异，随卷脚本实际积分并画图。$\omega<2\Delta$没有零温吸收；阈值处两态密度虽发散，相干因子同时压低，整体从0起；高频在仍满足脏近似窗口内趋近1。零频超流$\delta$项与虚部响应不在这条有限频率吸收曲线中。
\end{solution}
\begin{center}\begin{tikzpicture}\begin{axis}[width=11cm,height=5.7cm,xlabel={$\omega/\Delta$},ylabel={$\sigma_1/\sigma_N$},title={零温脏极限光吸收},xmin=0,ymin=0,ymax=1.05]
\addplot[thick,color=structurecolor] coordinates {(0,0) (0.17272727,0) (0.34545455,0) (0.51818182,0) (0.69090909,0) (0.86363636,0) (1.0363636,0) (1.2090909,0) (1.3818182,0) (1.5545455,0) (1.7272727,0) (1.9,0) (2,0) (2.06,0.045415756) (2.12,0.087636604) (2.18,0.12696447) (2.24,0.16366573) (2.3,0.19797619) (2.36,0.23010529) (2.42,0.26023955) (2.48,0.28854562) (2.54,0.31517274) (2.6,0.34025498) (2.66,0.36391301) (2.72,0.38625572) (2.78,0.40738163) (2.84,0.42738) (2.9,0.44633191) (2.96,0.46431117) (3.02,0.48138505) (3.08,0.49761502) (3.14,0.51305731) (3.2,0.52776348) (3.26,0.54178084) (3.32,0.55515291) (3.38,0.56791976) (3.44,0.58011835) (3.5,0.59178281) (3.56,0.60294469) (3.62,0.61363321) (3.68,0.62387547) (3.74,0.63369662) (3.8,0.64312003) (3.86,0.65216745) (3.92,0.66085913) (3.98,0.66921395) (4.04,0.67724955) (4.1,0.68498239) (4.16,0.69242788) (4.22,0.69960045) (4.28,0.7065136) (4.34,0.71318) (4.4,0.71961157) (4.46,0.72581946) (4.52,0.73181419) (4.58,0.73760565) (4.64,0.74320316) (4.7,0.74861548) (4.76,0.75385089) (4.82,0.75891722) (4.88,0.76382183) (4.94,0.76857172) (5,0.77317346) (5.06,0.77763331) (5.12,0.78195719) (5.18,0.78615068) (5.24,0.79021912) (5.3,0.79416753) (5.36,0.79800072) (5.42,0.80172321) (5.48,0.80533934) (5.54,0.80885321) (5.6,0.81226873) (5.66,0.81558962) (5.72,0.81881943) (5.78,0.82196154) (5.84,0.82501917) (5.9,0.82799538) (5.96,0.8308931) (6.02,0.83371515) (6.08,0.83646418) (6.14,0.83914276) (6.2,0.84175333) (6.26,0.84429821) (6.32,0.84677966) (6.38,0.84919979) (6.44,0.85156067) (6.5,0.85386426) (6.56,0.85611243) (6.62,0.85830699) (6.68,0.86044967) (6.74,0.86254213) (6.8,0.86458595) (6.86,0.86658268) (6.92,0.86853376) (6.98,0.87044063) (7.04,0.87230463) (7.1,0.87412707) (7.16,0.87590919) (7.22,0.87765221) (7.28,0.87935729) (7.34,0.88102554) (7.4,0.88265805) (7.46,0.88425583) (7.52,0.88581991) (7.58,0.88735123) (7.64,0.88885072) (7.7,0.89031929) (7.76,0.89175779) (7.82,0.89316705) (7.88,0.89454789) (7.94,0.89590107) (8,0.89722734) (8.5,0.90732454) (9,0.91595542) (9.5,0.92339547) (10,0.92985773) (10.5,0.93550917) (11,0.94048223) (11.5,0.94488308) (12,0.94879772) (12.5,0.95229638) (13,0.95543696) (13.5,0.9582675) (14,0.96082819) (14.5,0.96315281) (15,0.96527) (15.5,0.96720412) (16,0.96897603) (16.5,0.97060366) (17,0.97210251) (17.5,0.97348605) (18,0.97476599) (18.5,0.97595258) (19,0.97705484) (19.5,0.9780807) (20,0.97903717)};
\end{axis}\end{tikzpicture}\end{center}

\originalsection{作业8（2009年4月13日截止）}
\begin{exercise}[原题1、2(a)--(c)]\label{p:512-8-1}
近$T_c$时Gaussian制GL自由能$F=\Delta f\int\dd r[t|\psi|^2+|\psi|^4/2+\xi^2|(-\ii\nabla+2eA/\hbar c)\psi|^2]$，$t=(T-T_c)/T_c$；原题取$\hbar=1$。(1)求超导--正常自由能差与热力学场$H_c$。(2a)变分求GL方程。(2b)从正常态降场，利用线性化Landau级求$H_{c2}$。(2c)求London深度$\lambda(T)$，以$\xi(T)=\xi/\sqrt{-t}$与$\lambda$表示$H_c$，证明$H_{c2}>H_c$等价$\kappa>1/\sqrt2$。
\end{exercise}
\begin{solution}
均匀极小$|\psi|^2=-t$给$\delta f=-\Delta f\,t^2/2$，故$H_c^2/(8\pi)=\Delta f\,t^2/2$。变分$\xi^2(-\ii\nabla+2eA/\hbar c)^2\psi+t\psi+|\psi|^2\psi=0$。最低磁动能本征值$\xi^2(2eH/\hbar c)$，当与$-t$抵消时正常态首次失稳，$H_{c2}=(-t)\Phi_0/(2\pi\xi^2)=\Phi_0/[2\pi\xi(T)^2]$。

梯度项二阶$A$系数为$\Delta f\xi^2(2e/\hbar c)^2|\psi|^2$，电流变分配Maxwell得$\lambda^{-2}=32\pi e^2\Delta f\xi^2(-t)/(\hbar^2c^2)$。因此$H_c=\Phi_0/[2\sqrt2\pi\lambda(T)\xi(T)]$，$H_{c2}/H_c=\sqrt2\lambda/\xi(T)$，得到临界$\kappa$。此比较是GL近$T_c$判据，不能直接套到温度很低且GL参数强变的任意材料。
\end{solution}

\originalchapter{8.512作业9--12：配对与磁响应}
\psource{Patrick A. Lee，8.512 Spring 2009，公开作业9--12。}
\originalsection{作业9（2009年4月22日截止）}
\begin{exercise}[原题1(a)--(b)]\label{p:512-9-1}
取声子吸引$V_p=-V_0$在$|\omega|<\omega_D$，库仑排斥$N_\uparrow(0)V_c=\mu>0$在$|\omega|<E_F$，原题把频率作能量并取$\hbar=1$。(a)写有限温隙方程，近$T_c$试分段$\Delta_1$（低于$\omega_D$）与$\Delta_2$（高于$\omega_D$）；说明库仑项$A(\xi)$缓变且高能声子项小，导出$\mu^*=\mu/[1+\mu\log(E_F/\omega_D)]$及$1=[N_\uparrow V_0-\mu^*]\log(\omega_D/\kb T_c)$，估计典型金属参数。(b)仅考虑离子质量使$\omega_D$变化，求无库仑与有库仑的同位素系数。
\end{exercise}
\begin{solution}
完整方程$\Delta(\xi)=-N_\uparrow\int V(\xi-\xi')\Delta(\xi')\tanh(E_{\xi'}/2\kb T)/(2E_{\xi'})\dd\xi'$。近$T_c$取$E=|\xi'|$，两侧能量积分合并。宽库仑窗口$E_F$使在小$\xi$范围卷积变化缓慢，可以常数$A(0)$代替；高$\xi$的声子窗口只覆盖窄能段，其贡献无低能对数，故$\Delta_2\simeq A(0)$。令$\lambda=N_\uparrow V_0$、$L_1=\log(\omega_D/\kb T_c)$、$L_2=\log(E_F/\omega_D)$，分段方程为
\begin{equation}
\Delta_1=(\lambda-\mu)L_1\Delta_1-\mu L_2\Delta_2,\qquad\Delta_2=-\mu L_1\Delta_1-\mu L_2\Delta_2.
\end{equation}
第二式给$\Delta_2=-\mu L_1\Delta_1/(1+\mu L_2)$，回代得$1=(\lambda-\mu^*)L_1$。高能对数使有效排斥较小，例如假设$\mu=0.3,E_F/\omega_D=100$，$\mu^*=0.126$；这只是典型参数量级，不替代材料屏蔽势的计算。

(b)声子频率$\omega_D\propto M^{-1/2}$，所以$\delta\log\omega_D=-\delta\log M/2$。无排斥$T_c\propto\omega_D$，同位素指数$\alpha=-\partial\log T_c/\partial\log M=1/2$。有排斥$\partial\mu^*/\partial\log\omega_D=(\mu^*)^2$，$\log T_c=\log\omega_D-1/(\lambda-\mu^*)+\mathrm{const}$，因此$\alpha=\frac12[1-(\mu^*/(\lambda-\mu^*))^2]$。该结果固定$\lambda,E_F,\mu$，不包括结构或电子态随同位素变化。
\end{solution}
\begin{exercise}[原题2]\label{p:512-9-2}
两个类氢原子位于$R_a,R_b$，轨道$\phi_a,\phi_b$，重叠$l=\int\phi_a^*\phi_b\dd r$。定义$\Delta H=e_*^2/R_{ab}+e_*^2/r_{12}-e_*^2/r_{1b}-e_*^2/r_{2a}$，$V=\iint|\phi_a(1)\phi_b(2)|^2\Delta H$，$I=\iint\phi_a^*(1)\phi_b^*(2)\phi_b(1)\phi_a(2)\Delta H$。证明Heitler--London单重、三重能为$E_a+E_b+(V\pm I)/(1\pm|l|^2)$；若轨道实可写$l^2$。
\end{exercise}
\begin{solution}
空间态$\Psi_\pm=[\phi_a(1)\phi_b(2)\pm\phi_b(1)\phi_a(2)]/\sqrt{2(1\pm|l|^2)}$，分别配自旋反对称单重与对称三重，使总态反对称。原子哈密顿量作用给$E_a+E_b$；$\Delta H$的两对角矩阵元相等为$V$，两交叉矩阵元为$I$和$I^*$，实无流情况取$I$实。归一化相除得答案。括号只包$(V\pm I)$，不能把$E_a+E_b$也误除以重叠因子。
\end{solution}
\originalsection{作业10（2009年4月29日截止）}
\begin{exercise}[原题1(a)--(b)]\label{p:512-10-1}
(a)由线性响应证明纵向静态磁化率可用$S_{zz}(q,\omega)$的频谱积分表示，权重为$(1-\e^{-\beta\hbar\omega})/(\hbar\omega)$，取$q\to0$并明确磁矩单位。(b)若总磁化$M_z$与$H$对易，由$F=-\kb T\log\Tr\e^{-\beta(H-H_zM_z)}$导出$\chi=\beta\avg{M_z^2}$（零均值），证明与(a)一致。
\end{exercise}
\begin{solution}
(a)静态响应也可写虚时间积分$\chi=\int_0^\beta\avg{\delta M_z(-\ii\hbar\tau)\delta M_z(0)}\dd\tau$。插入能量态，$\int_0^\beta\e^{-\tau(E_m-E_n)}\dd\tau=(1-\e^{-\beta(E_m-E_n)})/(E_m-E_n)$，所以
\begin{equation}
\chi(q,0)=\int\frac{\dd\omega}{2\pi}\frac{1-\e^{-\beta\hbar\omega}}{\hbar\omega}\,S_{MM}(q,\omega).
\end{equation}
$S$若以自旋而非磁矩定义须乘$(g\mu_B)^2$，体积归一同样不能省；原题取$\hbar=1$即其表达。

(b)$[M_z,H]=0$时$\partial_{H_z}\log Z=\beta\avg{M_z}$，再求导$\chi=\beta(\avg{M_z^2}-\avg{M_z}^2)$。零場均值为0时得题示形式。谱只有$S_{MM}(0,\omega)=2\pi\avg{(\delta M_z)^2}\delta(\omega)$，权重在$\omega=0$的连续极限为$\beta$，两式一致。严格孤立守恒总量的有限频率响应与热库中等温静态导数可有极限顺序差，题目使用的是热平衡静态磁化率。
\end{solution}
\originalsection{作业11（2009年5月6日截止）}
\begin{exercise}[原题1(a)--(b)]\label{p:512-11-1}
用前题结果求三维Heisenberg反铁磁的低温沿Néel轴磁化率，证明$\chi_\parallel\propto T^2$。若有Ising易轴各向异性，说明指数激活形式并给激活隙$\Delta$。
\end{exercise}
\begin{solution}
两种反铁磁磁振子携带相反的$\Delta S_z=\pm1$，小纵场能量移$E_k\mp g\mu_BH$。每体积响应为
\begin{equation}
\chi_\parallel=2(g\mu_B)^2\int\frac{\dd^3k}{(2\pi)^3}\,\beta n_B(E_k)[1+n_B(E_k)].
\end{equation}
无各向异性$E_k=\hbar ck$，换变量$x=\beta\hbar ck$，积分$\int x^2\e^x/(\e^x-1)^2\dd x=\pi^2/3$，得$\chi_\parallel=(g\mu_B)^2(\kb T)^2/(3\hbar^3c^3)$，SI还按定义乘$\mu_0$。基态不能沿Néel轴连续极化，所以此纵向热响应趋0，横向响应则另有常数项。

有各向异性$E_k\simeq\sqrt{\Delta^2+\hbar^2c^2k^2}$，最低隙$\Delta=\sqrt{(\omega_0+\omega_A)^2-\omega_0^2}=\sqrt{\omega_A(2\omega_0+\omega_A)}$，其中题面$\omega$按能量计。低温积分进一步给$T^{1/2}\e^{-\Delta/\kb T}$前因子，因此原题的$\sim\e^{-\Delta/T}$应理解为主指数，而非无幂因子的严格等价。
\end{solution}
\begin{exercise}[原题2(a)--(b)]\label{p:512-11-2}
立方双子格反铁磁有易轴$-D(S^z)^2$，$D>0$，自旋波$E_k=\sqrt{(\omega_0+\omega_A)^2-\omega_0^2\gamma_k^2}$，$\omega_0=2zJS,\omega_A=2DS,\gamma_k=z^{-1}\sum_\delta\e^{\ii k\delta}$。(a)纵場证明两支$E_k\pm g\mu_BH_z$。(b)横場$H_y$如何改变谱？需先确定新的经典基态。
\end{exercise}
\begin{solution}
(a)磁振子的总$S_z$量子数为$\pm1$，Zeeman直接移能。纵場若超过稳定范围会发生自旋翻转，不能仍用旧Néel态。

(b)采用与题示$\omega_0$相容的$H=2J\sum_{\avg{ij}}\vct S_i\cdot\vct S_j-D\sum_i(S_i^z)^2-h\sum_iS_i^y$，$h=g\mu_BH_y$。两子格取$\vct S_A=S(0,s,c)$、$\vct S_B=S(0,s,-c)$，$s=\sin\theta,c=\cos\theta$。每站经典能$E/S^2=-zJ+(2zJ+D)s^2-hs/S-D$，极小$s=h/[2S(2zJ+D)]=h/(2\omega_0+\omega_A)$，要求$|s|<1$。

分别转到各子格局域轴后作HP展开，二次正常对角$A=\omega_0+\omega_A-(\omega_A/2)s^2$，正常子格杂化$B_k=\omega_0\gamma_ks^2$，反常子格配对$C_k=\omega_0\gamma_kc^2$，同站反常项$D_0=(\omega_A/2)s^2$。对子格对称与反对称模作玻色Bogoliubov变换，得到
\begin{equation}
E_{k,\pm}=\sqrt{(A\pm B_k)^2-(D_0\pm C_k)^2}.
\end{equation}
$h=0$恢复题示二重谱；横场修正从$h^2$开始且两支分开，并非简单的$E_k\pm h$。当$|s|\to1$达到饱和边界，旧低场分支需重新判稳定。随卷数值脚本核对$h=0$退化、正常/反常矩阵本征值与上式一致。
\end{solution}
\originalsection{作业12（2009年5月13日截止）}
\begin{exercise}[原题1--4(a)--(b)]\label{p:512-12-1}
接触排斥$U$的费米气在Stoner平均场有$\Delta=U(n_\uparrow-n_\downarrow)$，$\widetilde\epsilon_{k\uparrow}=\epsilon_k-\Delta/2$、$\widetilde\epsilon_{k\downarrow}=\epsilon_k+\Delta/2$，采用原题$\hbar=1$。(1)求自旋翻转横向响应泡图$\chi_\perp^0=\mu_B^2\Gamma_0$。(2)自洽响应证明$\chi_\perp=\mu_B^2\Gamma_0/(1-U\Gamma_0)$。(3)画单粒子--空穴连续谱在$(q,\omega)$中的支撑。(4a)证明$q=\omega=0$分母为0。(4b)小$q,\omega$展开求$\omega=Cq^2$及$C$的积分形式。
\end{exercise}
\begin{solution}
用正静态约定取
\begin{equation}
\Gamma_0(q,\omega)=\int\frac{\dd^dk}{(2\pi)^d}\frac{f(\epsilon_k-\Delta/2)-f(\epsilon_{k+q}+\Delta/2)}{\Delta+\epsilon_{k+q}-\epsilon_k-\omega-\ii0}.
\end{equation}
与原题把分子、分母同时反号的式相同，明确$k$不是两个不同的下标误排。外部横场加诱导交换场$U\delta S$，$\delta S=\Gamma_0(H+U\delta S)$，解出几何级数分母。

(3)把源题的连续色散专门取为三维各向同性抛物线$\epsilon_k=k^2/(2m)$，$T=0$，记多数与少数费米半径$a=k_{F\uparrow}>b=k_{F\downarrow}>0$，同一化学势要求$\Delta=(a^2-b^2)/(2m)$。正频率吸收只能由占据的$\uparrow$态到空的$\downarrow$态：$|k|\le a,\ |k+q|\ge b$，且
\begin{equation}
\omega=\Delta+\frac{q^2+2\vct k\cdot\vct q}{2m},\qquad q=|\vct q|.
\end{equation}
上边界取$\vct k$与$\vct q$同向、$|k|=a$，末态必在少数费米球外，故$\omega_+(q)=\Delta+q^2/(2m)+aq/m$。反向$\vct k=-a\hat q$的候选下边界$L(q)=\Delta+q^2/(2m)-aq/m$，仅在$|q-a|\ge b$时满足末态Pauli条件。若$a-b<q<a+b$，两个费米球的边界相交，可令初末态都在各自费米面上，连续谱下边界改为0。因此完整支撑是
\begin{equation}
\operatorname{Im}\chi^0_\perp(q,\omega)\ne0\quad\Longleftrightarrow\quad
\omega_-(q)<\omega<\omega_+(q),\qquad
\omega_-(q)=\begin{cases}
L(q),&0\le q\le a-b\ \text{或}\ q\ge a+b,\\
0,&a-b\le q\le a+b.
\end{cases}
\end{equation}
边界上响应可随相空间端点趋零，故用开区间表示非零内部；$q=0$连续谱退化到$\omega=\Delta$，仍有来自$b<|k|<a$的非零谱权重。两球相交的$q$区间里低频粒子--空穴态存在，不能把$\Delta\pm v_{F\uparrow}q$外推成全局下界。示意图以$a=2,b=1,m=1$取单位，此时$\Delta=3/2$、两处阈值为$q=1,3$；纵轴为正频率，阴影是允许区：
\begin{center}
\begin{tikzpicture}
\begin{axis}[width=.79\linewidth,height=5.1cm,xmin=0,xmax=4,ymin=0,ymax=18,
 xlabel={$q$},ylabel={$\omega$},axis lines=left,xtick={0,1,3,4},
 xticklabels={$0$,$a-b$,$a+b$,$4$},ytick={0,1.5},yticklabels={$0$,$\Delta$},
 legend style={at={(.97,.97)},anchor=north east,draw=none,fill=white,font=\small}]
\addplot[forget plot,fill=blue!15,draw=none,domain=0:4,samples=100]{1.5+x*x/2+2*x}\closedcycle;
\addplot[forget plot,fill=white,draw=none,domain=0:4,samples=160]{max(0,1.5+x*x/2-2*x)}\closedcycle;
\addplot[domain=0:4,samples=100,draw=blue!70!black,thick]{1.5+x*x/2+2*x};
\addplot[domain=0:4,samples=160,draw=red!70!black,thick]{max(0,1.5+x*x/2-2*x)};
\legend{$\omega_+$,$\omega_-$}
\end{axis}
\end{tikzpicture}
\end{center}
图的$\omega_-$折到0是末态Pauli限制的几何结果，不是令自旋分裂$\Delta$消失。小$q<a-b$时Goldstone模$\omega\propto q^2$在连续谱下方；更大$q$的集体模是否衰减还要解完整响应分母。

(4a)$\Gamma_0(0,0)=(n_\uparrow-n_\downarrow)/\Delta=1/U$，正是平均场自洽式。(4b)令$A=\partial_\omega\Gamma_0|_0=(n_\uparrow-n_\downarrow)/\Delta^2$，$B=(2d)^{-1}\nabla_q^2\Gamma_0|_0$，反演对称使一阶$q$消失，则$1-U\Gamma_0\simeq-U(A\omega+Bq^2)$，所以$C=-B/A$。抛物线可先把第二项换元为$k'=k+q$，再展开两个分母，得到
\begin{equation}
B=-\frac{n}{2m\Delta^2}+\frac1{d m^2\Delta^3}\int\frac{\dd^dk}{(2\pi)^d}k^2[f(\epsilon_k-\Delta/2)-f(\epsilon_k+\Delta/2)].
\end{equation}
故$C=n/[2m(n_\uparrow-n_\downarrow)]-\int k^2(f_\uparrow-f_\downarrow)/[d m^2\Delta(n_\uparrow-n_\downarrow)]$。恢复$\hbar$时每个$k$速度与$q^2$曲率带相应$\hbar^2$。符号正表示所选均匀态的长波稳定；若算出负刚度，应审查模型态是否不稳定，不能强行叫稳定自旋波。
\end{solution}

\originalchapter{8.512公开2004年补充作业}
\psource{8.512 Spring 2009作业页另列2004年10套；下表按原卷位置指向已经给出的同题完整题面与解答，再补独有问题。}
\begin{center}\begin{tabular}{lll}\toprule 2004套号&原题&本册定位\\\midrule
1&1、2&习题\ref{p:512-1-1}、\ref{p:512-1-2}\\
2&1、2&习题\ref{p:512-2-1}、\ref{p:512-2-2}（倒格标签改为带标签）\\
3&1、2&习题\ref{p:512-4-1}、\ref{p:512-4-2}\\
4&1&习题\ref{p:512-6-1}\\
6&1&习题\ref{p:512-9-1}\\
7&1&习题\ref{p:512-9-2}\\
8&1、2&习题\ref{p:512-10-1}、\ref{p:512-11-1}\\
9&1--4&习题\ref{p:512-12-1}中对应1--4\\\bottomrule\end{tabular}\end{center}
上述对应保留原套号与年份，不把一次解答重复统计成新增内容。2004套2的$G$标签是折叠带的另一表示，物理条件一致；2004套9与2009套12的个别自旋上下标底稿排版有混乱，均以所列明确分裂能与自洽关系核查。
\originalsection{2004作业4：核磁弛豫（3月11日截止）}
\begin{exercise}[原题2(a)--(b)]
核自旋为两态系统$I$，与电子在原点接触耦合$H_{\rm hf}=A\vct I\cdot\vct S(0)$，$S(r)=\Psi_\alpha^\dagger(r)\vct\sigma_{\alpha\beta}\Psi_\beta(r)$。(a)核自旋初始向上，证明弛豫率可用$\int\avg{S^+(t)S^-(0)}\cos\omega_nt\dd t$表示，原题约定前因子$2A^2/\hbar^2$；$\omega_n$远小于电子尺度可取零。(b)求自由费米气低温率并证明$1/(T_1T)\propto[N(E_F)]^2$，说明超导相干因子的角色。原书阅读指针不收录。
\end{exercise}
\begin{solution}
核翻转耦合来自$A(I^-S^++I^+S^-)/2$，Golden rule对电子热初末态求和，转为时间相关的Fourier变换；纵向回到平衡的$T_1$还合并上翻与下翻率。因原题自旋与$S^\pm$归一不同于常见$S=\sigma/2$，固定其前因子后使用所给定义，不能无条件复制另一教材的数值2。

自由电子$S^+(0)$在同一位置将一自旋态变成另一，谱权重为$\int\rho_\uparrow(E)\rho_\downarrow(E+\hbar\omega_n)f(E)[1-f(E+\hbar\omega_n)]\dd E$。小$\omega_n$及平滑态密度下$\int f(1-f)\dd E=\kb T$，故$1/T_1=\mathcal C(A^2/\hbar)\rho_\uparrow(E_F)\rho_\downarrow(E_F)\kb T$，$\mathcal C$由原自旋归一固定，Korringa温度与态密度结论不依该记号。超导态须将电子算符作Bogoliubov展开，核自旋通道的$uu'+vv'$型相干因子与能隙态密度共同决定低温行为；原题只要求阅读商业书了解改动，未给独立需证明的完整超导率题面。
\end{solution}
\originalsection{2004作业5：节点超导与响应（3月18日截止）}
\begin{exercise}[原题1(a)--(d)]
二维$d$波隙$\Delta(\phi)=\Delta_0\cos2\phi$，正常费米面圆形，原题取$\hbar=\kb=1$，能谱为
\begin{equation}
E(k)=\sqrt{v_F^2(|k|-k_F)^2+\Delta(\phi)^2}.
\end{equation}
(a)找4节点、线性化$E=\sqrt{v_F^2k_1^2+v_2^2k_2^2}$、求$v_2$及每节点每自旋态密度。(b)证明$\rho_s(T)/m=\rho_s(0)/m-(2\log2/\pi)(v_F/v_2)T$。(c)在极端II类单涡旋外$\xi_0\ll R\ll\lambda$，Doppler位移$v_F\cdot[\nabla\theta/2+eA/c]$使零能态密度非零，分别求沿$x$与$(1,1)$接近涡核的结果。(d)三角涡晶格中估计$c_V=\alpha\sqrt H\,T$及$\alpha$量级；原题所列实验论文只作参考链接。
\end{exercise}
\begin{solution}
(a)节点$\phi=\pi/4+n\pi/2$，近节点$\Delta\simeq\pm2\Delta_0\delta\phi=\pm(2\Delta_0/k_F)k_2$，故$v_2=2\Delta_0/k_F$。缩放$p_1=v_Fk_1,p_2=v_2k_2$，$\int\delta(E-|p|)\dd^2p/(2\pi)^2v_Fv_2=E/(2\pi v_Fv_2)$。

(b)热准粒子的反向电流给$\delta(\rho_s/m)=-2\sum_{\rm nodes}\int v_x^2[-f'(E)]\dd^2k/(2\pi)^2$，节点的$x$分量平方$v_F^2/2$；四节点并双自旋归一后系数为$-(2/\pi)(v_F/v_2)\int_0^\infty E[-f'(E)]\dd E$。积分分部给$\int f(E)\dd E=T\log2$，得到线性律。

(c)流动动量$\vct p_s\simeq\hat\phi/(2R)$，每节点的$D(0)=|v_F\hat n\cdot p_s|/(2\pi v_Fv_2)$，自旋另乘2。$R\parallel x$时四节点投影绝对值皆$1/\sqrt2$，总$D(0,R)=\sqrt2/(\pi v_2R)$；$R\parallel(1,1)$时两节点投影零、两节点为1，总$1/(\pi v_2R)$。角依赖是节点结构，不是标量正常态密度的简单平移。

(d)涡间距$r_H\sim\sqrt{\Phi_0/H}$，把$1/R$在一个胞内积分给平均$\overline D(0)\sim1/(v_2r_H)\propto\sqrt H$，因此费米型低温$C\propto\overline D(0)T$。用圆胞近似$R_c=\sqrt{\Phi_0/(\pi H)}$，角平均$\overline{\sum_n|\sin(\phi_n-\phi)|}=8/\pi$，得$\overline D\simeq8/(\pi^2v_2R_c)$，所以按该态密度与常用热归一$\alpha\simeq8\sqrt\pi\kb^2/(3v_2\sqrt{\Phi_0})$。真实三角格胞、核区和速度单位恢复改变数值常数，$\sqrt H\,T$是本题稳健的低场低温量级，不宣称这个粗略常数精确。
\end{solution}
\begin{exercise}[原题2]
列完美金属、无序金属、洁净超导、无序超导的实部横纵$K$，分别取静态$\omega=0,q\to0$与动态$q=0,\omega\to0$，共16个量，并给主要非零贡献。
\end{exercise}
\begin{solution}
采用SI核$j=-KA$，$D=ne^2/m,D_s(T)=n_s(T)e^2/m$，相位已作规范一致的纵向调整。表中洁净体系取单带Galilean不变连续模型、电流与总动量成正比且没有动量弛豫；无序体系取有限弹性散射时间$\tau$。温度低于$T_c$，先固定$T$再取表中静态或动态极限。16个实部的极限与先导项如下；零不表示整条频率函数恒零。
\begin{center}\begin{tabular}{lllll}\toprule 系统&静态$K_T$&静态$K_L$&动态$K_T$&动态$K_L$\\\midrule
完美金属&$-\chi_Dq^2/\mu_0$&0&$D$&$D$\\
无序金属&$-\chi_Dq^2/\mu_0$&0&$D(\omega\tau)^2$&$D(\omega\tau)^2$\\
洁净超导&$D_s(T)+O(q^2)$&0&$D+O(\omega^2)$&$D+O(\omega^2)$\\
无序超导&$D_s(T)+O(q^2)$&0&$D_s(T)+O(\omega^2)$&$D_s(T)+O(\omega^2)$\\\bottomrule\end{tabular}\end{center}
无序正常金属全复核$K=-\ii\omega D\tau/(1-\ii\omega\tau)$，实部$D(\omega\tau)^2/[1+(\omega\tau)^2]$，虚部首阶$-\omega\sigma_0$，故动态实极限为0。洁净Galilean体系因$\vct j\propto\vct P$且$\vct P$守恒，$q=0$动态核是$D$，即使$0<T<T_c$仍然如此；静态Meissner核却是$D_s(T)<D$。两极限不可互换。有限弛豫的无序超导在$\omega\to0$只保留凝聚体刚度$D_s(T)$，正常准粒子的正则电导给$-\ii\omega\sigma_{\rm reg}$并在实部极限消失；$T=0$时洁净Galilean体系$D_s(0)=D$。静态纵$A$是纯规范，若固定序参量相位而不让它响应，会得到错误非零值。
\end{solution}
\originalsection{2004作业7：易轴磁振子（4月13日截止）}
\begin{exercise}[原题2(a)--(c)，公开文件引Ashcroft--Mermin第723页题5]
各向异性铁磁$H=-\frac12\sum_{RR'}[J_z(R-R')S_z(R)S_z(R')+J_\perp(R-R')\vct S_\perp(R)\cdot\vct S_\perp(R')]$，$J_z>J_\perp>0$。(a)证明全极化与单自旋波仍为本征态，并给相对各向同性的能量抬升。(b)证明低温磁化偏离饱和按激活指数小。(c)解释二维无连续自旋序的红外论证为何失效。原PDF的抬升式两项都排为$J_z$，应由哈密顿量独立判断。
\end{exercise}
\begin{solution}
全上态横向交换翻转不能作用，故本征。单翻转可在格点间传播，纵向代价$S\sum_RJ_z(R)$，横向迁移给$-SJ_\perp(k)$，所以$E_k=S[J_z(0)-J_\perp(k)]$（此处$J(0)$表示Fourier零波矢）。相对各向同性$SJ_\perp(0)-SJ_\perp(k)$，抬升为$\Delta=S\sum_R[J_z(R)-J_\perp(R)]>0$，不是原底稿误排的零。

低温$E_k\simeq\Delta+Ak^2$，磁振子数$\int\e^{-E_k/T}\dd^dk\propto T^{d/2}\e^{-\Delta/T}$，得到激活磁化减小。二维连续对称模型$\int\dd^2k/k^2$红外发散；易轴打开隙且剩离散翻转对称，积分$\int\dd^2k/(\Delta+Ak^2)$红外有限，所以那条禁止论证不再成立，但仅“不再禁止”不是单凭此积分就证明某温度必有序。
\end{solution}
\originalsection{2004作业10：共振杂质的总电荷（5月6日截止）}
\begin{exercise}[原题1--6]
非相互作用$d$轨道$\epsilon_d$与传导带通过$V_k$杂化，$\Sigma(E)=\sum_k|V_k|^2/(E-\epsilon_k+\ii0)$，$G_d=[E-\epsilon_d-\Sigma]^{-1}$，$\Gamma=-\Sigma''$，相移由$\tan\delta=\Gamma/(\epsilon_d+\Sigma'-E)$连续选支。(1)用$\operatorname{Im}\log G_d$表达相移。(2)积分$\partial_E\log G_d$写为$G_d$和$\Sigma'$的式。(3)求$\delta(E_F)$与$n_d$的关系。(4)以$G_k-G_k^0$求传导电荷改变。(5)证明Friedel规则含总电荷。(6)估计窄共振$\Gamma\ll E_F$时忽略$\Sigma$能量依赖的误差。公开底稿个别式的负号与$1/\pi$归一需按Green函数定义重核。
\end{exercise}
\begin{solution}
(1)推迟$G_d$的相位在$[-\pi,0]$，使空共振在低能的相移为0可取$\delta=\pi+\operatorname{Im}\log G_d$，连续追踪支，不直接用$\arctan$主值跳变。(2)$\partial_E\log G_d=-G_d(1-\partial_E\Sigma)$，积分后
\begin{equation}
\delta(E_F)=\pi n_d+\operatorname{Im}\int_{-\infty}^{E_F}G_d(E)\partial_E\Sigma(E)\dd E,\qquad n_d=-\frac1\pi\operatorname{Im}\int^{E_F}G_d\dd E.
\end{equation}
这个符号从微分恒等式直接得出，不能沿用排版相反号却同时声称匹配总电荷。

(4)Dyson方程$G_k-G_k^0=|V_k|^2(G_k^0)^2G_d$，而$\partial_E\Sigma=-\sum_k|V_k|^2(G_k^0)^2$，故$\sum_k(G_k-G_k^0)=-G_d\partial_E\Sigma$。因此$\delta n_c=-\pi^{-1}\operatorname{Im}\int\sum_k(G_k-G_k^0)\dd E=\pi^{-1}\operatorname{Im}\int G_d\partial_E\Sigma\dd E$，(5)$\delta(E_F)=\pi(n_d+\delta n_c)$。

(6)若杂化带函数在能量$E_F$或带宽$D$上变化，$\partial_E\Sigma=O(\Gamma/D)$，共振窗口的积分修正相对为该小尺度。常$\Gamma,\Sigma'$的宽带近似使$\delta n_c=0$，才得到$\delta/\pi=n_d$。接近带边、伪隙或强能量依赖时窄峰本身不保证修正小，必须检查实际$\Sigma'$尺度。
\end{solution}

\originalchapter{8.513作业1--4：几何动力学与玻色场}
\psource{Xiao-Gang Wen，8.513 Fall 2021，公开作业1--4；多数题取$\hbar=e=c=1$，需要时明确恢复。}
\originalsection{作业1}
\begin{exercise}[原题1(a)--(c)]
系统满足$\hbar b_{ij}\dot\xi^j=\partial_iH$，$b_{ij}=\partial_ia_j-\partial_ja_i$可逆，逆矩阵为$b^{ij}$。(a)证明$\hbar\dot O=[O,H]_{\rm cl}$，$[A,B]_{\rm cl}=b^{ij}\partial_iA\partial_jB$。(b)证明反对称性。(c)证明无显含时哈密顿量守恒。
\end{exercise}
\begin{solution}
左乘逆矩阵得$\hbar\dot\xi^i=b^{ij}\partial_jH$，链式法则即(a)。$b^T=-b$使$(b^{-1})^T=-b^{-1}$，交换$A,B$并重命名指标得负号，故$[H,H]_{\rm cl}=0$。若$H$显含$t$则$\dot H=\partial_tH$，不能省这个条件。Jacobi恒等式还来自$b$闭合，见第\ref{ch:berry}章的一般辛结构。
\end{solution}
\begin{exercise}[原题2(a)--(c)]
二维强磁场的导引中心满足$\dot x^i=B^{-1}\epsilon^{ij}\partial_jV$，$\epsilon^{12}=1$。(a)找相空间作用量。(b)量子化并给哈密顿算符。(c)若$V=v[(x^1)^2+(x^2)^2]/2$，求能级。这是最低Landau级的有效导引中心理论。
\end{exercise}
\begin{solution}
(a)取$L=-(B/2)(x\dot y-y\dot x)-V$，其$b_{xy}=-B$，逆$b^{xy}=1/B$，给题示方程。磁场与电荷的方向以所给方程固定，不另套相反导引中心符号。(b)$[\hat x,\hat y]=\ii\hbar/B$，$\hat H=V(\hat x,\hat y)$，一般非二次$V$需说明算符排序，此题无歧义。(c)取$\hat Q=\sqrt{|B|}\hat x,\hat P=\operatorname{sgn}(B)\sqrt{|B|}\hat y$，$[Q,P]=\ii\hbar$，$H=v(Q^2+P^2)/(2|B|)$，所以$E_n=\hbar v(n+1/2)/|B|$，$v>0$时有界。被投影掉的回旋零点能是共同常数，不含在$V$的有效谱中。
\end{solution}
\begin{exercise}[原题3(a)--(b)]
(a)$a=a_x\dd x+a_y\dd y$，$b=\dd a=(\partial_xa_y-\partial_ya_x)\dd x\wedge\dd y$，证明圆盘的Stokes公式。(b)在三维原点外，$b=r^{-3}(x\dd y\wedge\dd z-y\dd x\wedge\dd z+z\dd x\wedge\dd y)$，证明$\dd b=0$并求半径$r$球面的积分。
\end{exercise}
\begin{solution}
(a)这是平面Green公式：分别把$\partial_xa_y$沿$x$积分、$-\partial_ya_x$沿$y$积分，区域相邻切片的内边界相消，只剩逆时针定向外边界$\oint a$。先对矩形证明，再将光滑圆盘分割逼近，要求$a$及一阶导数足够连续。(b)三项对应向量场$\vct r/r^3$的通量二形式，其散度$3r^{-3}-3r^{-5}r^2=0$于$r\ne0$。球面拉回为$\sin\theta\dd\theta\wedge\dd\phi$，积分$4\pi$。它闭合但在穿孔空间不全局恰当：若$b=\dd a$全局光滑，闭球面积分应为0，与$4\pi$矛盾。
\end{solution}
\begin{exercise}[原题4(a)--(b)]
量子自旋$S$有$H=-\vct B\cdot\vct S$，$\vct B=B\hat z$。取沿单位向量$n$的最大本征值相干态。(a)求$n$的相空间作用量。(b)求经典运动。
\end{exercise}
\begin{solution}
自旋$1/2$相干态$(\cos(\theta/2),\e^{\ii\phi}\sin(\theta/2))^T$的$\ii\avg{n|\partial_tn}=-(1-\cos\theta)\dot\phi/2$。自旋$S$可视为$2S$个对称自旋$1/2$的相干乘积，Berry项相加，$L=-\hbar S(1-\cos\theta)\dot\phi+BS\cos\theta$。变分给$\dot\theta=0,\dot\phi=-B/\hbar$，即$\dot n=n\times B/\hbar$。另一规范可加$\hbar S\dot\phi$总导数，却不改变局部方程；球面曲率通量$-4\pi S$需用两个相位片，量子条件$2S\in\mathbb Z$保证全局一致。
\end{solution}
\originalsection{作业2}
\begin{exercise}[原题1(a)--(c)]
带电粒子在$L_x\times L_y$环面，$H=(-\ii\nabla-A)^2/(2m)$，$A=(-By,0)$，$\hbar=e=c=1$。(a)$\psi'=\e^{\ii f}\psi$时求$H'$。(b)保留$x$周期，求使$A(y+L_y)$与$A(y)$一致的规范函数及$y$边界；证明总磁通量子化。(c)证明最低Landau级简并等于通量数，并说明与自旋球面“态数比Chern数多1”的区别。
\end{exercise}
\begin{solution}
(a)$A'=A+\nabla f$，$H'=\e^{\ii f}H\e^{-\ii f}$。(b)$f=-BL_yx$，$U=\e^{-\ii BL_yx}$，所以$\psi(x,y+L_y)=\e^{-\ii BL_yx}\psi(x,y)$。先绕$x$再绕$y$与反次序应同态，要求$\e^{-\ii BL_xL_y}=1$，即$BL_xL_y=2\pi n$。这正是线丛过渡函数一致性。

(c)Landau规范导引中心随$x$动量$2\pi l/L_x$移动，$y$方向周期识别$l\sim l+n$，有$|n|$个独立中心。若$B=0$则不是Landau投影问题。球面自旋相空间曲率通量绝对Chern数$2S$，量子自旋Hilbert维数$2S+1$；不同拓扑与量子化的端点修正使环面的等号不能照搬到球面。
\end{solution}
\begin{exercise}[原题2(a)--(d)]
三角格的四子格磁矩为$n_1=(1,1,1)/\sqrt3,n_2=(1,-1,-1)/\sqrt3,n_3=(-1,1,-1)/\sqrt3,n_4=(-1,-1,1)/\sqrt3$。以三角原矢$a_1,a_2$，格点$(n,m)$的类型按偶偶、奇偶、偶奇、奇奇为1、2、3、4。电子自旋始终平行局域矩，$\ket{n_i}=(\cos\theta_i/2,\e^{\ii\phi_i}\sin\theta_i/2)^T$，$t_{ij}=A\avg{n_i|n_j},A>0$。(a)求沿逆时针三角环的均匀磁通。(b)重新选自旋态相位，使跃迁仅使胞加倍。(c)求两带和隙。(d)关掉$a_1$水平键后求新带与隙。条件由原图转为向量与格点标签，不复用原示意图。
\end{exercise}
\begin{solution}
(a)任意不同两方向内积为$-1/3$，故$|t|=A/\sqrt3$。投影矩阵$P_i=(1+n_i\cdot\sigma)/2$给Bargmann积$\Tr(P_iP_jP_k)=[1+n_i\cdot n_j+n_j\cdot n_k+n_k\cdot n_i+\ii n_i\cdot(n_j\times n_k)]/4$。前三内积使实部为0，按原图逆时针$1\to2\to3$，虚部正，通量$\pi/2$模$2\pi$；相邻反向排布补偿空间方向，所有逆时针三角同号。

(b)原相干态规范下相位依四子格，选择附加相位$(\zeta_1,\zeta_2,\zeta_3,\zeta_4)=(0,\pi/4,-\pi/4,-\pi/2)$。令$t=A/\sqrt3$，沿$a_1$为$t$，沿$a_2$为$(-1)^nt$，沿$a_1-a_2$为$-\ii(-1)^nt$，都与$m$无关且只依$n$奇偶，胞仅沿$a_1$倍增。

(c)以$k_i=k\cdot a_i$、保留两子格内部位置相位的Bloch规范，矩阵可取$H=2t[\cos k_1\sigma_x+\cos(k_1-k_2)\sigma_y+\cos k_2\sigma_z]$，$E_\pm=\pm2t\sqrt{\cos^2k_1+\cos^2k_2+\cos^2(k_1-k_2)}$。最小平方和为$3/4$，可由令$k_1=-k_2=\pi/3$达到，并由三角恒等式或求驻点与边界核对全局最小；全隙$2\sqrt3t=2A$。

(d)去掉$\sigma_x$项，$E_\pm=\pm2t\sqrt{\cos^2k_2+\cos^2(k_1-k_2)}$，在$k_2=\pi/2,k_1=0$等点闭隙，故隙0。少一类键改变的是完整二维Bloch矩阵，不能只把原带宽乘$2/3$。
\end{solution}
\originalsection{作业3}
\begin{exercise}[原题1(a)--(c)，斜入射选做]
自旋$1/2$粒子$H=p^2/(2m)+g\sigma\cdot p+V(x)$，$g>0$。(a)$V=0$求给定$p$最低能与自旋态。(b)该低态束沿$+z$、$p>mg$撞击$z>0$无限墙，反射保持自旋，求反射能、动量与速度。(c)取以$(x,p)$为中心、宽$\xi$的Gaussian波包，$V$在$\xi$内缓变，求相空间作用量与方程。斜入射可选做，须保留墙平行动量。
\end{exercise}
\begin{solution}
(a)$\sigma\cdot\hat p$本征值$\pm1$，低支$\epsilon_-=p^2/(2m)-g|p|$，自旋反平行$p$。(b)初自旋沿$-z$，固定该自旋的能量为$p_z^2/(2m)-gp_z$。等能第二根$p'_z=2mg-p$，反射速度$v'=p'/m-g=-(p/m-g)$，能量不变。若$p>2mg$，$p'<0$而自旋仍负，反射落在该动量的高螺旋支；不能强迫反射后仍为最低支而违反边界自旋守恒。

(c)在孤立低支的绝热条件下$L=p\cdot\dot x+\hbar\mathcal A_p\cdot\dot p-\epsilon_-(p)-V(x)$，$\Omega_p=\hat p/(2|p|^2)$（本册$\mathcal A=\ii\avg{u|\nabla u}$）。$\dot p=-\nabla V$、$\dot x=p/m-g\hat p-\hbar\dot p\times\Omega_p$。$p=0$两支闭隙，不满足单支波包条件。对硬墙本身也不能以缓变单支方程代替(b)的两支边界散射；斜入射须保持$p_x,p_y$并叠加允许的反射自旋本征支来满足$\psi(z=0)=0$，具体地，保持平行波矢$q=p\sin\theta$，固定能量$E=p^2/(2m)-gp$。可用壳半径为$r_1=p$及$r_2=|p-2mg|$；$p>2mg$时第二壳属于高螺旋支，$mg<p<2mg$时属于低支内壳。每壳取满足向外法向群速度的$p_{z,j}=\pm\sqrt{r_j^2-q^2}$；若根虚，选$\operatorname{Im}p_z<0$使波在$z\to-\infty$衰减。以$(g\sigma\cdot p_j)u_j=[E-r_j^2/(2m)]u_j$求对应两分量自旋矢，边界系数由$[u_1,u_2](A_1,A_2)^T=-u_{\rm in}$解出。传播支概率$R_j=|A_j|^2|v_{z,j}|/v_{z,\rm in}$，倏逝支不带远处净流。这个两根加边界的构造覆盖斜入射，法向$q=0$退化为(b)的自旋守恒单支反射；壳退化点需取连续极限。
\end{solution}
\begin{exercise}[原题2(a)--(b)]
石墨烯两谷近似$H_K=\hbar v(k_x\sigma_x+k_y\sigma_y)$、$H_{K'}=\hbar v(k_x\sigma_y+k_y\sigma_x)$，$\mu=0$。(a)求热激发导带电子及价带空穴密度的温度幂律。(b)假设$\tau$不随$T$，由Boltzmann求电导，判断是否与密度同为$T^2$。
\end{exercise}
\begin{solution}
每谷每自旋$D(E)=|E|/(2\pi\hbar^2v^2)$。设总简并$g_d$，$n_e=n_h=g_d(\kb T)^2\pi/(24\hbar^2v^2)$，由$\int_0^\infty x/(\e^x+1)\dd x=\pi^2/12$。电导$\sigma=e^2\tau(v^2/2)\int D_{\mathrm{total}}(E)[-f'(E)]\dd E$，正负带均计，$\int_0^\infty E[-f']\dd E=\kb T\log2$，得$\sigma=g_de^2\tau\kb T\log2/(2\pi\hbar^2)$，仅线性$T$。原因是Dirac载流子的惯性尺度$|E|/v^2$也随温度改变，不能固定质量套Drude密度比例。$T\to0$的真实量子极限另含无序、接触与带间效应。
\end{solution}
\begin{exercise}[原题3(a)--(c)]
在上述两谷都加$\Delta\sigma_z$。(a)求两带Berry曲率峰、每谷总通量与宽$\delta k$。(b)零温K谷有导带电子密度$n_e^K$，K'谷有价带空穴密度$n_h^{K'}$，求小、大密度的Hall响应并说明连续依赖。(c)K谷导带电子初始$x=k=0$，受$E=E\hat x$，短时间$\tau$内曲率近常数，求位置。
\end{exercise}
\begin{solution}
令$\eta_K=+1,\eta_{K'}=-1$，$E_k=\sqrt{\Delta^2+\hbar^2v^2k^2}$，本册约定$\Omega_{\eta,\pm}=\mp\eta\Delta\hbar^2v^2/(2E_k^3)$。峰宽$|\Delta|/(\hbar v)$，每谷单带全平面积分$\Phi_\pm=\mp\eta\pi\operatorname{sgn}\Delta$，完整晶格两谷合计为整数且此同号质量使总Chern为0。

设各密度包含自旋简并$g_s$，$k_F^2=4\pi n/g_s$。定义$F(n)=\frac{\operatorname{sgn}\Delta}{2}[1-|\Delta|/\sqrt{\Delta^2+\hbar^2v^2k_F^2}]$。满价带的两个谷响应相消，新增电子和移除空穴后的曲率积分为$g_s[-F(n_e^K)+F(n_h^{K'})]$。用$\sigma_{xy}=-e^2\int f\Omega/\hbar$的方向约定，$\sigma_{xy}=g_s(e^2/h)[F(n_e^K)-F(n_h^{K'})]$。小密度$F\simeq\operatorname{sgn}\Delta\,\pi\hbar^2v^2n/(g_s\Delta^2)$，大密度各自趋$\operatorname{sgn}\Delta/2$。相同谷相反带的电子--空穴贡献相加，而本题相反谷会抵消；等密度时净Hall为0，不能遗漏K'手性负号。

(c)$k_x=-eEt/\hbar$，带底群速度$x$分量$\hbar v^2k_x/|\Delta|$，因此$x(\tau)=-eEv^2\tau^2/(2|\Delta|)$。异常速度$\dot y=-eE\Omega_{K,+}(0)/\hbar$，所以$y(\tau)=-eE\Omega_{K,+}(0)\tau/\hbar$。适用$eEv\tau\ll|\Delta|$且跨带跃迁小；速度与外力的线性关系包含横向几何项。
\end{solution}
\originalsection{作业4}
\begin{exercise}[原题1(a)--(c)]
一维线性色散玻色场在相距$L$的两固定墙间，$k_n=n\pi/L$，零点能$\frac12\sum\hbar ck_n$。(a)热核正规化乘$\e^{-ak_n}$，求小$a$展开并忽略线性及更高项。(b)格点正规化$E_{\rm latt}=\sum_{n=1}^{N}c\hbar a^{-1}\sin(k_na/2),N=L/a$，同样展开。(c)三墙总长$L$，中墙距左墙$l$，求总能及$l\ll L$时力，区分共同紫外项与有限力。
\end{exercise}
\begin{solution}
(a)几何级数导数$\sum n q^n=q/(1-q)^2$，$x=\pi a/L$给$E_{\mathrm{heat}}=\pi c\hbar\e^{-x}/[2L(1-\e^{-x})^2]=c\hbar L/(2\pi a^2)-\pi c\hbar/(24L)+O(a^2/L^3)$。

(b)有限正弦和为$\sum_{n=1}^N\sin(n\pi/2N)=[\cot(\pi/4N)+1]/2$，所以$E_{\mathrm{latt}}=2c\hbar L/(\pi a^2)+c\hbar/(2a)-\pi c\hbar/(24L)+O(a^2/L^3)$。两方案体积与边界发散不同，有限$1/L$项相同。形式$\sum n=-1/12$在此是扣掉规定发散后的解析正规化值，不是普通正项级数的和。

(c)分别加$E(l)+E(L-l)$，共同体积项只依总$L$，格点边界项只依墙数，故对中墙力仅有限项贡献：$F_l=-\pi c\hbar[1/l^2-1/(L-l)^2]/24$，$l\ll L$时$F=-\pi c\hbar/(24l^2)$，负号指向近墙。移动墙时保持同一微观截止，不用两个随长度变化的截止“消掉”发散。
\end{solution}
\begin{exercise}[原题2(a)--(e)]
格点玻色模型$H=\sum_i[2t\phi_i^\dagger\phi_i-t(\phi_i^\dagger\phi_{i+1}+\mathrm{h.c.})-\mu\phi_i^\dagger\phi_i+U\phi_i^{\dagger2}\phi_i^2]$，$t,U>0$。(a)证明未归一相干态重叠$\langle0|\e^{\sum\phi_i^*\hat\phi_i}\e^{\sum\widetilde\phi_i\hat\phi_i^\dagger}|0\rangle=\e^{\sum\phi_i^*\widetilde\phi_i}$。(b)求归一及共同本征值。(c)求$L=\langle\phi|\ii\partial_t-H|\phi\rangle$。(d)$\mu\ge0$时求均匀极小。(e)线性化运动求$\omega_k$，比较$\mu>0,=0,<0$的低能与动态指数。
\end{exercise}
\begin{solution}
(a)湮灭指数移过创建指数产生数对易子$\sum\phi_i^*\widetilde\phi_i$，余指数对真空作用为1。(b)归一因子$\e^{\sum|\phi_i|^2/2}$，且湮灭算符可移过指数得本征值$\phi_i$。(c)$\bar H=\sum[t|\phi_{i+1}-\phi_i|^2-\mu|\phi_i|^2+U|\phi_i|^4]$，Berry项$\ii\sum\phi_i^*\dot\phi_i$差一总导数后即原题形式。

(d)均匀跃迁项0，极小$|\bar\phi|^2=\mu/(2U)$，相位任意，$\mu=0$在原点。(e)取$\bar\phi$实，$\ii\dot{\delta\phi_i}=t(2\delta\phi_i-\delta\phi_{i+1}-\delta\phi_{i-1})+\mu(\delta\phi_i+\delta\phi_i^*)$，一纯复平面波通常不闭合，需$u\e^{\ii(ki-\omega t)}+v^*\e^{-\ii(ki-\omega t)}$的Bogoliubov双分量。$\epsilon_k=2t(1-\cos k)$，$\omega_k=\sqrt{\epsilon_k(\epsilon_k+2\mu)}$。$\mu>0$长波$\sqrt{2t\mu}|k|$，$z=1$；$\mu=0$为$tk^2$，$z=2$；$\mu<0$均匀真空稳定，$\omega_k=\epsilon_k+|\mu|$有隙。该相干态平均场刻画真空--超流入口，不证明一维存在通常长程凝聚，相位涨落仍需第\ref{ch:boson}章处理。
\end{solution}

\originalchapter{8.513作业5--8：量子临界与一维异常}
\psource{Xiao-Gang Wen，8.513 Fall 2021，公开作业5--8。}
\originalsection{作业5}
\begin{exercise}[原题1(a)--(b)]
自旋1链为
\begin{equation}
H=\sum_i[-J(S_i^xS_{i+1}^x+S_i^yS_{i+1}^y)-BS_i^z+V(S_i^z)^2],\quad J>0.
\end{equation}
试验态为各站$(\phi_i,1,\widetilde\phi_i)^T/\sqrt{1+|\phi_i|^2+|\widetilde\phi_i|^2}$的张量积。(a)大$V$的Mott态$\phi=\widetilde\phi=0$附近求二次相空间作用量。(b)从大$V$降低$V$，求$B=0$和$B\ne0$首次超流失稳的低能色散与动态指数。
\end{exercise}
\begin{solution}
$\avg{S^z}=|\phi|^2-|\widetilde\phi|^2$，$\avg{(S^z)^2}=|\phi|^2+|\widetilde\phi|^2$到二阶，$\avg{S^+}=\sqrt2(\phi^*+\widetilde\phi)$到一阶。Berry项为$\ii(\phi^*\dot\phi+\widetilde\phi^*\dot{\widetilde\phi})$差总导数。二次能$\sum[(V-B)|\phi|^2+(V+B)|\widetilde\phi|^2]-J\sum[(\phi_i^*+\widetilde\phi_i)(\phi_{i+1}+\widetilde\phi_{i+1}^*)+\mathrm{c.c.}]$。

Fourier后对$(\phi_k,\widetilde\phi_{-k}^*)$的动态矩阵为$\left(\begin{smallmatrix}V-B-2J\cos k&-2J\cos k\\2J\cos k&-V-B+2J\cos k\end{smallmatrix}\right)$，正激发两支$E_\pm=\sqrt{V(V-4J\cos k)}\pm B$。$B=0$在$V_c=4J$闭隙，$E\simeq2\sqrt2J|k|$，$z=1$。$B\ne0$首次失稳满足$V_c(V_c-4J)=B^2$，$V_c=2J+\sqrt{4J^2+B^2}$，低支$E\simeq V_cJk^2/|B|$，$z=2$。这是同一变分模型的粒子--空穴对称尖端与普通边界的区别；不将一维平均场相图当精确临界参数。
\end{solution}
\begin{exercise}[原题2(a)--(b)]
一维右行复费米子$E_k=\hbar vk,k=2\pi n/L$，基态填满$k\le0$。(a)保持总费米子数的低能多体谱，给前五个激发能级简并，与右行玻色子正整数模比较。(b)求热容并比较。
\end{exercise}
\begin{solution}
把费米海中各粒子从最高占据开始依次抬升，Pauli限制使抬升量成为非增整数序列，即整数分拆。激发能$n(2\pi\hbar v/L)$，$n=0,1,2,3,4,5$简并$1,1,2,3,5,7$；若“前五级”含基态则取前五项。玻色各模占据生成函数$\prod_{m\ge1}(1-q^m)^{-1}$给同谱，零模总电荷扇区已固定。

复费米子相对海的正能粒子与负能空穴均热激发，$E/L=2\int_0^\infty(\hbar vk)f(\hbar vk)\dd k/(2\pi)=\pi(\kb T)^2/(12\hbar v)$，用$\int x/(\e^x+1)=\pi^2/12$。因此$C/L=\pi\kb^2T/(6\hbar v)$，同一支手性玻色子的结果。Majorana不是复费米子，热容另减半。
\end{solution}
\originalsection{作业6}
\begin{exercise}[原题1(a)--(f)]
开链XY模型$H=-\sum_{i=1}^{N-1}(\sigma_i^x\sigma_{i+1}^x+\sigma_i^y\sigma_{i+1}^y+h\sigma_i^z)$，旋转$\prod_i\e^{\ii\theta\sigma_i^z/2}$为$U(1)$。(a)作Jordan--Wigner映射。(b)按原题改为偶数$N$的平移不变费米环，求$E_{\rm ring}(h,N)$。(c)找$\pm h_c$。(d)小$|h|$低能线性色散、中心荷与基态$1/N$项。(e)接近$h_c$速度如何消失，临界色散是什么？(f)中子给动量$k$并使自旋改变1，态$\sum_j\e^{\ii kj}\sigma_j^-\ket0$是否像$\sum\e^{\ii kj}c_j\ket0$一样是能量态？描述散射谱。
\end{exercise}
\begin{solution}
(a)$\sigma_j^z=1-2n_j$，XY键变$-2(c_j^\dagger c_{j+1}+\mathrm{h.c.})$，磁场变$2hn_j-h$。题面将场也求和到$N-1$，最后站无场，是边界缺项；(b)的平移环按题目指定改为所有站同场，不能声称与该开链每个有限尺寸能量相同。

(b)环取所选边界动量集$\mathcal K$，$\epsilon_k=-4\cos k+2h$，精确$E_{
m ring}=-hN+\sum_{k\in\mathcal K,\epsilon_k<0}\epsilon_k$，零能态须说明占据与宇称选择。体极限$|h|<2$时$k_F=\arccos(h/2)$，$E/N=-h+[2hk_F-4\sin k_F]/\pi$；$h>2$为空费米海$E/N=-h$，$h<-2$全满$E/N=h$，所以(c)$h_c=2$。

(d)两个费米点附近$E_{k_F+q}-E_F\simeq\pm vq$，$v=4\sin k_F$。左右复费米子构成非手性$c_R=c_L=1$的理论，通常称中心荷$c=1$；$c_{\rm tot}=2$。选CFT真空扇区时$E_0=Ne_\infty-\pi v/(6N)+\cdots$。一般有限环因粒子数舍入与边界扭曲另有零模$O(1/N)$项，必须用(b)的精确和核对，不能对任意周期动量海不加扇区条件套同一系数。

(e)$v=4\sqrt{1-h^2/4}\to0$，在$h=2$带底$\epsilon_k=2k^2+O(k^4)$，$z=2$，不是保持有限光速的Lorentz理论。(f)$\sigma_j^-=(\prod_{l<j}(1-2n_l))c_j^\dagger$或其共轭，依所选升降约定。该字符串作用于半填充海会叠加许多不同粒子--空穴激发，不是单个$k$费米本征态。动态谱为$\sum_n|\avg{n|S^-_k|0}|^2\delta(\omega-E_n+E_0)$，热力学极限形成连续谱与幂律边缘；小动量低能边缘与$v|q|$关联，但不等于一条裸$\epsilon_k$的$\delta$曲线。有限链当然仍由离散峰组成，连续谱是尺寸极限。
\end{solution}
\begin{exercise}[原题2]
环上两个手性复费米场$H=\int\dd x[c_1^\dagger v_1\ii\partial_xc_1+c_2^\dagger v_2\ii\partial_xc_2+M c_1^\dagger c_2+M^*c_2^\dagger c_1]$，标准反对易。分别$v_1v_2>0$与$<0$求谱，说明质量混合何时仍无隙，联系净手性中心荷。
\end{exercise}
\begin{solution}
Fourier$\e^{\ii kx}$下矩阵$\left(\begin{smallmatrix}-v_1k&M\\M^*&-v_2k\end{smallmatrix}\right)$，根$E_\pm=-(v_1+v_2)k/2\pm\sqrt{(v_1-v_2)^2k^2/4+|M|^2}$。若$v_1v_2>0$，行列式$v_1v_2k^2-|M|^2$仍在$k=\pm|M|/\sqrt{v_1v_2}$为0，只移开交叉，无全隙；净手性$c_-=\pm2$。若$v_1v_2<0$，无实零点，最小正能$2|M|\sqrt{-v_1v_2}/|v_1-v_2|>0$，净手性0。只看$k=0$的$\pm|M|$不能判断整个谱是否有隙；连续理论需交叉仍在其紫外有效窗口。
\end{solution}
\originalsection{作业7}
\begin{exercise}[原题1(a)]
手性玻色子的$E_n=n\,2\pi\hbar v/L$简并$p(n)\sim\e^{2\pi\sqrt{n/6}}/(4\sqrt3n)$。以此微正则态密度做鞍点求热容，并推广中心荷$c$的$D(E)\sim\exp[L\sqrt{c\pi(E/L)/(3\hbar v)}]$。区分$c$是手性总态计数还是非手性左右各一份。
\end{exercise}
\begin{solution}
主指数$\log D=\sqrt{\pi LE/(3\hbar v)}$，$\log Z$积分鞍点满足$\partial_E\log D=\beta$，所以$E_*=\pi L/(12\hbar v\beta^2)$，$C=\pi L\kb^2T/(6\hbar v)$。代数前因子产生有限尺寸次阶修正，在$\kb T\gg2\pi\hbar v/L$可略。将热能系数乘总手性计数$c_{\rm tot}$再用$S=2E/T$，得到题示指数的$c=c_{\rm tot}$。如果采用非手性教科书$c_R=c_L=c$，指数中应放$2c$，不可混用。
\end{solution}
\begin{exercise}[原题2(a)--(d)]
周期横场Ising链长$L\equiv0\pmod4$，$H=-\sum(J\sigma_i^x\sigma_{i+1}^x+h\sigma_i^z)$，$J>0$。(a)以$\lambda_i^{x,y}=(\prod_{j<i}\sigma_j^z)\sigma_i^{x,y}$作映射。(b)比较周期、反周期Majorana模型与自旋$Z=\prod\sigma^z$扇区，找正确FNP投影。(c)求两种费米边界真空宇称，单独考虑$k=0,\pi$。(d)$J>|h|$求最低两自旋态劈裂的指数形式和$\xi$；题面引用附论文未随公开作业附上，但可独立分析支点。
\end{exercise}
\begin{solution}
(a)内键$\sigma_i^x\sigma_{i+1}^x=-\ii\lambda_i^y\lambda_{i+1}^x$，边界多串$Z$，所以$c_{L+1}=-Zc_1$。(b)$FNP=\prod_i(\ii\lambda_i^y\lambda_i^x)=Z$，偶自旋扇区必须取反周期费米且投偶宇称；奇自旋取周期且投奇宇称。原题“even or maybe odd”是要确定这个配对，不能把两个完整费米空间都保留。

(c)反周期$k,-k$全部成对，BCS真空偶。周期的$k=0,\pi$配对幅度零，单模能$\xi_k=2(h-J\cos k)$，负则填。$-J<h<J$仅$k=0$填，真空奇；$h>J$两者空，$h<-J$两者满，皆偶。任一所需投影不匹配时最低允许态需加最便宜的准粒子。

(d)$E(k)=2\sqrt{J^2+h^2-2Jh\cos k}$，基态能在匹配扇区为$-\frac12\sum_kE(k)$。两动量格点错半步，Poisson求和差由第$L$阶及更高Fourier系数决定。最近复支点$k=\pm\ii\kappa$（负$h$可平移实部$\pi$），$\cosh\kappa=(J^2+h^2)/(2J|h|)$，$\kappa=\log(J/|h|)$。平方根支点的Fourier系数带代数因子，故$\Delta E=O(L^\mu)\e^{-L/\xi}$，$\xi^{-1}=\log(J/|h|)$，$h=0$是严格零劈裂极限。只求指数不需取得所引论文，也没有从该论文猜造额外问题。
\end{solution}
\begin{exercise}[原题3(a)--(c)]
反周期手性Majorana环$H=\sum_{k>0}vk\lambda_k^\dagger\lambda_k$，$k=2\pi(n+1/2)/L$，$\lambda_{-k}=\lambda_k^\dagger$，$\lambda_k\ket0=0$于$k>0$。(a)求虚时间有序$\avg{T\lambda_{k_1}(\tau)\lambda_{k_2}(0)}$，注意$\tau<0$负号。(b)$\lambda(x)=L^{-1/2}\sum\lambda_k\e^{\ii kx}$，证明Hermitian与局域反对易。(c)求环相关与无穷长极限。
\end{exercise}
\begin{solution}
(a)$\tau\ge0$只有$k_1>0,k_2=-k_1<0$非零，权重$[1-\operatorname{sgn}k_2]\e^{-v|\tau k_2|}\delta_{k_1+k_2}/2$；$\tau<0$先把算符交换，权重$-[1+\operatorname{sgn}k_2]\e^{-v|\tau k_2|}\delta_{k_1+k_2}/2$。(b)共轭后换$k\to-k$得自身，反对易为$L^{-1}\sum\e^{\ii k(x-y)}$，在反周期环上是带对应边界的$\delta$核，不同点为0。

(c)$\tau>0$求正模几何和得到$G=L^{-1}/[\e^{-\ii\pi(x+\ii v\tau)/L}-\e^{\ii\pi(x+\ii v\tau)/L}]$，即第\ref{ch:chiral}章的$\ii/[2L\sin(\pi z/L)]$。$L\to\infty$得$\ii/(2\pi z)=(-2\pi\ii)^{-1}/z$。$\tau<0$以时间排序的解析延续处理，不能把正模和直接当所有$\tau$收敛的级数。
\end{solution}
\originalsection{作业8}
\begin{exercise}[原题1(a)--(c)]
$U(1)$紧致相位低能Lagrangian为$L=\int_0^L\dd x[(\partial_t\theta)^2/(4\pi V_2)-V_1(\partial_x\theta)^2/(4\pi)+\rho_0\partial_t\theta]$，$\rho_0=|\bar\phi|^2/a$，$\hbar=1$。(a)零模$\theta=2\pi mx/L+\theta_0$求每$(m,N)$扇区最低能。(b)求总玻色子数。(c)平移引起相位$2\pi m\Delta x/L$，求总动量；也可取$\theta=2\pi m[x-x_0(t)]/L+\theta_0$由共轭动量求之，解释混合异常。
\end{exercise}
\begin{solution}
零模$p_0=L\dot\theta_0/(2\pi V_2)+\rho_0L$。紧致角要求$p_0=N\in\mathbb Z$，所以$\dot\theta_0=2\pi V_2(N-\rho_0L)/L$、$E_{m,N}=\pi V_2(N-\rho_0L)^2/L+\pi V_1m^2/L$，$N$就是该扇区的粒子数。固定$m$的最低态选择最接近$\rho_0L$的整数$N$；半整数位置可简并，不能擅自当$N$连续。

平移使态获$\e^{\ii N2\pi m\Delta x/L}$，相应$P=2\pi mN/L$按该主动平移约定。用$x-x_0$的参数则$p_{x_0}=-(2\pi m/L)N$，符号来自把图样位置与场主动平移相反定义；一致恢复同物理动量。改变$U(1)$扭曲$m$泵送平移量子数，即混合异常，Berry总导数在量子零模里有实际作用。
\end{solution}
\begin{exercise}[原题2(a)--(e)]
环上$H_0=\sum_i[-t(b_i^\dagger b_{i+1}+\mathrm{h.c.})+\mu n_i+Vn_in_{i+1}]$，整数粒子数$N_0$、平移$Tb_iT^{-1}=b_{i+1}$，基态$T\ket\Phi=\e^{\ii K_0}\ket\Phi$。定义$U_\theta=\exp[(\ii\theta/L)\sum_{i=0}^{L-1}in_i]$、$\ket{\Phi_\theta}=U_\theta\ket\Phi$。(a)$\theta=2\pi r$时求动量差。(b)求$U_\theta^{-1}H_{i,i+1}U_\theta$，说明跨边界。(c)求其二阶导数阶数。(d)利用$\pm2\pi$变分能均不低于基态证明每键能差$O(L^{-2})$。(e)推总能差$O(L^{-1})$与既约填充$p/q$的$q$个近简并低能态；讨论无理填充。题页“spin1/2”与标准玻色对易同时出现不相容，证明只用整数$n_i$和$[n_i,b_j]=-\delta_{ij}b_j$，对硬核或有限密度玻色模型分别注明界。
\end{exercise}
\begin{solution}
(a)$TU_\theta T^{-1}=U_\theta\e^{-\ii\theta N_0/L}\e^{\ii\theta n_0}$，整数$n_0$且$\theta=2\pi r$时最后因子1，所以$\Delta K=-\theta N_0/L$模$2\pi$，不同动量的态正交。

(b)内键$H_i^\theta=-t(\e^{\ii\theta/L}b_i^\dagger b_{i+1}+\mathrm{h.c.})+\mu n_i+Vn_in_{i+1}$。跨边界的纯规范变换相位实际为$\e^{-\ii\theta(L-1)/L}$，仅$\theta=2\pi r$时同$\e^{\ii\theta/L}$。要取任意$\theta$导数，可把内键函数解析延拓为均匀扭曲键；不能声称任意非整数通量都是全局纯规范。

(c)局域跃迁系数二阶导数为$O(L^{-2})$，在平移基态、有限密度$\rho$下$|\avg{b_i^\dagger b_{i+1}}|\le\sqrt{\avg{n_i}\avg{n_{i+1}}}=\rho$，所以每键能$\epsilon(\theta)$的二阶导数也$O(L^{-2})$。对普通无界玻色算符不是全Hilbert空间算符范数界，题目需要的是此期望界；硬核模型可直接用局域算符有界性。

(d)在$\theta=\pm2\pi$，扭曲态仍平移本征，能$L\epsilon(\pm2\pi)$。变分原理使两差$\epsilon(\pm2\pi)-\epsilon(0)\ge0$。Taylor正负相加消去一阶，和至多$C/L^2$，故各自也至多$C/L^2$，不必假定基态电流预先为0。(e)乘$L$得$O(1/L)$总差。既约$p/q$填充的$r=0,\ldots,q-1$扭曲动量相异，固定$q$时都有消失能差，排除唯一有隙对称基态。无理填充极限可产生任意多相异低能动量；若要求有限基态简并且保持平移，必须无隙。原题的“无论相互作用都无隙”还依赖局域、有限密度、这些对称与热力学极限条件，不能扩展到非局域或无限简并的任意模型。
\end{solution}

\originalchapter{8.513作业9--11：边缘谱与拓扑质量}
\psource{Xiao-Gang Wen，8.513 Fall 2021，公开作业9--11。}
\originalsection{作业9}
\begin{exercise}[原题1(a)--(c)]
石墨烯两谷手性相反。普通亚格势使两谷质量同号，Haldane质量使两谷分别$\pm\sqrt3t'$；只使K'质量跨过0可构造平凡与Chern绝缘体界面。循环置换Pauli矩阵后取$\hbar v=1,m(y)=y$，
\begin{equation}
H=-\ii(\sigma_z\partial_x-\sigma_y\partial_y)+y\sigma_x=\begin{pmatrix}-\ii\partial_x&\partial_y+y\\-\partial_y+y&\ii\partial_x\end{pmatrix}.
\end{equation}
(a)取$\Psi=\e^{\ii kx}\Psi_k(y)$，分别找只有上或下分量、Gaussian形式的解，去掉不可归一支。(b)求可归一边缘色散。(c)求并画其余带。
\end{exercise}
\begin{solution}
(a)只有上分量时$(-\partial_y+y)\psi_1=0$，$\psi_1\propto\e^{y^2/2}$不可归一，形式能$E=k$；只有下分量时$(\partial_y+y)\psi_2=0$，$\psi_2=\pi^{-1/4}\e^{-y^2/2}$可归一，(b)$E_0(k)=-k$。因此只有一支物理手性低能态。

(c)平方$H_k$得两对角$k^2-\partial_y^2+y^2+1$与$k^2-\partial_y^2+y^2-1$。上分量的振子$n-1$与下分量$n$配对，$n\ge1$时$E_{n,\pm}(k)=\pm\sqrt{k^2+2n}$。最低下分量$n=0$没有可归一上分量伙伴，故不组成$\pm|k|$两支。向量图如下，坐标为本题无量纲单位。
\begin{center}\begin{tikzpicture}\begin{axis}[width=10cm,height=5cm,axis lines=middle,xlabel={$k$},ylabel={$E$},xmin=-2.5,xmax=2.5,ymin=-3.2,ymax=3.2]
\addplot[structurecolor,thick,domain=-2.5:2.5]{-x};
\foreach \n in {1,2}{\addplot[gray,domain=-2.5:2.5,samples=80]{sqrt(x*x+2*\n)};\addplot[gray,domain=-2.5:2.5,samples=80]{-sqrt(x*x+2*\n)};}
\end{axis}\end{tikzpicture}\end{center}
线性质量无限增大是便于精确计算的模型，真实界面质量饱和时高能束缚带数及连续谱不同，但净低能手性仍受体拓扑差保护。
\end{solution}
\originalsection{作业10}
\begin{exercise}[原题1(a)--(c)]
开链Majorana模型$H=\ii J\sum_{i=1}^{L-1}\lambda_i^y\lambda_{i+1}^x+\ii h\sum_{i=1}^L\lambda_i^x\lambda_i^y$，$\{\lambda_i^\alpha,\lambda_j^\beta\}=2\delta_{ij}\delta_{\alpha\beta}$，$c_i=(\lambda_i^x+\ii\lambda_i^y)/2$。(a)$J=0,h\ne0$证明多体基态不简并。(b)$h=0,J\ne0$证明二重简并。(c)一般$J,h$求单粒子级，并核对两个极限。原题提示普通正弦波可用于部分边界处理，须核对真实Majorana开链端点。
\end{exercise}
\begin{solution}
(a)$\ii\lambda_i^x\lambda_i^y=2n_i-1$，每站独立，$h>0$全部空、$h<0$全部满，基态唯一，激发代价$2|h|$；$J=h=0$是另一个全空间简并点。(b)用$\widetilde c_i=(\lambda_i^y+\ii\lambda_{i+1}^x)/2$配对$L-1$个键，余$\lambda_1^x,\lambda_L^y$形成零能复费米子，所以两种占据同能。若物理设置固定总费米宇称，只可选其中一个，二重是在允许两宇称的多体空间中陈述。

(c)写$H=\ii\sum_{ij}a_iD_{ij}b_j$，$a_i=\lambda_i^x,b_i=\lambda_i^y$，$D_{ij}=h\delta_{ij}-J\delta_{i,j+1}$。正准粒子能$\epsilon_n=2s_n(D)$，$s_n$为奇异值。这给所有参数与有限$L$的完整、稳定计算法。$D^TD$对角前$L-1$个为$h^2+J^2$、最后为$h^2$，邻对角$-hJ$。对$hJ\ne0$，内部解$f_j=\sin kj$、
\begin{equation}
\epsilon(k)=2\sqrt{h^2+J^2-2hJ\cos k},\qquad h\sin[(L+1)k]=J\sin(Lk).
\end{equation}
因此一般$k$并非$n\pi/(L+1)$，这类简单正弦格点会漏掉末站特殊边界及拓扑端模。所有实根连同必要的虚$k=\ii\kappa$或$\pi+\ii\kappa$给恰$L$个正能（零模极限取连续）。$h=0$有一零奇异值、其余$|J|$，$J=0$全部$|h|$，与(a)(b)匹配。$|h/J|<1$大链最低能$\epsilon_0\simeq2|J|[1-(h/J)^2]|h/J|^L$，故一般有限链并非精确二重，而是指数近简并；无限链在$|h|=|J|$闭体隙。
\end{solution}
\originalsection{作业11}
\begin{exercise}[原题1(a)--(d)]
三维粒子数守恒、$T^2=-1$的自由费米体系，$M(k)=v(k)\cdot\sigma$在两零点$k_R,k_L$附近分别为$\delta k\cdot\sigma$与$-\delta k\cdot\sigma$。(a)写四分量连续Dirac哈密顿量。(b)转为实Majorana基，令$\varepsilon=\ii\sigma_2=\left(\begin{smallmatrix}0&1\\-1&0\end{smallmatrix}\right)$，
\begin{equation}
\gamma_1=\sigma_0\otimes\sigma_1\otimes\sigma_3,\quad\gamma_2=\varepsilon\otimes\varepsilon\otimes\sigma_3,\quad\gamma_3=\sigma_0\otimes\sigma_3\otimes\sigma_3,
\end{equation}
$Q=\varepsilon\otimes\sigma_0\otimes\sigma_0$、$T=\sigma_3\otimes\varepsilon\otimes\sigma_0$，找实反对称质量$A$满足$\{A,\gamma_i\}=\{A,T\}=0,[A,Q]=0$（$TQ$关系随之）。(c)解释加入空间动能与对称生成元后的Clifford扩张如何约为题页所用$R_0^1$质量空间。(d)若$A^2=-I$、$\{A,\rho\}=0$、$\rho=\operatorname{diag}(I_n,-I_n)$，证明$A=\left(\begin{smallmatrix}0&-O^T\\O&0\end{smallmatrix}\right)$、$O\in O(n)$，从而有两个连通分支。源题为稳定自由费米分类，不能用此直接分类任意相互作用相。
\end{exercise}
\begin{solution}
(a)以最后指标区分两谷，按原题Fourier$i\partial\leftrightarrow k$约定$H=\sigma_1\otimes\sigma_3\ii\partial_x+\sigma_2\otimes\sigma_3\ii\partial_y+\sigma_3\otimes\sigma_3\ii\partial_z$。若使用$\e^{\ii kx}$定义，$i\partial\leftrightarrow-k$且两谷标签相应反换，谱与分类不变。

(b)先把虚数显式提出，$\sigma_2=-\ii\varepsilon$，再在实化新增第一指标上将$-\ii$换为$\varepsilon$、实1换$\sigma_0$，得到原题三个实对称$\gamma_i$。直接Pauli乘法给$\{\gamma_i,\gamma_j\}=2\delta_{ij}$，$Q^2=T^2=-I,QT=-TQ$。把$A$展开在$\{I,\sigma_1,\sigma_3,\varepsilon\}^{\otimes3}$的64个实独立矩阵中，逐个施加反对称、$Q$对易及四个反对易条件，只剩
\begin{equation}
A=m\,\varepsilon\otimes\sigma_0\otimes\sigma_1.
\end{equation}
每个条件的矩阵积仍是该正交基的一个带符号元素，故不存在靠不同基元相消而遗漏的线性组合。可直接核对这个$A$的全部条件及$A^2=-m^2I$；随卷矩阵脚本又从64维线性条件独立核对零空间维数1。

(c)先固定记号，课程第16--19讲第29页的$R_q^p$是满足$A^T=-A,A^2=-I$的实质量矩阵空间；它与$p$个平方$+I$、$q$个平方$-I$的两两反对易正交生成元$\rho_i$逐个反对易。正平方生成元为$\gamma_1,\gamma_2,\gamma_3$，负平方生成元为$T,TQ$：从$Q^2=T^2=-I$、$QT=-TQ$可验$(TQ)^2=-I,\{T,TQ\}=0$。$\gamma_i$与$Q$对易、与$T$反对易，因此也与$TQ$反对易。条件$[A,Q]=0,\{A,T\}=0$推出$\{A,TQ\}=0$。故归一质量的扩张确实是$R_2^3$，并非把$Q$错计为第三个负平方生成元。

说明两次约化的具体矩阵操作。设$(\widetilde A,\widetilde\rho_1,\ldots)$代表$R_q^p$，令$\varepsilon=\left(\begin{smallmatrix}0&1\\-1&0\end{smallmatrix}\right)$，则
\begin{equation}
A=\widetilde A\otimes\sigma_3,\quad
\rho_i=\widetilde\rho_i\otimes\sigma_3,\quad
\rho_+=I\otimes\sigma_1,\quad
\rho_-=I\otimes\varepsilon
\end{equation}
定义$R_{q+1}^{p+1}$中的一组矩阵：$\rho_+^2=I,\rho_-^2=-I$，$\sigma_3$与$\sigma_1,\varepsilon$均反对易，故所有要求逐条保持。反过来，在稳定扩带后把指定的一正一负生成元选为$\rho_+,\rho_-$的标准实表示，其余生成元和$A$必有$\widetilde X\otimes\sigma_3$的形式；这样移去这一正一负对。第一次从$R_2^3$移到$R_1^2$，第二次从$R_1^2$移到$R_0^1$。最后只余一个正平方$\rho$与实反对称$A$，正是(d)的问题，$R_2^3\simeq R_1^2\simeq R_0^1$。这里$\simeq$指稳定质量空间等价，允许增加平凡带，并不声称原8维矩阵空间本身等于任意$n$的$O(n)$；也没有借此宣称证明一般Bott周期性。

(d)反对易使$A$对角块为0，反对称使两个非对角块互为负转置，写$A=\left(\begin{smallmatrix}0&-O^T\\O&0\end{smallmatrix}\right)$。$A^2=-\operatorname{diag}(O^TO,OO^T)=-I$要求$O$正交，反之每个正交矩阵均给合法$A$。$\det O=\pm1$沿连续路不变；各定号分支均连通（$n\ge2$可分解连续平面旋转，$n=1$各是单点），所以$\pi_0=\mathbb Z_2$。这给三维AII强拓扑的两个稳定自由相；仅用未正则化Dirac质量的正负不能判断哪一支是某个晶格模型的绝对平凡相，须给完整紫外晶格参照。
\end{solution}
\originalsection{学期论文的公开要求}
作业8及作业页说明论文约四页，可为短综述或一个实际研究问题，主题示例包括磁序、自旋波、Berry输运、Chern绝缘体、石墨烯、Weyl半金属、分数量子Hall、共形场、Haldane链、Majorana与量子临界。公开底稿给12月9日截止、建议11月30日前提交草稿，这是2021课的历史安排。论文选题属于开放任务，未公开标准答案，本册不把编者虚构论文计入原课程交付题数，也不代学生发送邮件或上传Canvas。

\originalchapter{来源、覆盖与阅读边界}
本册以实际公开材料为范围：8.511的24份摘要对应第1--6章，8.512的25讲主题对应第7--12章，8.513的四组幻灯片对应第13--18章。相同物理主题合并推导，课程内部进度仍保留在来源清单。讲义是概念重编，不能把章节数当成原讲数，也没有把英文PDF整页贴入正文。

8.511：1--4讲见第\ref{ch:crystal}章，5--10见第\ref{ch:bands}章，11--14见第\ref{ch:transport}章，15--18见第\ref{ch:interaction}及\ref{ch:phonon}章，19--21见第\ref{ch:phonon}章，22--24见第\ref{ch:bcs}章。8.512：1--4讲见第\ref{ch:response}章，5--8见第\ref{ch:disorder}章，9--13见第\ref{ch:scresponse}章，14--19见第\ref{ch:magnet}章，20--22见第\ref{ch:kondo}章，23--25见第\ref{ch:fermi}章。8.513：1--6讲见第\ref{ch:berry}章，7--15见第\ref{ch:boson}--\ref{ch:chiral}章，16--19见第\ref{ch:spt}章，20--24见第\ref{ch:topological}章。

原始文件登记共92份、547个PDF物理页，包括摘要、讲义、当学期及2004补充作业。不同文件名的Kondo讲义同文，幻灯片含重复动画页，这个页数不等于不重复教学页数。第三方权利例外、可取得材料与商业书题号引用分开登记；本稿不收录例外图片、商业教材整书、未公开测验卷、学生论文或课堂视频逐字稿。严格一般Bott周期性与高阶范畴分类的证明未纳入；已说明调用与物理应用范围。


\originalsection{考核覆盖与缺项}
按每套原印刷顶层编号计，8.511为19个位置（含一个阅读指令），8.512当期22个、公开2004补充25个、8.513为23个，合计89个原位置。2004补充中15个位置与当期题相同，保留定位并回指完整题解；63个中文习题环境合并书写相关原题，但没有以合并环境数代替原题数。72个不重复的可获计算题位置均给现有公开要求的中文题面和解答。

84个原位置（含重复）全完成，3个混合位置完成公开可核部分，1个仅商业书指针缺题面，1个为阅读要求：8.511 PS2题1开头的Marder第38页2.5(b)未取得，公开补充(a)、(b)已解；PS7题1仅Marder第17章题6与热流提示，未补造；PS3题1(e)算出自由费米能，商业书KKR图比较未完成；PS4题1(e)给所指定二维拟合点能量，未取得完整三维APW曲线核对。PS8题2是阅读Marder附录C，正文已建立二次量子化先修，不计虚构题。原卷缺失的8.511测验、8.513论文实例与答案不在交付中。

原PDF均曾实际下载并核对页数、字节与SHA-256，公开仓库提供官方原件链接及元数据，不重分发混有第三方图的原件。中文PDF的图是原生TikZ/PGFPlots；源码ZIP不含原英文PDF、字体、凭据或运行缓存。8.512第一套抽取文字失败的页面已实际看图恢复；原题中的响应符号、宇称边界和印刷误植在解答处明确核查。数值曲线是固定种子与有限样本的实际计算，不能等同热力学极限证明。

\begin{thebibliography}{9}
\bibitem{lee1} Patrick A. Lee，MIT OCW，8.511 Theory of Solids I，Fall 2004。\url{https://ocw.mit.edu/courses/8-511-theory-of-solids-i-fall-2004/}
\bibitem{lee2} Patrick A. Lee；Mark Spencer Rudner（公开页署名讲义1--10），MIT OCW，8.512 Theory of Solids II，Spring 2009。\url{https://ocw.mit.edu/courses/8-512-theory-of-solids-ii-spring-2009/}
\bibitem{wen} Xiao-Gang Wen，MIT OCW，8.513 Modern Quantum Many-body Physics for Condensed Matter Systems，Fall 2021。\url{https://ocw.mit.edu/courses/8-513-modern-quantum-many-body-physics-for-condensed-matter-systems-fall-2021/}
\bibitem{license} MIT OCW使用条款及逐件例外。\url{https://ocw.mit.edu/pages/privacy-and-terms-of-use/}。基于OCW的中文改编：CC BY-NC-SA 4.0，\url{https://creativecommons.org/licenses/by-nc-sa/4.0/}。
\end{thebibliography}

\end{document}
